\documentclass[10pt]{amsart}
\usepackage{amsmath,amssymb,amsthm,mathtools,mathrsfs}
\usepackage[a4paper,hmargin=0.76in,vmargin=0.60in]{geometry}
\usepackage{booktabs}
\usepackage{array}
\usepackage{longtable}
\usepackage{tabularx}
\usepackage[hypertexnames=false]{hyperref}
\hypersetup{hidelinks}
\usepackage[final,expansion=false]{microtype}
\usepackage{url}
\emergencystretch=2.5em

\newtheorem{theorem}{Theorem}[section]
\newtheorem{proposition}[theorem]{Proposition}
\newtheorem{lemma}[theorem]{Lemma}
\newtheorem{corollary}[theorem]{Corollary}
\theoremstyle{definition}
\newtheorem{definition}[theorem]{Definition}
\newtheorem{problem}[theorem]{Problem}
\theoremstyle{remark}
\newtheorem{remark}[theorem]{Remark}

\newcommand{\C}{\mathbb C}
\newcommand{\R}{\mathbb R}
\newcommand{\N}{\mathbb N}
\newcommand{\Z}{\mathbb Z}
\newcommand{\Q}{\mathbb Q}
\newcommand{\Rel}{\operatorname{Re}}
\newcommand{\Iml}{\operatorname{Im}}
\newcommand{\supp}{\operatorname{supp}}
\newcommand{\TV}{\operatorname{TV}}
\newcommand{\Kp}{\mathcal K^{\mathrm{prim}}_{2T}}
\newcommand{\Kw}{\mathcal K_{2T}}
\newcommand{\PT}{\mathcal P_T}
\newcommand{\IT}{I_T}
\newcommand{\lamP}{\lambda^{\mathrm{prim}}_T}
\newcommand{\lamW}{\lambda^{\mathrm W}_T}
\newcommand{\lamA}{\lambda_{\mathrm A}(T)}
\newcommand{\sT}{\sigma_T}
\newcommand{\gh}[1]{\widehat g_{#1}}
\newcommand{\nuh}{\widehat\nu}
\newcommand{\wh}{\widehat w}
\newcommand{\Fh}{\widehat F}
\newcommand{\cF}{c_F}
\newcommand{\Lev}{\mathfrak E}
\newcommand{\Wcl}{\mathcal W_T}
\newcommand{\NT}{\mathcal N_T}
\newcommand{\Aconst}{\mathsf A_T}
\newcommand{\Ainf}{\mathcal A_\infty}
\newcommand{\PiT}{\Pi_T}
\newcommand{\mT}{\mathfrak m_T}
\newcommand{\laminf}{\lambda^{\infty}_{\min}}
\newcommand{\laminfP}{\lambda^{\infty,\mathrm P}_{\min}}
\newcolumntype{Y}{>{\raggedright\arraybackslash}X}

\providecommand{\anPa}{\Phi_a}
\providecommand{\anPb}{\Phi_b}
\providecommand{\aneu}{\mathrm{eu}_d}
\providecommand{\anrd}{\mathrm{rd}}
\providecommand{\auAz}{\mathcal A_0}
\providecommand{\auNm}{\operatorname{Nm}}
\providecommand{\auO}{\mathcal O}
\providecommand{\auP}{\mathcal P}
\providecommand{\auPl}{\Phi}
\providecommand{\auS}{\mathcal S}
\providecommand{\auU}{\mathrm U}
\providecommand{\pbLM}{\mathrm{LM}}
\providecommand{\pbW}{\mathbb W}
\providecommand{\pbeB}{\mathsf B}
\providecommand{\pbel}{\ell}
\providecommand{\pbels}{\tilde\ell}
\providecommand{\pbeq}{\mathsf q}
\providecommand{\pbip}[2]{\langle#1,#2\rangle}
\providecommand{\pbmD}{\mathsf D}
\providecommand{\pbmPhi}{\varphi}
\providecommand{\pbmdec}{\mathsf d}
\providecommand{\pbmnuS}{\nu_*}
\providecommand{\pbnat}{\natural}
\providecommand{\pboF}{\mathsf F}
\providecommand{\pboR}{R^0}
\providecommand{\pboU}{U^0}
\providecommand{\pboV}{\operatorname{vol}\mathrm U}
\providecommand{\pbone}{\mathbf 1}
\providecommand{\pbrpsi}{r_\psi}
\newenvironment{quotedthm}[2]{\par\medskip\noindent\begingroup
  \textbf{#1}#2.\ \itshape\ignorespaces}{\endgroup\par\medskip}
\newenvironment{quoteddef}[2]{\par\medskip\noindent\begingroup
  \textbf{#1}#2.\ \ignorespaces}{\endgroup\par\medskip}
\usepackage{xurl}
\begin{document}
\addtolength{\abovedisplayskip}{-3.4pt}
\addtolength{\belowdisplayskip}{-3.4pt}
\addtolength{\belowdisplayshortskip}{-2pt}
\title[The Spanning Theorem: local computations]{The Spanning Theorem: sieve systems at the real and finite places and Siegel's step}
\author{David Alejandro Trejo Pizzo}
\date{}
\begin{abstract}
In the proof of the Spanning Theorem of this series the main term of the second moment of a
Poincar\'e map, on a unitary group over a CM field, is obtained from the open cell of a relative
trace formula by Siegel's formula followed by repeated Poisson summation. We prove the local
statements on which this rests. At every real place, for the Bessel function of a discrete series
with a large parameter, and at every finite set of finite places, for the characteristic function
of a congruence subgroup twisted by a unipotent character, we construct a \emph{sieve system}: a
family of forms on pairs of sequences of numbers, linked by integration over one entry, positive
at the top, of prescribed mass at the bottom and decaying in the frequency. The construction uses
a theory of correlations of a Bessel function along two frames which is common to all places. We
also prove Siegel's step, which bounds sums of squares of orbital integrals over the rational
points of a sphere by the top forms, with the constant of Siegel's formula.
\end{abstract}
\maketitle
\setcounter{tocdepth}{1}\tableofcontents

\section{Introduction}\label{B2:intro}

The main paper~\cite{P46} of this series deduces the Riemann hypothesis for Dirichlet
$L$-functions from three theorems on integral algebraic modular forms on a chain of theta lifts
between unitary groups over a CM field $E$. The second of them, the Spanning Theorem, rests on an
estimate for the second moment of a Poincar\'e map $\auP$ from algebraic modular forms on a
definite unitary group $\auU(X_0)$ of rank $k$ to automorphic forms on the unitary group $G$ of
the space $Y_0=X_0\perp\mathbb H^{\varrho}\perp\langle-N\rangle$ ({\cite[Proposition~14.32]{P46}}).
This estimate is proved in~\cite{P51} by a relative trace formula, and its main term must be
obtained with the exact constant. In~\cite{P51} the main term is deduced from three statements
({\cite[Theorem~3.51]{P51}}): the existence of a \emph{sieve system} at every real place and at
every suitable finite set of finite places, and an inequality called \emph{Siegel's step}. We prove
these three statements here, as Theorems~\ref{pfr:R}, \ref{pff:F}
and~\ref{pff:S}.

Let $e_a,e_a^*$ $(a<\varrho)$ be the hyperbolic pairs of $\mathbb H^{\varrho}$ and $z$ the vector
of norm $-N$; the \emph{frame} is $(f_0,\dots,f_\varrho)=(e_0,\dots,e_{\varrho-1},z)$, and
$n=\varrho+1$. Let $U\subset G$ be the unipotent group of the elements that act unitriangularly
on the span of the $e_a$ and trivially on the quotient of its orthogonal by it, with its generic
character, and let $G^\flat$ be the unitary group of $X^+=X_0\perp Ez$. A \emph{head} is a
sequence $p=(p_0,p_1,\dots)$ of invertible numbers; it defines the Weyl element $w_p$ with
$w_pe_a=p_ae_a^*$. An element $\gamma$ of the open cell is $u\,w_p\,\sigma\,u'$ with $u,u'\in U$,
a head $p$ of length $\varrho$ and $\sigma\in G^\flat$, and the number $t=h(\sigma z,z)$ is its
\emph{last entry}; we regard $(p,t)$ as a head of length $n$. In the relative trace formula
of~\cite{P51} the contribution of the rational elements $\gamma=w_p\sigma$ is a product of orbital
integrals: at a real place $\tau$ the integral $\int W^\flat_\tau(y)\otimes
\overline{W^\flat_\tau(\gamma^{-1}y)}\,dy$ of a Bessel function $W^\flat_\tau$ on the real group,
obtained from the coefficient of the lowest $K$-type of a discrete series and cut off at a fixed
distance from the Bessel subgroup; at the finite places a Kloosterman sum, which is the orbital
integral of the function $\chi_M$ equal to the character on $U$ times the principal congruence
subgroup of level $M$ and to zero elsewhere.

A sieve system ({\cite[Definition~3.44]{P51}}) lives at a \emph{place} $(C,dc,\chi,\iota)$: a
locally compact ring with a Haar measure, a character $\chi$ of its additive group and a map
$\iota$ on its units. It is a family of functions $J_i$, $0\le i\le n$, on pairs of heads of
length $i$ with entries in $C^\times$, with a \emph{mass} $\mathsf E$, \emph{weights} $\Lambda_i$
and a \emph{decay} $\beta$, such that: the integral of $J_{i+1}((p,c),(q,c))$ over $c\in C$ is
$J_i(p,q)$, and its Fourier transform at a unit $\xi$ is at most $\beta(\xi)\Lambda_i(p)
\Lambda_i(q)$ in absolute value; the top form $J_n$ is continuous and, for a fixed last entry,
hermitian and positive semidefinite in the heads; $|J_0|\le\mathsf E$; and $|J_i(p,q)|\le
\beta(\xi)\Lambda_i(p)\Lambda_i(q)$ when the last entries of $p$ and $q$ differ, with
$\xi=\iota(p_{i-1})-\iota(q_{i-1})$ a unit. Given such systems, $n$ successive Poisson summations
over the entries of the heads bound the sum of the top forms over rational heads by the mass
times an explicit constant, up to a small error; this is the sieve of~\cite{P51}
({\cite[Theorem~3.50]{P51}}). The mass is
therefore the quantity that has to be exact.

Theorem~\ref{pfr:R} concerns a real place $\tau$. Under hypotheses (H0)--(H5) on the sizes and on
the Bessel functions, which hold for the data of~\cite{P46} and the Bessel functions
of~\cite{P51} (\S\ref{prp:ss-H}), there is a sieve system $(J^\tau_i)$ at the place $(\C,
\text{Lebesgue measure},c\mapsto e^{c-\bar c},w\mapsto w^{-1})$ whose top form is the Gram form
\[
J^\tau_n(x,x')=\bigl\langle\mathsf N_\sigma(x),\mathsf N_\sigma(x')\bigr\rangle,\qquad
\mathsf N_\sigma(x)=c(x)\int_{G_\tau}W^\flat_\tau(y)\otimes\overline{W^\flat_\tau\bigl((w_x\sigma)^{-1}
y\bigr)}\,dy ,
\]
of the orbital integrals with an explicit normalising factor $c(x)\ge0$, for $\sigma\in G^\flat$
at $\tau$ and heads $x,x'$ of length $n$ with the same last entry $h(\sigma z,z)$. Its mass is $(1+b^{-10})/\deg_\tau$, where
$\deg_\tau$ is the formal degree of the discrete series at $\tau$, and its decay is
\[
\beta^\infty(\xi)=\Theta\,b^{2000n^9}\,|\xi|^{-(k-n^7)}\min(1,|\xi|^{-1})^6
\]
with an explicit constant $\Theta$; the weights are explicit functions of the partial products of
the entries. Here $b$ is one of the scales of~\cite{P46}, and $n\approx\sqrt x$, $k\approx x^5$ in
terms of another, much smaller one.

Theorem~\ref{pff:F} concerns a finite set $T$ of finite places of $E$, stable under conjugation
and containing the places above $MN$, a set $\mathsf A$ of the primes above $p_0$ at which the
level is imposed, and a non-zero vector of the span of the frame modulo each prime of
$\mathsf A$. It gives a group $\mathsf Z$ between the principal congruence subgroup and the
stabiliser of the lattice at $T$, containing the integral elements whose reductions fix the
given vectors, and a sieve system $(J^T_i)$ at the place $(E_T,dx,\psi_T,w\mapsto M^{-1}w^{-1})$,
built from the translates of $\chi_M$ by $\mathsf Z$, whose mass is
\[
(1+b^{-6})\,b^{-6k\,\#\mathsf A},
\]
with a decay which, for a frequency $\xi$, is essentially the power $(k-n^7)/2$ of the absolute
value of $\xi^{-1}$ in $E_T$. Thus every vector that is required to be fixed gains a factor
$b^{-6k}$ in the mass; in~\cite{P51} this is what makes an induction on the set $\mathsf A$
possible.

Theorem~\ref{pff:S} is Siegel's step. Let $(c_j,c_j',s_j,\sigma_j)$ be finitely many distinct
quadruples of classes $c_j,c_j'$ of the genus of the lattice of $X_0$ with its level, elements
$s_j$ of the finite group through which the level acts that fix the given vectors, and rational
elements $\sigma_j$ of $G^\flat$ that are integral outside $T$ relative to the classes and whose
last entries $t_j$ are larger than $N$ at every infinite place. Let $\mathcal O_{c,s,c'}(\gamma)$
be the sum of the contributions to the relative trace formula of the terms in
$U(E^+)\gamma U(E^+)$, let $\mathcal T$ be the finite set of the possible entries, and suppose
that $T$ contains the places at which the possible minors of the frame matrices are not units.
Then
\[
\mathsf r_*^{\,2}\sum_{j}\Bigl\|\sum_{p\in\mathcal T^\varrho}\mathcal O_{c_j,s_j,c'_j}(w_p\sigma_j)
\Bigr\|^2\le(1+b^{-6})\,\mathsf a_*\,\Rel\sum_{t\in\{t_j\}}\ \sum_{p,q\in\mathcal T^\varrho}
\prod_\tau J^\tau_n\bigl([p,t]_\tau,[q,t]_\tau\bigr)\,J^T_n\bigl((p,t),(q,t)\bigr),
\]
where $[p,t]_\tau$ is the head $(p,t)$ at the place $\tau$ in normalised coordinates,
$\mathsf r_*$ is a normalising constant of the orbital integrals, and
$\mathsf a_*$ is, up to the orders of two finite groups, the number of classes times the
constant of Siegel's formula. The three theorems are stated with exactly the constants that are
used in~\cite{P51}, where the product of $\mathsf a_*$, of the masses and of the constants of the
Poisson summations is shown to be the main term of the second moment.

We turn to the proofs. The forms below the top are defined by the integration that a sieve system requires: $J_i$ is
$J_n$ integrated over the last $n-i$ entries. The group $U$ is a tower of Heisenberg groups, one
for each vector $e_a$ of the frame, which we call the \emph{floors}, and Fourier inversion along
these groups (Lemma~\ref{pfr:descent}) turns $J_i$ into an integral, over a \emph{position}
$x\in G$ and over the floors $i$ and above, of the product of two \emph{correlations} of the
Bessel function, that is, of integrals of $\langle W(g),W(gx)\rangle$ over translates of $g$, one
for each of the two frames given by the heads $p$ and $q$. At the bottom $J_0$ is the integral of
$|\langle W,R(x)W\rangle|^2$ over all positions, the square of the $L^2$ norm of a matrix
coefficient of the right regular representation; at a real place this is the inverse of the
formal degree, up to the effect of the cut, and at the finite places it is the integral of
$|\int_{\mathsf Z}\chi_M(ygy^{-1})\,dy|^2$ over $g$. This explains the two masses. What has to be
proved is that these integrals converge, that the positions far from the identity contribute
little to the mass, and that a non-zero frequency, or two heads with different entries, produce
the decay.

Sections~\ref{pfr:sec} and~\ref{pfs:sec} do this for an abstract \emph{place datum}, which
covers the real place and the finite places at once. The basic notion is the \emph{depth} of a
vector $e_a$ of the frame at a pair of points $(g,gx)$: the inverse of the volume of the set of
the coordinates along $e_a^*$ of the vectors that are small at both points. We prove a bound for
the integrals over a ball of positions, the one-frame bound (Theorem~\ref{pfr:oneframe}); a gain
by a power of the inverse depth at a position with a deep vector (Theorem~\ref{pfr:deep}); a
gain, for a frequency $\xi$, by a power of the modulus of $\xi^{-1}$ whose exponent is of the
order of $k$ (Theorem~\ref{pfr:gain-char}); and the same gain for two heads whose entries differ
(Theorem~\ref{pfr:gain-heads}). The last two rest on Plancherel's
inequality on the quadric of the vectors of given norm and on Parseval's identity along a
Heisenberg group (\S\ref{pfs:ss-planch}). The positions are then cut into dyadic shells
according to their distance from the identity; a remote position has a deep vector, the gain
from the depth pays for the volume of the shell, and the shells can be summed
(Theorem~\ref{pfr:shells}). Sizes are measured throughout by a \emph{gauge}, a filtration of the
ring by compact sets with a multiplicative loss, so that the same proofs apply to $\C$ and to
$E_T$.

At a real place (Sections~\ref{pfr:sec} and~\ref{prp:sec}) it remains to verify the axioms of a
gauge and of a place datum, to compare the measures on the floors with Haar measure
(Lemma~\ref{pfr:haar}), to prove that far positions are deep (Lemma~\ref{pfr:deepreal}), and to
evaluate the mass: the positions near the identity give $1/\deg_\tau$ by the square
integrability of the Bessel function and the reproducing property, and the others are small
(Proposition~\ref{pfr:mass}). At the finite places (Sections~\ref{pff:sec} and~\ref{pfg:sec}) the
Bessel function is the vector of the translates of $\chi_M$, the depth is the index of an ideal,
and the mass of the full stabiliser of the lattice is at most $1+b^{-6}$ times the part that
comes from the positions in the principal congruence subgroup
(Corollary~\ref{pff:massK}). The gain $b^{-6k}$ for a fixed vector is obtained by comparing
$\mathsf Z$ with the full stabiliser: the stabiliser of a vector modulo a prime above $p_0$ has
index $p_0^{\,r}-1$, with $r$ the rank of $Y_0$, which makes the volume of $\mathsf Z$ small,
while its mass integral remains controlled by its main part (Theorem~\ref{pff:fixed}). The top
form at $T$ is a density on the quadric of the vectors of norm $-N$ of $X^+\otimes E_T$ with a
prescribed last entry (Proposition~\ref{pff:top}), constructed with the fibre measures of a
unimodular hermitian lattice (Proposition~\ref{pfg:fibres}).

For Siegel's step, the elements $\sigma$ with a given last entry $t$ correspond, through the
vector $\sigma z$, to the classes and to the vectors $v\in E^k$ with $v^*v=\alpha$,
$\alpha=(t\bar t-N^2)/N$. The square of the sum in the theorem is the value at such a point of a
non-negative function on the adelic sphere, which is the product of the Gram form of the orbital
integrals at the real places and of a pairing of orbital integrals at $T$
(\S\ref{pff:ss-siegel}). We then use Siegel's formula for the number of representations of a
totally positive number by the genus of the hermitian lattice with its level
(Theorem~\ref{pff:siegelformula}): for a non-negative locally constant function $\Phi$ on the
finite adelic points of $X_0$, invariant under the level,
\[
\sum_{c}\ \sum_{v\in E^k,\ v^*v=\alpha}\Phi(h_c^{-1}v)=\#\mathcal C\cdot|\operatorname{disc}E|^{1/4}
\prod_{\tau}\frac{\pi^k\,\tau(\alpha)^{k-1}}{(k-1)!}\cdot\delta(\Phi,\alpha),
\]
where $c$ runs over the set $\mathcal C$ of classes, with representatives $h_c$, and
$\delta(\Phi,\alpha)$ is the local density at the finite places. Since all terms are
non-negative, the formula bounds the sum over any subset of the points
(Proposition~\ref{pff:siegelsum}); the density outside $T$ is at most $1+b^{-6}$ times a product of volumes
(Lemma~\ref{pff:offT}), and the density at $T$ is the integral against the fibre measure which
defines $J^T_n$. For quadratic forms the formula is Siegel's~\cite{pff-Siegel35}, for hermitian forms
Braun's~\cite{pff-Braun41}; since every constant matters here, we derive it in
Proposition~\ref{pfg:siegel} from Weil's form of the Siegel--Weil
theorem~\cite{WeilSiegel} for the pair $(\auU(1,1),\auU(X_0))$, which is in the range of absolute
convergence because $X_0$ is anisotropic of rank $k\ge3$, checking that the measures are Tamagawa
measures and computing the archimedean density.

The notation of the data is that of~\cite{P46}. The setting, the Bessel functions with their
properties, the description of the open cell and of the orbit sums, the set of entries and the
notion of sieve system are those of~\cite{P51}, from which we quote in particular
{\cite[Theorems~3.16 and~3.22]{P51}}, {\cite[Propositions~3.18, 3.41 and~3.42]{P51}} and {\cite[Definition~3.44]{P51}}, as well as lemmas on Haar measure in the
coordinates of the open cell ({\cite[Lemmas~3.6, 8.80 and~8.81]{P51}}), the structure of the level ({\cite[Lemma~2.1]{P51}}) and Witt's theorem in the form of {\cite[Lemma~7.12]{P51}}. None of
these statements depends on the three theorems proved here: these are stated together in~\cite{P51} as {\cite[Theorem~3.51]{P51}} and are used there only for
the main term of the second moment ({\cite[Theorem~3.53]{P51}}), which is not used in this paper. When
we refer to the sieve, to the evaluation of the constants or to the main term
of~\cite{P51} ({\cite[Theorems~3.50 and~3.53]{P51}}, {\cite[Lemma~3.52]{P51}}), it is to compare constants, not to use them. In the other direction, the proof of {\cite[Lemma~3.6]{P51}} in~\cite{P51} uses Lemma~\ref{prp:heis} of this paper, whose proof does not refer to~\cite{P51}. So the two papers quote each other
without a circle. The statements quoted are reproduced in Appendix~\ref{B2:quoted}. The main
published result used is Weil's theorem, which we use as stated.

Section~\ref{pfr:sec} sets up the local theory of pairs of frames: places, floors and place data,
correlations and forms, depth and gauges, the statements of the shared estimates and the sum over
the shells; it then states Theorem~\ref{pfr:R} and proves it from local statements at the real
place. Section~\ref{pfs:sec} proves the shared estimates. Section~\ref{prp:sec} proves the local
statements at the real place and verifies the hypotheses of Theorem~\ref{pfr:R}.
Section~\ref{pff:sec} treats the finite places: the datum at $T$, the mass integral, the gain for
fixed vectors, the decay, the exact statements and the proofs of Theorems~\ref{pff:F}
and~\ref{pff:S}. Section~\ref{pfg:sec} proves the lemmas used there on gauges, on Fourier analysis
and fibre measures on hermitian lattices, and Siegel's formula.


\providecommand{\panot}[1]{\par\smallskip{\raggedright\noindent\textbf{[Not proved here:} #1\textbf{]}\par}\smallskip}
\providecommand{\pbdeg}{\operatorname{deg}}
\providecommand{\pbvol}{\operatorname{vol}}
\providecommand{\pbmTh}{\Theta}
\providecommand{\pfrG}{\mathsf G}
\providecommand{\pfrH}{\mathsf H}
\providecommand{\pfrU}{\mathsf U}
\providecommand{\pfrR}{\mathsf R}
\providecommand{\pfrM}{\mathsf M}
\providecommand{\pfrm}{\mathsf m}
\providecommand{\pfrK}{\mathsf K}
\providecommand{\pfrP}{\mathcal P}
\providecommand{\pfrCl}{\mathcal C}
\providecommand{\pfrI}{\mathcal I}
\providecommand{\pfrF}{\mathsf F}
\providecommand{\pfrB}{\mathsf B}
\providecommand{\pfrmod}{\operatorname{mod}}
\providecommand{\pfrdep}{\operatorname{dp}}
\providecommand{\pfrwt}{\wp}
\providecommand{\pfrkap}{\kappa}
\providecommand{\pfrS}{\mathsf S}
\providecommand{\pfrL}{\Lambda}
\providecommand{\pfrN}{\mathsf N}
\providecommand{\pfrX}{\Xi}
\providecommand{\pfrip}[2]{\langle#1,#2\rangle}
\providecommand{\pfrhw}{\mathsf w}
\providecommand{\pfrell}{\ell}

\section{Two frames: the local theory, and the sieve system at a real place}\label{pfr:sec}

This section concerns part \textup{(R)} of {\cite[Theorem~3.51]{P51}}, and the local theory of
\emph{pairs of frames} that parts \textup{(R)} and \textup{(F)} of that theorem share.

\emph{What is proved.}
\begin{itemize}
\item[(I)] \emph{The exact statement of \textup{(R)}.} Theorem~\ref{pfr:R} states \textup{(R)} with
its hypotheses. Its constants are \emph{exactly} those of {\cite[Theorem~3.51\,(R)]{P51}} as used in
{\cite[Theorem~3.50]{P51}}, {\cite[Lemma~3.52]{P51}} and {\cite[Theorem~3.53]{P51}}: mass $(1+b^{-10})/\pbdeg_\tau$,
weights $\Lambda^\infty$ with the radius $b^{60n^2}$ and the exponent $n^6$, decay $\beta^\infty$ with
the constant $\pbmTh\,b^{2000n^9}$ and the exponent $k-n^7$; the top form is the Gram form of the
normalised orbital integrals described in {\cite[\S3.3.7]{P51}}. The comparison is made term by term
in Remark~\ref{pfr:rem-match}. Two points have to be added to the statement printed in
{\cite[Theorem~3.51]{P51}}: the identification \textup{(R-b)} of the top form, which Siegel's step uses,
and the hypotheses on the data under which \textup{(R)} holds (Remark~\ref{pfr:rem-match}).
\item[(II)] \emph{The reduction of \textup{(R)} to four local estimates is proved in full}
(Theorem~\ref{pfr:R}, \S\ref{pfr:ss-Rproof}): positivity, continuity and descent of the forms, the
mass, the decay in the frequency and the decay for unequal entries.
\item[(III)] \emph{The mass is proved} (Proposition~\ref{pfr:mass}), from the bound for the remote
positions (Proposition~\ref{pfr:remote}) and the comparison of the floor measure with Haar measure
(Lemma~\ref{pfr:haar}), both proved in \S\ref{prp:sec}.
\item[(IV)] \emph{The shared theory.} The objects (place data, floors, correlations, forms,
majorants, depths, gauges, shells) are defined in \S\ref{pfr:ss-setting}--\S\ref{pfr:ss-gauge}. The
sum over the shells (Theorem~\ref{pfr:shells}) is proved in full from the one-frame bound
(Theorem~\ref{pfr:oneframe}) and the gains from a deep vector. Theorems~\ref{pfr:nogain} and~\ref{pfr:deep}
are proved (\S\ref{pfr:ss-gainsproof}), and so is Lemma~\ref{pfr:depth-rare}, with all the facts on
floors, quadrics and sections they use (\S\ref{pfr:ss-floors}, \S\ref{pfr:ss-section}); Lemma~\ref{pfr:descent}
is proved (\S\ref{pfr:ss-descent}). Theorems~\ref{pfr:oneframe}, \ref{pfr:gain-char} and \ref{pfr:gain-heads}
are proved in \S\ref{pfs:sec}, which also supplies the facts on sections and on the skew Haar measure used below.
\item[(V)] \emph{The real place.} Lemmas~\ref{pfr:realgauge}, \ref{pfr:deepreal} and
Propositions~\ref{pfr:conv}, \ref{pfr:top}, \ref{pfr:formbound}, \ref{pfr:remote} are proved in
\S\ref{prp:sec}, and so is Lemma~\ref{pfr:haar} (using {\cite[Lemmas~8.80 and~8.81]{P51}}).
\end{itemize}

\subsection{Places, floors and place data}\label{pfr:ss-setting}

\subsubsection*{Places}
A \emph{place} is a locally compact commutative ring $A$ with an involution $c\mapsto\bar c$, a Haar
measure $dc$, a continuous unitary character $\psi$ of $(A,+)$, a \emph{modulus}
$\pfrmod\colon A\to[0,\infty)$ (the factor by which multiplication by $c$ changes $dc$) and a
constant $P\ge1$ (the constant of Plancherel's inequality: axiom~(P) of \S\ref{pfs:ss-planch}). The two places used
are:
\begin{itemize}
\item the \emph{real place} $A=\C$, $\bar c$ complex conjugation, $dc$ Lebesgue measure,
$\psi(c)=e^{c-\bar c}$, $\pfrmod(c)=|c|^2$, $P=\pi^2$;
\item the \emph{finite place} $A=E_T=\prod_{w\in T}E_w$ of {\cite[\S3.3.3]{P51}}, with the measure
giving $\auO_T$ volume $1$, $\psi=\psi_T$, $\pfrmod(c)=|c|_T$; it is treated in
\textup{B-frames-finite}.
\end{itemize}
Each place also carries an element $\tau_0$ with $\tau_0+\bar\tau_0=1$ and a skew unit $\delta$
($\bar\delta=-\delta$): $\tau_0=\frac12$, $\delta=\frac i2$ at the real place, and the elements of
Lemma~\ref{pfg:psi}(b) at the finite place. Put $A^+=\{c:\bar c=c\}$ and $A_{\rm sk}=\{c:\bar c=-c\}$. We fix
the Haar measure $d\beta_{\rm sk}$ of $A_{\rm sk}$ of Lemma~\ref{pfs:skewhaar} (the one for which
$(\beta_1,\beta_2)\mapsto\tau_0\delta^{-1}\beta_1+\beta_2$ carries $d\beta_{\rm sk}\otimes d\beta_{\rm sk}$ to $dc$), and the
Haar measure $dh$ of $A^+$: Lebesgue measure of $\R$ at the real place, the measure giving $A^+\cap\auO_T$ volume $1$
at the finite place. The \emph{skew constant} $\kappa_A$ of the place is defined by: the isomorphism
$A^+\times A_{\rm sk}\to A$, $(h,b)\mapsto\tau_0h+b$, carries $dh\otimes d\beta_{\rm sk}$ to $\kappa_A\,dc$. At the
real place, $\tau_0\delta^{-1}=-i$, so $(it_1,it_2)\mapsto t_1+it_2$ is measure preserving and $d\beta_{\rm sk}$ is the
Lebesgue measure $dt$ of $\beta=it$; then $(h,it)\mapsto\frac h2+it$ has Jacobian $\frac12$, and $\kappa_\C=2$. At the
finite place $\kappa_{E_T}=1$, by Lemma~\ref{pfg:duality}(d) (whose proof shows that $\beta\mapsto\delta\beta$ carries
the measure of $A^+$ to that of $A_{\rm sk}$). At both places multiplication by $c\bar c$ ($c$ a unit) multiplies $dh$
by $\pfrmod(c)$ (at the finite place by Lemma~\ref{pfg:duality}(d), as $|c\bar c|_T^{1/2}=|c|_T$).

Hermitian forms follow convention (Cv1) of {\cite[\S2.1]{P51}} throughout this section: $h$ is linear in the
first variable, so the coordinate of a vector $x$ along $e_a^*$ is $h(x,e_a)$ and along $e_a$ it is $h(x,e_a^*)$.

\subsubsection*{The space and its floors}
Let $\varrho\ge1$, $n=\varrho+1$, and let $V$ be a free $A$-module with a basis
\[
e_0,\dots,e_{\varrho-1};\quad z;\quad (x_i)_{i\in\iota};\quad e^*_{\varrho-1},\dots,e^*_0
\]
and a hermitian form $h$ for which the $(e_a,e^*_a)$ are hyperbolic pairs
($h(e_a,e^*_b)=\delta_{ab}$, $h(e_a,e_b)=h(e_a^*,e_b^*)=0$), orthogonal to $z$ and to the $x_i$,
$h(z,z)=-1$ at the real place, and the $x_i$ are an orthonormal family orthogonal to $z$ (at the
finite place the middle block is the given unimodular lattice; see \textup{B-frames-finite}). Put
$k=\#\iota$ and $r=\operatorname{rank}V=2\varrho+k+1$. The \emph{frame} is $f_a=e_a$ $(a<\varrho)$,
$f_\varrho=z$. Let $\pfrG=\auU(V,h)$.

The \emph{floor} of $e_a$ and of $e_a^*$ is $a$, the floor of $z$ is $\varrho$ and the floor of the
$x_i$ is $n=\varrho+1$. For $l\ge0$ let $\pfrG^{(l)}$ be the subgroup of the $g\in\pfrG$ that fix the
basis vectors of floor $<l$; so $\pfrG^{(0)}=\pfrG$, $\pfrG^{(\varrho)}=\pfrG^\flat$ fixes the
$e_a,e_a^*$, and $\pfrH=\pfrG^{(n)}$ fixes them and $z$. Let $\pfrR^{(l)}$ be the group of the $g\in\pfrG$
that change each $f_a$, $a<l$, by a combination of the $f_b$, $b<a$ (that is, $gf_a\in f_a+\langle
f_b:b<a\rangle$), let $\pfrU=\pfrU_{(<\varrho+1)}$ be the Bessel unipotent group (the block upper
unitriangular elements for the flag of the frame, with middle block $1$), and for $u\in\pfrU$ let
$\sigma(u)$ be the sum of the entries of $u$ just above the diagonal of its block on the frame, as in
{\cite[(4)]{P51}}. For $l\le\varrho$ the \emph{Heisenberg group of floor $l$}, $\pfrN_l\subseteq
\pfrG^{(l)}\cap\pfrU$, is the unipotent radical of the stabiliser in $\pfrG^{(l)}$ of the line
$Af_l$. For $l<\varrho$ its elements are the $n(\omega,\beta)$, with $\omega$ in the span $V_{l+1}$ of the
basis vectors of floor $>l$ and $\beta\in A$ with $\beta+\bar\beta=-h(\omega,\omega)$, acting by $f_l\mapsto f_l$,
$v\mapsto v-h(v,\omega)f_l$ ($v\in V_{l+1}$), $f_l^*\mapsto f_l^*+\omega+\beta f_l$, and as the identity on the basis
vectors of floor $<l$ (\S\ref{pfr:ss-floors}). Writing $\beta=b-\tau_0h(\omega,\omega)$ with $b\in A_{\rm sk}$, the
measure $\eta_l$ is the image of $d\omega\otimes d\beta_{\rm sk}(b)$ under $(\omega,b)\mapsto n(\omega,b-\tau_0h(\omega,\omega))$,
where $d\omega$ is the Haar measure of $V_{l+1}\cong A^{d_l}$ in the coordinates of the basis ($d_l=\operatorname{rank}
V_{l+1}$; Lebesgue measure at the real place, the measure giving the lattice volume $1$ at the finite place); it is a
Haar measure of $\pfrN_l$ ((N2) of \S\ref{pfr:ss-floors}). At the real place, with $b=it$, this is the image of
$d\omega\,dt$. For $l\ge\varrho$ we put $\pfrN_l=1$, $\eta_l=\delta_1$.

For a sequence $p=(p_0,p_1,\dots)$ of units of $A$ and $j\le\varrho$ the \emph{Weyl element}
$w_p=w_{j,p}$ is $e_a\mapsto p_ae^*_a$, $e_a^*\mapsto\bar p_a^{\,-1}e_a$ $(a<j)$, and the identity on the
other basis vectors; and $w_{q}^{-1}w_p$ is the \emph{Weyl ratio}, the diagonal element with entries
$p_a/q_a$ at $e_a$ and $\overline{q_a/p_a}$ at $e_a^*$ $(a<j)$.

\subsubsection*{Place data}
A \emph{place datum} $\pfrF$ consists of
\begin{itemize}
\item a measure $\theta_\varrho$ on $\pfrG^{(\varrho)}$ (the \emph{top floor measure}) whose image under
$\sigma\mapsto\sigma^{-1}z$ is $c_\varrho\mu_\varrho$ for a constant $0<c_\varrho\le2$, where $\mu_\varrho$ is the
quadric measure of $z$ defined below;
\item a non-zero measure $\mu_\pfrH$ carried by $\pfrH$ and invariant under right translation by
$\pfrH$;
\item a unit $M$ of $A$ with $\bar M=M$ (the \emph{level});
\item an open set $\pfrCl$ of pairs of vectors containing $0$, with compact closure (the \emph{close
pairs});
\item a compact set $\pfrK\subseteq\pfrG$ (the \emph{box}) and a measurable set $\pfrP\subseteq\pfrG$
(the \emph{positions});
\end{itemize}
such that $\mu_\pfrH\{h:\exists u\in\pfrU,\ uhg\in\pfrK\}\le1$ for every $g\in\pfrG$. Its
\emph{floor measures} are
\[
\theta_l=\lambda_l*\eta_l\ (l<\varrho),\qquad\theta_\varrho,\qquad\theta_{n}=\mu_\pfrH,\qquad
\theta_l=\delta_1\ (l>n),
\]
where $\lambda_l$ is the image of the measure $\pfrmod(c)^{d_l}\,dc$ on $A^\times$ under $c\mapsto w_{l,c}^{-1}$
(the law of the inverse Weyl elements of floor $l$; $w_{l,c}$ is the Weyl factor $e_l\mapsto ce_l^*$,
$e_l^*\mapsto\bar c^{-1}e_l$, the identity on the other basis vectors), so that
\[
\int F\,d\theta_l=\iint F\bigl(w_{l,c}^{-1}\,n\bigr)\,\pfrmod(c)^{d_l}\,dc\,d\eta_l(n)\qquad(l<\varrho,\ F\ge0\text{ Borel}).
\]
\emph{Quadric measures.} For $l\le\varrho$ let $V_l$ be the span of the basis vectors of floor $\ge l$, with the Haar
measure $dy$ in the coordinates of the basis (Lebesgue measure at the real place; the measure giving the lattice
spanned by these basis vectors volume $1$ at the finite place). The \emph{quadric measure} $\mu_l$ of $f_l$ is the
fibre, at the value $h(f_l,f_l)$, of the disintegration of $dy$ along $y\mapsto h(y,y)$ (with $dh$ on $A^+$):
\begin{equation}\label{pfr:eq-slab}
\int\varphi\,d\mu_l=\lim_{B\to0}\frac1{\pbvol(B)}\int_{\{y\in V_l:\ h(y,y)-h(f_l,f_l)\in B\}}\varphi(y)\,dy,
\end{equation}
for $\varphi$ continuous with compact support (locally constant with compact support at the finite place), where
$B=(-\varepsilon,\varepsilon)$, $\varepsilon\to0$, at the real place and $B=\varpi_1^{\ell}(A^+\cap\auO_T)$, $\ell\to\infty$, at the
finite place ($\varpi_1$ of Lemma~\ref{pfg:psi}(c)). The limit exists and defines a positive Radon measure carried by
the quadric $\{h(y,y)=h(f_l,f_l)\}$: at the finite place this is Proposition~\ref{pfg:fibres}(a),(b) (the lattice of
$V_l$ is unimodular of rank $\ge3$); for $l<\varrho$ it also follows from Lemma~\ref{pfr:floorquad} below, and for
$l=\varrho$ at the real place from the computation in its proof. As $dy$ and the sets in
\eqref{pfr:eq-slab} are invariant under $\auU(V_l)$, \emph{$\mu_l$ is invariant under $\auU(V_l)\supseteq\pfrG^{(l)}$}
(restricted to $V_l$). This is the measure called the disintegration of $dy$ in Lemma~\ref{pfr:block}.
The constants of the floor measures are $c_l$ for $l\le\varrho$: for $l<\varrho$, $c_l=\kappa_A$ (Lemma~\ref{pfr:floorquad});
$c_\varrho$ is the constant of the top floor measure. At the real place $c_l=2$ for all $l\le\varrho$
(\S\ref{pfr:ss-realforms}); at the finite place $c_l=1$ for all $l\le\varrho$ (\S\ref{pff:ss-large}). In all cases
$0<c_l\le2$, and the image of $\theta_l$ under $\sigma\mapsto\sigma^{-1}f_l$ is $c_l\mu_l$ ($l\le\varrho$).

For measures $\nu_a$ on $\pfrG$ write
$\nu_{[l,l+i)}=\nu_{l+i-1}*\dots*\nu_l$ (convolution on $\pfrG$). Put
\[
\pfrM=\bigl((\eta*\theta)\bigr)_{[0,\varrho+2)},\qquad\pfrm_j=\bigl((\eta*\theta)\bigr)_{[0,j)},\qquad
\theta_{\ge l}=\theta_{[l,\varrho+2)},
\]
where $(\eta*\theta)_a=\eta_a*\theta_a$: the \emph{group measure} $\pfrM$, the measure $\pfrm_j$ of the
floors below $j$, and the measure $\theta_{\ge l}$ that represents the quotient
$(\pfrU\cap\pfrG^{(l)})\backslash\pfrG^{(l)}$. We always assume that $\pfrM$ is a Haar measure of
$\pfrG$ (at the real place this is Lemma~\ref{pfr:haar}; at the finite place see \S\ref{pff:ss-large}).

\begin{lemma}[Floor measures and quadrics]\label{pfr:floorquad}
Let $l<\varrho$. For every $\varphi$ as in \eqref{pfr:eq-slab} the limit \eqref{pfr:eq-slab} exists, and
\[
\int\varphi(\sigma^{-1}f_l)\,d\theta_l(\sigma)=\kappa_A\int\varphi\,d\mu_l .
\]
That is, the image of $\theta_l$ under $\sigma\mapsto\sigma^{-1}f_l$ is $c_l\mu_l$ with $c_l=\kappa_A$: $c_l=2$ at the
real place and $c_l=1$ at the finite place.
\end{lemma}
\begin{proof}
\emph{Coordinates.} In the basis of $V_l$ (the vectors $f_l$, $f_l^*$ and the basis of $V_{l+1}$) write
$y=sf_l+\omega'+tf_l^*$ with $s,t\in A$,
$\omega'\in V_{l+1}$; then $dy=ds\,d\omega'\,dt$ and, by (Cv1), $h(y,y)=s\bar t+\bar st+h(\omega',\omega')$. The non-units
$t$ form a null set (axiom (P) of \S\ref{pfs:ss-planch}); for a unit $t$ write $s=t(\tau_0h+b)$ with $h\in A^+$,
$b\in A_{\rm sk}$, and $h'=t\bar t\,h$. Then $s\bar t+\bar st=t\bar t(\tau_0h+b+\bar\tau_0h-b)=h'$, and, for fixed
$t,\omega'$, the map $(h',b)\mapsto s$ carries $dh'\otimes d\beta_{\rm sk}$ to $\kappa_A\,ds$: indeed $(h,b)\mapsto\tau_0h+b$
carries $dh\otimes d\beta_{\rm sk}$ to $\kappa_A\,du$, $u\mapsto tu$ multiplies $du$ by $\pfrmod(t)$, and $h\mapsto t\bar t\,h$
multiplies $dh$ by $\pfrmod(t)$ (\S\ref{pfr:ss-setting}). Write $Y(t,\omega',b,h')=sf_l+\omega'+tf_l^*$ for the
point with these coordinates, $s=\tau_0h'\bar t^{-1}+tb$. Hence, with $c_0=h(f_l,f_l)=0$,
\[
\begin{gathered}\frac1{\pbvol(B)}\int_{\{h(y,y)\in B\}}\varphi\,dy=\frac1{\kappa_A}\iint F_B(t,\omega')\,dt\,d\omega',\\
F_B(t,\omega')=\frac1{\pbvol(B)}\iint\mathbf 1\bigl[h'+h(\omega',\omega')\in B\bigr]\,\varphi\bigl(Y(t,\omega',b,h')\bigr)\,
db\,dh'.\end{gathered}
\]
\emph{A bound.} Let $\varphi$ vanish outside the set of the $y$ whose coordinates lie in a compact set $R_0$
($R_0=\{|c|\le R\}$, resp.\ $R_0=\varpi_1^{-a}\auO_T$), and let $B\subseteq B_0$, $B_0=(-1,1)$ resp.\ $A^+\cap\auO_T$.
For fixed $t,\omega',h'$ the set of the $b$ with $Y(t,\omega',b,h')$ in the support of $\varphi$ has
$d\beta_{\rm sk}$-measure at most $C_0\pfrmod(t)^{-1/2}$ with $C_0$ depending only on $R_0$: at the real place $s=
\frac{h'}{2\bar t}+tb$ with $b=ix$, and $\frac{h'}{2\bar t}=\frac{h't}{2|t|^2}$ is orthogonal to $tix$ in $\C=\R^2$, so
$|s|\le R$ forces $|x|\le R/|t|$; at the finite place the set is empty or a coset of
$\{b:tb\in\varpi_1^{-a}\auO_T\}$, of measure at most $|\varpi_1^at|_T^{-1/2}$ by Lemma~\ref{pfg:duality}(d). As the
$h'$ with $h'+h(\omega',\omega')\in B$ form a set of measure $\pbvol(B)$,
\[
|F_B(t,\omega')|\le\sup|\varphi|\cdot C_0\,\pfrmod(t)^{-1/2}\,\mathbf 1[t\in R_0,\ \omega'\in R_0^{d_l}],
\]
and $\int_{R_0}\pfrmod(t)^{-1/2}dt<\infty$ (at the real place $\int_{|t|\le R}|t|^{-1}dt=2\pi R$; at the finite place
the integral is a finite product over $w\in T$ of the sums $\sum_{j\ge-a'}(1-q_w^{-1})q_w^{-j}q_w^{j/2}$).
\emph{The limit.} Fix $t,\omega'$ and put $h'_0=-h(\omega',\omega')$ and $G(h')=\int\varphi(Y(t,\omega',b,h'))\,db$. The
function $G$ is continuous at $h'_0$: at the real place by dominated convergence (the $b$ with non-zero integrand lie in
a fixed bounded set, by the bound above); at the finite place $\varphi$ is invariant under translation by a lattice
$\varpi_1^{m}L$ ($L$ the lattice of $V_l$), and for $h'$ so close to $h'_0$ that $\tau_0(h'-h'_0)\bar t^{-1}\in\varpi_1^m\auO_T$
the points $Y(t,\omega',b,h')$ and $Y(t,\omega',b,h'_0)$ differ by $\tau_0(h'-h'_0)\bar t^{-1}f_l\in\varpi_1^mL$ for every $b$, so
$G(h')=G(h'_0)$. As $F_B(t,\omega')$ is the mean of $G$ over
$h'_0+B$, $F_B(t,\omega')\to G(h'_0)$. By dominated convergence the limit \eqref{pfr:eq-slab} exists and
\[
\kappa_A\int\varphi\,d\mu_l=\iiint\varphi\bigl(Y(t,\omega',b,-h(\omega',\omega'))\bigr)\,db\,d\omega'\,dt .
\]
\emph{The image of $\theta_l$.} For $\sigma=w_{l,c}^{-1}n$ with $n=n(\omega,\beta)$, $\beta=b-\tau_0h(\omega,\omega)$, we have
$n^{-1}=n(-\omega,\bar\beta)$ (as $\beta+\bar\beta=-h(\omega,\omega)$) and $\sigma^{-1}f_l=n^{-1}(cf_l^*)=c(f_l^*-\omega+\bar\beta f_l)$. Its
coordinates are $t=c$, $\omega'=-c\omega$, $s=c\bar\beta$; with $h'=-h(\omega',\omega')=-c\bar c\,h(\omega,\omega)$ we get
$s=\tau_0h'\bar c^{-1}+cb'$ for $b'=\bar\beta+\tau_0h(\omega,\omega)=-b+(\tau_0-\bar\tau_0)h(\omega,\omega)\in A_{\rm sk}$. So
$\sigma^{-1}f_l=Y(c,-c\omega,b',-h(\omega',\omega'))$. For fixed $c$, $\omega\mapsto-c\omega$ multiplies $d\omega$ by
$\pfrmod(c)^{d_l}$, and for fixed $c,\omega$, $b\mapsto b'$ preserves $d\beta_{\rm sk}$. By the definition of $\theta_l$,
\[
\int\varphi(\sigma^{-1}f_l)\,d\theta_l(\sigma)=\iiint\varphi\bigl(Y(c,-c\omega,b',\cdot)\bigr)\pfrmod(c)^{d_l}\,dc\,d\omega\,db=
\iiint\varphi\bigl(Y(t,\omega',b',-h(\omega',\omega'))\bigr)\,db'\,d\omega'\,dt,
\]
which is $\kappa_A\int\varphi\,d\mu_l$. The values of $\kappa_A$ are those of \S\ref{pfr:ss-setting}.
\end{proof}

A point $g$ is \emph{near at $j$} if $\pfrR^{(j)}g$ meets $\pfrK$ (for $j\ge\varrho+2$ we put $\pfrR^{(j)}=\pfrU$). A function $W\colon\pfrG\to\C^I$ ($I$
a finite set) is a \emph{cut Bessel function} for $\pfrF$ if it is bounded and measurable,
$W(ug)=\psi(M^{-1}\sigma(u))W(g)$ for $u\in\pfrU$, and $W(g)\ne0$ only if $ug\in\pfrK$ for some
$u\in\pfrU$.

\subsubsection*{Correlations and forms}
For a kernel $\mathsf k$ on $\pfrG\times\pfrG$ the \emph{correlation from floor $l$ on} is
\[
\mathsf k_l(g,y)=\int_\pfrG\mathsf k(xg,xy)\,d\theta_{\ge l}(x).
\]
For a cut Bessel function $W$ and $d\in\pfrG$ put $\mathsf k^{d}(g,y)=\pfrip{W(g)}{W(dy)}$. Fix $j\le
\varrho$ and sequences of units $p,q$; write $w=w_{j,p}$ and $d=w_{j,q}^{-1}w_{j,p}$. The
\emph{head density} is
\[
\pfrkap_j(p,q)=\prod_{a<j}\pfrmod(p_aq_a)^{d_a/2},\qquad d_a=k+1+2(\varrho-a-1).
\]
For a position $x$, a point $g$ and a bounded measurable weight $\omega\colon\pfrG\to\C$ put
\[
\Phi_j(x;\omega;g)=\iint\omega(\sigma\tau^{-1})\,\mathsf k^1_{j+1}(\sigma g,\sigma gx)\,
\overline{\mathsf k^{d}_{j+1}(\tau w^{-1}g,\tau w^{-1}gx)}\,d\theta_j(\tau)\,d\theta_j(\sigma)
\]
The
\emph{form of floor $j$} is
\[
\pfrI_j(p,q;\omega)=\pfrkap_j(p,q)\int_{\pfrG\times\pfrG}\Phi_j(x;\omega;g)\,d\pfrM(x)\,d\pfrm_j(g),
\]
whenever the integral converges absolutely. Its \emph{integrand norm} and \emph{majorant} at the
position $x$ are
\[
\pfrF_j(x;\omega)=\int|\Phi_j(x;\omega;g)|\,d\pfrm_j(g),\qquad
\pfrF_j^+(x)=\int\pfrB^1_j(g,gx)\,\pfrB^{d}_j(w^{-1}g,w^{-1}gx)\,d\pfrm_j(g),
\]
where $\pfrB^{d}_j(g,y)=\int|\mathsf k^{d}_{j+1}(\sigma g,\sigma y)|\,d\theta_j(\sigma)$ is the
\emph{pivot norm}.

\begin{lemma}\label{pfr:L1}
Let $W$ be a cut Bessel function, $j\le\varrho$, $p,q$ sequences of units and $|\omega|\le1$. Then
$\pfrF_j(x;\omega)\le\pfrF^+_j(x)$ for every $x$, and, whenever $\int\pfrF_j^+\,d\pfrM<\infty$, the
form $\pfrI_j(p,q;\omega)$ is defined and
$|\pfrI_j(p,q;\omega)|\le\pfrkap_j(p,q)\int\pfrF_j(x;\omega)\,d\pfrM(x)$.
\end{lemma}
\begin{proof}
For fixed $g$, $|\Phi_j(x;\omega;g)|\le\iint|\mathsf k^1_{j+1}(\sigma g,\sigma gx)|\,
|\mathsf k^{d}_{j+1}(\tau w^{-1}g,\tau w^{-1}gx)|\,d\theta_j(\tau)\,d\theta_j(\sigma)$, and by
Tonelli's theorem the right side is $\pfrB^1_j(g,gx)\pfrB^{d}_j(w^{-1}g,w^{-1}gx)$. Integrating in $g$
gives the first claim. The second is the triangle inequality for the integral defining $\pfrI_j$, whose
integrand is then absolutely integrable by the first claim and Tonelli's theorem.
\end{proof}

\subsection{Depth, weights and gauges}\label{pfr:ss-gauge}

\subsubsection*{Stable vectors and depth}
Let $Y^{(l)}$ be the module of the vectors orthogonal to $e_0,\dots,e_{l-1}$ modulo their span. A
vector $\omega$ is \emph{stable at $(g,g')$ for $l$} if it is orthogonal to the $e_b$, $b<l$, and
$(g^{-1}\omega,g'^{-1}\omega')\in\pfrCl$ for some $\omega'\equiv\omega$ in $Y^{(l)}$. For $a<\varrho$ and
points $g$, $gx$ let
\[
\mathfrak d_a(g,x)=\{h(\omega,e_a):\ \omega\text{ stable at }(g,gx)
\text{ for }a\}\subseteq A\ \text{(the $e_a^*$-coordinates)},\qquad\pfrdep_a(g,x)=\pbvol(\mathfrak d_a(g,x))^{-1},
\]
the \emph{depth} of $e_a$ at the two points. At the real place $\mathfrak d_a$ is a disc and
$\pfrdep_a(g,x)=\pi^{-1}\inf_\zeta(|\zeta|_g^2+|\zeta|_{gx}^2)$, the infimum over the duals $\zeta$ of
$e_a$ in $Y^{(a)}$ (Lemma~\ref{pfr:realgauge}); at the finite places it is the index of an ideal.
Let $X'$ be the span of the orthonormal vectors $x_i$ left after two blocks of $2m+2$ of them are set
aside ($m$ the number of further hyperbolic pairs of the place datum; $m=0$ at the real place), and, at
the real place, also the vector $x_{i_0}$ of the boosts (\S\ref{pfr:ss-realdata}), and
$s=\operatorname{rank}X'-(3n+1+m)$, the \emph{volume exponent}. The \emph{depth weight} of length $j$ is
$\pfrwt_j(g,x)=\prod_{a<j}\pfrdep_a(g,x)^{-s}$.

\subsubsection*{Second moments}
For a kernel $\mathsf k$, a position $x$ and $j\le n$ let
\[
\mathsf a_j(g)=\pfrwt_j(g,x)\,|\mathsf k_{j+1}(g,gx)|^2,\qquad
\mathsf b_j(g)=\int\bigl(\mathsf a_j+\mathsf b_{j+1}\bigr)(\sigma g)\,d\theta_j(\sigma)\ (j<n),\quad
\mathsf b_j=0\ (j\ge n),
\]
and $\mathsf S_j(\mathsf k;x)=\int\mathsf b_j\,d\theta_{[0,j)}$, the \emph{weighted second moment} from
floor $j$ on.

\subsubsection*{Gauges}
A \emph{gauge} $\Gamma$ for the place datum $\pfrF$ is a filtration of $A$ by compact sets
$\Gamma(0)\subseteq\Gamma(1)\subseteq\dots$ (``the numbers of size $i$'': $|c|\le L^i$ at the real place
for a scale $L$, $L^ic$ integral at the finite places), together with a
number $\Lambda=\Lambda_\Gamma\in[2,\infty)$, the \emph{loss}, such that:
\begin{enumerate}
\item[(G0)] (\emph{size rules}) $\Gamma(i)\Gamma(i')\subseteq\Gamma(i+i')$; a sum of at most $r+1$ elements of
$\Gamma(i)$ lies in $\Gamma(i+1)$; $\Gamma(i)$ is stable under $c\mapsto-c$ and $c\mapsto\bar c$ and contains $0,1$;
the entries of the Gram matrix $J$ of the basis and of $J^{-1}$ lie in $\Gamma(1)$. A vector or matrix
\emph{has size $i$} if its entries lie in $\Gamma(i)$. Consequences used below: if $X,Y$ have sizes $i,i'$,
then $XY$ has size $i+i'+1$; an isometry $k$ of size $i$ has $k^{-1}=J^{-1}k^*J$ of size $i+4$; an
$l\times l$ determinant of entries of size $i$ ($l\le r+1$) has size $l(i+1)$ (Laplace expansion and
induction on $l$);
\item[(G1)] $m\le9n$, $P\ge1$, $M\in\Gamma(1)$, $\tau_0\delta^{-1}\in\Gamma(1)$ for a fixed element $\tau_0$ of trace one and a skew unit $\delta$, the basis and the elements of $\pfrK$ have entries in
$\Gamma(1)$;
\item[(G2)] for a close pair $(\omega_1,\omega_2)$ the vectors $\omega_1$ and $M^{-1}(\omega_1-\omega_2)$
have entries in $\Gamma(1)$; and if $M^{-1}\in\Gamma(1)$, for every $i$ there is a unit $c$ with
$c^{-1}\in\Gamma(i+2)$ such that $(c\omega_1,c\omega_2)$ is close for all vectors $\omega_1,\omega_2$ with
entries in $\Gamma(i)$;
\item[(G3)] $\pfrmod(u)\le\Lambda^i$ for units $u\in\Gamma(i)$, and $\pbvol\Gamma(i)\le\Lambda^{i+1}$;
\item[(G4)] (\emph{quadrics of a block}) Lebesgue measure on $A^{2m+2}$ disintegrates along the
standard hermitian form into measures each of which gives mass at most $\Lambda^{(2m+2)(i+1)}$ to the
vectors with entries in $\Gamma(i)$;
\item[(G5)] (\emph{common vectors}) for $g\in\pfrK$, a position $x\in\pfrP$ and $\varrho$ linear forms
$v$ on $X'$, let $\mathcal W$ be the set of the $\omega\in X'$ killed by $v$ with $(g^{-1}\omega,
(gx)^{-1}\omega)\in\pfrCl$, and let $\mathcal D\subseteq A$ be the set of the differences
$M^{-1}\bigl(h(\omega_1,\omega_1')-h(\omega_2,\omega_2')\bigr)$ for close pairs $(\omega_1,\omega_2)$,
$(\omega_1',\omega_2')$; there is a measurable $\phi\colon X'\to[0,\infty]$ such that for every $u\in X'$
\[
\pbvol\{h(\omega,u):\omega\in\mathcal W\}^{-1}\le\phi(u),\qquad
\pbvol\{c\in A:h(\omega,cu)\in\mathcal D\ \forall\omega\in\mathcal W\}\le\phi(u),
\]
and $\int_{\Gamma(i)^{X'}}\phi^s\le\Lambda^{99n(i+1)}$, the integral over the $u\in X'$ whose coordinates in the
orthonormal basis of $X'$ lie in $\Gamma(i)$, for the Haar measure in these coordinates.
\end{enumerate}
In (G5) the form $h$ is that of (Cv1): for $u=\sum_iu_ix_i$ and $\omega=\sum_i\omega_ix_i$ in $X'$,
$h(\omega,u)=\sum_i\omega_i\bar u_i=u^*\omega$, and $h(\omega,cu)=\bar c\,h(\omega,u)$. The set $\mathcal D$ is stable
under conjugation ($\bar M=M$, $\overline{h(\omega,\omega')}=h(\omega',\omega)$, and the roles of the two pairs of close
pairs may be exchanged), so the second set in (G5) is also $\{c:h(cu,\omega)\in\mathcal D\ \forall\omega\in\mathcal W\}$.
For a vector $y$ we write $y_{X'}=\sum_ih(y,x_i)x_i$ ($x_i$ running over the orthonormal basis of $X'$) for its
orthogonal projection on $X'$; its coordinates $h(y,x_i)$ are coordinates of $y$, and $h(\omega,y)=h(\omega,y_{X'})$
for $\omega\in X'$.
A sequence $p$ is a \emph{head of length $j$} for $\Gamma$ if $(p_0\cdots p_{l-1})^{-1}\in\Gamma(1)$ for
$l\le j$. The \emph{position ball} of $\Gamma$ is the set of the positions $x\in\pfrP$ such that
$M^{-1}(x-1)$ has entries in $\Gamma(1)$. Put $C_\Gamma=\Lambda^{10^6n^3}$ and
\[
\mathsf Q_\Gamma(x)=\mathsf Q_{\Gamma,j,p,q,W}(x)=\Bigl(\prod_{a<j}\mathsf c_a(p_a^{-1})\Bigr)^2C_\Gamma^3\,
\mathsf S_0(\mathsf k^1;x)\,\mathsf S_j(\mathsf k^{d};x),\qquad\mathsf c_l(g)=\pfrmod(g)^{d_l+1/2},
\]
the \emph{weighted square product}; we say that a number $f\ge0$ \emph{has the gain $e$ at $x$} if
$f^2\le e^2\mathsf Q_\Gamma(x)$. Put also
\[
\mathsf N(W)=\Bigl(\prod_{l<\varrho}\mathsf c_l(M)\Bigr)P^{\,\varrho\operatorname{rank}X'}\iint
|\pfrip{W(g)}{W(y)}|^2\,d\theta_{\ge0}(g)\,d\theta_{\ge0}(y),
\]
the \emph{shell normalisation}. A point $g$ is \emph{active} at $x$ (for $j,p,q,W$) if
$\pfrip{W(g)}{W(gx)}\ne0$ and $w_p^{-1}g$, $w_q^{-1}gx$ are near at $j$; the position $x$ \emph{has a
deep vector of depth $D$} if every active point $g$ has an $a$ with $j\le a<\varrho$ and
$\pfrdep_a(g,x)\ge D$. Finally, for $V\ge0$, \emph{the mass of floor $j$ is at most $V$} if
$\theta_j\{\sigma:\sigma g\text{ near at }j+1\}\le V$ for every $g$.

\subsection{The shared estimates}\label{pfr:ss-shared}

Throughout this subsection $\pfrF$ is a place datum whose group measure $\pfrM$ is a Haar measure,
$\Gamma$ a gauge with loss $\Lambda$, $W$ a cut Bessel function, $j\le\varrho$, $p,q$ sequences of units,
and $\mathsf Q=\mathsf Q_{\Gamma,j,p,q,W}$, $\pfrF^+_j$, $\mathsf N(W)$ are as in \S\ref{pfr:ss-gauge}.
These are the statements that \textup{B-frames-finite} also uses.

\begin{lemma}[Descent]\label{pfr:descent}
Let $i<\varrho$ and assume $\int\pfrF^+_i\,d\pfrM<\infty$ (for the given $p,q$). For $c\in A^\times$ let
$p^c$, $q^c$ be $p$, $q$ with the entry of index $i$ replaced by $c$. Then for every $\xi\in A$ the
function $c\mapsto\psi(\xi c)\pfrI_{i+1}(p^c,q^c;1)$ is integrable on $A$ (the complement of $A^\times$
has measure $0$), and
\[
\int_A\psi(\xi c)\,\pfrI_{i+1}(p^c,q^c;1)\,dc=\pfrI_i\bigl(p,q;\ g\mapsto\psi(\xi\,h(gf_i,f_i))\bigr).
\]
\end{lemma}
(Proof in \S\ref{pfr:ss-descent}.)

\begin{theorem}[One frame]\label{pfr:oneframe}
Let $q$ be a head of length $j$ for $\Gamma$, $\Gamma'$ a second gauge, $S$ a subset of the position
ball of $\Gamma'$, and $f\ge0$ measurable with the gain $e$ at every $x\in S$. Then
\[
\pfrkap_j(p,q)\int_Sf\,d\pfrM\le e\,\Lambda^{1500000n^3+1}\,\Lambda_{\Gamma'}^{25000n^2}\,\mathsf N(W).
\]
\end{theorem}
(Proof in \S\ref{pfs:ss-oneframe}.)

\begin{theorem}[One position, no gain]\label{pfr:nogain}
Let $p$ be a head of length $j$ for $\Gamma$, $x\in\pfrP$, and let the mass of floor $j$ be at most $V$.
Then $\pfrF^+_j(x)^2\le V^2\,\mathsf Q(x)$.
\end{theorem}

\begin{theorem}[The gain from a deep vector]\label{pfr:deep}
Let $p$ be a head of length $j$ for $\Gamma$, $x\in\pfrP$, let the mass of floor $j$ be at most $V\ge1$,
and let $0<\varsigma$ be such that $x$ has a deep vector of depth $\varsigma^{-1}$. Then
$\pfrF^+_j(x)^2\le\varsigma^{\,s}\,V^2\,\mathsf Q(x)$. If $j=0$, then moreover
$\pfrF^+_0(x)^2\le\varsigma^{\,2s}\,\mathsf Q(x)$ (at floor $0$ the two frames coincide and both pivots gain).
\end{theorem}
(Proofs in \S\ref{pfr:ss-gainsproof}.)

\begin{theorem}[The gain from a frequency]\label{pfr:gain-char}
Let $p$ be a head of length $j$ for $\Gamma$, $x\in\pfrP$, $\xi$ a unit and $\eta\in A$ with
$\psi(\eta c)=\psi(\xi c)$ for all $c\in\Gamma(200n^2)$. Then
\[
\pfrF_j\bigl(x;\ g\mapsto\psi(\eta\,h(gf_j,f_j))\bigr)^2\le\Bigl(\frac{P}{\pfrmod(\xi)}\Bigr)^{\operatorname{rank}X'}
\mathsf Q(x).
\]
\end{theorem}

\begin{theorem}[The gain from two heads]\label{pfr:gain-heads}
Let $j=i+1$, $p$ a head of length $j$ for $\Gamma$, $x\in\pfrP$, and assume that
$\xi=M^{-1}(p_i^{-1}-q_i^{-1})$ is a unit. Then
$\pfrF_j(x;1)^2\le(P/\pfrmod(\xi))^{\operatorname{rank}X'}\,\mathsf Q(x)$.
\end{theorem}
(Proofs in \S\ref{pfs:ss-gains}; the constant $P$ is that of axiom~(P) of \S\ref{pfs:ss-planch}.)

\begin{lemma}[Deep vectors are rare]\label{pfr:depth-rare}
For $x\in\pfrP$, $a<\varrho$ and every $g$,
$\int_{\{\sigma:\ \sigma g\text{ near at }a+1\}}\pfrdep_a(\sigma g,x)^s\,d\theta_a(\sigma)\le\Lambda^{10^4n}$.
\end{lemma}
We use two facts on the floor measures, stated exactly.

\begin{lemma}[Charts of a floor]\label{pfr:chart}
Let $a\le\varrho$ and let $c\ge0$ be measurable with $c(\xi g)=c(g)$ for $\xi\in\pfrR^{(a+1)}$. Suppose that for
every $g\in\pfrK$ there is a measurable $\Psi_g\ge0$ on vectors with $c(\sigma g)\le\Psi_g(\sigma^{-1}f_a)$ for
$\sigma\in\pfrG^{(a)}$ and $\int_{\Gamma(7)}\Psi_g\,d\mu_a\le C$, where $\mu_a$ is the measure of the quadric
of $f_a$ and $\Gamma(7)$ the vectors with entries of size $7$. Then
$\int_{\{\sigma:\sigma g\text{ near at }a+1\}}c(\sigma g)\,d\theta_a(\sigma)\le c_aC\le2C$ for every $g$ (so
$\le C$ at the finite place, where $c_a=1$).
\end{lemma}

\begin{lemma}[From the quadric to $X'$]\label{pfr:quadricX}
Let $a\le\varrho$. For measurable $\phi\ge0$ on $X'$ with $\int_{\Gamma(11)^{X'}}\phi^s\le B$, the function
$\Psi(v)=\phi(v_{X'})^s$, $v_{X'}$ the orthogonal projection of $v$ on $X'$ (\S\ref{pfr:ss-gauge}), satisfies
$\int_{\Gamma(7)}\Psi\,d\mu_a\le\Lambda^{99n\cdot13}B$.
\end{lemma}
(Lemma~\ref{pfr:chart} is proved in \S\ref{pfr:ss-floors}.)
(Lemma~\ref{pfr:quadricX} is proved in \S\ref{pfr:ss-section}.)

\begin{proof}[Proof of Lemma~\ref{pfr:depth-rare}]
Let $c(g)=\pfrdep_a(g,x)^s$; it is $\pfrR^{(a+1)}$-invariant by (C3) of \S\ref{pfr:ss-gainsproof}. Let
$g\in\pfrK$ and let $\phi$ be the function of gauge axiom (G5) for $g$, $x$ and $v=0$, and $\mathcal W$ its set
of common vectors. For $\sigma\in\pfrG^{(a)}$ and $\omega\in\mathcal W$, the vector $\sigma\omega$ is orthogonal to the
$e_b$, $b<a$ ($\omega\in X'$ is orthogonal to them and $\sigma$ fixes them), and $(\,(\sigma g)^{-1}\sigma\omega,
(\sigma gx)^{-1}\sigma\omega)=(g^{-1}\omega,(gx)^{-1}\omega)$ is close; so $\sigma\omega$ is stable at $(\sigma g,\sigma gx)$,
and its $e_a^*$-coordinate is $h(\sigma\omega,f_a)=h(\omega,\sigma^{-1}f_a)=h(\omega,u)$ with $u=(\sigma^{-1}f_a)_{X'}$. Hence
$\mathfrak d_a(\sigma g,x)\supseteq\{h(\omega,u):\omega\in\mathcal W\}$, and $\pfrdep_a(\sigma g,x)\le\phi(u)$ by the first
inequality of (G5). By Lemma~\ref{pfr:quadricX} with $B=\Lambda^{99n\cdot12}$ (axiom (G5) at size $11$) and
Lemma~\ref{pfr:chart} with $\Psi_g(v)=\phi(v_{X'})^s$, the integral is at most
$c_a\Lambda^{99n\cdot25}\le2\Lambda^{2475n}\le\Lambda^{10^4n}$ (as $\Lambda\ge2$).
\end{proof}

\subsubsection*{Shells}
\begin{definition}\label{pfr:shell}
Let $H\in[1,\infty)$. A set $S\subseteq\pfrG$ is a \emph{shell of height $H$} (for $\Gamma,j,p,q,W$)
if
\begin{enumerate}
\item[(a)] $S$ is covered by the position balls of a family of at most $2H^{9n}$ gauges $\Gamma'$ with
$\Lambda_{\Gamma'}\le\Lambda H^2$;
\item[(b)] for every $\varsigma\in(0,1)$ with $\Lambda^{1200n^2}\le\varsigma^{400n^2}H$, every $x\in S$
has a deep vector of depth $\varsigma^{-2}$.
\end{enumerate}
\end{definition}
By (a) and the definition of the position balls, a shell lies in $\pfrP$.

\begin{lemma}[One shell]\label{pfr:oneshell}
Let $S$ be a shell of height $H$, $q$ a head of length $j$, and $f\ge0$ measurable with the gain $e$
at every point of $S$. Then
\[
\pfrkap_j(p,q)\int_Sf\,d\pfrM\le e\,\Lambda^{2\cdot10^6n^3}H^{10^5n^2}\,\mathsf N(W).
\]
\end{lemma}
\begin{proof}
Let $\mathcal G$ be the family of (a). By Theorem~\ref{pfr:oneframe} applied to $S\cap(\text{position
ball of }\Gamma')$ for each $\Gamma'\in\mathcal G$, and $\Lambda_{\Gamma'}\le\Lambda H^2$,
\[
\pfrkap_j\int_{S}f\,d\pfrM\le\sum_{\Gamma'\in\mathcal G}e\,\Lambda^{1500000n^3+1}(\Lambda H^2)^{25000n^2}
\mathsf N(W)\le2H^{9n}\,e\,\Lambda^{1500000n^3+25000n^2+1}H^{50000n^2}\mathsf N(W).
\]
As $2\le\Lambda$ and $H\ge1$, the right side is at most $e\,\Lambda^{2\cdot10^6n^3}H^{10^5n^2}\mathsf N(W)$.
\end{proof}

\begin{lemma}[A remote shell]\label{pfr:remoteshell}
Let $S$ be a shell of height $H>\Lambda^{2400n^2}$, let $p,q$ be heads of length $j$, let the mass of
floor $j$ be at most $V\ge1$, and let $\tau\in(0,1)$ with $\tau^{800n^2}=H^{-1}$. Then
\[
\pfrkap_j(p,q)\int_S\pfrF^+_j\,d\pfrM\le V\,\tau^{\,s}\,\Lambda^{2\cdot10^6n^3}H^{10^5n^2}\,\mathsf N(W).
\]
\end{lemma}
\begin{proof}
$\tau<1$ because $H>1$. Moreover $\Lambda^{2400n^2}\le H=\tau^{800n^2}H^2$, so
$\Lambda^{1200n^2}\le\tau^{400n^2}H$, and by (b) every $x\in S$ has a deep vector of depth $\tau^{-2}$.
By Theorem~\ref{pfr:deep} with $\varsigma=\tau^2$, $\pfrF^+_j(x)$ has the gain $V\tau^{s}$ at every
$x\in S$; Lemma~\ref{pfr:oneshell} concludes.
\end{proof}

\begin{lemma}[A few factors of the gain pay for the height]\label{pfr:payheight}
Let $S$ be a shell of height $H$, $p,q$ heads of length $j$, the mass of floor $j$ at most $V\ge1$,
$e\le1$, $r'\ge0$ with $9\cdot10^7n^4+r'\le s$, and $f\ge0$ measurable with $f\le\pfrF_j^+$ everywhere
and with the gain $Ve^{s}$ on $S$. Then
\[
\pfrkap_j(p,q)\int_Sf\,d\pfrM\le H^{-1}\,V\,e^{r'}\,\Lambda^{3\cdot10^8n^4}\,\mathsf N(W).
\]
\end{lemma}
\begin{proof}
It suffices to find $\mathsf g\le e$ with
\begin{equation}\label{pfr:eq-pay}
\mathsf g^{9\cdot10^7n^4}H^{1+10^5n^2}\le\Lambda^{29\cdot10^7n^4},\qquad
\pfrkap_j\int_Sf\,d\pfrM\le V\mathsf g^{s}\Lambda^{2\cdot10^6n^3}H^{10^5n^2}\mathsf N(W):
\end{equation}
indeed $\mathsf g^s\le\mathsf g^{9\cdot10^7n^4}e^{r'}$ (as $\mathsf g\le e\le1$), and then the right side
of the second inequality is at most $H^{-1}Ve^{r'}\Lambda^{29\cdot10^7n^4+2\cdot10^6n^3}\mathsf N(W)$, and
$29\cdot10^7n^4+2\cdot10^6n^3\le3\cdot10^8n^4$.

If $H\le\Lambda^{2400n^2}$ take $\mathsf g=e$: the first inequality holds since $e\le1$ and
$H^{1+10^5n^2}\le\Lambda^{2400n^2(1+10^5n^2)}\le\Lambda^{29\cdot10^7n^4}$, and the second is
Lemma~\ref{pfr:oneshell}. If $H>\Lambda^{2400n^2}$, let $\tau^{800n^2}=H^{-1}$ and
$\mathsf g=\min(e,\tau)$. Then $\mathsf g^{9\cdot10^7n^4}H^{1+10^5n^2}\le(\tau^{800n^2})^{112500n^2}
H^{112500n^2}=1$. If $e\le\tau$ the second inequality is Lemma~\ref{pfr:oneshell}; if $\tau<e$ it is
Lemma~\ref{pfr:remoteshell}, since $f\le\pfrF^+_j$.
\end{proof}

\begin{theorem}[The sum over the shells]\label{pfr:shells}
Let $S_i$ $(i\ge0)$ be shells of heights $2^i$ that cover the support of a measurable $f\ge0$. Let $p,q$
be heads of length $j\le\varrho$ for $\Gamma$, let the mass of floor $j$ be at most $V\ge1$, let $f\le
\pfrF_j^+$ everywhere, and let $e\ge0$ and $r'$ with $9\cdot10^7n^4+r'\le s$ be such that, if $e<1$, $f$
has the gain $Ve^{s}$ at every $x\in\pfrP$. Then
\[
\pfrkap_j(p,q)\int_\pfrG f\,d\pfrM\le V\,\min(1,e)^{r'}\,\Lambda^{4\cdot10^8n^4}\,\mathsf N(W).
\]
\end{theorem}
\begin{proof}
Let $e'=\min(1,e)$. On each $S_i\subseteq\pfrP$, $f$ has the gain $Ve'^{\,s}$: if $e<1$ by hypothesis,
and if $e\ge1$ by $f\le\pfrF^+_j$ and Theorem~\ref{pfr:nogain} (as $V\ge1$). By
Lemma~\ref{pfr:payheight}, $\pfrkap_j\int_{S_i}f\le2^{-i}Ve'^{\,r'}\Lambda^{3\cdot10^8n^4}\mathsf N(W)$. Summing
over $i$ (the $S_i$ cover the support of $f$) gives the factor $\sum_i2^{-i}=2\le\Lambda^{10^8n^4}$.
\end{proof}

\subsection{The real place: data and the statement of \textup{(R)}}\label{pfr:ss-realdata}

\subsubsection*{The group}
Fix $\tau\in\auPl$. As in {\cite[\S3.3.7]{P51}} we work in the normalised basis of
{\cite[\S3.2.1]{P51}}: $\pfrG=G_\tau$ is the unitary group of the hermitian form on $\C^r$ with $\varrho$
hyperbolic pairs $e_a,e_a^*$, an orthonormal family $(x_i)_{i\in\iota}$, $\#\iota=k$, and $z$ with
$h(z,z)=-1$; the frame is $(e_0,\dots,e_{\varrho-1},z)$, $\pfrU$ is the Bessel unipotent group and
$\pfrH=\auU(k)$ the fixer of the $e_a,e_a^*,z$. For $\gamma\in\pfrG$ let $Z(\gamma)=(h(\gamma f_b,f_a))_{a,b}$
be its \emph{frame matrix}, with diagonal entries $Z_{aa}(\gamma)=h(\gamma f_a,f_a)$. For a sequence
$x=(x_0,\dots,x_\varrho)$ of non-zero complex numbers with $|x_\varrho|>1$ let $w_x=w_{\varrho,x}$ and let
$\sigma_x\in\pfrG^{(\varrho)}$ be the boost in the plane of $z$ and $x_{i_0}$ ($i_0\in\iota$ fixed) with
$h(\sigma_xz,z)=x_\varrho$ (for $|x_\varrho|\le1$ we let $\sigma_x$ be the boost of parameter
$x_\varrho/\min(1,|x_\varrho|)$; this case does not occur in the sieve, as the last entries of the
open terms satisfy $|\tau(t)|>N$, that is $|x_\varrho|>1$ after normalisation).

\subsubsection*{The cut Bessel function and its orbital integrals}
Let $\mathbb W$ be a finite-dimensional unitary representation of $\pfrH$, of dimension $\dim\mathbb W=\#I$,
irreducible, and let $W,W^\flat\colon\pfrG\to\mathbb W$ be continuous Bessel functions of type $\mathbb W$: they transform on the left by the
character $u\mapsto e^{\sigma(u)-\overline{\sigma(u)}}$ under $\pfrU$ and by the representation under $\pfrH$, as in
{\cite[Definition~2.16]{P51}}. Write $\pfrip{A}{R(x)B}=\int_{\pfrU\backslash\pfrG}\pfrip{A(y)}{B(yx)}\,dy$ for
the matrix coefficients in $L^2(\pfrU\backslash\pfrG)$ (the invariant measure normalised as in
{\cite[\S3.2.1]{P51}}). For $\sigma\in\pfrG^{(\varrho)}$ and a sequence $x$ as above the \emph{normalised
orbital integral} is
\[
\pfrN_\sigma(x)=\frac{\sqrt{c_G}\,\bigl((|x_\varrho|^2-1)^+\bigr)^{(k-1)/2}}{\#I}\prod_{a<\varrho}|x_a|^{k+1+2(\varrho-a-1)}
\int_\pfrG W^\flat(y)\otimes\overline{W^\flat\bigl((w_x\sigma)^{-1}y\bigr)}\,dy\ \in\mathbb W\otimes\overline{\mathbb W},
\]
with $c_G=\pbvol\auU(k)^2/(\pbvol\auU(r)\pbvol\auU(k-1))$ as in {\cite[\S3.3.4]{P51}}, and
$\pfrX(x)=\pfrN_{\sigma_x}(x)$. (The integrand has compact support, by {\cite[Proposition~3.41]{P51}}
and the support of $W^\flat$, as in {\cite[\S3.3.2]{P51}}.)

\subsubsection*{Hypotheses}
Let $b\ge2$, $\rho=b^{6n+2}$, $f_W\in\N$ and $\pbdeg_\tau>0$. We assume:
\begin{enumerate}
\item[(H0)] (\emph{sizes}) $8r\le\rho$; $k=s_0+3n+6$ with $s_0>10^9n^3$; $n^7\le k$; $\varrho\ge10^8$;
\item[(H1)] (\emph{cut pair}) $\int_{\pfrU\backslash\pfrG}\|W\|^2=\#I$, $\int_{\pfrU\backslash\pfrG}\|W^\flat\|^2\le\#I$,
$\int_{\pfrU\backslash\pfrG}\|W^\flat-W\|^2\le b^{-2r^3f_W}\#I$, and $W^\flat(y)\ne0$ only if
$\|uy\|\le\rho$ for some $u\in\pfrU$ (operator norm);
\item[(H2)] (\emph{formal degree}) $\int_\pfrG|\pfrip{W}{R(x)W}|^2\,dx\le(\#I)^2/\pbdeg_\tau$;
\item[(H3)] $\pbdeg_\tau\,\pbmTh\le b^{4r^2}$, with $\pbmTh=\pi^{(2\varrho+1)k}\pbvol\auU(k)/\pbvol\auU(r)$;
\item[(H4)] (\emph{balls}) $\pbdeg_\tau\cdot\pbvol\{g\in\pfrG:\|g\|\le T\}\le T^{9r^2}b^{99r^2}$ for $T\ge1$;
\item[(H5)] $f_W\ge4000n$.
\end{enumerate}
The Haar measure $dx$ of $\pfrG$ in (H2), (H4) and in $\pfrN_\sigma$ is the one fixed in
{\cite[\S3.2.1]{P51}}; (H2) says that $\pbdeg_\tau$ is at most the formal degree of the discrete series
generated by $W$ for that measure.

\begin{theorem}[{\cite[Theorem~3.51]{P51}}\,\textup{(R)}]\label{pfr:R}
Assume \textup{(H0)}--\textup{(H5)}. There are functions $J_i=J^\tau_i$, $0\le i\le n$, on pairs of
sequences of non-zero complex numbers such that:
\begin{enumerate}
\item[(R-a)] $(J_i)_{0\le i\le n}$ is a sieve system ({\cite[Definition~3.44]{P51}}) at the place
$(\C,\text{Lebesgue},c\mapsto e^{c-\bar c},w\mapsto w^{-1})$ with mass $(1+b^{-10})/\pbdeg_\tau$,
weights $\Lambda^\infty_i$ and decay $\beta^\infty$ of {\cite[(38)]{P51}}:
\[
\Lambda^\infty_i(p)=\prod_{l=1}^i\bigl[|\delta_l|\le b^{60n^2}\bigr]\max(1,|\delta_l|^{-1})^{n^6},\quad
\delta_l=p_0\cdots p_{l-1};\qquad\beta^\infty(\xi)=\pbmTh\,b^{2000n^9}\,\pfrell_{k-n^7}(|\xi|^{-1}),
\]
$\pfrell_K(t)=t^K\min(1,t)^6$;
\item[(R-b)] $J_n(x,x')=\pfrip{\pfrX(x)}{\pfrX(x')}$; and for every $\sigma\in\pfrG^{(\varrho)}$ and all $x,x'$
with $x_\varrho=x'_\varrho=h(\sigma z,z)$, $J_n(x,x')=\pfrip{\pfrN_\sigma(x)}{\pfrN_\sigma(x')}$.
\end{enumerate}
\end{theorem}

\begin{remark}[Comparison with {\cite[Theorem~3.51]{P51}}]\label{pfr:rem-match}
(a) \emph{The constants agree exactly.} The mass $(1+b^{-10})/\pbdeg_\tau$, the weights
$\Lambda^\infty$ (radius $b^{60n^2}$, exponent $n^6$, product over the partial products
$\delta_l=p_0\cdots p_{l-1}$, $1\le l\le i$) and the decay $\beta^\infty$ (constant $\pbmTh\,b^{2000n^9}$,
the same $\pbmTh$ as in {\cite[\S3.3.4]{P51}}, exponent $k-n^7$, the factor $\min(1,t)^6$) are those of
{\cite[(38)]{P51}} and {\cite[Theorem~3.51\,(R)]{P51}}, which is all that {\cite[Theorem~3.50]{P51}},
{\cite[Lemma~3.52]{P51}} and {\cite[Theorem~3.53]{P51}} use from \textup{(R)}. The constant $c_G$ in
$\pfrN_\sigma$ is that of {\cite[\S3.3.4]{P51}}: $c_G=\frac{\pbvol\auU(k)}{\pbvol\auU(r)}\cdot
\frac{\pbvol\auU(k)}{\pbvol\auU(k-1)}=\frac{\pbvol\auU(k)}{\pbvol\auU(r)}\cdot\frac{2\pi^k}{(k-1)!}$. The
normalisation of $\pfrN_\sigma$ is that of {\cite[\S3.3.7]{P51}}. The indexing of
{\cite[Definition~3.44]{P51}} matches: (S1) is required for $i\le\varrho$ (integration over the entry of
index $i$ of the forms $J_{i+1}$), (S5) for $1\le i\le\varrho$ with $\xi=p_{i-1}^{-1}-q_{i-1}^{-1}$.

(b) \emph{What must be added to {\cite[Theorem~3.51\,(R)]{P51}}.} First, the identification (R-b) of the
top form with the Gram form of the normalised orbital integrals: Siegel's step \textup{(S)} is about
the forms $J_n^\tau$, and it is through (R-b) that they are related to the orbit sums of
{\cite[Proposition~3.42]{P51}}. Second, the hypotheses (H0)--(H5) on the data and the Bessel functions;
they are properties of the Bessel functions of {\cite[Theorem~3.22]{P51}} and of the data, and have to be
checked there. Among them, (H0) asks $\varrho\ge10^8$ and $k\ge n^7$, $k>10^9n^3+3n+6$, which is
stronger than the condition $n\ge100$ listed after {\cite[Theorem~3.53]{P51}}, but holds for the data: by
{\cite[Lemma~13.4]{P46}} (S1), (S2), (S7), $\varrho=m/2-1\ge\sqrt x/2-1$ with $x\ge128^{100}$, $k\ge x^5/2$ and
$n\le\sqrt x$, so $\varrho\ge10^8$, $k\ge x^5/2\ge n^{10}/2\ge n^7$ and $k-3n-6\ge n^{10}/2-3n-6>10^9n^3$ (these three conditions are also stated directly in (S7)(a)). The remaining condition of (H0), $8r\le\rho$, holds since $8r\le8nr\le b\le b^{6n+2}=\rho$ by {\cite[Lemma~13.4\,(S7)(b)]{P46}}.
(c) The functions $J_i$, $i\le\varrho$, are the normalised forms of floor $i$:
$J_i(p,q)=\mathsf c\cdot\pfrI_i(p,q;1)$ for the place datum of \S\ref{pfr:ss-realforms} and the constant
$\mathsf c$ of \eqref{pfr:eq-c}; this is the description of {\cite[\S3.3.7]{P51}} (``$J_i$ is $J_n$
integrated over the last $n-i$ entries'').
\end{remark}

\subsection{The real place: the place datum, the forms and the local estimates}\label{pfr:ss-realforms}

\subsubsection*{The place datum at $\C$}
The \emph{place datum at $\tau$} is: level $M=1$; close pairs the $(\omega_1,\omega_2)$ with
$|\omega_1|^2+|\omega_2|^2<1$ (Euclidean norm of the coordinates); box $\pfrK=\{g:\|g\|\le\rho\}$; all
of $\pfrG$ as positions; $\mu_\pfrH$ the Haar probability measure of $\pfrH=\auU(k)$; $\lambda_l$ ($l<\varrho$)
the image of the measure $|c|^{2d_l}\,dc$ on $\C^\times$ (the \emph{cell density} of floor $l$;
$2d_l=2(k-1)+4(\varrho-l)$) under $c\mapsto w_{l,c}^{-1}$, as in \S\ref{pfr:ss-setting} ($\pfrmod(c)^{d_l}=|c|^{2d_l}$);
and $\theta_\varrho=\lambda_\varrho*\mu_\pfrH$, where $\lambda_\varrho$ is the image of the measure
$\vartheta(c)\,dc$, $\vartheta(c)=\frac{2\pi^k}{(k-1)!}\bigl((|c|^2-1)^+\bigr)^{k-1}$, under $c\mapsto\sigma_c^{-1}$ (the
boosts of \S\ref{pfr:ss-realdata}); that is, $\int F\,d\theta_\varrho=\iint F(\sigma_c^{-1}k)\,\vartheta(c)\,dc\,d\mu_\pfrH(k)$.
The constant $\frac{2\pi^k}{(k-1)!}=\frac{\pbvol\auU(k)}{\pbvol\auU(k-1)}$ is the one in $c_G$
(Remark~\ref{pfr:rem-match}(a)), and it is forced by the identity $N(p^c)N(q^c)=\mathsf c\lambda^2\vartheta(c)\pfrkap_\varrho(p,q)$
of the proof of Proposition~\ref{pfr:top}. The mass condition holds since $\mu_\pfrH$ is a probability measure.

\emph{The constants $c_l$ at $\C$.} For $l<\varrho$, $c_l=\kappa_\C=2$ by Lemma~\ref{pfr:floorquad}. For $l=\varrho$,
$c_\varrho=2$: write $y=y_zz+y_X\in V_\varrho$ with $y_X$ in the span $X_\iota\cong\C^k$ of the $x_i$, so $h(y,y)=|y_X|^2-|y_z|^2$.
For fixed $y_z$, polar coordinates give $dy_X=\mathsf m_k(R)\,dR\,dS$, with $R=|y_X|^2$,
$\mathsf m_k(R)=\pi^kR^{k-1}/(k-1)!$ (the derivative of the volume $\pi^kR^k/k!$ of the ball) and $dS$ the
$\auU(k)$-invariant probability measure on the sphere of radius $\sqrt R$. As $R\mapsto\mathsf m_k(R)\int\varphi\,dS$ is
continuous, \eqref{pfr:eq-slab} gives $\int\varphi\,d\mu_\varrho=\int_{|y_z|>1}\mathsf m_k(|y_z|^2-1)\int\varphi(y_zz+y_X)\,
dS(y_X)\,dy_z$, the inner integral over the sphere of radius $(|y_z|^2-1)^{1/2}$ (dominated convergence; $\varphi$ has
compact support). On the other side, for $\sigma=\sigma_c^{-1}k$, $\sigma^{-1}z=k^{-1}\sigma_cz=-cz+k^{-1}v_c$ with $|v_c|^2=
|c|^2-1$ (proof of Proposition~\ref{pfr:top}), and $k^{-1}v_c$ is distributed by $dS$ when $k$ is distributed by
$\mu_\pfrH$. So the image of $\theta_\varrho$ is $\int\vartheta(c)\int\varphi(-cz+y_X)\,dS\,dc=2\int\varphi\,d\mu_\varrho$
(substitute $y_z=-c$, and $\vartheta(c)=2\mathsf m_k(|c|^2-1)$ for $|c|>1$). Thus $c_l=2$ for every $l\le\varrho$ at the real place. The gauge of scale
$L\ge\rho$, $\Gamma_L$, has $\Gamma_L(i)=\{c:|c|\le L^i\}$ and loss $\Lambda=L^2$. All of this is the
specialisation of \S\ref{pfr:ss-setting}--\S\ref{pfr:ss-gauge}; that the gauge axioms (G0)--(G5)
hold for $\Gamma_L$ under (H0) is Lemma~\ref{pfr:realgauge}.

\begin{lemma}[Haar measure]\label{pfr:haar}
The group measure $\pfrM$ of the place datum at $\tau$ is $\lambda$ times the Haar measure $dx$ of
$\pfrG$ for a number $0<\lambda<\infty$ with $\lambda\,\pbvol\auU(k)/\pbvol\auU(r)\le1$; and for every
measurable $\phi\ge0$ on $\pfrG$ invariant under left translation by $\pfrU$,
$\int\phi\,d\theta_{\ge0}=\lambda\int_{\pfrU\backslash\pfrG}\phi$.
\end{lemma}
(Proof in \S\ref{prp:ss-haar}.)

\begin{lemma}[The gauge at $\C$]\label{pfr:realgauge}
Under \textup{(H0)}, for every $L\ge\rho$ the filtration $\Gamma_L$ with loss $L^2$ is a gauge for the
place datum at $\tau$, and its volume exponent is $s=k-(3n+1)-2\cdot2-1=s_0$ ($x_{i_0}$ and two blocks of rank
$2m+2=2$ are set aside, $m=0$); by (H0), $s\ge10^{10}n^4$. The depth of $e_a$ at
$(g,gx)$ is $\pi^{-1}\inf_\zeta(|g^{-1}\zeta|^2+|(gx)^{-1}\zeta|^2)$ over the duals $\zeta$ of $e_a$ in
$Y^{(a)}$.
\end{lemma}
(Proof in \S\ref{prp:ss-gauge}.)

\subsubsection*{The forms}
Put
\begin{equation}\label{pfr:eq-c}
\mathsf c=\frac{\pbvol\auU(k)}{\pbvol\auU(r)}\cdot\frac1{(\lambda\,\#I)^2},\qquad
J_i(p,q)=\mathsf c\,\pfrI_i(p,q;1)\ (0\le i\le\varrho),\qquad J_n(p,q)=\pfrip{\pfrX(p)}{\pfrX(q)},
\end{equation}
where $\pfrI_i$ is the form of floor $i$ of the place datum at $\tau$ for the cut Bessel function
$W^\flat$ (that $W^\flat$ is a cut Bessel function for this datum is (H1) with $M=1$).

\begin{proposition}[Absolute convergence and support]\label{pfr:conv}
Under \textup{(H0)}--\textup{(H1)}, for $j\le\varrho$ and all $p,q$, $\int\pfrF^+_j\,d\pfrM<\infty$.
If $|p_0\cdots p_{l-1}|>\rho^{9n}$ or $|q_0\cdots q_{l-1}|>\rho^{9n}$ for some $l\le j$, then
$\pfrF^+_j=0$, so $\pfrI_j(p,q;\omega)=0$ for $|\omega|\le1$.
\end{proposition}

\begin{proposition}[The top floor]\label{pfr:top}
Under \textup{(H0)}--\textup{(H1)}: $\pfrX$ is continuous; for $\sigma\in\pfrG^{(\varrho)}$ and
$x_\varrho=x'_\varrho=h(\sigma z,z)$, $\pfrip{\pfrX(x)}{\pfrX(x')}=\pfrip{\pfrN_\sigma(x)}{\pfrN_\sigma(x')}$; and
for $\xi\in\C$ the function $c\mapsto e^{\xi c-\overline{\xi c}}\pfrip{\pfrX(p^c)}{\pfrX(q^c)}$ ($p^c,q^c$:
the entry of index $\varrho$ replaced by $c$) is integrable, with integral
$\mathsf c\,\pfrI_\varrho\bigl(p,q;\ g\mapsto e^{\xi Z_{zz}(g)-\overline{\xi Z_{zz}(g)}}\bigr)$.
\end{proposition}
(Proposition~\ref{pfr:conv} is proved in \S\ref{prp:ss-sum}, Proposition~\ref{pfr:top} in \S\ref{prp:ss-top}.)

\begin{proposition}[The bound for a form]\label{pfr:formbound}
Under \textup{(H0)}--\textup{(H1)}, let $i\le\varrho$, $p,q$ sequences, $\nu\in\C^\times$, and let
$\omega$ be the weight $g\mapsto e^{\nu Z_{ii}(g)-\overline{\nu Z_{ii}(g)}}$, or, if $i\ge1$ and
$\nu=p_{i-1}^{-1}-q_{i-1}^{-1}$, the weight $\omega=1$. Then
\[
\#I\cdot\mathsf c\,|\pfrI_i(p,q;\omega)|\le\pbmTh\,b^{2000n^9}\,\pfrell_{k-n^7}(|\nu|^{-1})\,
\Lambda^\infty_i(p)\Lambda^\infty_i(q).
\]
\end{proposition}
(Proof in \S\ref{prp:ss-sum}.)

\begin{proposition}[Remote positions]\label{pfr:remote}
Under \textup{(H0)}--\textup{(H1)}, let $\mathcal B$ be the set of the $x\in\pfrG$ all of whose entries
of $x-1$ have modulus at most $(b^{6r}\rho^{11n})^{99n}$ (the \emph{position ball}). Then
\[
\mathsf c\,\pfrkap_0\int_{\pfrG\smallsetminus\mathcal B}\pfrF^+_0\,d\pfrM\le\frac{\pbmTh\,\rho^{10^9n^4}\rho^{24k}
(b^{6r})^{10^8n^4}}{(b^{6r})^{s_0}}.
\]
\end{proposition}
(Proof in \S\ref{prp:ss-sum}.)

\begin{lemma}[Far means deep]\label{pfr:deepreal}
Let $p,q$ be heads of length $j\le\varrho$ for $\Gamma_L$, $L\ge\rho$, let $T\ge L^{11n}$ and
$T^{-1}\le\tau'/2$ with $0<\tau'\le1$. If some entry of $x-1$ has modulus $>T^{99n}$, then $x$ has a deep
vector of depth $\tau'^{-2}$.
\end{lemma}
(Proof in \S\ref{prp:ss-deep}.)

\subsection{The mass}\label{pfr:ss-mass}

\begin{proposition}[The mass at a real place]\label{pfr:mass}
Assume \textup{(H0)}--\textup{(H5)}. Then $|J_0(p,q)|\le(1+b^{-10})/\pbdeg_\tau$ for all $p,q$.
\end{proposition}

We first collect three lemmas. Write $\mathsf C(x)=\pfrip{W^\flat}{R(x)W^\flat}$ and
$\mathsf C_0(x)=\pfrip{W}{R(x)W}$ for the matrix coefficients.

\begin{lemma}\label{pfr:F0}
For every $x$ and all $p,q$, $\pfrF_0(x;1)=\lambda^2|\mathsf C(x)|^2$, where $\lambda$ is the number of
Lemma~\ref{pfr:haar}.
\end{lemma}
\begin{proof}
At floor $0$ the Weyl elements and the Weyl ratio are $1$, $\pfrm_0=\delta_1$, the weight is $1$, and the
two factors of $\Phi_0(x;1;1)$ are conjugate: $\Phi_0(x;1;1)=|\mathsf k^1_0(1,x)|^2$, where $\mathsf
k^1_0(1,x)=\int\pfrip{W^\flat(y)}{W^\flat(yx)}\,d\theta_{\ge0}(y)$ (the integration over $\theta_0$ in $\Phi_0$ and
the correlation from floor $1$ on combine to the correlation from floor $0$ on). The function
$y\mapsto\pfrip{W^\flat(y)}{W^\flat(yx)}$ is invariant under left translation by $\pfrU$, because the two
factors $\psi(\sigma(u))$ cancel. By Lemma~\ref{pfr:haar} (applied to the positive and negative parts of
the real and imaginary parts, each integrable by (H1) and Cauchy--Schwarz), $\mathsf k^1_0(1,x)=\lambda\,
\mathsf C(x)$.
\end{proof}

\begin{lemma}\label{pfr:coeffdiff}
For every $x$, $|\mathsf C(x)-\mathsf C_0(x)|\le2b^{-r^3f_W}\,\#I$.
\end{lemma}
\begin{proof}
$\mathsf C(x)-\mathsf C_0(x)=\pfrip{W^\flat-W}{R(x)W^\flat}+\pfrip{W}{R(x)(W^\flat-W)}$. The difference of two
Bessel functions of the same type is a Bessel function, right translation preserves the norm of
$L^2(\pfrU\backslash\pfrG)$, and by Cauchy--Schwarz and (H1) each term has modulus at most
$(b^{-2r^3f_W}\#I)^{1/2}(\#I)^{1/2}$.
\end{proof}

\begin{lemma}\label{pfr:ballnorm}
Let $\mathcal B$ be the position ball of Proposition~\ref{pfr:remote}. Under \textup{(H0)},
$\|x\|\le b^{700nr}$ for $x\in\mathcal B$, and
$\int_{\mathcal B}|\mathsf C(x)|^2dx\le(1+b^{-12})^2(\#I)^2/\pbdeg_\tau$ under \textup{(H1)}, \textup{(H2)},
\textup{(H4)}, \textup{(H5)}.
\end{lemma}
\begin{proof}
Let $R_0=(b^{6r}\rho^{11n})^{99n}$. For $x\in\mathcal B$ the entries of $x-1$ are at most $R_0$, so
$\|x\|\le1+r^2R_0\le\rho^2R_0$ (as $r^2+1\le\rho^2$ by $8r\le\rho$). With $\rho=b^{6n+2}\le b^{7n}$ (as $n\ge2$),
$\rho^2R_0\le b^{14n+99n(6r+77n^2)}\le b^{700nr}$, because $14n+7623n^3\le106nr$ when $r\ge10^9n^3$ (which
follows from (H0): $r=2\varrho+k+1\ge k>10^9n^3$).

For the second claim, let $\alpha=(\#I)\pbdeg_\tau^{-1/2}$, so that $\int_{\mathcal B}|\mathsf C_0|^2\le\alpha^2$ by
(H2). Let $T=b^{700nr}$ and $V_0=T^{9r^2}b^{99r^2}$; by the first claim and (H4),
$\pbvol(\mathcal B)\le V_0/\pbdeg_\tau$. By Lemma~\ref{pfr:coeffdiff},
\[
\int_{\mathcal B}(2b^{-r^3f_W}\#I)^2\,dx\le(2b^{-r^3f_W})^2V_0\cdot\alpha^2\le(b^{-12}\alpha)^2,
\]
since $(2b^{-r^3f_W})^2V_0\le b^{2+6300nr^3+99r^2-2r^3f_W}\le b^{-24}$: indeed $2+6300nr^3+99r^2+24\le
2r^3\cdot4000n\le2r^3f_W$ by (H5) and $r\ge1$, $n\ge1$. By Minkowski's inequality in $L^2(\mathcal B)$
applied to $|\mathsf C|\le|\mathsf C_0|+2b^{-r^3f_W}\#I$,
$\bigl(\int_{\mathcal B}|\mathsf C|^2\bigr)^{1/2}\le\alpha+b^{-12}\alpha$.
\end{proof}

\begin{proof}[Proof of Proposition~\ref{pfr:mass}]
By Proposition~\ref{pfr:conv}, Lemma~\ref{pfr:L1} and $\pfrkap_0=1$,
$|J_0(p,q)|\le\mathsf c\int\pfrF_0(x;1)\,d\pfrM(x)$. Split the integral into $\mathcal B$ and its
complement. On $\mathcal B$, by Lemmas~\ref{pfr:haar}, \ref{pfr:F0} and~\ref{pfr:ballnorm},
\[
\mathsf c\int_{\mathcal B}\pfrF_0\,d\pfrM=\mathsf c\,\lambda\cdot\lambda^2\int_{\mathcal B}|\mathsf C|^2dx\le
\frac{\pbvol\auU(k)}{\pbvol\auU(r)}\frac{\lambda^3}{\lambda^2(\#I)^2}\cdot\frac{(1+b^{-12})^2(\#I)^2}{\pbdeg_\tau}
\le\frac{(1+b^{-12})^2}{\pbdeg_\tau},
\]
using $\lambda\,\pbvol\auU(k)/\pbvol\auU(r)\le1$. On the complement, $\pfrF_0\le\pfrF_0^+$
(Lemma~\ref{pfr:L1}), and by Proposition~\ref{pfr:remote} and (H3) the contribution is at most
\[
\pbmTh\,\rho^{10^9n^4}\rho^{24k}(b^{6r})^{10^8n^4-s_0}\le\frac{b^{4r^2}}{\pbdeg_\tau}\,b^{7n(10^9n^4+24k)+6r\cdot
10^8n^4-6rs_0}\le\frac{b^{-12}}{\pbdeg_\tau}.
\]
For the last step write $r=s_0+5n+5$ and $k=s_0+3n+6$ (by (H0) and $r=2\varrho+k+1$), and note
$s_0\ge10^{10}n^4$: by (H0), $s_0=k-3n-6\ge n^7-3n-6$ and $n\ge10^8$. We need $E+12\le6rs_0$ with
$E=4r^2+7\cdot10^9n^5+168nk+6\cdot10^8n^4r$. Now $6rs_0-4r^2=2r(s_0-10n-10)\ge1.9\,rs_0$, while
$6\cdot10^8n^4r\le0.06\,rs_0$, $7\cdot10^9n^5\le0.7\,n\cdot10^{10}n^4\le0.7\,rs_0$ (as $r\ge s_0\ge n$),
and $168nk+12\le169\,nr\le0.01\,rs_0$. Altogether
\[
|J_0(p,q)|\le\frac{(1+b^{-12})^2+b^{-12}}{\pbdeg_\tau}\le\frac{1+b^{-10}}{\pbdeg_\tau},
\]
since $(1+u^{12})^2+u^{12}=1+(3+u^{12})u^2\cdot u^{10}\le1+u^{10}$ for $0<u\le\frac12$.
\end{proof}

\subsection{Proof of Theorem~\ref{pfr:R}}\label{pfr:ss-Rproof}

\begin{proof}[Proof of Theorem~\ref{pfr:R} from Propositions~\ref{pfr:conv}--\ref{pfr:formbound}, Lemma~\ref{pfr:descent} and Proposition~\ref{pfr:mass}]
Let $J_i$ be as in \eqref{pfr:eq-c}. Heads are sequences of non-zero complex numbers; in (S1) the entry
$c$ ranges over $\C$, and $c=0$ (where the head is not defined) is a null set.

(S1) Let $i\le\varrho$, $p,q$ heads, $\xi\in\C$. If $i<\varrho$, by Proposition~\ref{pfr:conv} and
Lemma~\ref{pfr:descent} (with $\psi(c)=e^{c-\bar c}$), $c\mapsto e^{\xi c-\overline{\xi c}}J_{i+1}(p^c,q^c)$ is
integrable and its integral is $\mathsf c\,\pfrI_i(p,q;\omega_\xi)$ with $\omega_\xi(g)=e^{\xi Z_{ii}(g)-
\overline{\xi Z_{ii}(g)}}$. If $i=\varrho$ the same holds by Proposition~\ref{pfr:top}. For $\xi=0$,
$\omega_0=1$ and the integral is $J_i(p,q)$. For $\xi\ne0$, by Proposition~\ref{pfr:formbound} and
$\#I\ge1$,
$|\mathsf c\,\pfrI_i(p,q;\omega_\xi)|\le\pbmTh\,b^{2000n^9}\pfrell_{k-n^7}(|\xi|^{-1})\Lambda^\infty_i(p)\Lambda^\infty_i(q)
=\beta^\infty(\xi)\Lambda^\infty_i(p)\Lambda^\infty_i(q)$.

(S2) For heads $p^{(1)},\dots,p^{(\kappa)}$ (with any common last entry) the matrix
$(\pfrip{\pfrX(p^{(a)})}{\pfrX(p^{(b)})})_{a,b}$ is a Gram matrix, hence hermitian and positive
semidefinite.

(S3) $J_n$ is continuous because $\pfrX$ is (Proposition~\ref{pfr:top}).

(S4) is Proposition~\ref{pfr:mass}.

(S5) Let $1\le i\le\varrho$ and $\xi=p_{i-1}^{-1}-q_{i-1}^{-1}\ne0$. By Proposition~\ref{pfr:formbound}
with $\omega=1$, $|J_i(p,q)|\le\beta^\infty(\xi)\Lambda^\infty_i(p)\Lambda^\infty_i(q)$.

(R-b) is the definition of $J_n$ and the second claim of Proposition~\ref{pfr:top}.
\end{proof}

\subsection{Proofs of Theorems~\ref{pfr:nogain} and~\ref{pfr:deep}}\label{pfr:ss-gainsproof}

We keep the assumptions of \S\ref{pfr:ss-shared}. Two facts on the floors are used without proof in
this subsection; they are stated exactly.

\begin{lemma}[Weyl invariance of the lower floors]\label{pfr:weylinv}
Let $j\le\varrho$, $p$ a sequence of units and $\Phi\ge0$ measurable and invariant under left
translation by $\pfrG^{(j)}$. Then $\int\Phi(w_{j,p}^{-1}g)\,d\pfrm_j(g)=\int\Phi(g)\,d\pfrm_j(g)$.
\end{lemma}

\begin{lemma}[The inverse weight on the floors below a head]\label{pfr:invweight}
Let $p$ be a head of length $j\le\varrho$ for $\Gamma$, $x\in\pfrP$, $j\le b'$, and let $\Phi\ge0$ be
measurable, invariant under left translation by $\pfrR^{(b')}$ and zero at the points not near at $b'$.
Then
\[
\int\Phi(g)\,\pfrwt^{-}_j(w_{j,p}g)\,d\pfrm_j(g)\le\Bigl(\prod_{a<j}\mathsf c_a(p_a^{-1})\Bigr)C_\Gamma
\int\Phi\,d\theta_{[0,j)},
\]
where $\pfrwt^-_j(g)=\pfrwt_j(g,x)^{-1}$ if $g$ is near at $j$, and $0$ otherwise.
\end{lemma}
(Proofs in \S\ref{pfr:ss-floors}.)

\subsubsection*{Elementary facts on correlations}
Let $\mathsf k=\mathsf k^{d'}$ for a cut Bessel function $W$ and an element $d'$ commuting with
$\pfrG^{(j)}$ (for us $d'=1$ or the Weyl ratio $w_q^{-1}w_p$, which is diagonal on the hyperbolic
pairs of index $<j$ and the identity elsewhere). We use:
\begin{enumerate}
\item[(C1)] $\mathsf k_l(g,y)=\int\mathsf k_{l+1}(\sigma g,\sigma y)\,d\theta_l(\sigma)$, by the
definition $\theta_{\ge l}=\theta_{\ge l+1}*\theta_l$ and Fubini's theorem (the integrals converge
absolutely: $W$ is bounded and, by the cut condition and the compactness of $\pfrK$, the integrand is
supported in a set of finite measure); hence $|\mathsf k_l(g,y)|\le\pfrB_l(g,y)$, the pivot norm
$\int|\mathsf k_{l+1}(\sigma g,\sigma y)|d\theta_l(\sigma)$.
\item[(C2)] If $\mathsf k_l(g,y)\ne0$ then $g$ is near at $l$: some $x'g$, $x'\in\pfrG^{(l)}$, has
$W(x'g)\ne0$, hence $ux'g\in\pfrK$ for some $u\in\pfrU$; and the element $ux'$ lies in $\pfrR^{(l)}$
after the factor of $u$ that moves the $f_a$, $a\ge l$, is absorbed (precisely: $\pfrR^{(l)}\supseteq
\pfrU\pfrG^{(l)}$).
\item[(C3)] The depths $\pfrdep_a(g,x)$, $a<l$, and therefore $\pfrwt_l(g,x)$, are invariant under
$g\mapsto\xi g$ for $\xi\in\pfrR^{(l)}$, in particular for $\xi\in\pfrG^{(l)}$: the stable vectors for
$a$ at $(\xi g,\xi gx)$ are the $\xi\omega$ with $\omega$ stable at $(g,gx)$, and $\xi$ preserves
$\langle e_b:b\le a\rangle$ and the coordinate on $e^*_a$ modulo that span.
\item[(C4)] $|\mathsf k_{l+1}(ng,n'y)|=|\mathsf k_{l+1}(g,y)|$ for $n,n'\in\pfrU_{(<l)}$ (the equivariance of
the correlation by a character, (K1) of Lemma~\ref{pfs:kernels}); the function $\mathsf a_l$ is invariant under
$\pfrR^{(l+1)}$ ((K3)) and $\mathsf b_l$ under $\pfrR^{(l)}$ (Lemma~\ref{pfr:floorright}, descending induction) on the
left, and both vanish off the points near at $l+1$, resp.\ $l$, by (C2).
\end{enumerate}

\begin{lemma}[A weighted square on one floor]\label{pfr:floorsq}
For $l\le\varrho$ and all $g$:
\[
\int|\mathsf k_{l+1}(\sigma g,\sigma gx)|^2\,d\theta_l(\sigma)\le\mathsf b_l(g)\,\pfrwt_l(g,x)^{-1}.
\]
If moreover the mass of floor $l$ is at most $V$, then $\pfrB_l(g,gx)^2\le V\int|\mathsf k_{l+1}(\sigma g,\sigma
gx)|^2\,d\theta_l(\sigma)$.
\end{lemma}
\begin{proof}
$\theta_l$ is carried by $\pfrG^{(l)}$, on which $\pfrwt_l(\sigma g,x)=\pfrwt_l(g,x)$ by (C3); so the left
side is $\pfrwt_l(g,x)^{-1}\int\mathsf a_l(\sigma g)\,d\theta_l(\sigma)\le\pfrwt_l(g,x)^{-1}\mathsf b_l(g)$ (the
weight is finite and non-zero, as each depth is). For the second claim, by (C2) the integrand of
$\pfrB_l(g,gx)$ vanishes unless $\sigma g$ is near at $l+1$, and by Cauchy--Schwarz on that set, of
$\theta_l$-measure at most $V$, $\pfrB_l^2\le V\int|\mathsf k_{l+1}|^2$.
\end{proof}

\begin{lemma}[Separating the frames]\label{pfr:separate}
Let $p$ be a head of length $j\le\varrho$, $x\in\pfrP$, $K\ge0$ and $f\ge0$ measurable on $\pfrG$ with
\[
f(g)^2\le K\,C_\Gamma\,\frac{\mathsf b^{(1)}_j(g)}{\pfrwt_j(g,x)}\cdot\frac{\mathsf b^{(d)}_j(w_p^{-1}g)}{\pfrwt_j(w_p^{-1}g,x)}
\quad\text{wherever }f(g)\ne0,
\]
where $\mathsf b^{(1)}$, $\mathsf b^{(d)}$ are the functions $\mathsf b$ of the kernels $\mathsf k^1$,
$\mathsf k^{d}$. Then $\bigl(\int f\,d\pfrm_j\bigr)^2\le K\,\mathsf Q(x)$.
\end{lemma}
\begin{proof}
Put $B=\mathsf b^{(1)}_j$, $B'=\mathsf b^{(d)}_j$; by (C4) both are $\pfrR^{(j)}$-invariant and vanish
off the points near at $j$, so we may replace $B(g)/\pfrwt_j(g,x)$ by $B(g)\pfrwt^-_j(g)$, and likewise
for $B'$. Crossing the weights, $f(g)^2\le\bigl(B'(w_p^{-1}g)\pfrwt^-_j(g)\bigr)\cdot KC_\Gamma\bigl(B(g)
\pfrwt^-_j(w_p^{-1}g)\bigr)$, and by Cauchy--Schwarz
\[
\Bigl(\int f\,d\pfrm_j\Bigr)^2\le\int B'(w_p^{-1}g)\pfrwt^-_j(g)\,d\pfrm_j\cdot KC_\Gamma\int B(g)\pfrwt^-_j(w_p^{-1}g)\,
d\pfrm_j .
\]
In the first integral substitute $g\mapsto w_pg$ by Lemma~\ref{pfr:weylinv}: the function
$g\mapsto B'(g)\pfrwt^-_j(w_pg)$ is $\pfrG^{(j)}$-invariant (as $w_p$ commutes with $\pfrG^{(j)}$ and by
(C3), (C4)); then Lemma~\ref{pfr:invweight} bounds it by $D\int B'\,d\theta_{[0,j)}$ with
$D=\prod_{a<j}\mathsf c_a(p_a^{-1})C_\Gamma$. In the second, $w_p^{-1}=w_{\bar p}$ (indeed $w_p^{-1}e_a=\bar p_ae_a^*$ and $w_p^{-1}e_a^*=p_a^{-1}e_a$),
$\bar p$ is again a head for $\Gamma$ (the filtration is stable under the involution) and
$\mathsf c_a(\bar p_a^{-1})=\mathsf c_a(p_a^{-1})$; so Lemma~\ref{pfr:invweight} for the head $\bar p$
gives the bound $D\int B\,d\theta_{[0,j)}$. Hence
$(\int f)^2\le KC_\Gamma D^2\,\mathsf S_j(\mathsf k^1;x)\,\mathsf S_j(\mathsf k^d;x)$. Finally
$\mathsf S_j(\mathsf k^1;x)\le\mathsf S_0(\mathsf k^1;x)$ (unfolding $\mathsf b_0$ down to floor $j$ gives
$\mathsf S_0=\mathsf S_j+\sum_{l<j}\int\mathsf a_l\,d\theta_{[0,l+1)}\ge\mathsf S_j$), and $KC_\Gamma D^2\mathsf S_0
\mathsf S_j=K\mathsf Q(x)$ by definition.
\end{proof}

\begin{proof}[Proof of Theorem~\ref{pfr:nogain}]
Let $y$ be a point and $z=w_p^{-1}y$. By Lemma~\ref{pfr:floorsq} for both kernels,
\[
\bigl(\pfrB^1_j(y,yx)\pfrB^{d}_j(z,zx)\bigr)^2\le V^2\frac{\mathsf b^{(1)}_j(y)}{\pfrwt_j(y,x)}\cdot\frac{\mathsf
b^{(d)}_j(z)}{\pfrwt_j(z,x)},
\]
and $C_\Gamma\ge1$. Lemma~\ref{pfr:separate} with $K=V^2$ and $f(y)=\pfrB^1_j(y,yx)\pfrB^d_j(z,zx)$ gives the
claim, since $\pfrF^+_j(x)=\int f\,d\pfrm_j$.
\end{proof}

\subsubsection*{The cut at the first deep vector}
For a kernel $\mathsf k$, a position $x$, $D>0$ and $a<\varrho$ let $\mathsf d_a(y)=|\mathsf k_{a+1}(y,
yx)|$ if $\pfrdep_a(y,x)\ge D$, and $0$ otherwise.

\begin{lemma}\label{pfr:cut}
Let $l\le\varrho$ and $g$ be such that for every $\sigma\in\pfrG^{(l)}$ with $\mathsf k(\sigma g,\sigma gx)\ne0$
there is $a$ with $l\le a<\varrho$ and $\pfrdep_a(\sigma g,x)\ge D$. Then
\[
\pfrB_l(g,gx)\le\sum_{l\le a<\varrho}\int\mathsf d_a(\sigma g)\,d\theta_{[l,a+1)}(\sigma).
\]
\end{lemma}
\begin{proof}
Descending induction on $l$. If $l=\varrho$ no $a$ is available, so $\mathsf k(\sigma g,\sigma gx)=0$ for
all $\sigma\in\pfrG^{(l)}$, hence $\mathsf k_{l+1}(\sigma g,\sigma gx)=0$ for $\sigma\in\pfrG^{(l)}$ and
$\pfrB_l=0$ ($\theta_l$ is carried by $\pfrG^{(l)}$). Let $l<\varrho$ and $\sigma\in\pfrG^{(l)}$. If
$\pfrdep_l(\sigma g,x)\ge D$, then $|\mathsf k_{l+1}(\sigma g,\sigma gx)|=\mathsf d_l(\sigma g)$, which is the
term $a=l$. Otherwise, for $\tau\in\pfrG^{(l+1)}$, $\pfrdep_l(\tau\sigma g,x)=\pfrdep_l(\sigma g,x)<D$ by (C3),
so the hypothesis for $\tau\sigma\in\pfrG^{(l)}$ gives a deep vector with $a\ge l+1$; by (C1) and the
induction hypothesis at $(l+1,\sigma g)$,
$|\mathsf k_{l+1}(\sigma g,\sigma gx)|\le\pfrB_{l+1}(\sigma g,\sigma gx)\le\sum_{l<a<\varrho}\int\mathsf
d_a(\tau\sigma g)\,d\theta_{[l+1,a+1)}(\tau)$. Integrating over $\sigma$ against $\theta_l$ and using
$\theta_{[l,a+1)}=\theta_{[l+1,a+1)}*\theta_l$ gives the claim.
\end{proof}

\begin{lemma}\label{pfr:invweight-floors}
For $x\in\pfrP$, $l\le a<\varrho$ and every $g$,
$\int\pfrwt^-_{a+1}(\sigma g)\,d\theta_{[l,a+1)}(\sigma)\le\Lambda^{10^4n(a-l+1)}\,\pfrwt^-_l(g)$.
\end{lemma}
\begin{proof}
By induction it suffices to show $\int\pfrwt^-_{b+1}(\sigma g)\,d\theta_b(\sigma)\le\Lambda^{10^4n}\pfrwt^-_b(g)$
for every $g$ and $b<\varrho$. For $\sigma\in\pfrG^{(b)}$ with $\sigma g$ near at $b+1$, $g$ is near at $b$
(as $\pfrR^{(b+1)}\pfrG^{(b)}\subseteq\pfrR^{(b)}$) and $\pfrwt^-_{b+1}(\sigma g)=\pfrdep_b(\sigma
g,x)^{s}\pfrwt^-_b(g)$ by (C3); if $\sigma g$ is not near at $b+1$ the integrand is $0$. Lemma~\ref{pfr:depth-rare}
concludes.
\end{proof}

\begin{lemma}[The gain on one pivot]\label{pfr:pivotdeep}
Let $x\in\pfrP$, $l\le\varrho$, $\varepsilon>0$, and let $g$ satisfy the hypothesis of
Lemma~\ref{pfr:cut} with $D=\varepsilon^{-1}$. Then
$\pfrB_l(g,gx)^2\le\varepsilon^{s}\,C_\Gamma\,\mathsf b_l(g)\,\pfrwt_l(g,x)^{-1}$.
\end{lemma}
\begin{proof}
Fix $a$ with $l\le a<\varrho$. Where $\mathsf d_a(y)\ne0$ we have $1\le\varepsilon\,\pfrdep_a(y,x)$, $y$ near
at $a+1$ by (C2), and so $\mathsf d_a(y)^2\le\varepsilon^s\pfrdep_a(y,x)^s|\mathsf k_{a+1}(y,yx)|^2=
\varepsilon^s\pfrwt^-_{a+1}(y)\,\mathsf a_a(y)$. By Cauchy--Schwarz, Lemma~\ref{pfr:invweight-floors} and
$\int\mathsf a_a(\sigma g)\,d\theta_{[l,a+1)}(\sigma)\le\mathsf b_l(g)$ (unfold the recursion of $\mathsf b$),
\[
\Bigl(\int\mathsf d_a(\sigma g)\,d\theta_{[l,a+1)}\Bigr)^2\le\varepsilon^s\Lambda^{10^4n\varrho}\,\pfrwt_l(g,x)^{-1}\mathsf b_l(g).
\]
By Lemma~\ref{pfr:cut} and the triangle inequality, $\pfrB_l^2\le\varrho^2\cdot\varepsilon^s\Lambda^{10^4n\varrho}
\pfrwt_l^{-1}\mathsf b_l$, and $\varrho^2\Lambda^{10^4n\varrho}\le\Lambda^{2\varrho+10^4n\varrho}\le C_\Gamma$.
\end{proof}

\begin{proof}[Proof of Theorem~\ref{pfr:deep}]
Let $y$ be a point with $\pfrB^1_j(y,yx)\pfrB^d_j(z,zx)\ne0$, $z=w_p^{-1}y$. The second factor is
bounded by Lemma~\ref{pfr:floorsq}: $\pfrB^d_j(z,zx)^2\le V\mathsf b^{(d)}_j(z)/\pfrwt_j(z,x)$. As it is not
zero, there is $\tau\in\pfrG^{(j)}$ with $W(\tau z)\ne0$ and $W(\tau dzx)\ne0$, $d=w_q^{-1}w_p$; by the cut
condition and since $\tau$ commutes with the Weyl elements, $w_p^{-1}(\rho y)$ and $w_q^{-1}(\rho yx)$ are near
at $j$ for every $\rho\in\pfrG^{(j)}$. Hence for $\rho\in\pfrG^{(j)}$ with $\mathsf k^1(\rho y,\rho yx)\ne0$, the
point $\rho y$ is active at $x$, and by hypothesis it has a deep vector of depth $\varsigma^{-1}$ with index
$\ge j$. Lemma~\ref{pfr:pivotdeep} with $l=j$, $\varepsilon=\varsigma$ gives $\pfrB^1_j(y,yx)^2\le\varsigma^sC_\Gamma
\mathsf b^{(1)}_j(y)/\pfrwt_j(y,x)$. Multiply, use $V\le V^2$, and apply Lemma~\ref{pfr:separate} with
$K=\varsigma^sV^2$.

\emph{The case $j=0$.} Here $\pfrm_0=\delta_1$, $w_p=1$, $d=1$ and the two kernels coincide, so
$\pfrF^+_0(x)=\pfrB^1_0(1,x)^2$. Every point is near at $0$ (as $\pfrR^{(0)}=\pfrG$), so every $\rho\in\pfrG$ with
$\mathsf k^1(\rho,\rho x)\ne0$ is active at $x$ and has a deep vector of depth $\varsigma^{-1}$.
Lemma~\ref{pfr:pivotdeep} with $l=0$, $g=1$, $\varepsilon=\varsigma$ gives
$\pfrB^1_0(1,x)^2\le\varsigma^sC_\Gamma\,\mathsf b^{(1)}_0(1)$ (the weight $\pfrwt_0$ is $1$), and
$\mathsf b^{(1)}_0(1)=\mathsf S_0(\mathsf k^1;x)$ since $\theta_{[0,0)}=\delta_1$. For $j=0$,
$\mathsf Q(x)=C_\Gamma^3\,\mathsf S_0(\mathsf k^1;x)^2$; hence
$\pfrF^+_0(x)^2\le\varsigma^{2s}C_\Gamma^2\,\mathsf S_0(\mathsf k^1;x)^2\le\varsigma^{2s}\,\mathsf Q(x)$, as $C_\Gamma\ge1$.
\end{proof}

\subsection{Facts on the floors: proofs of Lemmas~\ref{pfr:weylinv}, \ref{pfr:invweight}, \ref{pfr:chart}}\label{pfr:ss-floors}

Three structural facts on place data are used; they are stated exactly.

\begin{lemma}[The upper floors]\label{pfr:upperinv}
For $0\le j\le\varrho+2$ let $\rho_j=\bigl((\eta*\theta)\bigr)_{[j,\varrho+2)}$. Then $\pfrM=\rho_j*\pfrm_j$,
$\rho_j$ is carried by $\pfrG^{(j)}$, and $\rho_j*\delta_t=\rho_j$ for every $t\in\pfrG^{(j)}$.
\end{lemma}

\begin{lemma}[One Heisenberg floor after a head]\label{pfr:section}
Let $a<\varrho$, $p$ a head of length $a+1$ for $\Gamma$, $x\in\pfrP$, $y\in\pfrK$ and $\nu\in\pfrU_{(<a)}$. Then
\[
\int\mathbf 1[w_{a+1,p}n\nu y\text{ near at }a+1]\;\pfrdep_a(w_{a+1,p}n\nu y,x)^s\,d\eta_a(n)\le\mathsf c_a(p_a^{-1})
\,\Lambda^{10^6n^2}.
\]
\end{lemma}

\begin{lemma}[Near points on a floor]\label{pfr:nearsize}
Let $b\le\varrho$, $g\in\pfrK$ and $\sigma\in\pfrG^{(b)}$ with $\sigma g$ near at $b+1$. Then $\sigma^{-1}f_b$ has
entries of size $7$ for $\Gamma$.
\end{lemma}
(Lemmas~\ref{pfr:upperinv} and~\ref{pfr:nearsize} are proved below; Lemma~\ref{pfr:section} in \S\ref{pfr:ss-section}.)

\begin{proof}[Proof of Lemma~\ref{pfr:weylinv}]
Write $w=w_{j,p}^{-1}$, $\rho=\rho_j$, $\nu=\pfrm_j$. Since $w$ commutes with $\pfrG^{(j)}$, which carries $\rho$,
$\rho*\delta_w=\delta_w*\rho$. Choose $f>0$ measurable with $\int f\,d\pfrM<\infty$ ($\pfrM$ is a Haar measure,
hence $\sigma$-finite), and put $\mathrm{av}(g)=\int f(yg)\,d\rho(y)$; by Lemma~\ref{pfr:upperinv},
$\mathrm{av}(tg)=\mathrm{av}(g)$ for $t\in\pfrG^{(j)}$, and $\mathrm{av}>0$ everywhere as $\rho\ne0$ ($\pfrM\ne0$).
Let $\mu$ be any $\sigma$-finite measure with $\pfrM=\rho*\mu$. Then $\mathrm{av}<\infty$ $\mu$-a.e., and for
$\Phi\ge0$ measurable and $\pfrG^{(j)}$-invariant,
\[
\int\frac{\Phi}{\mathrm{av}}f\,d\pfrM=\int\!\!\int\frac{\Phi(yg)}{\mathrm{av}(yg)}f(yg)\,d\rho(y)\,d\mu(g)=
\int\frac{\Phi(g)}{\mathrm{av}(g)}\,\mathrm{av}(g)\,d\mu(g)=\int\Phi\,d\mu,
\]
since $\rho$ is carried by $\pfrG^{(j)}$. The left side does not depend on $\mu$. Apply this to $\mu=\nu$ and to
$\mu=\delta_w*\nu$, which is admissible because $\rho*\delta_w*\nu=\delta_w*\rho*\nu=\delta_w*\pfrM=\pfrM$
(left invariance of the Haar measure $\pfrM$). Hence $\int\Phi\,d(\delta_w*\nu)=\int\Phi\,d\nu$, that is
$\int\Phi(wg)\,d\nu(g)=\int\Phi\,d\nu$.
\end{proof}

\begin{proof}[Proof of Lemma~\ref{pfr:invweight}]
\emph{Step 1: a pointwise bound.} We show: for $g$ near at $j$,
\begin{equation}\label{pfr:eq-iw}
\int\pfrwt^-_j(w_{j,p}\nu g)\,d\eta_{[0,j)}(\nu)\le\prod_{a<j}\mathsf c_a(p_a^{-1})\cdot C_\Gamma,
\end{equation}
where $\eta_{[0,j)}=\eta_{j-1}*\dots*\eta_0$ is carried by $\pfrU_{(<j)}$. Let $\xi g\in\pfrK$ with
$\xi\in\pfrR^{(j)}$ and write $\xi=\gamma n_0$ with $\gamma\in\pfrG^{(j)}$, $n_0\in\pfrU_{(<j)}$. The measure
$\eta_{[0,j)}$ is invariant under right translation by $n_0$ and under conjugation by $\gamma$ (which
normalises each Heisenberg group of floor $<j$ and preserves its Haar measure), $\gamma$ commutes with
$w_{j,p}$, and $\pfrwt^-_j$ is $\pfrG^{(j)}$-invariant by (C3); so we may replace $g$ by $y=\xi g\in\pfrK$. (The right
invariance: by the last step of the proof of Lemma~\ref{pfr:rearr}, applied repeatedly, $\eta_{[0,j)}=
\eta_0*\dots*\eta_{j-1}$. As $\pfrU_{(<j)}=\pfrN_0\cdots\pfrN_{j-1}$, it suffices to treat $n\in\pfrN_a$, $a<j$.
Write $\eta_{[0,j)}=\eta_0*\dots*\eta_{a-1}*\eta_a*A$ with $A=\eta_{a+1}*\dots*\eta_{j-1}$, carried by
$\pfrG^{(a+1)}$ since $\pfrN_b\subseteq\pfrG^{(b)}\subseteq\pfrG^{(a+1)}$ for $b>a$. For $x\in\pfrG^{(a+1)}$,
$n_x:=xnx^{-1}\in\pfrN_a$ by (N2), so for $\Phi\ge0$, $\iint\Phi(zxn)\,d\eta_a(z)\,dA(x)=\iint\Phi(zn_xx)\,d\eta_a(z)\,dA(x)
=\iint\Phi(zx)\,d\eta_a(z)\,dA(x)$ by the invariance $\eta_a*\delta_{n_x}=\eta_a$ of (N2). Thus
$\eta_a*A*\delta_n=\eta_a*A$, and so $\eta_{[0,j)}*\delta_n=\eta_{[0,j)}$.) We
prove by induction on $a\le j$ that $I_a:=\int\pfrwt^-_a(w_{a,p}\nu y)\,d\eta_{[0,a)}(\nu)\le\prod_{l<a}\mathsf
c_l(p_l^{-1})\cdot\Lambda^{10^6n^2a}$. For $a=0$ the integrand is $\le1$ ($\pfrwt_0=1$). For the step, write
$\eta_{[0,a+1)}=\eta_a*\eta_{[0,a)}$ and $\nu'=n\nu$. For $n$ in the Heisenberg group of floor $a$, the
weight of length $a$ at $w_{a+1,p}n\nu y$ equals that at $w_{a,p}\nu y$ ($w_{a+1,p}=w_{a,p}\cdot$(the factor of
floor $a$), and both the factor of floor $a$ and $n$ lie in $\pfrG^{(a)}$, under which the depths of index
$<a$ are invariant); if the point is near at $a+1$ it is near at $a$, and
$\pfrwt^-_{a+1}=\pfrwt^-_a\cdot\pfrdep_a^{s}$. Hence, by Lemma~\ref{pfr:section},
$\int\pfrwt^-_{a+1}(w_{a+1,p}n\nu y)\,d\eta_a(n)\le\pfrwt^-_a(w_{a,p}\nu y)\,\mathsf c_a(p_a^{-1})\Lambda^{10^6n^2}$, and
integrating over $\nu$ gives the step. Finally $\Lambda^{10^6n^2j}\le\Lambda^{10^6n^3}=C_\Gamma$.

\emph{Step 2.} $\pfrm_j=\eta_{[0,j)}*\theta_{[0,j)}$ up to the order of the factors; more precisely, since the
Heisenberg groups of floors $<j$ are normalised by the floor measures of the lower floors, we
write $\int\Psi\,d\pfrm_j=\int\!\int\Psi(\nu y)\,d\eta_{[0,j)}(\nu)\,d\theta_{[0,j)}(y)$ (the ``ascending''
convolution of the Heisenberg measures followed by that of the floors). For fixed $y$,
$\Phi(\nu y)=\Phi(y)$ for $\nu\in\pfrU_{(<j)}\subseteq\pfrR^{(b')}$; if $y$ is not near at $b'$, $\Phi(y)=0$;
otherwise $y$ is near at $j$ ($j\le b'$) and \eqref{pfr:eq-iw} applies. Integrating over $y$ gives the claim.
\end{proof}
(The rearrangement is Lemma~\ref{pfr:rearr}.)

\begin{lemma}[Right translation on one floor]\label{pfr:floorright}
Let $b\le\varrho+1$, $\xi\in\pfrR^{(b)}$ and $\Psi\ge0$ measurable on $\pfrG$. Assume either
(a) $\Psi(\zeta\sigma)=\Psi(\sigma)$ for all $\zeta\in\pfrR^{(b+1)}$ and $\sigma\in\pfrG$; or
(b) $\xi\in\pfrG^{(b)}$ and $\Psi(\zeta\sigma)=\Psi(\sigma)$ for all $\zeta\in\pfrR^{(b+1)}\cap\pfrG^{(b)}$ and
$\sigma\in\pfrG^{(b)}$.
Then $\int\Psi(\sigma\xi)\,d\theta_b(\sigma)=\int\Psi(\sigma)\,d\theta_b(\sigma)$.
\end{lemma}
\begin{proof}[Proof of Lemma~\ref{pfr:chart}]
Write $b=a$ and $\Phi=c\cdot\mathbf 1[\text{near at }b+1]$, which is $\pfrR^{(b+1)}$-invariant on the left. If
$g$ is not near at $b$, then no $\sigma g$, $\sigma\in\pfrG^{(b)}$, is near at $b+1$ (as
$\pfrR^{(b+1)}\pfrG^{(b)}\subseteq\pfrR^{(b)}$), and the integral is $0$ since $\theta_b$ is carried by $\pfrG^{(b)}$.
Otherwise let $\xi\in\pfrR^{(b)}$ with $g'=\xi g\in\pfrK$. By Lemma~\ref{pfr:floorright}(a) applied to
$\Psi(\sigma)=\Phi(\sigma g)$, $\int\Phi(\sigma g)\,d\theta_b(\sigma)=\int\Phi(\sigma\xi g)\,d\theta_b(\sigma)=\int\Phi(\sigma g')\,d\theta_b(\sigma)$. For
$\sigma\in\pfrG^{(b)}$ with $\sigma g'$ near at $b+1$, $\sigma^{-1}f_b\in\Gamma(7)$ by Lemma~\ref{pfr:nearsize}, and
$\Phi(\sigma g')\le\Psi_{g'}(\sigma^{-1}f_b)$. As $\theta_b$ maps to $c_b\mu_b$ under $\sigma\mapsto\sigma^{-1}f_b$, the
integral is at most $c_b\int_{\Gamma(7)}\Psi_{g'}\,d\mu_b\le c_bC$, with $c_b\le2$ (\S\ref{pfr:ss-setting}).
\end{proof}

\subsubsection*{The Heisenberg groups and the floors}
We record the elementary group structure behind \S\ref{pfr:ss-setting}. Let $l\le\varrho$, $V_l$ the span of the
basis vectors of floor $\ge l$, and $\pfrN_l$ the Heisenberg group of floor $l$: the isometries of $V_l$
(extended by the identity) of the form $n_{\omega,\beta}\colon f_l\mapsto f_l$, $v\mapsto v-h(v,\omega)f_l$ for $v\perp
f_l,f_l^*$ in $V_{l+1}$-part, $f^*_l\mapsto f_l^*+\omega+\beta f_l$, with $\beta+\bar\beta=-h(\omega,\omega)$ (for $l=\varrho$
we have $f_\varrho=z$ and $\pfrN_\varrho=1$ in our normalisation, the top floor being handled by $\theta_\varrho$ alone).
\begin{enumerate}
\item[(N1)] \emph{$\pfrR^{(l+1)}\cap\pfrG^{(l)}=\pfrG^{(l+1)}\pfrN_l$.} An element $t$ of the left side fixes the basis
vectors of floor $<l$, and $tf_l\in f_l+\langle f_b:b<l\rangle$ is orthogonal to the $e_b^*$, $b<l$, hence $tf_l=f_l$.
So $t$ stabilises $f_l$, hence $f_l^\perp\cap V_l$ and the flag $Af_l\subset f_l^\perp\cap V_l\subset V_l$; the
stabiliser of this flag in $\auU(V_l)$ fixing $f_l$ is $\pfrN_l\rtimes\auU(V_{l+1})$, and $\auU(V_{l+1})=\pfrG^{(l+1)}$.
\item[(N2)] \emph{$\eta_l$ is a Haar measure of $\pfrN_l$}, so $\eta_l*\delta_n=\delta_n*\eta_l=\eta_l$ for $n\in\pfrN_l$;
and for $u\in\pfrG^{(l+1)}$, $u\,n_{\omega,\beta}u^{-1}=n_{u\omega,\beta}$, a linear change of the coordinate
$\omega$ of determinant of modulus $1$, which preserves $\eta_l$; hence $\delta_u*\eta_l=\eta_l*\delta_u$.
\item[(N3)] $\pfrN_a$ for $a<l$ is normalised by $\pfrG^{(l)}$ with the same formula, so the ascending convolution
$\eta_0*\eta_1*\dots*\eta_{l-1}$ commutes with $\delta_u$ for $u\in\pfrG^{(l)}$.
\end{enumerate}

\begin{lemma}\label{pfr:rearr}
$\pfrm_j=(\eta_0*\dots*\eta_{j-1})*\theta_{[0,j)}$.
\end{lemma}
\begin{proof}
Induction on $j$; $j=0$ is trivial. As $\theta_j$ is carried by $\pfrG^{(j)}$, (N3) gives $\theta_j*(\eta_0*\dots*\eta_{j-1})=
(\eta_0*\dots*\eta_{j-1})*\theta_j$ (both sides integrate $\Phi$ to $\iint\Phi(us)\,d\eta(u)\,d\theta_j(s)$ after
the substitution $u\mapsto sus^{-1}$). Hence $\pfrm_{j+1}=\eta_j*\theta_j*\pfrm_j=\eta_j*\theta_j*(\eta_0*\dots*\eta_{j-1})*
\theta_{[0,j)}=\eta_j*(\eta_0*\dots*\eta_{j-1})*\theta_j*\theta_{[0,j)}$. Finally $\eta_j*(\eta_0*\dots*\eta_{j-1})=
\eta_0*\dots*\eta_{j}$: by (N3), conjugation by $n\in\pfrN_j\subseteq\pfrG^{(j)}$ preserves each $\eta_a$, $a<j$
($\omega\mapsto n\omega$ has modulus $|\det n|=1$), hence preserves $\nu'=\eta_0*\dots*\eta_{j-1}$; so $\iint\Phi(nx)\,d\nu'(x)\,
d\eta_j(n)=\iint\Phi(xn)\,d\nu'(x)\,d\eta_j(n)$, that is $\eta_j*\nu'=\nu'*\eta_j$ (details in {\cite[\S3.1.6(i)]{P51}}).
\end{proof}

\begin{proof}[Proof of Lemma~\ref{pfr:floorright}]
In case (b) put $\gamma=\xi$. In case (a) write $\xi=\gamma\nu$ with $\gamma\in\pfrG^{(b)}$ and
$\nu\in\pfrU_{(<b)}$: an element of $\pfrR^{(b)}$ preserves the flag $\langle e_0\rangle\subset\dots\subset\langle
e_0,\dots,e_{b-1}\rangle$ and acts unitriangularly on it, so it lies in the parabolic subgroup of this flag with trivial
torus part, which is $\pfrU_{(<b)}\rtimes\auU(V_b)$; its Levi component $\gamma$ lies in $\auU(V_b)=\pfrG^{(b)}$. For
$\sigma\in\pfrG^{(b)}$, $\sigma\gamma\nu=(\sigma\gamma\nu(\sigma\gamma)^{-1})\,\sigma\gamma$ and $\sigma\gamma\nu(\sigma\gamma)^{-1}\in\pfrU_{(<b)}
\subseteq\pfrR^{(b+1)}$ (conjugation by $\pfrG^{(b)}$ preserves $\pfrU_{(<b)}$), so $\Psi(\sigma\xi)=\Psi(\sigma\gamma)$.
In both cases $\Psi$ restricted to $\pfrG^{(b)}$ is invariant on the left under $\pfrR^{(b+1)}\cap\pfrG^{(b)}$, and
only this is used below.
If $b<n$: $\Psi$ restricted to $\pfrG^{(b)}$ is invariant under $\pfrG^{(b+1)}\pfrN_b=\pfrR^{(b+1)}\cap\pfrG^{(b)}$, which
is the stabiliser of $f_b$ in $\pfrG^{(b)}$ by (N1); by Lemma~\ref{pfs:section}, $\Psi(\sigma)=\Psi_0(\sigma^{-1}f_b)$ on
$\pfrG^{(b)}$ for a Borel function $\Psi_0$ on $V_b$, and $\mu_b$ is carried by the orbit of $f_b$ (Remark~\ref{pfs:witt}:
no transitivity on the whole quadric is needed). Then $\Psi(\sigma\gamma)=\Psi_0(\gamma^{-1}\sigma^{-1}f_b)$, and since
$\theta_b$ maps to $c_b\mu_b$, a multiple of the quadric measure $\mu_b$, which is $\pfrG^{(b)}$-invariant,
$\int\Psi(\sigma\gamma)\,d\theta_b=c_b\int\Psi_0(\gamma^{-1}v)\,d\mu_b(v)=c_b\int\Psi_0\,d\mu_b=\int\Psi\,d\theta_b$. If $b=n$, $\theta_n=\mu_\pfrH$
is right invariant under $\pfrH=\pfrG^{(n)}$. If $b>n$, $\pfrG^{(b)}=1$.
\end{proof}

\begin{proof}[Proof of Lemma~\ref{pfr:upperinv}]
The factorisation is associativity of convolution, and the support claim holds factor by factor ($\eta_a$,
$\theta_a$ are carried by $\pfrG^{(a)}\subseteq\pfrG^{(j)}$ for $a\ge j$). Invariance by descending induction on
$j$; for $j\ge\varrho+2$, $\pfrG^{(j)}=1$. Let $\rho_j=\rho'*\theta_j$ with $\rho'=\rho_{j+1}*\eta_j$. For
$t=un\in\pfrR^{(j+1)}\cap\pfrG^{(j)}$ ((N1)), $\rho'*\delta_t=\rho_{j+1}*\delta_u*\eta_j*\delta_n=\rho_{j+1}*\eta_j=\rho'$
by (N2) and the induction hypothesis. Let $f\ge0$ and $\mathrm{av}(y)=\int f(xy)\,d\rho'(x)$; then $\mathrm{av}(ty)=
\mathrm{av}(y)$ for such $t$, and for $s\in\pfrG^{(j)}$, by Lemma~\ref{pfr:floorright}(b) (with $\xi=s$),
$\int f\,d(\rho_j*\delta_s)=\int\mathrm{av}(ys)\,d\theta_j(y)=\int\mathrm{av}(y)\,d\theta_j(y)=\int f\,d\rho_j$.
\end{proof}

\begin{proof}[Proof of Lemma~\ref{pfr:nearsize}]
Let $\xi\in\pfrR^{(b+1)}$ with $k:=\xi\sigma g\in\pfrK$. Then $\sigma^{-1}\xi^{-1}=g\,k^{-1}$, a product of an element of
$\pfrK$ and the inverse of one; by (G1) and the closure properties of the filtration it has entries of
size $\le7$ ($k^{-1}=J^{-1}k^*J$ has size $5$ by (G0)).
Now $\xi^{-1}f_b=f_b+\sum_{c<b}\lambda_cf_c$ and $\sigma^{-1}f_c=f_c$ for $c<b$; so the coordinates of
$\sigma^{-1}f_b$ of floor $\ge b$ are those of the column $\sigma^{-1}\xi^{-1}f_b$, and those of floor $<b$ vanish
($\sigma^{-1}$ preserves $V_b$). Hence $\sigma^{-1}f_b\in\Gamma(7)$.
\end{proof}

\subsection{Quadrics, sections and the proof of Lemmas~\ref{pfr:quadricX}, \ref{pfr:section}}\label{pfr:ss-section}

\subsubsection*{Coordinates}
For $b\le\varrho$ let $V_b$ be the span of the basis vectors of floor $\ge b$ and let its \emph{floor basis}
consist of $f_b,f_b^*$, the remaining hyperbolic pairs of $V_b$, and the orthonormal vectors $x_i$. In
these coordinates $h$ is the standard form; the coordinates of $v$ on the $x_i$ of $X'$ are the $h(v,x_i)$, the
coordinates of $v_{X'}$ (\S\ref{pfr:ss-gauge}). A vector of size $j$ has floor coordinates of size $j+4$ (they are entries of
$B^*Jv$ with $B$ the basis matrix of size $1$, by (G0)). Recall (G3): $\pbvol\Gamma(i)\le\Lambda^{i+1}$, and
(G4): Lebesgue measure on $A^{2m+2}$ disintegrates along the standard form, $dy=\int\kappa_\beta\,d\beta$,
with $\kappa_\beta(\Gamma(i)^{2m+2})\le\Lambda^{(2m+2)(i+1)}$ for every value $\beta$.

\begin{lemma}[A solved block]\label{pfr:block}
Let $\varrho_b$ be Lebesgue measure on $V_b$ or the quadric measure $\mu_b$ (the disintegration of
Lebesgue measure of $V_b$ along $y\mapsto h(y,y)$, at the value $h(f_b,f_b)$, given by \eqref{pfr:eq-slab}). Fix a block $\beta$ of $2m+2$ of
the orthonormal coordinates, disjoint from those of $X'$ (possible by $2(2m+2)\le\#\iota$). Then, in floor
coordinates $y=(y',y_\beta)$, $\varrho_b=\int\kappa'_{y'}\,dy'$ with Lebesgue measure $dy'$ in the coordinates off
the block and measures $\kappa'_{y'}$ on the block with $\kappa'_{y'}(\Gamma(i)^{2m+2})\le\Lambda^{(2m+2)(i+2)}$ for
every $i$ and every $y'$.
\end{lemma}
\begin{proof}
For Lebesgue measure take $\kappa'_{y'}$ Lebesgue on the block, of mass $\pbvol\Gamma(i)^{2m+2}\le\Lambda^{(2m+2)(i+1)}$ on
$\Gamma(i)^{2m+2}$, and $\Lambda^{(2m+2)(i+1)}\le\Lambda^{(2m+2)(i+2)}$. For $\mu_b$ put $L_i=\Lambda^{(2m+2)(i+2)}$ and let
$U_i$ be the interior of $\Gamma(i+1)^{2m+2}$; then $\Gamma(i)^{2m+2}\subseteq U_i$ (at the real place
$\Gamma(i)=\{|c|\le L^i\}$ with $L\ge r+1\ge2$, so $\Gamma(i)\subseteq\{|c|<L^{i+1}\}$; at the finite places $\Gamma(i)$,
defined by an integrality condition, is open) and $\bigcup_iU_i=A^{2m+2}$. Let $t_0=h(f_b,f_b)$.

\emph{Step 1: $\mu_b(E'\times U_i)\le L_i\,\pbvol(E')$ for every Borel set $E'$ of the coordinates off the block.} Let
$\varphi'\ge0$ and $\chi$ with $0\le\chi\le\mathbf 1_{U_i}$ be test functions as in \eqref{pfr:eq-slab} in the two groups
of coordinates. We have $h(y,y)=h(y',y')+h(y_\beta,y_\beta)$ with the standard form on the block, and by (G4)
$dy_\beta=\int\kappa_s\,ds$. By Tonelli, for each $y'$ the substitution $s=t-h(y',y')$ (a translation, preserving $dh$)
and $\kappa_s(\chi)\le\kappa_s(\Gamma(i+1)^{2m+2})\le L_i$,
\[
\int_{\{h(y,y)-t_0\in B\}}\varphi'(y')\chi(y_\beta)\,dy=\int\varphi'(y')\int_{\{s:\ s+h(y',y')-t_0\in B\}}\kappa_s(\chi)\,ds\,dy'
\le\pbvol(B)\,L_i\int\varphi'\,dy'.
\]
Dividing by $\pbvol(B)$ and letting $B\to0$, \eqref{pfr:eq-slab} gives $\int\varphi'\otimes\chi\,d\mu_b\le L_i\int\varphi'dy'$.
Taking the supremum over $\chi\uparrow\mathbf 1_{U_i}$ ($U_i$ is open; monotone convergence), then the infimum over
$\varphi'\downarrow\mathbf 1_{E'}$ for $E'$ compact (dominated convergence, $\mu_b$ being Radon), gives the claim for
compact $E'$. The Borel measure $\pi_i(E')=\mu_b(E'\times U_i)$ is thus finite on compact sets, hence inner regular,
and the claim follows for every Borel $E'$.

\emph{Step 2: disintegration.} Let $w(y)=\sum_i2^{-i}(1+L_i)^{-1}\mathbf 1_{U_i}(y_\beta)$, a Borel function with
$0<w\le1$. By Step 1 the image $\pi_w$ of $w\mu_b$ under $y\mapsto y'$ satisfies $\pi_w(E')\le\sum_i2^{-i}\pbvol(E')
\le2\pbvol(E')$, so $\pi_w=g\,dy'$ with $0\le g\le2$. By the disintegration theorem (all spaces are locally compact
and second countable), $w\mu_b=\int\lambda_{y'}\,d\pi_w(y')$ with probability measures $\lambda_{y'}$ on the fibre
$\{y'\}\times A^{2m+2}$, identified with the block. Put $\kappa'_{y'}=g(y')\,w^{-1}\lambda_{y'}$; then
$\mu_b=\int\kappa'_{y'}\,dy'$, and for every Borel $E'$, $\int_{E'}\kappa'_{y'}(U_i)\,dy'=\pi_i(E')\le L_i\pbvol(E')$, so
$\kappa'_{y'}(U_i)\le L_i$ for almost every $y'$. Replacing $\kappa'_{y'}$ by $0$ on the null set of the $y'$ where
this fails for some $i$ (countably many conditions), $\kappa'_{y'}(\Gamma(i)^{2m+2})\le\kappa'_{y'}(U_i)\le L_i$ for
every $i$ and $y'$.
\end{proof}

\begin{proof}[Proof of Lemma~\ref{pfr:quadricX}]
More generally, for $\varrho_b$ as in Lemma~\ref{pfr:block} and $\int_{\Gamma(j+4)^{X'}}\phi\le B$ we show
$\int_{\Gamma(j)}\phi(v_{X'})\,\allowbreak d\varrho_b(v)\le\Lambda^{99n(j+6)}B$; the case $j=7$, $\phi\mapsto\phi^s$ is the lemma. A
vector of size $j$ has floor coordinates in $\Gamma(j+4)$, so the integral is at most
$\int_{\Gamma(j+4)}\phi(y_{X'})\,d\kappa'_{y'}(y_\beta)\,dy'$ over $y'$ in $\Gamma(j+4)$. Integrate first over the block
($\le\Lambda^{(2m+2)(j+6)}$), then over the coordinates of $X'$ ($\le B$), then over the remaining $d$ coordinates
($\le(\Lambda^{j+5})^d$). The number of coordinates of $V_b$ off $X'$ is at most $2+2(m+\varrho)+4(m+1)\le99n-(2m+2)$
(by $m\le9n$); so the product is at most $\Lambda^{(2m+2)(j+6)+(99n-(2m+2))(j+5)}B\le\Lambda^{99n(j+6)}B$.
\end{proof}

\subsubsection*{The Heisenberg group of floor $a$ in coordinates}
For $a<\varrho$ the elements of $\pfrN_a$ are $n=n(\omega,\beta)$ with $\omega\in V_{a+1}$ ($r_a:=d_a=\dim V_{a+1}$
coordinates) and $\beta=b-\tau_0h(\omega,\omega)$, $b\in A_{\rm sk}$; as defined in \S\ref{pfr:ss-setting}, $\eta_a$ is
the image of $d\omega\otimes d\beta_{\rm sk}(b)$, where $d\beta_{\rm sk}$ is the Haar measure of the skew elements for which
$(\beta_1,\beta_2)\mapsto\tau_0\delta^{-1}\beta_1+\beta_2$ is a measure-preserving isomorphism of $A_{\rm sk}^2$ onto $A$
(Lemma~\ref{pfs:skewhaar}; $\tau_0\delta^{-1}\in\Gamma(1)$ by (G1)). As $n^{-1}=n(-\omega,\bar\beta)$,
$n^{-1}e_a^*=e_a^*-\omega+\mu e_a$ with $\mu=\bar\beta=-b-\bar\tau_0h(\omega,\omega)$; in the decomposition $\mu=\tau_0(\mu+\bar\mu)
+(\mu-\tau_0(\mu+\bar\mu))$ its skew part is $-b+(\tau_0-\bar\tau_0)h(\omega,\omega)$, that is $-b$ up to a term depending only
on $\omega$. Below, ``the $\beta$-integral'' is the integral in $b$ against $d\beta_{\rm sk}$.

\begin{lemma}[Skew sets]\label{pfr:skew}
If $S$ is a set of skew elements and $u$ a unit with $u^{-1}(\beta-\beta')\in\Gamma(j+1)$ for $\beta,\beta'\in S$, then
$\pbvol_{\rm sk}(S)\le\pfrmod(u)^{1/2}\Lambda^{j+2}$.
\end{lemma}
\begin{proof}
The image of $S\times S$ under $(\beta_1,\beta_2)\mapsto\tau_0\delta^{-1}\beta_1+\beta_2$ has measure $\pbvol_{\rm sk}(S)^2$, and,
after fixing a point $(\beta^0_1,\beta^0_2)$ of $S\times S$, all its points $c$ satisfy
$u^{-1}(c-c^0)\in\Gamma(1+(j+1)+1)$ by (G0). Hence $\pbvol_{\rm sk}(S)^2\le\pfrmod(u)\pbvol\Gamma(j+3)\le\pfrmod(u)\Lambda^{j+4}
\le(\pfrmod(u)^{1/2}\Lambda^{j+2})^2$.
\end{proof}

\begin{lemma}[Integration over $\pfrN_a$]\label{pfr:heis}
Let $u$ be a unit, $\Psi\ge0$ measurable on $V_{a+1}$ with $\int_{\Gamma(j)}\Psi\le I$, $C\ge0$, and $\Phi\ge0$
measurable on $\pfrG$ such that whenever $\Phi(n(\omega,\beta))\ne0$: $u^{-1}(n^{-1}-1)e_a^*$ is congruent modulo
$\langle e_b:b<a\rangle$ to a vector of size $j$, and $\Phi(n)\le C\Psi(u^{-1}\omega)$. Then
$\int\Phi\,d\eta_a\le C\,\mathsf c_a(u)\,\Lambda^{j+2}I$, where $\mathsf c_a(u)=\pfrmod(u)^{d_a+1/2}$.
\end{lemma}
\begin{proof}
Substitute $\omega=uv$ (Jacobian $\pfrmod(u)^{d_a}$). For fixed $v$, on the set $S$ of the skew coordinates $b$
(with $n=n(\omega,b-\tau_0h(\omega,\omega))$) with $\Phi(n)\ne0$, the vector $u^{-1}(n^{-1}-1)e_a^*=u^{-1}(-\omega+\mu e_a)$ has,
modulo the $e_{b'}$ ($b'<a$), size $j$: its $\omega$-part gives $v\in\Gamma(j)$ and its $e_a$-coordinate gives
$u^{-1}\mu=-u^{-1}(b+c_v)\in\Gamma(j)$ for the number $c_v=\bar\tau_0h(\omega,\omega)$ depending only on $v$; so differences
satisfy $u^{-1}(b-b')\in\Gamma(j+1)$, and by Lemma~\ref{pfr:skew} the $b$-integral is at most $\pfrmod(u)^{1/2}\Lambda^{j+2}C\Psi(v)\mathbf 1[v\in\Gamma(j)]$. Integrating over $v$ gives the
claim.
\end{proof}

For a unit $c$ the \emph{section} of the quadric of $f_a$ ($h(y,y)=0$ in $V_a$) at $h(y,e_a)=c$ (the
$e_a^*$-coordinate) is parametrised by
$n\mapsto y_c(n)=c\,n^{-1}e_a^*$, $n\in\pfrN_a$; its \emph{section measure} is $\pfrmod(c)^{d_a}$ times the image of
$\eta_a$. By the definition of $\theta_a=\lambda_a*\eta_a$ ($\lambda_a$ the image of $\pfrmod(c)^{d_a}dc$ under
$c\mapsto w_{a,c}^{-1}$, and $(w_{a,c}^{-1}n)^{-1}e_a=c\,n^{-1}e_a^*$), the image of $\theta_a$ under $\sigma\mapsto\sigma^{-1}e_a$ is
$\int(\text{section measure at }c)\,dc$, which is $c_a\mu_a$ by Lemma~\ref{pfr:floorquad}.

\begin{lemma}[Integration over a section]\label{pfr:sectionint}
If $\int_{\Gamma(j+5)^{X'}}\phi\le B$, then $\int_{\Gamma(j)}\phi(y_{X'})\,d(\text{section at }c)(y)\le
\pfrmod(c^{-1})^{1/2}\Lambda^{j+3}\Lambda^{99n(j+7)}B$.
\end{lemma}
\begin{proof}
$y_c(n)=c(e_a^*-\omega+\mu e_a)$, so $y_c(n)_{X'}=(-c\,\omega)_{X'}$ and $c(n^{-1}-1)e_a^*=y_c(n)-ce_a^*$, which has size
$j+1$ when $y_c(n)$ has size $j$ (then $c=h(e_a,y)$-coordinate has size $j$ and (G0)); hence $(-c)(n^{-1}-1)e_a^*$ has size $j+1$ too ((G0): $\Gamma(i)=-\Gamma(i)$). Apply Lemma~\ref{pfr:heis} with
$u=(-c)^{-1}$, $j+1$ in place of $j$, $C=1$, $\Psi(v)=\phi(v_{X'})$ on $V_{a+1}$, and $I=\Lambda^{99n(j+7)}B$ from the
general form of Lemma~\ref{pfr:quadricX} for Lebesgue measure on $V_{a+1}$ (with $j+1$ in place of $j$). The result is
$\mathsf c_a(c^{-1})\Lambda^{j+3}\Lambda^{99n(j+7)}B$; multiplying by $\pfrmod(c)^{d_a}$ leaves $\pfrmod(c^{-1})^{1/2}$.
\end{proof}

\subsubsection*{Sizes at a head}
\begin{lemma}[Entries at a head]\label{pfr:headentries}
Let $j\le\varrho$, $p$ a head of length $j$, $g\in\pfrG$, $s\in\pfrR^{(j)}$ with $sg\in\pfrK$ and $\eta\in\pfrR^{(j)}$ with
$\eta w_p^{-1}g\in\pfrK$, and $a<j$. Then $p_a$, $p_a^{-1}$ and the entries of the column $se_a^*$ have size $10n$.
\end{lemma}
\begin{proof}
Let $k_1=sg$, $k_2=\eta w_p^{-1}g$, of size $1$ by (G1); then $\gamma:=k_1k_2^{-1}=sw_p\eta^{-1}$ has size $1+5+1=7$ by
(G0). \emph{Products.} For $u\in\pfrR^{(j)}$, $uf_c\in f_c+\langle f_{c'}:c'<c\rangle$ for $c<j$, and the same holds for
$u^{-1}$. Hence, for $l\le j$, the leading $l\times l$ block of the frame matrix of $\gamma$ is, since
$h(sw_p\eta^{-1}f_c,f_{c'})=h(w_p\eta^{-1}f_c,s^{-1}f_{c'})$, a product $\mathsf L\,D_l\,\mathsf R$ with unitriangular
$\mathsf L,\mathsf R$ and $D_l=(h(w_pf_c,f_{c'}))_{c,c'<l}=\operatorname{diag}(p_c)$ up to conjugation of the entries
(as $w_pe_c=p_ce_c^*$); so its determinant is $P_l=p_0\cdots p_{l-1}$ or $\bar P_l$. The entries of $Z(\gamma)$ are entries of $J\gamma$, of size
$9$; so $P_l$ (or $\bar P_l$, the same size) has size $10l\le10\varrho$. As $P_a^{-1},P_{a+1}^{-1}\in\Gamma(1)$ (head),
$p_a=P_a^{-1}P_{a+1}$ and $p_a^{-1}=P_{a+1}^{-1}P_a$ have size $10\varrho+1\le10n$.
\emph{Columns} (Cramer's rule). Fix a basis index $i$. Let $Z'$ be the $(a+1)\times(a+1)$ block of $\gamma$ with rows
$e_0^*,\dots,e_{a-1}^*,i$ and columns $e_0,\dots,e_a$; it has size $7$. Since $\eta^{-1}e_d\in e_d+\langle e_c:c<d\rangle$,
$Z'=\mathsf L'\mathsf C$ with $\mathsf C$ the unitriangular block of $\eta^{-1}$ on $e_0,\dots,e_a$ and $\mathsf L'$ the block
of $sw_p$ with the same rows and columns $e_0,\dots,e_a$; as $w_pe_d=p_de_d^*$, $\mathsf L'_{bd}=s_{R(b),e_d^*}p_d$. For
$b\le d\le a$ with $b<a$, $s_{e_b^*,e_d^*}=h(se_d^*,e_b)=h(e_d^*,s^{-1}e_b)=\delta_{bd}$ (as $s^{-1}e_b\in e_b+\langle e_c:c<b\rangle$); so $\mathsf
L'$ is lower triangular with diagonal $p_0,\dots,p_{a-1},s_{i,e_a^*}p_a$. Hence $\det Z'=P_{a+1}s_{i,e^*_a}$, of size
$(a+1)\cdot8\le8\varrho$, and $s_{i,e_a^*}=P_{a+1}^{-1}\det Z'$ has size $8\varrho+1\le10n$.
\end{proof}

\begin{proof}[Proof of Lemma~\ref{pfr:section}]
Write $w=w_{a+1,p}$ and let $\phi$ be the function of (G5) for the point $y\in\pfrK$, the position $x$, and the
$\varrho$ linear forms $\omega\mapsto(\nu\omega)_{e_b}$ on $X'$ (the $e_b$-coordinates, $b<\varrho$), with common vectors
$\mathcal W$.

\emph{The depth is bounded by $\phi$.} Let $n\in\pfrN_a$ and $\omega\in\mathcal W$. The vector $\omega'=wn\nu\omega$ is stable at
$(wn\nu y,wn\nu yx)$ for $a$: the pair $((wn\nu y)^{-1}\omega',(wn\nu yx)^{-1}\omega')=(y^{-1}\omega,(yx)^{-1}\omega)$ is close,
and for $b<a$ the $e_b^*$-coordinate of $\omega'$ is $p_b(n\nu\omega)_{e_b}=p_b(\nu\omega)_{e_b}=0$ ($n$ fixes the coordinates
of floor $<a$; $\omega$ is killed by the forms), i.e. $\omega'\perp e_b$. Its $e_a^*$-coordinate is $h(\omega',e_a)=
h(\omega,(wn\nu)^{-1}e_a)$, and $(wn\nu)^{-1}e_a=\nu^{-1}n^{-1}\bar p_ae_a^*$ is congruent modulo $\langle e_b:b<a\rangle$ to
the section point $y_{\bar p_a}(n)$ (as $\nu^{-1}\in\pfrU_{(<a)}$ and $y_{\bar p_a}(n)$ has no $e^*_b$-coordinates, $b<a$);
congruent vectors have the same projection on $X'$ (the $e_b$ are orthogonal to $X'$). So, with $u=y_{\bar p_a}(n)_{X'}$,
the $e_a^*$-coordinate of $\omega'$ is $h(\omega,u)$, $\mathfrak d_a\supseteq\{h(\omega,u):\omega\in\mathcal W\}$, and
$\pfrdep_a(wn\nu y,x)\le\phi(u)$ by the first inequality of (G5).

\emph{Size.} If $wn\nu y$ is near at $a+1$, apply Lemma~\ref{pfr:headentries} with $j=a+1$, the head $\bar p$
(for which $w_{\bar p}^{-1}=w$, \S\ref{pfr:ss-gainsproof}), $g=n\nu y$, $s=(n\nu)^{-1}\in\pfrR^{(a+1)}$ ($sg=y\in\pfrK$) and
the $\eta$ given by nearness: $\bar p_a$ and the column $(n\nu)^{-1}e_a^*$ have size $10n$, so
$\bar p_a(n\nu)^{-1}e_a^*$ has size $20n$, and so has $y_{\bar p_a}(n)$, which agrees with it in the coordinates of
floor $\ge a$ and vanishes in the others.

\emph{Integration.} Hence the integrand is at most $\Psi(y_{\bar p_a}(n))$ with $\Psi(v)=\mathbf 1[v\in\Gamma(20n)]\,
\phi(v_{X'})^s$. By the definition of the section measure and Lemma~\ref{pfr:sectionint} (with $j=20n$ and
$B=\Lambda^{99n(20n+6)}$ from (G5)),
\[
\int\Psi(y_{\bar p_a}(n))\,d\eta_a(n)=\pfrmod(p_a^{-1})^{d_a}\int\Psi\,d(\text{section at }\bar p_a)\le\pfrmod(p_a^{-1})^{d_a+1/2}
\Lambda^{20n+3+99n(20n+7)+99n(20n+6)},
\]
and the exponent is at most $10^6n^2$.
\end{proof}

\subsection{Proof of Lemma~\ref{pfr:descent}}\label{pfr:ss-descent}

Write $a=i$, $w=w_{a,p}$, $d=w_{a,q}^{-1}w_{a,p}$, and for $c\in A^\times$ let $w_c=w_{a+1,c}$ restricted to the pair
$(e_a,e_a^*)$: $w_ce_a=ce_a^*$, $w_ce_a^*=\bar c^{-1}e_a$. For points $h,x$ and an element $v$ put
\[
\Pi(x;v;h)=\mathsf k^1_{a+1}(h,hx)\,\overline{\mathsf k^{d}_{a+1}(v^{-1}h,v^{-1}hx)} .
\]
\emph{Facts.} (i) $w_{a+1,p^c}=w\,w_c$, the Weyl ratio of $(p^c,q^c)$ at length $a+1$ is $d$ (the entries $c/c$
are $1$), and $\pfrkap_{a+1}(p^c,q^c)=\pfrkap_a(p,q)\pfrmod(c)^{d_a}$. (ii) By (C1) of \S\ref{pfr:ss-gainsproof},
applied to both factors (the elements $w,w_c$ commute with $\pfrG^{(a+1)}$, which carries $\theta_{a+1}$),
$\Phi_{a+1}(x;1;g')=\Pi(x;ww_c;g')$ for the pair $(p^c,q^c)$, and by definition
$\Phi_a(x;\omega;g)=\iint\omega(\sigma\tau^{-1})\,\Pi(x;\sigma w\tau^{-1};\sigma g)\,d\theta_a(\tau)\,d\theta_a(\sigma)$. (iii) For
$t$ in $\pfrR^{(a+1)}\cap\pfrG^{(a)}$, $tf_a=f_a$ ((N1)), so $Z_{aa}(tyt')=Z_{aa}(y)$; and $\mathsf k^{1}_{a+1}(th,thx)=
\mathsf k^{1}_{a+1}(h,hx)$, $\mathsf k^{d}_{a+1}(th,thx)=\mathsf k^{d}_{a+1}(h,hx)$ by (K2), second claim, of Lemma~\ref{pfs:kernels} (applied with
$l=a$, $\zeta=t$; for the factor coming from $\pfrN_a$ one also sees directly that the character factors
$\psi(M^{-1}(\sigma(dtd^{-1})-\sigma(t)))$ are $1$ since $dtd^{-1}=t$, $d$ being trivial on the floors $\ge a$ where $t$
acts). (iv) $\theta_a$ is the image of $\pfrmod(c)^{d_a}dc\times\eta_a$ under $(c,n)\mapsto w_c^{-1}n$
(the definition of the floor measures in \S\ref{pfr:ss-setting}, at every place), and $Z_{aa}(w_c)=c$.

\begin{proof}[Proof of Lemma~\ref{pfr:descent}]
Let $\omega=\varphi\circ Z_{aa}$ with $\varphi(c)=\psi(\xi c)$, $|\varphi|\le1$. Put
$F(x,g,\sigma,\tau)=\omega(\sigma\tau^{-1})\Pi(x;\sigma w\tau^{-1};\sigma g)$ on $\pfrG^4$ with the measure
$\pfrM\otimes\pfrm_a\otimes\theta_a\otimes\theta_a$, and $E(x,g',c)=\varphi(c)\Pi(x;ww_c;g')$ with the measure
$\pfrM\otimes\pfrm_{a+1}\otimes\pfrmod(c)^{d_a}dc$.

\emph{$|F|$ is integrable.} By Tonelli and $|\omega|\le1$, $\int|F|\le\int\pfrF_a^+\,d\pfrM<\infty$: the integral over
$(\sigma,\tau)$ of $|\Pi|$ is the product of the two pivot norms, as in Lemma~\ref{pfr:L1}.

\emph{$F$ and $E$ have the same law.} Fix $x,g$ and $\sigma\in\pfrG^{(a)}$. The function $\tau\mapsto F(x,g,\sigma,\tau)$
is invariant under $\tau\mapsto t\tau$ for $t\in\pfrR^{(a+1)}\cap\pfrG^{(a)}$, by (iii) (on the first factor
$\omega(\sigma\tau^{-1}t^{-1})$, and on the second because $(\sigma w\tau^{-1}t^{-1})^{-1}\sigma g=t\tau w^{-1}g$). By
Lemma~\ref{pfr:floorright}(b) (with $\xi=\sigma\in\pfrG^{(a)}$) and (iv), for every measurable $f\ge0$,
\[
\int f(F(x,g,\sigma,\tau))\,d\theta_a(\tau)=\int f(F(x,g,\sigma,\tau\sigma))\,d\theta_a(\tau)=\iint
f\bigl(F(x,g,\sigma,w_c^{-1}n\sigma)\bigr)\,d\eta_a(n)\,\pfrmod(c)^{d_a}dc .
\]
Now $\sigma w(w_c^{-1}n\sigma)^{-1}=\sigma w\sigma^{-1}n^{-1}w_c=wn^{-1}w_c$ ($\sigma$ commutes with $w$), so
$\omega(\sigma\tau^{-1})=\varphi(Z_{aa}(n^{-1}w_c))=\varphi(c)$ by (iii), (iv); and
$\Pi(x;wn^{-1}w_c;\sigma g)=\Pi(x;ww_c;n\sigma g)$ by (iii) applied to both correlations ($n$ commutes with $w$).
Hence $F(x,g,\sigma,w_c^{-1}n\sigma)=E(x,n\sigma g,c)$. Integrating over $\sigma\sim\theta_a$, $g\sim\pfrm_a$ and $x\sim\pfrM$,
and using $\pfrm_{a+1}=(\eta_a*\theta_a)*\pfrm_a$, we get $\int f\circ F=\int f\circ E$ for all $f\ge0$.

\emph{Conclusion.} Hence $E$ is integrable, $\int E=\int F$, and by Fubini's theorem (in the order $c$ last)
$c\mapsto\pfrmod(c)^{d_a}\varphi(c)\iint\Pi(x;ww_c;g')\,d\pfrM\,d\pfrm_{a+1}$ is integrable with integral $\int F$. By
(i), (ii), the integrand is $\psi(\xi c)\pfrI_{a+1}(p^c,q^c;1)/\pfrkap_a(p,q)$, and $\pfrkap_a\int F=\pfrI_a(p,q;\omega)$ by
definition.
\end{proof}



\providecommand{\panot}[1]{\par\smallskip{\raggedright\noindent\textbf{[Not proved here:} #1\textbf{]}\par}\smallskip}
\providecommand{\pbvol}{\operatorname{vol}}
\providecommand{\pfrG}{\mathsf G}
\providecommand{\pfrH}{\mathsf H}
\providecommand{\pfrU}{\mathsf U}
\providecommand{\pfrR}{\mathsf R}
\providecommand{\pfrM}{\mathsf M}
\providecommand{\pfrm}{\mathsf m}
\providecommand{\pfrK}{\mathsf K}
\providecommand{\pfrP}{\mathcal P}
\providecommand{\pfrI}{\mathcal I}
\providecommand{\pfrF}{\mathsf F}
\providecommand{\pfrB}{\mathsf B}
\providecommand{\pfrmod}{\operatorname{mod}}
\providecommand{\pfrdep}{\operatorname{dp}}
\providecommand{\pfrwt}{\wp}
\providecommand{\pfrkap}{\kappa}
\providecommand{\pfrN}{\mathsf N}
\providecommand{\pfrip}[2]{\langle#1,#2\rangle}
\providecommand{\pfsA}{\mathsf A}
\providecommand{\pfsBall}{\mathcal B}
\providecommand{\pfscm}{\operatorname{cm}}
\providecommand{\pfscut}{\operatorname{cut}}
\providecommand{\pfsCB}{\mathcal C\!\mathcal B}
\providecommand{\pfsc}{\mathsf c}
\providecommand{\pfsrk}{\operatorname{rk}}
\providecommand{\pfsfc}{\mathsf f}
\providecommand{\pfsO}{\mathcal O}
\providecommand{\pfsst}{\operatorname{St}}

\section{Two frames: the shared estimates}\label{pfs:sec}

This section proves the three statements of \S\ref{pfr:ss-shared} stated in
\S\ref{pfr:sec}: the one-frame bound (Theorem~\ref{pfr:oneframe}), the gain from a frequency
(Theorem~\ref{pfr:gain-char}) and the gain from two heads (Theorem~\ref{pfr:gain-heads}). It also
supplies the three facts that \S\ref{pfr:sec} used as standard (the action on the quadric, measurable
sections, the Haar measure on the skew elements), in \S\ref{pfs:ss-conv}.

\emph{What is proved.} All three theorems are proved in full, from the statements of \S\ref{pfr:sec}
that are proved there and the gauge axioms (G0)--(G5), together with one property of the place that
\S\ref{pfr:sec} names (the constant $P$) but does not state; it is stated here as axiom~(P)
(\S\ref{pfs:ss-planch}). Fact (C4) of \S\ref{pfr:ss-gainsproof} is proved here as (K1)--(K3), and we
recall the convention $\pfrR^{(j)}=\pfrU$ ($j\ge\varrho+2$) in the definition of nearness
(Remark~\ref{pfs:rem-corr}).

Throughout, the assumptions of \S\ref{pfr:ss-shared} are in force: $\pfrF$ is a place datum whose group
measure $\pfrM$ is a Haar measure, $\Gamma$ is a gauge with loss $\Lambda$, $W$ is a cut Bessel function,
$n=\varrho+1$, and $\pfsrk:=\operatorname{rank}X'$.

\subsection{Conventions, kernels, and three standard facts}\label{pfs:ss-conv}

\subsubsection*{Conventions}
\begin{enumerate}
\item[(i)] The hermitian form $h$ is linear in its first and conjugate linear in its second variable
(convention (Cv1) of {\cite[\S2.1]{P51}});
this is the reading under which the maps $v\mapsto v-h(v,\omega)f_l$ of \S\ref{pfr:ss-floors} are
linear. The coordinate of a vector $x$ along $e_a^*$ is $h(x,e_a)$, and along $e_a$ it is
$h(x,e_a^*)$. The scalar product $\pfrip{\cdot}{\cdot}$ on $\C^I$ is conjugate linear in its first
variable. Since $\psi$ is unitary, $\overline{\psi(c)}=\psi(-c)$.
\item[(ii)] For $l\ge0$ put $\pfrU_{(<l)}=\pfrN_0\pfrN_1\cdots\pfrN_{l-1}$ and
$\pfsA_l=\eta_0*\eta_1*\dots*\eta_{l-1}$ (the ascending convolution of \S\ref{pfr:ss-floors};
$\pfsA_0=\delta_1$, and $\pfsA_{l}=\pfsA_{\varrho}$ for $l\ge\varrho$ since $\eta_l=\delta_1$ for $l\ge\varrho$).
By Lemma~\ref{pfr:rearr} and its proof, $\pfsA_l$ is a Haar measure of the nilpotent (hence unimodular)
group $\pfrU_{(<l)}$, so $\delta_\nu*\pfsA_l=\pfsA_l*\delta_\nu=\pfsA_l$ for $\nu\in\pfrU_{(<l)}$; by (N3) of
\S\ref{pfr:ss-floors}, $\delta_\sigma*\pfsA_l=\pfsA_l*\delta_\sigma$ for $\sigma\in\pfrG^{(l)}$; and
$\pfsA_{l+1}=\pfsA_l*\eta_l=\eta_l*\pfsA_l$ (associativity, and (N3) for $\eta_l$, which is carried by
$\pfrN_l\subseteq\pfrG^{(l)}$).
\item[(iii)] The facts on $\pfrU_{(<l)}$ used in \S\ref{pfr:ss-floors}: it fixes the $e_b$, $b<l$; it acts on
$V_l$ by $v\mapsto v+\sum_{b<l}\lambda_b(v)e_b$ (each $\pfrN_b$, $b<l$, acts on $V_{b+1}\supseteq V_l$ by
$v\mapsto v-h(v,\omega)e_b$ and fixes the $e_c$, $c\le b$, and moves the $e_c$, $c>b$, by multiples of $e_b$);
it is normalised by $\pfrG^{(l)}$; and $\pfrR^{(l)}=\pfrG^{(l)}\pfrU_{(<l)}$, by (N1) applied on each floor
$b<l$. In particular, $\pfrU_{(<l)}$ acts trivially on $Y^{(l)}$.
\item[(iv)] \emph{The Heisenberg group of floor $a<\varrho$.} Write $n(\omega,\beta)$ for its elements as in
\S\ref{pfr:ss-floors}, with $\omega\in V_{a+1}$ orthogonal to $f_a,f_a^*$. Directly from the formulas:
$n(\omega,\beta)n(\omega',\beta')=n(\omega+\omega',\beta'')$ for some $\beta''$ (the vector part is additive);
$\gamma n(\omega,\beta)\gamma^{-1}=n(\gamma\omega,\beta)$ for $\gamma\in\pfrG^{(a+1)}$ ((N2)); and the entry of
$n(\omega,\beta)$ just above the diagonal of its frame block, at $(a,a+1)$, is the coefficient of $f_a$ in
$n(\omega,\beta)f_{a+1}=f_{a+1}-h(f_{a+1},\omega)f_a$, so that
\begin{equation}\label{pfs:eq-sigma}
\sigma(n(\omega,\beta))=-h(f_{a+1},\omega),\qquad
\sigma\bigl(\gamma n(\omega,\beta)\gamma^{-1}\bigr)=-h(\gamma^{-1}f_{a+1},\omega)\quad(\gamma\in\pfrG^{(a+1)}).
\end{equation}
\item[(v)] The size rules (G0) are used without comment; ``size $i$'' always refers to $\Gamma$ unless
another gauge is named. A vector of $V_b$ ``of size $i$'' is one whose entries lie in $\Gamma(i)$.
\end{enumerate}

\subsubsection*{Kernels}
A \emph{kernel at floor $j$} is $\mathsf k=\mathsf k^{d'}$, $\mathsf k^{d'}(g,y)=\pfrip{W(g)}{W(d'y)}$, where
$d'$ is diagonal in the basis, of the form $e_a\mapsto\lambda_ae_a$, $e_a^*\mapsto\bar\lambda_a^{-1}e_a^*$ ($a<j$),
and the identity on the basis vectors of floor $\ge j$; such a $d'$ is an isometry, normalises $\pfrU$, and
commutes with $\pfrG^{(j)}$. Examples: $d'=1$, and the Weyl ratio $w_{j,q}^{-1}w_{j,p}$ of two sequences of units.
A kernel at floor $j$ is a kernel at every floor $j'\ge j$. The correlations $\mathsf k_l$ are those of \S\ref{pfr:ss-setting};
facts (C1) and (C2) of \S\ref{pfr:ss-gainsproof} hold for them. We use three further facts, in the form in
which the argument needs them.

\begin{lemma}\label{pfs:kernels}
Let $\mathsf k$ be a kernel at floor $j$.
\begin{enumerate}
\item[(K1)] \textup{(Equivariance.)} For $l\le n$, $\nu,\nu'\in\pfrU_{(<l)}$ and all $g,y$,
\[
\mathsf k_{l+1}(\nu g,\nu'y)=\psi\bigl(M^{-1}(\sigma(d'\nu'd'^{-1})-\sigma(\nu))\bigr)\,\mathsf k_{l+1}(g,y).
\]
\item[(K2)] \textup{(Invariance.)} For $j\le l$, $\gamma\in\pfrG^{(l)}$ and all $h,h'$:
$\mathsf k_l(\gamma h,\gamma h')=\mathsf k_l(h,h')$. For $j\le l\le\varrho$ and $\zeta\in\pfrR^{(l+1)}\cap\pfrG^{(l)}$:
$\mathsf k_{l+1}(\zeta h,\zeta h')=\mathsf k_{l+1}(h,h')$.
\item[(K3)] For $l\le\varrho$ with $j\le l$, and every position $x$, the function
$\mathsf a_l(g)=\pfrwt_l(g,x)|\mathsf k_{l+1}(g,gx)|^2$ of \S\ref{pfr:ss-gauge} satisfies $\mathsf a_l(\xi g)=\mathsf
a_l(g)$ for $\xi\in\pfrR^{(l+1)}$; for every $l$, $|\mathsf k_{l+1}(g,\nu y)|=|\mathsf k_{l+1}(g,y)|$ for
$\nu\in\pfrU_{(<l)}$.
\end{enumerate}
\end{lemma}
\begin{proof}
(K1) The measure $\theta_{\ge l+1}$ is carried by $\pfrG^{(l+1)}$. For $x\in\pfrG^{(l+1)}$, $x\nu x^{-1}$ lies in
$\pfrU_{(<l)}\subseteq\pfrU$ and has the same entries as $\nu$ just above the diagonal of the frame block: these are
the entries $(a,a+1)$ with $a<l$, so $a+1\le l$, that is $h(\nu f_{a+1},f^\vee_a)$ for the dual
vector $f_a^\vee$ of $f_a$, and $x$ fixes $f_{a+1}$ and $f_a^\vee$ (both of floor $\le l<l+1$). So $\sigma(x\nu x^{-1})=\sigma(\nu)$, and
$W(x\nu g)=W((x\nu x^{-1})xg)=\psi(M^{-1}\sigma(\nu))W(xg)$. Likewise $d'x\nu'x^{-1}d'^{-1}\in\pfrU$ ($d'$ is diagonal
on the frame, so it normalises $\pfrU$ and multiplies each entry above the diagonal by a scalar that does not
depend on $x$), $\sigma(d'x\nu'x^{-1}d'^{-1})=\sigma(d'\nu'd'^{-1})$, and $W(d'x\nu'y)=\psi(M^{-1}\sigma(d'\nu'd'^{-1}))
W(d'xy)$. Insert both in $\mathsf k(x\nu g,x\nu'y)$, use convention (i) and integrate over $x$.

(K2) Descending induction on $l$. For $l>n$, $\pfrG^{(l)}=1$. For $l=n$, $\gamma\in\pfrH$ and
$\theta_{\ge n}=\theta_{\ge n+1}*\theta_n$ with $\theta_{\ge n+1}=\delta_1$ and $\theta_n=\mu_\pfrH$; so
$\mathsf k_n(\gamma h,\gamma h')=\int\mathsf k(x\gamma h,x\gamma h')\,d\mu_\pfrH(x)=\mathsf k_n(h,h')$ by the right
invariance of $\mu_\pfrH$. Let $j\le l\le\varrho$ and assume the first claim at $l+1$. \emph{Second claim at
$l$.} By (N1), $\zeta=\gamma'n_0$ with $\gamma'\in\pfrG^{(l+1)}$ and $n_0\in\pfrN_l$. By the first claim at $l+1$,
$\gamma'$ may be dropped. For $n_0$: by (C1), $\mathsf k_{l+1}(n_0h,n_0h')=\int\mathsf k_{l+2}(\sigma n_0h,\sigma
n_0h')\,d\theta_{l+1}(\sigma)$, and for $\sigma\in\pfrG^{(l+1)}$ the element $s=\sigma n_0\sigma^{-1}$ lies in
$\pfrN_l\subseteq\pfrU_{(<l+1)}$; by (K1) at level $l+2$ (note $l+1\le n$), $\mathsf k_{l+2}(s\sigma h,s\sigma h')=
\psi(M^{-1}(\sigma(d'sd'^{-1})-\sigma(s)))\mathsf k_{l+2}(\sigma h,\sigma h')$, and $d'sd'^{-1}=s$ as $s\in\pfrG^{(l)}
\subseteq\pfrG^{(j)}$. So the factor is $1$. \emph{First claim at $l$.} By (C1), $\mathsf k_l(\gamma h,\gamma h')=
\int\Psi(\sigma\gamma)\,d\theta_l(\sigma)$ with $\Psi(\sigma)=\mathsf k_{l+1}(\sigma h,\sigma h')$, which is integrable
for $\theta_l$ and, by the second claim, satisfies $\Psi(\zeta\sigma)=\Psi(\sigma)$ for $\zeta\in\pfrR^{(l+1)}\cap
\pfrG^{(l)}$. Lemma~\ref{pfr:floorright}(b) with $\xi=\gamma\in\pfrG^{(l)}$, applied to $\varphi\circ\Psi$ for every
Borel $\varphi\ge0$ on $\C$, shows that $\sigma\mapsto\Psi(\sigma\gamma)$ and $\Psi$ have the same law under
$\theta_l$; hence the same integral.

(K3) The second claim is (K1) with $\nu=1$ and $|\psi|=1$. For the first, write $\xi=\gamma\nu$ with
$\gamma\in\pfrG^{(l)}$, $\nu\in\pfrU_{(<l)}$ (convention (iii)); as $\nu\in\pfrR^{(l+1)}$, also $\gamma\in\pfrR^{(l+1)}
\cap\pfrG^{(l)}$. By (K2), $\mathsf k_{l+1}(\gamma\nu g,\gamma\nu gx)=\mathsf k_{l+1}(\nu g,\nu gx)$, and by (K1) its
modulus is $|\mathsf k_{l+1}(g,gx)|$. The weight $\pfrwt_l$ is invariant under $\pfrR^{(l)}\supseteq\pfrR^{(l+1)}$ by
(C3).
\end{proof}

\begin{remark}[Two points of \S\ref{pfr:sec} recalled]\label{pfs:rem-corr}
(a) Fact (C4) of \S\ref{pfr:ss-gainsproof} is stated in the form that holds and that the proofs in
\S\ref{pfr:ss-gainsproof} use, namely (K1)--(K3): invariance of $|\mathsf k_{l+1}|$ under $\pfrU_{(<l)}$, and of
$\mathsf a_l$ under $\pfrR^{(l+1)}$ (the functions $\mathsf b_l$ are then $\pfrR^{(l)}$-invariant by
Lemma~\ref{pfr:floorright}, by descending induction). (b) A kernel
$\mathsf k_{\varrho+2}=\mathsf k$ is non-zero only at points $g$ with $\pfrU g\cap\pfrK\ne\emptyset$; so in (C2) for
$l=\varrho+2$, and wherever \S\ref{pfr:sec} speaks of points ``near at $\varrho+2$'', the group $\pfrR^{(\varrho+2)}$ is
to be read as $\pfrU$ (this convention is part of the definition of nearness in \S\ref{pfr:ss-setting}). We use this only in Lemma~\ref{pfs:climb} below, where
we argue directly.
\end{remark}

\subsubsection*{Three standard facts}
\begin{lemma}[The Haar measure on the skew elements]\label{pfs:skewhaar}
Let $\tau_0\in A$ with $\tau_0+\bar\tau_0=1$ and let $\delta$ be a skew unit. The map
$\Sigma(\beta_1,\beta_2)=\tau_0\delta^{-1}\beta_1+\beta_2$ is an isomorphism of topological groups from
$A_{\rm sk}^2$ onto $A$ ($A_{\rm sk}=\{\beta:\beta+\bar\beta=0\}$), and there is exactly one Haar measure
$d\beta_{\rm sk}$ on $A_{\rm sk}$ such that $\Sigma$ carries $d\beta_{\rm sk}\otimes d\beta_{\rm sk}$ to $dc$.
\end{lemma}
\begin{proof}
For skew $\beta_1$, $r=\delta^{-1}\beta_1$ is hermitian ($\bar r=(-\delta^{-1})(-\beta_1)$). If
$c=\tau_0r+\beta_2$, then $c+\bar c=(\tau_0+\bar\tau_0)r=r$. Hence $\Sigma$ has the continuous inverse
$c\mapsto(\delta(c+\bar c),\,c-\tau_0(c+\bar c))$ (the second entry is skew: its sum with its conjugate is
$(c+\bar c)(1-\tau_0-\bar\tau_0)=0$). Let $\nu_0$ be a Haar measure of $A_{\rm sk}$. The image of $\nu_0\otimes\nu_0$
under the isomorphism $\Sigma$ is a Haar measure of $A$, so it is $\lambda\,dc$ with $\lambda>0$; then
$\nu=\lambda^{-1/2}\nu_0$ is the required measure, and it is unique since any other is $c\nu$ with $c^2=1$,
$c>0$.
\end{proof}
At the real place ($\tau_0=\frac12$, $\delta=\frac i2$, as fixed in \S\ref{pfr:ss-setting}), $\Sigma(it_1,it_2)=t_1+it_2$, so
$d\beta_{\rm sk}$ is the Lebesgue measure $dt$ of $\beta=it$. At the finite place $d\beta_{\rm sk}$ gives the integral skew elements
volume $1$ (Lemma~\ref{pfg:duality}(d)).

\begin{lemma}[Orbits and measurable sections]\label{pfs:section}
Let $b\le\varrho$, let $\pfsst_b=\{\sigma\in\pfrG^{(b)}:\sigma f_b=f_b\}$ and $\pfsO_b=\{\sigma^{-1}f_b:\sigma\in\pfrG^{(b)}\}
\subseteq V_b$. Then $\pfsst_b=\pfrR^{(b+1)}\cap\pfrG^{(b)}$; $\pfsO_b$ is a Borel set; and there is a Borel map
$s\colon\pfsO_b\to\pfrG^{(b)}$ with $s(v)^{-1}f_b=v$. Consequently, if $\Psi$ is a Borel function on $\pfrG^{(b)}$
(with values in $\C$ or $[0,\infty]$) with $\Psi(\zeta\sigma)=\Psi(\sigma)$ for $\zeta\in\pfsst_b$, then
$\Psi(\sigma)=\Psi_0(\sigma^{-1}f_b)$ for the Borel function $\Psi_0=\Psi\circ s$ on $\pfsO_b$, extended by $0$
to $V_b$. Moreover the quadric measure $\mu_b$ is carried by $\pfsO_b$.
\end{lemma}
\begin{proof}
The equality $\pfsst_b=\pfrR^{(b+1)}\cap\pfrG^{(b)}$ is shown in (N1) of \S\ref{pfr:ss-floors} (one inclusion is
proved there; conversely an element of $\pfrG^{(b)}$ fixing $f_b$ fixes all $f_c$, $c\le b$, hence lies in
$\pfrR^{(b+1)}$). The group $\pfrG^{(b)}$ is a closed subgroup of $\mathrm{GL}_r(A)$, a Polish group ($A$ is $\C$ or
a finite product of local fields), and $\pfsst_b$ is a closed subgroup. By the theorem on Borel transversals of
cosets of a closed subgroup of a Polish group (Kechris, \emph{Classical descriptive set theory}, Theorem~12.17,
applied to the left cosets of $\pfsst_b$ after the inversion $\sigma\mapsto\sigma^{-1}$) there is a Borel set $T\subseteq\pfrG^{(b)}$ meeting each coset $\pfsst_b\sigma$ in exactly one
point. The continuous map $o(\sigma)=\sigma^{-1}f_b$ is constant on these cosets ($o(\zeta\sigma)=\sigma^{-1}\zeta^{-1}
f_b=o(\sigma)$) and separates them ($o(\sigma)=o(\sigma')$ gives $\sigma'\sigma^{-1}\in\pfsst_b$); so $o|_T$ is a
continuous injection with image $\pfsO_b$. By the Lusin--Souslin theorem (Kechris, op.\ cit., Theorem~15.1) the
image $\pfsO_b$ is Borel and $s=(o|_T)^{-1}$ is Borel. For $\Psi$ as stated, $s(\sigma^{-1}f_b)\in\pfsst_b\sigma$,
so $\Psi(s(\sigma^{-1}f_b))=\Psi(\sigma)$. Finally, the image of $\theta_b$ under $o$ is $c_b\mu_b$ with $c_b>0$
(definition of a place datum, \S\ref{pfr:ss-setting}), and $\theta_b$ is carried by $\pfrG^{(b)}$.
\end{proof}

\begin{remark}[Witt's theorem]\label{pfs:witt}
The proof of Lemma~\ref{pfr:floorright} (and of Lemma~\ref{pfr:upperinv}, which uses it) invokes the transitivity of
$\pfrG^{(b)}$ on the quadric of $f_b$ in order to write a $\pfsst_b$-invariant function as a function of $\sigma^{-1}f_b$.
Transitivity is not needed: Lemma~\ref{pfs:section} gives $\Psi=\Psi_0(\sigma^{-1}f_b)$ on $\pfrG^{(b)}$ with
$\Psi_0$ Borel, and the integrals in that proof are integrals against $\mu_b$, which is carried by the orbit
$\pfsO_b$; the remaining step, the $\pfrG^{(b)}$-invariance of $\mu_b$, is used there as stated. For the record:
at the real place $\pfsO_b$ is the whole quadric minus $\{0\}$, by Witt's extension theorem for hermitian forms
over $\C$ (e.g.\ Scharlau, \emph{Quadratic and Hermitian Forms}), applied to
the map $f_b\mapsto v$ between the lines $Af_b$ and $Av$ for $v\ne0$ with $h(v,v)=h(f_b,f_b)$.
\end{remark}

\subsection{Plancherel's inequality and the quadrics}\label{pfs:ss-planch}

\subsubsection*{The constant of a place}
The constant $P$ of a place (\S\ref{pfr:ss-setting}) is the constant of the following axiom, which
we add to the definition of a place, together with two properties of the modulus that every place of
\S\ref{pfr:ss-setting} has.
\begin{enumerate}
\item[(P)] There is a continuous integrable function $G\colon A\to[0,\infty)$ with $G(0)=1$ whose
Fourier transform $\widehat G(z)=\int_AG(v)\psi(\bar vz)\,dv$ satisfies $\int_A|\widehat G(z)|\,dz\le P$.
Moreover $\pfrmod(\bar c)=\pfrmod(c)$, $\pfrmod$ is multiplicative, $d(cv)=\pfrmod(c)\,dv$ for units $c$, and
the non-units of $A$ form a null set.
\end{enumerate}
For $d\ge1$ and $v,w\in A^d$ write $\pfrip vw=\sum_i\bar v_iw_i$.

\begin{proposition}[Plancherel's inequality]\label{pfs:planch}
Let $\eta$ be a unit and $f\in L^1(A^d)$. Then
\[
\int_{A^d}\Bigl|\int_{A^d}\psi(\eta\pfrip vw)f(w)\,dw\Bigr|^2dv\le\Bigl(\frac P{\pfrmod(\eta)}\Bigr)^d\int_{A^d}|f|^2 .
\]
\end{proposition}
\begin{proof}
We may assume $f\in L^2$. Write $\widehat f(v)=\int\psi(\eta\pfrip vw)f(w)\,dw$ and, for $x\in A$,
$I(x)=\int|\widehat f(v)|^2\prod_iG(\bar xv_i)\,dv$. \emph{Expansion.} The function $(v,w,w')\mapsto\prod_iG(\bar x
v_i)|f(w)||f(w')|$ is integrable, so by Fubini's theorem, with $z=w'-w$,
\[
I(x)=\iint\overline{f(w)}f(w')\prod_i\widehat G_x(z_i)\,dw\,dw',\qquad\widehat G_x(t)=\int_AG(\bar xv)\psi(\eta\bar vt)\,dv .
\]
For a unit $x$ substitute $v=\bar x^{-1}v'$: $dv=\pfrmod(x)^{-1}dv'$ and $\eta\bar vt=\eta x^{-1}\bar v't$, so
$\widehat G_x(t)=\pfrmod(x)^{-1}\widehat G(\eta x^{-1}t)$ and $\int|\widehat G_x|\le\pfrmod(x)^{-1}\pfrmod(\eta
x^{-1})^{-1}P=P/\pfrmod(\eta)$. By Cauchy--Schwarz, $\int|f(w)||f(w+z)|\,dw\le\int|f|^2$ for every $z$; hence,
substituting $w'=w+z$,
\[
I(x)\le\int\prod_i|\widehat G_x(z_i)|\int|f(w)||f(w+z)|\,dw\,dz\le\Bigl(\frac P{\pfrmod(\eta)}\Bigr)^d\int|f|^2 .
\]
\emph{Limit.} The units are conull, hence dense, so there are units $x_m\to0$. As $G$ is continuous with
$G(0)=1$, $\prod_iG(\bar x_mv_i)\to1$ for every $v$, and Fatou's lemma gives $\int|\widehat f|^2\le\liminf
I(x_m)$.
\end{proof}

\subsubsection*{The Fourier kernel between two quadrics}
Fix $b\le\varrho$ and use the floor coordinates of $V_b$ (\S\ref{pfr:ss-section}): in them $h$ is the standard
form, a vector of size $i$ has floor coordinates in $\Gamma(i+4)$, and for $y,v\in V_b$ the product $h(y,v)$ is
$\sum_{c}(y_{c}\bar v_{c^*}+y_{c^*}\bar v_c)+\sum_{t}y_t\bar v_t$, the first sum over the hyperbolic pairs of $V_b$,
the second over the orthonormal coordinates. Fix two disjoint blocks $\beta,\beta'$ of $2m+2$ orthonormal
coordinates each, disjoint from those of $X'$ (by the definition of $X'$ in \S\ref{pfr:ss-gauge}). Recall that
$V_b$ has at most $99n$ floor coordinates off $X'$ (proof of Lemma~\ref{pfr:quadricX}).

\begin{proposition}[Plancherel on a quadric]\label{pfs:quadplanch}
Let $\rho$ be Lebesgue measure on $V_b$ or the quadric measure $\mu_b$, let $i,j\ge0$ and let $\xi$ be a unit.
For every Borel $\varphi\in L^1(\mu_b)$ vanishing off the vectors of size $j$,
\[
\int_{\{v\text{ of size }i\}}\Bigl|\int\psi(\xi h(y,v))\varphi(y)\,d\mu_b(y)\Bigr|^2d\rho(v)\le
\Lambda^{99n(i+6)}\Lambda^{99n(j+6)}\Bigl(\frac P{\pfrmod(\xi)}\Bigr)^{\pfsrk}\int|\varphi|^2\,d\mu_b .
\]
\end{proposition}
\begin{proof}
By Lemma~\ref{pfr:block}, applied to $\mu_b$ with the block $\beta$ and to $\rho$ with the block $\beta'$, write
$y=(y_o,y_X,y_\beta)$ and $v=(v_o,v_X,v_{\beta'})$, where the subscript $X$ denotes the coordinates of $X'$ and
$o$ the remaining (outer) ones; then $d\mu_b(y)=dy_o\,dy_X\,\kappa_{(y_o,y_X)}(dy_\beta)$ and
$d\rho(v)=dv_o\,dv_X\,\kappa'_{(v_o,v_X)}(dv_{\beta'})$ with Lebesgue measures in the outer and $X'$ coordinates,
and fibre measures giving mass at most $\Lambda^{(2m+2)(l+2)}$ to the vectors with entries in $\Gamma(l)$. Let
$R=\Gamma(j+4)$, $R'=\Gamma(i+4)$, and let $N_o$, $N_o'$ be the numbers of outer coordinates of $y$ and of $v$.

\emph{The phase.} The coordinates of $\beta'$ are outer coordinates of $y$ (as $\beta'$ is disjoint from $\beta$ and
from $X'$). Hence $h(y,v)=\pfrip{v_X}{y_X}+B_v(y_o)+a_{v_o}(y)$, where $\pfrip{v_X}{y_X}=\sum_{t\in X'}\bar v_ty_t$ is
the pairing of Proposition~\ref{pfs:planch}, $B_v(y_o)=\sum_{t\in\beta'}y_t\bar v_t$
depends on $y$ only through $y_o$, and $a_{v_o}(y)$ collects the terms containing a coordinate of $v_o$. So
$\psi(\xi h(y,v))=\psi(\xi B_v(y_o))\,\psi(\xi\pfrip{v_X}{y_X})\,\psi(\xi a_{v_o}(y))$.

\emph{Cauchy--Schwarz in $y_o$.} Put $F_{v_o}(y_o,y_X)=\int\psi(\xi a_{v_o}(y))\varphi(y)\,\kappa_{(y_o,y_X)}(dy_\beta)$.
As $\varphi$ vanishes unless $y_o\in R^{N_o}$, and $|\psi(\xi B_v)|=1$,
\[
T(v):=\Bigl|\int\psi(\xi h(y,v))\varphi\,d\mu_b\Bigr|^2\le\pbvol(R)^{N_o}\,H(v_o,v_X),\quad
H(v_o,v_X)=\int\Bigl|\int\psi(\xi\pfrip{v_X}{y_X})F_{v_o}(y_o,y_X)\,dy_X\Bigr|^2dy_o .
\]
\emph{Plancherel in $y_X$, Cauchy--Schwarz in the block.} By Proposition~\ref{pfs:planch} for each $y_o$ (the
function $y_X\mapsto F_{v_o}(y_o,y_X)$ is integrable for almost every $y_o$), and then
$|F_{v_o}|^2\le\kappa_{(y_o,y_X)}(R^{2m+2})\int|\varphi|^2d\kappa\le\Lambda^{(2m+2)(j+6)}\int|\varphi|^2d\kappa$ (as
$\varphi$ vanishes unless $y_\beta\in R^{2m+2}$),
\[
\int_{A^{X'}}H(v_o,v_X)\,dv_X\le\Bigl(\frac P{\pfrmod(\xi)}\Bigr)^{\pfsrk}\Lambda^{(2m+2)(j+6)}\int|\varphi|^2\,d\mu_b
\quad\text{for every }v_o .
\]
\emph{Integration in $v$.} A vector $v$ of size $i$ has floor coordinates in $R'$. As $T(v)\le\pbvol(R)^{N_o}H(v_o,v_X)$
does not depend on $v_{\beta'}$, integrating first over $v_{\beta'}$ (mass $\le\Lambda^{(2m+2)(i+6)}$), then over
$v_X\in A^{X'}$, then over $v_o\in R'^{N'_o}$ gives the bound
\[
\pbvol(R')^{N'_o}\Lambda^{(2m+2)(i+6)}\cdot\pbvol(R)^{N_o}\Lambda^{(2m+2)(j+6)}\Bigl(\frac P{\pfrmod(\xi)}\Bigr)^{\pfsrk}\int|\varphi|^2d\mu_b .
\]
By (G3), $\pbvol(R)\le\Lambda^{j+5}$, $\pbvol(R')\le\Lambda^{i+5}$, and $N_o+2m+2$, $N_o'+2m+2$ are at most the number
$99n$ of coordinates off $X'$; as $\Lambda\ge1$, $(i+5)N_o'+(i+6)(2m+2)\le99n(i+6)$ and likewise for $j$, so the
constant is at most $\Lambda^{99n(i+6)}\Lambda^{99n(j+6)}$.
\end{proof}

\subsubsection*{The Fourier transform of a correlation along a quadric}
\begin{lemma}[Near points, general size]\label{pfs:nearsize}
Let $b\le\varrho$, let $h'\in\pfrG$ have size $i$, and let $\rho\in\pfrG^{(b)}$ with $\rho h'$ near at $b+1$. Then
$\rho^{-1}f_b$ has size $i+6$.
\end{lemma}
\begin{proof}
As in the proof of Lemma~\ref{pfr:nearsize}: let $\xi\in\pfrR^{(b+1)}$ with $k=\xi\rho h'\in\pfrK$; then
$\rho^{-1}\xi^{-1}=h'k^{-1}$ has size $i+5+1$ ($k^{-1}$ has size $5$ by (G0), (G1)); $\xi^{-1}f_b=f_b+\sum_{c<b}\lambda_cf_c$
and $\rho^{-1}$ fixes the $f_c$, $c<b$, and preserves $V_b$; so the coordinates of $\rho^{-1}f_b$ are those of
floor $\ge b$ of the column $h'k^{-1}f_b$, and vanish on the floors $<b$.
\end{proof}

\begin{proposition}[Plancherel for a correlation]\label{pfs:corrplanch}
Let $b\le\varrho$, let $\mathsf k$ be a kernel at floor $j_0\le b$, let $t\in\pfrG^{(b)}$ and $q'\in\pfrG$ with
$t^{-1}q'\in\pfrU_{(<b)}$, let $g,y\in\pfrG$ with $q'g$ of size $i$, let $\xi$ be a unit, $v_0\in V_b$, $i'\ge0$, and
let $\rho$ be Lebesgue measure on $V_b$ or $\mu_b$. Put
$\mathcal F(v)=\int\psi(\xi h(t\sigma^{-1}f_b,v+v_0))\,\mathsf k_{b+1}(\sigma g,\sigma y)\,d\theta_b(\sigma)$. Then
\[
\int_{\{v\text{ of size }i'\}}|\mathcal F(v)|^2\,d\rho(v)\le2\Lambda^{99n(i'+6)}\Lambda^{99n(i+12)}\Bigl(\frac
P{\pfrmod(\xi)}\Bigr)^{\pfsrk}\int|\mathsf k_{b+1}(\sigma g,\sigma y)|^2\,d\theta_b(\sigma).
\]
\end{proposition}
\begin{proof}
Let $\kappa(\sigma)=\mathsf k_{b+1}(\sigma g,\sigma y)$; it is bounded, Borel, and integrable for $\theta_b$ by
(C1). \emph{Support.} Let $\sigma\in\pfrG^{(b)}$ with $\kappa(\sigma)\ne0$. By (C2), $\sigma g$ is near at $b+1$. The
element $\sigma(t^{-1}q')\sigma^{-1}$ lies in $\pfrU_{(<b)}\subseteq\pfrR^{(b+1)}$, so
$(\sigma t^{-1})(q'g)=\sigma(t^{-1}q')\sigma^{-1}\cdot\sigma g$ is near at $b+1$, and by Lemma~\ref{pfs:nearsize} (with
$\rho=\sigma t^{-1}$, $h'=q'g$), $t\sigma^{-1}f_b$ has size $i+6$. \emph{Factorisation.} By (K2), $\kappa(\zeta\sigma)=
\kappa(\sigma)$ for $\zeta\in\pfrR^{(b+1)}\cap\pfrG^{(b)}=\pfsst_b$, so by Lemma~\ref{pfs:section}, $\kappa(\sigma)=\kappa_0(
\sigma^{-1}f_b)$ on $\pfrG^{(b)}$ with $\kappa_0$ Borel. Put $\varphi(w)=\mathbf 1[w\text{ of size }i+6]\,\kappa_0(t^{-1}w)\,
\psi(\xi h(w,v_0))$. The map $\sigma\mapsto t\sigma^{-1}f_b=(\sigma t^{-1})^{-1}f_b$ carries $\theta_b$ to $c_b\mu_b$ (the
image of $\theta_b$ under $\sigma\mapsto\sigma^{-1}f_b$ is $c_b\mu_b$ with $0<c_b\le2$, \S\ref{pfr:ss-setting} and
Lemma~\ref{pfr:floorquad}, and $\mu_b$ is $\pfrG^{(b)}$-invariant). Hence, by the
support property, for $\theta_b$-almost every $\sigma$, $\kappa(\sigma)\psi(\xi h(t\sigma^{-1}f_b,v_0))=\varphi(t\sigma^{-1}
f_b)$, so $\mathcal F(v)=c_b\int\psi(\xi h(w,v))\varphi(w)\,d\mu_b(w)$ (additivity of $h$ in the second variable), and
$c_b\int|\varphi|^2d\mu_b=\int|\kappa|^2d\theta_b$, $\varphi\in L^1(\mu_b)$. Proposition~\ref{pfs:quadplanch} with $j=i+6$
bounds the left side by $c_b^2\Lambda^{99n(i'+6)}\Lambda^{99n(i+12)}(P/\pfrmod(\xi))^{\pfsrk}\int|\varphi|^2d\mu_b$, and
$c_b^2\int|\varphi|^2d\mu_b=c_b\int|\kappa|^2d\theta_b\le2\int|\kappa|^2d\theta_b$.
\end{proof}

\subsubsection*{Correlations along a Heisenberg group}
\begin{lemma}\label{pfs:heiscorr}
Let $a<\varrho$, let $\mathsf k$ be a kernel at floor $a+1$, $\varepsilon\in\{0,1\}$ and $\xi\in A$ with
\begin{equation}\label{pfs:eq-twist}
M^{-1}\bigl(\sigma(d'sd'^{-1})-\varepsilon\,\sigma(s)\bigr)=\xi\,\sigma(s)\qquad(s\in\pfrN_a).
\end{equation}
Then for $n=n(\omega,\beta)\in\pfrN_a$ and all $g,y$,
\[
\mathsf k_{a+1}(n^\varepsilon g,ny)=\int\psi\bigl(-\xi\,h(\sigma^{-1}f_{a+1},\omega)\bigr)\,\mathsf k_{a+2}(\sigma g,\sigma y)\,d\theta_{a+1}(\sigma).
\]
\end{lemma}
\begin{proof}
By (C1), $\mathsf k_{a+1}(n^\varepsilon g,ny)=\int\mathsf k_{a+2}(\sigma n^\varepsilon g,\sigma ny)\,d\theta_{a+1}(\sigma)$. For
$\sigma\in\pfrG^{(a+1)}$ put $s=\sigma n\sigma^{-1}\in\pfrN_a\subseteq\pfrU_{(<a+1)}$; then $\sigma n^\varepsilon g=s^\varepsilon\sigma g$
and $\sigma ny=s\sigma y$. By (K1) at level $a+2$ (as $a+1\le n$),
$\mathsf k_{a+2}(s^\varepsilon\sigma g,s\sigma y)=\psi(M^{-1}(\sigma(d'sd'^{-1})-\sigma(s^\varepsilon)))\mathsf k_{a+2}(\sigma g,\sigma y)$, and
$\sigma(s^\varepsilon)=\varepsilon\sigma(s)$. By \eqref{pfs:eq-twist} and \eqref{pfs:eq-sigma} the phase is
$\psi(\xi\sigma(s))=\psi(-\xi h(\sigma^{-1}f_{a+1},\omega))$.
\end{proof}

\begin{proposition}[Parseval along a Heisenberg group]\label{pfs:parseval}
Let $a<\varrho$, $\mathsf k$ a kernel at floor $a+1$, $\varepsilon\in\{0,1\}$, $\xi$ and $u$ units satisfying
\eqref{pfs:eq-twist}, $C\ge0$, $i,j\ge0$, $g,y\in\pfrG$, $S\subseteq\pfrN_a$, and $\Phi\ge0$ Borel on $\pfrG$ such
that:
\begin{enumerate}
\item[(a)] if $n\in\pfrN_a$ and $\Phi(n)\ne0$, then $n\in S$ and $\Phi(n)\le C|\mathsf k_{a+1}(n^\varepsilon g,ny)|^2$;
\item[(b)] for every $n_0\in S$ there are $t_0\in\pfrR^{(a+1)}\cap\pfrG^{(a)}$ and $q\in\pfrU_{(<a+1)}$ such that
$t_0qg$ has size $i$ and, for every $n\in S$, some vector $v$ congruent to $(t_0n_0(t_0n)^{-1}-1)e_a^*$ modulo
$\langle e_c:c<a\rangle$ has $u^{-1}v$ of size $j$.
\end{enumerate}
Then
\[
\int\Phi\,d\eta_a\le C\Bigl(\frac P{\pfrmod(\xi u)}\Bigr)^{\pfsrk}\pfsc_a(u)\,\Lambda^{100n(i+j+31)}
\int|\mathsf k_{a+2}(\sigma g,\sigma y)|^2\,d\theta_{a+1}(\sigma),
\]
with $\pfsc_a(u)=\pfrmod(u)^{d_a+1/2}$ as in \S\ref{pfr:ss-gauge}.
\end{proposition}
\begin{proof}
If $\int\Phi\,d\eta_a=0$ there is nothing to prove; otherwise, as $\eta_a$ is carried by $\pfrN_a$, there is
$n_0\in\pfrN_a$ with $\Phi(n_0)\ne0$, so $n_0\in S$. Take $t_0,q$ as in (b) and write $t_0=tr$ with
$t\in\pfrG^{(a+1)}$, $r\in\pfrN_a$ ((N1)). Let $z$ be the vector part of $n_0$, and for $n'=n(\omega',\beta')\in\pfrN_a$ put
\[
N(n')=r^{-1}\,(t^{-1}n't)\,r\,n_0\in\pfrN_a ,
\]
whose vector part is $t^{-1}\omega'+z$ (convention (iv)). \emph{Substitution.} The measure $\eta_a$ is invariant
under left and right translation by elements of $\pfrN_a$ and under conjugation by $t\in\pfrG^{(a+1)}$ ((N2)); so
$\int\Phi\,d\eta_a=\int\Phi(N(n'))\,d\eta_a(n')$.

\emph{The correlation at $N(n')$.} By Lemma~\ref{pfs:heiscorr} and $h(\sigma^{-1}f_{a+1},t^{-1}\omega'+z)=h(t\sigma^{-1}f_{a+1},
\omega'+tz)$,
\[
\mathsf k_{a+1}(N(n')^\varepsilon g,N(n')y)=\int\psi\bigl(\xi'\,h(t\sigma^{-1}f_{a+1},u^{-1}\omega'+v_0)\bigr)\,
\mathsf k_{a+2}(\sigma g,\sigma y)\,d\theta_{a+1}(\sigma)=:\mathcal F(u^{-1}\omega'),
\]
with $\xi'=-\xi\bar u$ and $v_0=u^{-1}tz$ (as $h(w,\omega')=\bar u\,h(w,u^{-1}\omega')$ by convention (i)).

\emph{The lift.} If $\Phi(N(n'))\ne0$, then $N(n')\in S$, and (b) with $n=N(n')$ gives a vector $v\equiv(t_0n_0(t_0
N(n'))^{-1}-1)e_a^*$ with $u^{-1}v$ of size $j$; and $t_0n_0(t_0N(n'))^{-1}=trn_0\,n_0^{-1}r^{-1}t^{-1}n'^{-1}tr\,r^{-1}t^{-1}=
n'^{-1}$. Also $\Phi(N(n'))\le C|\mathcal F(u^{-1}\omega')|^2$ by (a).

\emph{Integration.} Apply Lemma~\ref{pfr:heis} to $n'\mapsto\Phi(N(n'))$ with the unit $u$, the size $j$, the
constant $C$ and $\Psi=|\mathcal F|^2$ on $V_{a+1}$; the integral $I$ of $\Psi$ over the vectors of size $j$
is bounded by Proposition~\ref{pfs:corrplanch} (with $b=a+1$, the kernel $\mathsf k$ at floor $a+1$, $t$, $q'=t_0q$, so that
$t^{-1}q'=rq\in\pfrU_{(<a+1)}$ and $q'g=t_0qg$ has size $i$, the unit $\xi'$, $v_0$, $i'=j$ and Lebesgue measure).
As $\pfrmod(\xi')=\pfrmod(\xi u)$, the result is
\[
\int\Phi\,d\eta_a\le 2C\,\pfsc_a(u)\Lambda^{j+2}\Lambda^{99n(j+6)}\Lambda^{99n(i+12)}\Bigl(\frac P{\pfrmod(\xi u)}\Bigr)^{\pfsrk}
\int|\mathsf k_{a+2}(\sigma g,\sigma y)|^2\,d\theta_{a+1}(\sigma),
\]
and, as $2\le\Lambda$, $j+3+99n(j+6)+99n(i+12)\le100n(i+j+31)$.
\end{proof}

\subsection{Depth along the Heisenberg groups}\label{pfs:ss-depth}

Let $\pfsBall=\pfsBall_\Gamma$ be the position ball of $\Gamma$ (the $x\in\pfrP$ with $M^{-1}(x-1)$ of size $1$),
and for $a\le\varrho$ and $g\in\pfrG$ let
\[
\pfsCB_a(g)=\{z\in\pfrG:\ g^{-1}\nu z\in\pfsBall\text{ for some }\nu\in\pfrU_{(<a)}\}.
\]

\begin{lemma}[Two points of the position ball]\label{pfs:twopos}
If $g\in\pfrK$ and $x,y\in\pfsBall$, then $M^{-1}(gx(gy)^{-1}-1)$ has size $14$.
\end{lemma}
\begin{proof}
$y=M\cdot M^{-1}(y-1)+1$ has size $3$ ((G1): $M\in\Gamma(1)$), so $gy$ has size $5$ and $(gy)^{-1}$ size $9$. The
matrix $g(M^{-1}(x-1)-M^{-1}(y-1))$ has size $1+2+1=4$, and $M^{-1}(gx(gy)^{-1}-1)=g\,M^{-1}(x-y)\,(gy)^{-1}$ has size
$4+9+1$.
\end{proof}

\begin{lemma}[The lift of a Heisenberg element]\label{pfs:lift}
Let $a<\varrho$, $g\in\pfrK$, $y\in\pfrG$ with $g^{-1}y\in\pfsBall$, and $n\in\pfrN_a$ with $ny\in\pfsCB_a(g)$. Then
some vector $v$ congruent to $(n^{-1}-1)e_a^*$ modulo $\langle e_c:c<a\rangle$ has $M^{-1}v$ of size $17$.
\end{lemma}
\begin{proof}
Let $\nu\in\pfrU_{(<a)}$ with $g^{-1}\nu ny\in\pfsBall$, and $\mu=n^{-1}\nu^{-1}n\in\pfrU_{(<a)}$ (as $n\in\pfrG^{(a)}$
normalises $\pfrU_{(<a)}$). Then $g(g^{-1}y)\bigl(g(g^{-1}\nu ny)\bigr)^{-1}=n^{-1}\nu^{-1}=\mu n^{-1}$, so by
Lemma~\ref{pfs:twopos} the matrix $M^{-1}(\mu n^{-1}-1)$ has size $14$, and so has its column $M^{-1}v$ at $e_a^*$, where
$v=(\mu n^{-1}-1)e_a^*$. As $n^{-1}e_a^*\in V_a$ and $\mu$ acts on
$V_a$ by adding combinations of the $e_c$, $c<a$ (convention (iii)), $v\equiv(n^{-1}-1)e_a^*$. Size $14$ implies
size $17$.
\end{proof}

\begin{lemma}[A Heisenberg identity]\label{pfs:heisid}
Let $a<\varrho$, $n=n(\omega,\beta)\in\pfrN_a$, let $\zeta\equiv\zeta'$ modulo $\langle e_c:c<a\rangle$, and let $o\in X'$.
Then $h(\zeta,o)-h(n^{-1}\zeta',o)=h(\zeta,e_a)\,h(\omega,o)$.
\end{lemma}
\begin{proof}
$o\in V_{a+1}$ is orthogonal to $f_a,f_a^*$, so $no=o-h(o,\omega)e_a$ and $h(n^{-1}\zeta',o)=h(\zeta',no)=h(\zeta',o)-
\overline{h(o,\omega)}\,h(\zeta',e_a)$. Here $h(\zeta',o)=h(\zeta,o)$ and $h(\zeta',e_a)=h(\zeta,e_a)$, as $\zeta-\zeta'$ is a
combination of the $e_c$, $c<a$, which are orthogonal to $o$ and to $e_a$; and $\overline{h(o,\omega)}=h(\omega,o)$.
\end{proof}

\begin{lemma}[The inverse depth on a Heisenberg group]\label{pfs:invdepthN}
Let $a<\varrho$, $g\in\pfrK$ and $y\in\pfrG$ with $g^{-1}y\in\pfsBall$. Then
\[
\int\mathbf 1[ny\in\pfsCB_a(g)]\;\pfrdep_a(g,g^{-1}ny)^{-s}\,d\eta_a(n)\le\pfsc_a(M)\,\Lambda^{4600n}.
\]
\end{lemma}
\begin{proof}
\emph{The inverse depth is bounded by a common-vector function.} Let $\phi$ be the function of (G5) for the point
$g\in\pfrK$, the position $g^{-1}y\in\pfrP$ and $v=0$, with set of common vectors $\mathcal W$ (the $o\in X'$ with
$(g^{-1}o,y^{-1}o)$ close) and set of differences $\mathcal D$. Let $n=n(\omega,\beta)$ and $u=(M^{-1}\omega)_{X'}$. We claim
$\pfrdep_a(g,g^{-1}ny)^{-1}=\pbvol\mathfrak d_a(g,g^{-1}ny)\le\phi(u)$. Let $c\in\mathfrak d_a(g,g^{-1}ny)$: $c=h(\zeta,e_a)$
(the $e_a^*$-coordinate) for a vector $\zeta$ stable at $(g,ny)$ for $a$, with $(g^{-1}\zeta,(ny)^{-1}\zeta')$ close for some
$\zeta'\equiv\zeta$ modulo $\langle e_c:c<a\rangle$. For $o\in\mathcal W$ the pair $(g^{-1}o,y^{-1}o)$ is close, so
\[
\mathcal D\ni M^{-1}\bigl(h(g^{-1}\zeta,g^{-1}o)-h((ny)^{-1}\zeta',y^{-1}o)\bigr)=M^{-1}\bigl(h(\zeta,o)-h(n^{-1}\zeta',o)\bigr)=
c\,h(M^{-1}\omega,o)=h(cu,o),
\]
by Lemma~\ref{pfs:heisid}, $\bar M=M$, and $h(\omega',o)=h(\omega'_{X'},o)$ for $o\in X'$ (here $\omega'_{X'}$ is read as the
vector $\omega'_{X'}=\sum_th(\omega',x_t)x_t$ of $X'$, \S\ref{pfr:ss-gauge}). As $\mathcal D$ is stable under conjugation
(\S\ref{pfr:ss-gauge}) and $\overline{h(cu,o)}=h(o,cu)$, also $h(o,cu)\in\mathcal D$ for every $o\in\mathcal W$. So $c$ lies in the
set $\{c:h(o,cu)\in\mathcal D\ \forall o\in\mathcal W\}$ of (G5), whose volume is at most $\phi(u)$.

\emph{Integration.} If the integrand is non-zero at $n$, Lemma~\ref{pfs:lift} gives a vector $v\equiv(n^{-1}-1)e_a^*$
with $M^{-1}v$ of size $17$, and the integrand is at most $\Psi(M^{-1}\omega)$ with $\Psi(w)=\phi(w_{X'})^s$ on $V_{a+1}$.
By the general form of Lemma~\ref{pfr:quadricX} (proved there for Lebesgue measure) and (G5) at size $21$,
$\int_{\{w\text{ of size }17\}}\Psi\le\Lambda^{99n\cdot23}\Lambda^{99n\cdot22}$. Lemma~\ref{pfr:heis} with $u=M$, $j=17$, $C=1$ bounds the
integral by $\pfsc_a(M)\Lambda^{19}\Lambda^{99n\cdot45}\le\pfsc_a(M)\Lambda^{4600n}$.
\end{proof}

\begin{lemma}[The depth weight over the unipotent group]\label{pfs:depthweight}
Let $l\le\varrho$, and let $g\in\pfrG$ be near at $l$. Then for every $y\in\pfrG$,
\[
\int\mathbf 1[g^{-1}\nu y\in\pfsBall]\;\pfrwt_l(g,g^{-1}\nu y)\,d\pfsA_l(\nu)\le\prod_{a<l}\bigl(\pfsc_a(M)\Lambda^{4600n}\bigr).
\]
\end{lemma}
\begin{proof}
\emph{Reduction to $g\in\pfrK$.} Let $\xi g\in\pfrK$ with $\xi\in\pfrR^{(l)}$, $\xi=\gamma\nu_0$ with $\gamma\in\pfrG^{(l)}$,
$\nu_0\in\pfrU_{(<l)}$. Substitute $\nu\mapsto\nu_0^{-1}\gamma^{-1}\nu\gamma$, which preserves $\pfsA_l$ (convention (ii)). Then
$g^{-1}\nu_0^{-1}\gamma^{-1}\nu\gamma y=(\xi g)^{-1}\nu(\gamma y)$, and for every $x$, $\pfrwt_l(g,x)=\pfrwt_l(\xi g,x)$ by (C3)
(the second point being $gx$, resp.\ $\xi gx$). So the integral equals the one for $(\xi g,\gamma y)$; we may assume
$g\in\pfrK$.

\emph{Induction on the length.} Let $g\in\pfrK$ and $J_a(y)$ the integral with $a$ in place of $l$, $a\le l$. We show
$J_a(y)\le C_a:=\prod_{b<a}\bigl(\pfsc_b(M)\Lambda^{4600n}\bigr)$ for all $y$. For $a=0$, $\pfsA_0=\delta_1$ and $\pfrwt_0=1$. Let $a<l$ and
suppose $J_a\le C_a$ everywhere. If $J_{a+1}(y)\ne0$, there is $\tau\in\pfrU_{(<a+1)}$ with $g^{-1}\tau y\in\pfsBall$
($\pfsA_{a+1}$ is carried by $\pfrU_{(<a+1)}$). By the right invariance of $\pfsA_{a+1}$ under $\tau$ and
$\pfsA_{a+1}=\pfsA_a*\eta_a$ (convention (ii)),
\[
J_{a+1}(y)=\iint\mathbf 1[g^{-1}\nu nz\in\pfsBall]\,\pfrwt_{a+1}(g,g^{-1}\nu nz)\,d\pfsA_a(\nu)\,d\eta_a(n),\qquad z=\tau y .
\]
For fixed $n$ and $z'=nz$: $\pfrwt_{a+1}(g,g^{-1}\nu z')=\pfrwt_a(g,g^{-1}\nu z')\,\pfrdep_a(g,g^{-1}\nu z')^{-s}$, and
$\pfrdep_a(g,g^{-1}\nu z')=\pfrdep_a(g,g^{-1}z')$ for $\nu\in\pfrU_{(<a)}$: the second point is $\nu z'$, and the vectors
stable at $(g,\nu z')$ and at $(g,z')$ for $a$ are the same, because $(\nu z')^{-1}\zeta'=z'^{-1}(\nu^{-1}\zeta')$ and
$\nu^{-1}\zeta'\equiv\zeta'$ modulo $\langle e_c:c<a\rangle$ for $\zeta'$ orthogonal to the $e_c$, $c<a$ (convention (iii)). Moreover
the inner integral vanishes unless $z'\in\pfsCB_a(g)$. Hence the inner integral is at most
$J_a(z')\,\mathbf 1[z'\in\pfsCB_a(g)]\,\pfrdep_a(g,g^{-1}z')^{-s}\le C_a\,\mathbf 1[nz\in\pfsCB_a(g)]\,\pfrdep_a(g,g^{-1}nz)^{-s}$,
and Lemma~\ref{pfs:invdepthN} (for the point $z=\tau y$, with $g^{-1}z\in\pfsBall$) gives $J_{a+1}(y)\le C_a\pfsc_a(M)
\Lambda^{4600n}=C_{a+1}$.
\end{proof}

\subsection{One frame: proof of Theorem~\ref{pfr:oneframe}}\label{pfs:ss-oneframe}

Throughout this subsection $\Gamma$ is any gauge for $\pfrF$ (it will be applied to the gauge $\Gamma'$ of
Theorem~\ref{pfr:oneframe}), $\pfsBall$ its position ball, and $\mathsf k$ a kernel at floor $j\le\varrho$. Put
\[
\pfsfc_l=\pfsc_l(M)\,P^{\pfsrk}\Lambda^{4900n}\ (l<\varrho),\qquad\pfsfc_\varrho=1,\qquad
\|\mathsf k\|^2=\iint|\mathsf k(g,y)|^2\,d\theta_{\ge0}(y)\,d\theta_{\ge0}(g).
\]
For $l\le\varrho+1$ and $g,y\in\pfrG$ let
\[
\pfscm_l(g,y)=\pfsA_l\{\nu:\ g^{-1}\nu y\in\pfsBall\},\qquad
\pfscut_l(g,y)=\begin{cases}|\mathsf k_{l+1}(g,y)|^2&\text{if }\pfscm_l(g,y)>0,\\0&\text{otherwise}.\end{cases}
\]
Both are Borel functions of $(g,y)$ (Tonelli). If $\pfscm_l(g,y)>0$ then $y\in\pfsCB_l(g)$, since $\pfsA_l$ is carried by
$\pfrU_{(<l)}$.

\begin{lemma}[Translation rules]\label{pfs:cmrules}
For $l\le\varrho+1$: (a) $\pfscm_l(\sigma g,\sigma y)=\pfscm_l(g,y)$ for $\sigma\in\pfrG^{(l)}$; (b) $\pfscm_l(g,\nu y)=
\pfscm_l(g,y)$ for $\nu\in\pfrU_{(<l)}$; (c) $\int\pfscm_l(g,ny)\,d\eta_l(n)=\pfscm_{l+1}(g,y)$ for $l\le\varrho$.
\end{lemma}
\begin{proof}
(a) $\pfscm_l(\sigma g,\sigma y)=\int\mathbf 1_\pfsBall(g^{-1}\sigma^{-1}\nu\sigma y)\,d\pfsA_l(\nu)$, and $\nu\mapsto\sigma^{-1}\nu\sigma$
preserves $\pfsA_l$ (convention (ii)). (b) Right invariance of $\pfsA_l$ under $\nu$. (c) $\pfsA_{l+1}=\pfsA_l*\eta_l$ and Tonelli.
\end{proof}

\begin{lemma}[A factorisation of the Haar measure]\label{pfs:factor}
For $l\le\varrho+1$ let $\rho_l=\eta_l*\eta_{l+1}*\dots*\eta_\varrho*\theta_{\ge0}$ ($\rho_{\varrho+1}=\theta_{\ge0}$). Then
$\rho_l$ is s-finite, $\rho_l=\eta_l*\rho_{l+1}$ for $l\le\varrho$, and $\pfrM=\pfsA_l*\rho_l$. If $\Phi\ge0$ is Borel with
$\Phi(\nu y)=\Phi(y)$ for all $\nu\in\pfrU_{(<l)}$ and $y$, then $\int\Phi(\sigma y)\,d\rho_l(y)=\int\Phi\,d\rho_l$ for every
$\sigma\in\pfrG^{(l)}$.
\end{lemma}
\begin{proof}
By Lemma~\ref{pfr:rearr} with $j=\varrho+2$ (and $\eta_\varrho=\eta_{\varrho+1}=\delta_1$), $\pfrM=\pfsA_{\varrho+1}*\theta_{[0,\varrho+2)}=
\pfsA_{\varrho+1}*\theta_{\ge0}$; and $\pfsA_{\varrho+1}=\pfsA_l*\eta_l*\dots*\eta_\varrho$, which gives $\pfrM=\pfsA_l*\rho_l$ by associativity.
Convolutions of finitely many s-finite measures are s-finite. For the invariance we use the averaging argument of
the proof of Lemma~\ref{pfr:weylinv}. Choose a Borel $f>0$ with $\int f\,d\pfrM<\infty$ and put
$\mathrm{av}(y)=\int f(\nu y)\,d\pfsA_l(\nu)\in(0,\infty]$; $\mathrm{av}(\nu y)=\mathrm{av}(y)$ for $\nu\in\pfrU_{(<l)}$
by right invariance of $\pfsA_l$. For every s-finite $\mu$ with $\pfrM=\pfsA_l*\mu$ we have $\int\mathrm{av}\,d\mu=\int f\,d\pfrM<\infty$,
so $\mathrm{av}<\infty$ $\mu$-a.e., and
\[
\int\frac{\Phi f}{\mathrm{av}}\,d\pfrM=\iint\frac{\Phi(\nu y)f(\nu y)}{\mathrm{av}(\nu y)}\,d\pfsA_l(\nu)\,d\mu(y)=\int\Phi\,d\mu .
\]
The left side does not depend on $\mu$. For $\sigma\in\pfrG^{(l)}$, $\mu=\delta_\sigma*\rho_l$ is admissible:
$\pfsA_l*\delta_\sigma*\rho_l=\delta_\sigma*\pfsA_l*\rho_l=\delta_\sigma*\pfrM=\pfrM$ by convention (ii) and left invariance. Hence
$\int\Phi(\sigma y)\,d\rho_l(y)=\int\Phi\,d(\delta_\sigma*\rho_l)=\int\Phi\,d\rho_l$.
\end{proof}

\begin{lemma}[One step up]\label{pfs:climb}
Let $j\le l\le\varrho$ and $g,y\in\pfrG$. Then
\[
\int\pfscut_l(g,ny)\,d\eta_l(n)\le\pfsfc_l\int\pfscut_{l+1}(\sigma g,\sigma y)\,d\theta_{l+1}(\sigma).
\]
\end{lemma}
\begin{proof}
If $\pfscm_{l+1}(g,y)=0$, then by Lemma~\ref{pfs:cmrules}(c) $\pfscm_l(g,ny)=0$ for $\eta_l$-almost every $n$, and the left
side is $0$. Otherwise, by Lemma~\ref{pfs:cmrules}(a), $\pfscm_{l+1}(\sigma g,\sigma y)=\pfscm_{l+1}(g,y)>0$ for
$\sigma\in\pfrG^{(l+1)}$, which carries $\theta_{l+1}$; so the right side is $\pfsfc_l\int|\mathsf k_{l+2}(\sigma g,\sigma y)|^2\,
d\theta_{l+1}(\sigma)$, and $\pfscut_l(g,ny)\le|\mathsf k_{l+1}(g,ny)|^2$.

\emph{The top floor, $l=\varrho$.} Here $\eta_\varrho=\delta_1$ and $\pfsfc_\varrho=1$. As $\theta_{\ge\varrho+2}=\delta_1$ and
$\theta_{\varrho+1}=\mu_\pfrH$, (C1) gives $\mathsf k_{\varrho+1}(g,y)=\int\pfrip{W(\sigma g)}{W(d'\sigma y)}\,d\mu_\pfrH(\sigma)$; the integrand vanishes
unless $W(\sigma g)\ne0$, hence unless $u\sigma g\in\pfrK$ for some $u\in\pfrU$, a set of $\mu_\pfrH$-mass at most $1$ by the
definition of a place datum. By Cauchy--Schwarz, $|\mathsf k_{\varrho+1}(g,y)|^2\le\int|\mathsf k_{\varrho+2}(\sigma g,\sigma y)|^2\,
d\theta_{\varrho+1}(\sigma)$.

\emph{Below the top floor, $l<\varrho$.} We apply Proposition~\ref{pfs:parseval} with $a=l$, the kernel $\mathsf k$ (a kernel at
floor $l+1$), $\varepsilon=0$, $\xi=M^{-1}$, $u=M$, $C=1$, $i=1$, $j=17$, the points $g$ and $y$, and
\[
S=\{n\in\pfrN_l:\ ny\in\pfsCB_l(g)\text{ and }g\text{ is near at }l+1\},\qquad\Phi(n)=\pfscut_l(g,ny).
\]
\eqref{pfs:eq-twist}: for $s\in\pfrN_l\subseteq\pfrG^{(j)}$, $d'sd'^{-1}=s$, so $M^{-1}\sigma(d'sd'^{-1})=M^{-1}\sigma(s)$. (a): if
$\Phi(n)\ne0$, then $\pfscm_l(g,ny)>0$, so $ny\in\pfsCB_l(g)$, and $\mathsf k_{l+1}(g,ny)\ne0$, so $g$ is near at $l+1$ by (C2);
and $\Phi(n)\le|\mathsf k_{l+1}(n^0g,ny)|^2$. (b): let $n_0\in S$ and $\nu_0\in\pfrU_{(<l)}$ with $g^{-1}\nu_0n_0y\in\pfsBall$. As
$g$ is near at $l+1$, $\xi_0g\in\pfrK$ for some $\xi_0\in\pfrR^{(l+1)}\subseteq\pfrR^{(l)}$; write $\xi_0=t\nu$ with $t\in\pfrG^{(l)}$,
$\nu\in\pfrU_{(<l)}$ (convention (iii)). Then $t=\xi_0\nu^{-1}\in\pfrR^{(l+1)}\cap\pfrG^{(l)}$, $q=\nu\in\pfrU_{(<l+1)}$, and
$tqg=\xi_0g\in\pfrK$ has size $1$. Let $n\in S$, with $\mu\in\pfrU_{(<l)}$ such that $g^{-1}\mu ny\in\pfsBall$. Put $\delta=nn_0^{-1}
\in\pfrN_l$, $g'=t\nu g\in\pfrK$, $y'=t\nu\nu_0n_0y$ and $n'=t\delta t^{-1}$. Then $n'\in\pfrN_l$ (write $t=s'r'$ with
$s'\in\pfrG^{(l+1)}$, $r'\in\pfrN_l$ by (N1), and use (N2)); $g'^{-1}y'=g^{-1}\nu_0n_0y\in\pfsBall$; and
$n'y'\in\pfsCB_l(g')$, because $\kappa=t\nu\mu\,\delta(\nu\nu_0)^{-1}\delta^{-1}t^{-1}$ lies in $\pfrU_{(<l)}$ ($\delta$ and $t$ lie in
$\pfrG^{(l)}$, which normalises $\pfrU_{(<l)}$) and $g'^{-1}\kappa n'y'=g^{-1}\mu ny\in\pfsBall$. By Lemma~\ref{pfs:lift} (for
$g'$, $y'$, $n'$) some $v\equiv(n'^{-1}-1)e_l^*$ has $M^{-1}v$ of size $17$; and $n'^{-1}=tn_0n^{-1}t^{-1}=tn_0(tn)^{-1}$.

Proposition~\ref{pfs:parseval} gives the bound $(P/\pfrmod(M^{-1}M))^{\pfsrk}\pfsc_l(M)\Lambda^{100n\cdot49}\int|\mathsf k_{l+2}
(\sigma g,\sigma y)|^2d\theta_{l+1}(\sigma)$, and $100n\cdot49=4900n$.
\end{proof}

\begin{lemma}[All floors above $l$]\label{pfs:allfloors}
For $j\le l\le\varrho+1$,
\[
\Bigl(\prod_{a<l}\pfsfc_a\Bigr)\iint\pfscut_l(g,y)\,d\rho_l(y)\,d\theta_{[0,l+1)}(g)\le\Bigl(\prod_{a<\varrho}\pfsfc_a\Bigr)\|\mathsf k\|^2 .
\]
\end{lemma}
\begin{proof}
Descending induction on $l$. For $l=\varrho+1$: $\rho_{\varrho+1}=\theta_{[0,\varrho+2)}=\theta_{\ge0}$, $\pfscut_{\varrho+1}\le|\mathsf k_{\varrho+2}|^2=|\mathsf k|^2$
($\theta_{\ge\varrho+2}=\delta_1$), and $\prod_{a\le\varrho}\pfsfc_a=\prod_{a<\varrho}\pfsfc_a$. Let $j\le l\le\varrho$. For fixed $g$, by
$\rho_l=\eta_l*\rho_{l+1}$, Lemma~\ref{pfs:climb} and Tonelli,
\[
\int\pfscut_l(g,y)\,d\rho_l(y)=\iint\pfscut_l(g,ny)\,d\eta_l(n)\,d\rho_{l+1}(y)\le\pfsfc_l\iint\pfscut_{l+1}(\sigma g,\sigma y)\,
d\rho_{l+1}(y)\,d\theta_{l+1}(\sigma).
\]
For $\sigma\in\pfrG^{(l+1)}$ the function $\Phi(y)=\pfscut_{l+1}(\sigma g,y)$ satisfies $\Phi(\nu y)=\Phi(y)$ for $\nu\in\pfrU_{(<l+1)}$,
by Lemma~\ref{pfs:cmrules}(b) and (K3); so by Lemma~\ref{pfs:factor} the inner integral is $\int\pfscut_{l+1}(\sigma g,y)\,d\rho_{l+1}(y)$.
Integrating over $g$ against $\theta_{[0,l+1)}$ and using $\theta_{[0,l+2)}=\theta_{l+1}*\theta_{[0,l+1)}$, the left side of the
claim at $l$ is at most the left side at $l+1$.
\end{proof}

\begin{proposition}[One frame, one pivot]\label{pfs:pivot}
For $j\le l\le\varrho$ (and $\mathsf k$ a kernel at floor $j$),
\[
\int_{\pfsBall}\int\mathsf a_l(g;x)\,d\theta_{[0,l+1)}(g)\,d\pfrM(x)\le\Bigl(\prod_{a<\varrho}\pfsfc_a\Bigr)\|\mathsf k\|^2,
\qquad\mathsf a_l(g;x)=\pfrwt_l(g,x)\,|\mathsf k_{l+1}(g,gx)|^2 .
\]
\end{proposition}
\begin{proof}
By Tonelli it suffices to bound, for fixed $g$, $\int\mathbf 1_\pfsBall(x)\mathsf a_l(g;x)\,d\pfrM(x)$ by
$(\prod_{a<l}\pfsfc_a)\int\pfscut_l(g,y)\,d\rho_l(y)$, and then to apply Lemma~\ref{pfs:allfloors}. By the left invariance of
$\pfrM$ under $g^{-1}$ and $\pfrM=\pfsA_l*\rho_l$,
\[
\int\mathbf 1_\pfsBall(x)\mathsf a_l(g;x)\,d\pfrM(x)=\iint\mathbf 1_\pfsBall(g^{-1}\nu y)\,\pfrwt_l(g,g^{-1}\nu y)\,|\mathsf k_{l+1}(g,\nu y)|^2
\,d\pfsA_l(\nu)\,d\rho_l(y).
\]
By (K3), $|\mathsf k_{l+1}(g,\nu y)|=|\mathsf k_{l+1}(g,y)|$ for $\nu\in\pfrU_{(<l)}$. Fix $y$. If $\mathsf k_{l+1}(g,y)=0$, the inner
integral vanishes. If $\pfscm_l(g,y)=0$, the indicator vanishes for $\pfsA_l$-almost every $\nu$. Otherwise $g$ is near at $l+1$,
hence at $l$ (C2), $\pfscut_l(g,y)=|\mathsf k_{l+1}(g,y)|^2$, and by Lemma~\ref{pfs:depthweight} the inner integral is at most
$\prod_{a<l}\bigl(\pfsc_a(M)\Lambda^{4600n}\bigr)\cdot\pfscut_l(g,y)\le\prod_{a<l}\pfsfc_a\cdot\pfscut_l(g,y)$, as $P\ge1$.
\end{proof}

\begin{proposition}[One frame]\label{pfs:frame}
For a kernel $\mathsf k$ at floor $j\le\varrho$,
$\int_{\pfsBall}\mathsf S_j(\mathsf k;x)\,d\pfrM(x)\le n\,(\prod_{a<\varrho}\pfsfc_a)\,\|\mathsf k\|^2$.
\end{proposition}
\begin{proof}
Unfolding the recursion of \S\ref{pfr:ss-gauge} with $\theta_{[0,l+1)}=\theta_l*\theta_{[0,l)}$ gives $\int\mathsf b_l\,d\theta_{[0,l)}=
\int\mathsf a_l\,d\theta_{[0,l+1)}+\int\mathsf b_{l+1}\,d\theta_{[0,l+1)}$ for $l\le\varrho$, and $\mathsf b_n=0$; so
$\mathsf S_j(\mathsf k;x)=\sum_{l=j}^{\varrho}\int\mathsf a_l(g;x)\,d\theta_{[0,l+1)}(g)$, a sum of at most $n$ terms, each bounded
by Proposition~\ref{pfs:pivot}.
\end{proof}

\begin{lemma}[Translation by a Weyl ratio]\label{pfs:weylratio}
Let $j\le\varrho$, $p,q$ sequences of units, $d=w_{j,q}^{-1}w_{j,p}$, and $\Phi\ge0$ Borel on $\pfrG$. Then
\[
\int\Phi(dy)\,d\theta_{\ge0}(y)=\prod_{a<j}\pfrmod(p_a/q_a)^{d_a+1}\int\Phi\,d\theta_{\ge0}.
\]
\end{lemma}
\begin{proof}
Induction on $j$; $j=0$ is trivial. Let $j<\varrho$, write $d_l=w_{l,q}^{-1}w_{l,p}$, and for a unit $c$ let $w^{(j)}_c$ be the
Weyl factor of floor $j$: $e_j\mapsto ce_j^*$, $e_j^*\mapsto\bar c^{-1}e_j$, the identity elsewhere. By (iv) of
\S\ref{pfr:ss-descent}, $\theta_j=\lambda_j*\eta_j$ with $\lambda_j$ the image of $\pfrmod(c)^{d_j}dc$ under $c\mapsto(w^{(j)}_c)^{-1}$;
these elements commute with $\pfrG^{(j+1)}$, which carries $\theta_{\ge j+1}$. Hence $\theta_{\ge0}=\lambda_j*\mu$ with
$\mu=\theta_{\ge j+1}*\eta_j*\theta_{[0,j)}$, and for Borel $\Psi\ge0$,
$\int\Psi\,d\theta_{\ge0}=\int\pfrmod(c)^{d_j}\int\Psi((w^{(j)}_c)^{-1}y)\,d\mu(y)\,dc$. With $r=p_j/q_j$ one checks on $e_j,e_j^*$
that $d_{j+1}(w^{(j)}_{rc})^{-1}=d_j(w^{(j)}_c)^{-1}$ (both send $e_j^*\mapsto c^{-1}e_j$ and $e_j\mapsto\bar ce_j^*$, and agree with
$d_j$ elsewhere). Substituting $c=rc'$ ($dc=\pfrmod(r)\,dc'$; the non-units are null),
\[
\int\Phi(d_{j+1}y)\,d\theta_{\ge0}=\pfrmod(r)^{d_j+1}\int\pfrmod(c')^{d_j}\int\Phi(d_j(w^{(j)}_{c'})^{-1}y)\,d\mu\,dc'=
\pfrmod(r)^{d_j+1}\int\Phi(d_jy)\,d\theta_{\ge0},
\]
and the induction hypothesis concludes.
\end{proof}

\begin{proof}[Proof of Theorem~\ref{pfr:oneframe}]
Let $\pfsBall'$ be the position ball of $\Gamma'$, $\Lambda'=\Lambda_{\Gamma'}$, and $\pfsfc'_l$ the constants above for $\Gamma'$. Let
$d=w_{j,q}^{-1}w_{j,p}$, $D=(\prod_{a<j}\pfsc_a(p_a^{-1}))^2$, $J=\prod_{a<j}\pfrmod(p_a/q_a)^{d_a+1}$ and
$E_1=\int_{\pfsBall'}\mathsf S_0(\mathsf k^1;x)\,d\pfrM$, $E_2=\int_{\pfsBall'}\mathsf S_j(\mathsf k^d;x)\,d\pfrM$. Let
$\|W\|^2=\iint|\pfrip{W(g)}{W(y)}|^2\,d\theta_{\ge0}(g)\,d\theta_{\ge0}(y)$, so that $\mathsf N(W)=\prod_{l<\varrho}\pfsc_l(M)
P^{\varrho\pfsrk}\|W\|^2$.

\emph{Cauchy--Schwarz.} On $S\subseteq\pfsBall'$, $f^2\le e^2\mathsf Q(x)=e^2DC_\Gamma^3\mathsf S_0(\mathsf k^1;x)\mathsf S_j(\mathsf k^d;x)$, so
$(\int_Sf\,d\pfrM)^2\le e^2DC_\Gamma^3E_1E_2$. \emph{The two frames.} By Proposition~\ref{pfs:frame} for $\Gamma'$ and the kernels $\mathsf k^1$
(at floor $0$) and $\mathsf k^d$ (at floor $j$), $E_1\le nV\|\mathsf k^1\|^2$ and $E_2\le nV\|\mathsf k^d\|^2$ with $V=\prod_{l<\varrho}\pfsfc'_l$;
here $\|\mathsf k^1\|^2=\|W\|^2$ and, by Lemma~\ref{pfs:weylratio} applied to $y\mapsto|\pfrip{W(g)}{W(y)}|^2$ for each $g$,
$\|\mathsf k^d\|^2=J\|W\|^2$. \emph{The head.} For each $a<j$,
$\pfrmod(p_aq_a)^{d_a}\cdot\pfrmod(p_a)^{-(2d_a+1)}\cdot\pfrmod(p_a/q_a)^{d_a+1}=\pfrmod(q_a)^{-1}$; hence
$\pfrkap_j(p,q)^2DJ=\pfrmod((q_0\cdots q_{j-1})^{-1})\le\Lambda$ by the head condition on $q$ and (G3). Altogether
\[
\pfrkap_j(p,q)^2\Bigl(\int_Sf\,d\pfrM\Bigr)^2\le\Lambda\,e^2\,C_\Gamma^3\,(nV)^2\,\|W\|^4 .
\]
\emph{Constants.} $n\le2^n\le\Lambda'^n$, and $V=\prod_{l<\varrho}\pfsc_l(M)P^{\varrho\pfsrk}\Lambda'^{4900n\varrho}$; as $n+4900n\varrho\le25000n^2$,
$nV\|W\|^2\le\Lambda'^{25000n^2}\mathsf N(W)$. Taking square roots, with $C_\Gamma^{3/2}=\Lambda^{1500000n^3}$ and $\Lambda^{1/2}\le\Lambda$,
\[
\pfrkap_j(p,q)\int_Sf\,d\pfrM\le e\,\Lambda^{1500000n^3+1}\Lambda'^{25000n^2}\mathsf N(W).\qedhere
\]
\end{proof}

\subsection{The gains from a frequency and from two heads}\label{pfs:ss-gains}

\subsubsection*{Sizes at two near points}
\begin{lemma}[An element of the unipotent group]\label{pfs:unisize}
Let $j\le\varrho$, $d\ge1$, and $s\in\pfrU_{(<j)}$ such that the columns $se_a^*$, $a<j$, have size $d$. Then $s$ has size
$j(d+1)+4$ and $s^{-1}$ has size $j(d+1)+8$.
\end{lemma}
\begin{proof}
Let $C_{ab}=h(se_a^*,e_b)$ ($a,b<j$), the $e_b^*$-coordinate of $se_a^*$; it has size $d$. As $s^{-1}e_b\in e_b+\langle e_c:c<b\rangle$
(convention (iii)), $C_{ab}=h(e_a^*,s^{-1}e_b)=\delta_{ab}$ for $b\le a$; so the matrix $\bar C=(\bar C_{ab})$ is unitriangular,
and the entries of $\bar C^{-1}$ are $(j-1)\times(j-1)$ minors, of size $(j-1)(d+1)$ by (G0). Note $h(e_c,se_a^*)=\bar C_{ac}$.
\emph{Columns at $e_b$, $b<j$.} $se_b=e_b+\sum_{c<b}\mu_ce_c$, and $h(se_b,se_a^*)=h(e_b,e_a^*)=\delta_{ab}$ for $a<j$ gives
$\bar C\mu'=\mathbf e_b$ with $\mu'=(\mu_0,\dots,\mu_{b-1},1,0,\dots)$; so the $\mu_c$ are entries of $\bar C^{-1}$. \emph{Columns at basis
vectors $v$ of floor $\ge j$.} $sv=v+\sum_{c<j}\lambda_ce_c$ (convention (iii)), and $h(sv,se_a^*)=h(v,e_a^*)=0$ gives
$\bar C\lambda=-(h(v,se_a^*))_{a<j}$. Each $h(v,se_a^*)$ has size $d+3$ (products with the Gram matrix, of size $1$, and a sum of
at most $r$ terms), so each $\lambda_c$ has size $(j-1)(d+1)+(d+3)+1+1=j(d+1)+4$. \emph{Columns at $e_a^*$, $a<j$}: size $d$.
So $s$ has size $j(d+1)+4$, and $s^{-1}=J^{-1}s^*J$ has size $j(d+1)+8$ by (G0).
\end{proof}

\begin{lemma}[One translate for two near points]\label{pfs:twosizes}
Let $j\le\varrho$, $p$ a head of length $j$ for $\Gamma$, and $g\in\pfrG$ such that $g$ and $w_p^{-1}g$ are near at $j$. Then there
is $t\in\pfrG^{(j)}$ such that $tg$ has size $60n^2$ and $tw_p^{-1}g$ has size $99n^2$.
\end{lemma}
\begin{proof}
Let $\xi g\in\pfrK$ with $\xi\in\pfrR^{(j)}$, $\xi=t\nu$ with $t\in\pfrG^{(j)}$, $\nu\in\pfrU_{(<j)}$, and $s=t\nu t^{-1}\in\pfrU_{(<j)}
\subseteq\pfrR^{(j)}$; then $s(tg)=\xi g\in\pfrK$. The element $t$ commutes with $w_p$ ($w_p$ acts on the pairs of index $<j$, which $t$
fixes, and is the identity on $V_j$, which $t$ preserves), so $w_p^{-1}(tg)=t\,w_p^{-1}g$ is near at $j$ ($t\in\pfrR^{(j)}$). By
Lemma~\ref{pfr:headentries} at the point $tg$, $p_a$, $p_a^{-1}$ and the columns $se_a^*$ ($a<j$) have size $10n$. By
Lemma~\ref{pfs:unisize}, $s^{-1}$ has size $j(10n+1)+8\le20n^2$ (as $j<n$ and $n\ge2$), so $tg=s^{-1}(\xi g)$ has size $20n^2+2\le60n^2$.
The entries of $w_p^{-1}$ are $0$, $1$, $\bar p_a$ and $p_a^{-1}$, of size $10n$, so $tw_p^{-1}g=w_p^{-1}(tg)$ has size
$10n+60n^2+1\le99n^2$.
\end{proof}

\subsubsection*{The gain from a frequency}
\begin{proof}[Proof of Theorem~\ref{pfr:gain-char}]
Write $w=w_{j,p}$, $d=w_{j,q}^{-1}w_{j,p}$, $\omega(\gamma)=\psi(\eta h(\gamma f_j,f_j))$, $K=(P/\pfrmod(\xi))^{\pfsrk}$, and for
$\sigma\in\pfrG^{(j)}$ write $y_\sigma=\sigma^{-1}f_j$. Fix $g$ and put
\[
f(\sigma)=\mathsf k^1_{j+1}(\sigma g,\sigma gx),\qquad f'(\tau)=\mathsf k^d_{j+1}(\tau w^{-1}g,\tau w^{-1}gx).
\]
As $\sigma$ is an isometry, $h(\sigma\tau^{-1}f_j,f_j)=h(y_\tau,y_\sigma)$, so $\Phi_j(x;\omega;g)=\iint\psi(\eta h(y_\tau,y_\sigma))f(\sigma)
\overline{f'(\tau)}\,d\theta_j(\tau)\,d\theta_j(\sigma)$. By Lemma~\ref{pfr:separate} (with $f$ there the function $|\Phi_j(x;\omega;\cdot)|$,
whose $\pfrm_j$-integral is $\pfrF_j(x;\omega)$) and the first claim of Lemma~\ref{pfr:floorsq} for the two kernels (which bounds
$\int|f|^2d\theta_j$ by $\mathsf b^{(1)}_j(g)/\pfrwt_j(g,x)$ and $\int|f'|^2d\theta_j$ by $\mathsf b^{(d)}_j(w^{-1}g)/\pfrwt_j(w^{-1}g,x)$), it suffices to
show, for every $g$ with $\Phi_j(x;\omega;g)\ne0$,
\begin{equation}\label{pfs:eq-char}
|\Phi_j(x;\omega;g)|^2\le K\,C_\Gamma\int|f|^2\,d\theta_j\int|f'|^2\,d\theta_j .
\end{equation}
\emph{A common translate.} As $\Phi_j(x;\omega;g)\ne0$ and $\theta_j$ is carried by $\pfrG^{(j)}$, there are $\sigma_0,\tau_0\in\pfrG^{(j)}$
with $f(\sigma_0)\ne0\ne f'(\tau_0)$; by (C2), $\sigma_0g$ and $\tau_0w^{-1}g$ are near at $j+1$, so $g$ and $w^{-1}g$ are near at $j$
($\pfrR^{(j+1)}\pfrG^{(j)}\subseteq\pfrR^{(j)}$). Let $t$ be given by Lemma~\ref{pfs:twosizes}. \emph{Supports.} If $f(\sigma)\ne0$ ($\sigma\in
\pfrG^{(j)}$), then $\sigma g=(\sigma t^{-1})(tg)$ is near at $j+1$, and by Lemma~\ref{pfs:nearsize}, $ty_\sigma=(\sigma t^{-1})^{-1}f_j$ has size
$60n^2+6$. Likewise $f'(\tau)\ne0$ implies that $ty_\tau$ has size $99n^2+6$. Let $S$ and $S'$ be the sets of these $\sigma$, resp.\ $\tau$.
\emph{The character.} For $\sigma\in S$, $\tau\in S'$, $h(y_\tau,y_\sigma)=h(ty_\tau,ty_\sigma)$ has size $(60n^2+6)+(99n^2+6)+5\le200n^2$ (by
(G0); $n\ge2$), so $\psi(\eta h(y_\tau,y_\sigma))=\psi(\xi h(ty_\tau,ty_\sigma))$. Hence, with
\[
G(v)=\int\psi\bigl(-\xi\,h(ty_\tau,v)\bigr)f'(\tau)\,d\theta_j(\tau),
\]
the inner integral equals $f(\sigma)\overline{G(ty_\sigma)}$ for $\theta_j$-almost every $\sigma$ (both sides vanish for $\sigma\notin S$, and
$\overline{\psi(c)}=\psi(-c)$), and $\Phi_j(x;\omega;g)=\int_Sf(\sigma)\overline{G(ty_\sigma)}\,d\theta_j(\sigma)$. \emph{Cauchy--Schwarz and
Plancherel.} The map $\sigma\mapsto ty_\sigma$ carries $\theta_j$ to $c_j\mu_j$, $c_j\le2$; so
\[
|\Phi_j(x;\omega;g)|^2\le2\int|f|^2\,d\theta_j\cdot\int_{\{v\text{ of size }60n^2+6\}}|G(v)|^2\,d\mu_j(v),
\]
and Proposition~\ref{pfs:corrplanch} (with $b=j$, the kernel $\mathsf k^d$ at floor $j$, $t$, $q'=t$, the points $w^{-1}g$ and $w^{-1}gx$,
$i=99n^2$, the unit $-\xi$, $v_0=0$, $i'=60n^2+6$, $\rho=\mu_j$) bounds the last integral by
$2\Lambda^{99n(60n^2+12)}\Lambda^{99n(99n^2+12)}K\int|f'|^2d\theta_j$. As $4\le\Lambda^2$ and $99n(159n^2+24)+2\le10^6n^3$, \eqref{pfs:eq-char} follows.
\end{proof}

\subsubsection*{The gain from two heads}
\begin{lemma}\label{pfs:col}
Let $i<\varrho$, $t,t'\in\pfrG^{(i+1)}$, $n_0,n\in\pfrG$ and $h\in\pfrG$ such that $tn_0h$ and $t'nh$ have size $s$. Then the column
$(tn_0(tn)^{-1}-1)e_i^*$ has size $2s+6$.
\end{lemma}
\begin{proof}
$t^{-1}$ and $t'^{-1}$ fix $e_i^*$, so $tn_0(tn)^{-1}e_i^*=tn_0n^{-1}e_i^*=(tn_0h)(t'nh)^{-1}e_i^*$; the matrix $(tn_0h)(t'nh)^{-1}-1$ has size
$s+(s+4)+1+1$ by (G0).
\end{proof}

\begin{lemma}[Parseval along the last vector of two heads]\label{pfs:secbound}
Let $j=i+1\le\varrho$, $p$ a head of length $j$ for $\Gamma$, $x$ a position, $\xi=M^{-1}(p_i^{-1}-q_i^{-1})$ a unit, $w=w_{j,p}$,
$d=w_{j,q}^{-1}w_{j,p}$. Then for every $h\in\pfrG$,
\[
\int\mathbf 1[wnh\text{ near at }j]\,|\mathsf k^d_{i+1}(nh,nhx)|^2\,d\eta_i(n)\le\Bigl(\frac P{\pfrmod(\xi)}\Bigr)^{\pfsrk}\pfsc_i(p_i^{-1})\,C_\Gamma
\int|\mathsf k^d_{i+2}(\sigma h,\sigma hx)|^2\,d\theta_{i+1}(\sigma).
\]
\end{lemma}
\begin{proof}
Apply Proposition~\ref{pfs:parseval} with $a=i$, the kernel $\mathsf k^d$ (at floor $i+1$), $\varepsilon=1$, the units $-\xi p_i$ and
$u=p_i^{-1}$, $C=1$, sizes $99n^2$ and $999n^2$, the points $h$, $hx$, the function $\Phi(n)=\mathbf 1[wnh\text{ near at }j]|\mathsf k^d_{i+1}(nh,
nhx)|^2$ and $S=\{n\in\pfrN_i:\ wnh\text{ and }nh\text{ are near at }j\}$.
\eqref{pfs:eq-twist}: $d$ multiplies $f_i$ by $p_i/q_i$ and fixes $f_{i+1}$, so for $s=n(\omega,\beta)\in\pfrN_i$,
$dsd^{-1}f_{i+1}=f_{i+1}-(p_i/q_i)h(f_{i+1},\omega)f_i$, that is $\sigma(dsd^{-1})=(p_i/q_i)\sigma(s)$, and
$M^{-1}(p_i/q_i-1)\sigma(s)=-\xi p_i\,\sigma(s)$. (a): if $\Phi(n)\ne0$, then $wnh$ is near at $j$ and, by (C2), so is $nh$; and
$\Phi(n)\le|\mathsf k^d_{i+1}(n^1h,nhx)|^2$. (b): let $n_0\in S$. Lemma~\ref{pfs:twosizes} at the point $wn_0h$ (near at $j$, with
$w^{-1}(wn_0h)=n_0h$ near at $j$) gives $t\in\pfrG^{(j)}$ with $tn_0h$ of size $99n^2$; take $t_0=t\in\pfrR^{(i+1)}\cap\pfrG^{(i)}$ and
$q=n_0\in\pfrU_{(<i+1)}$. For $n\in S$, Lemma~\ref{pfs:twosizes} likewise gives $t'\in\pfrG^{(j)}$ with $t'nh$ of size $99n^2$, and by
Lemma~\ref{pfs:col} the vector $v=(tn_0(tn)^{-1}-1)e_i^*$ has size $198n^2+6$. By Lemma~\ref{pfr:headentries} at the point $wnh$ ($wnh$ and
$w^{-1}(wnh)=nh$ are near at $j$), $p_i$ has size $10n$, so $u^{-1}v=p_iv$ has size $10n+198n^2+7\le999n^2$.

Proposition~\ref{pfs:parseval} gives the bound with the constant $(P/\pfrmod(-\xi p_ip_i^{-1}))^{\pfsrk}\pfsc_i(p_i^{-1})\Lambda^{100n(1098n^2+31)}$,
and $100n(1098n^2+31)\le10^6n^3$.
\end{proof}

\begin{proof}[Proof of Theorem~\ref{pfr:gain-heads}]
Let $j=i+1$, $w=w_{j,p}$, $d=w_{j,q}^{-1}w_{j,p}$, $K=(P/\pfrmod(\xi))^{\pfsrk}$, $\varphi_1(g)=|\mathsf k^1_{j}(g,gx)|$, $\varphi_2(g)=|\mathsf k^d_{j}(g,
gx)|$, and $\pfrwt^-_l$ as in Lemma~\ref{pfr:invweight}. For $\omega=1$ the double integral defining $\Phi_j$ factorises, and by (C1)
$|\Phi_j(x;1;g)|=\varphi_1(g)\varphi_2(w^{-1}g)$; so $\pfrF_j(x;1)=\int\varphi_1(g)\varphi_2(w^{-1}g)\,d\pfrm_j(g)$.

\emph{Crossing the weights.} Put $u(h)=\mathbf 1[wh\text{ near at }j]\,\pfrwt_{j}(h,x)\,\pfrwt^-_i(w_{i,p}h)\,\varphi_2(h)^2$. We claim
\[
\bigl(\varphi_1(g)\varphi_2(w^{-1}g)\bigr)^2\le\bigl(\pfrwt_i(g,x)\varphi_1(g)^2\bigr)\,\pfrwt^-_j(w^{-1}g)\cdot u(w^{-1}g).
\]
The left side vanishes unless $g$ and $w^{-1}g$ are near at $j$ ((C2)). Then $\pfrwt^-_j(w^{-1}g)=\pfrwt_j(w^{-1}g,x)^{-1}$; and
$w_{i,p}w^{-1}$ is the inverse of the Weyl factor of floor $i$, an element of $\pfrG^{(i)}\subseteq\pfrR^{(i)}$, under which
$\pfrwt^-_i$ is invariant ((C3) and the invariance of nearness), so $\pfrwt^-_i(w_{i,p}w^{-1}g)=\pfrwt^-_i(g)=\pfrwt_i(g,x)^{-1}$. The
weights are finite and non-zero, and the right side is the left side.

\emph{Cauchy--Schwarz.} $\pfrF_j(x;1)^2\le I_1I_2$ with $I_1=\int\pfrwt_i(g,x)\varphi_1(g)^2\pfrwt^-_j(w^{-1}g)\,d\pfrm_j(g)$ and
$I_2=\int u(w^{-1}g)\,d\pfrm_j(g)$.

\emph{The first integral.} $w^{-1}=w_{\bar p}$ with the head $\bar p$ (proof of Lemma~\ref{pfr:separate}), and the function
$\pfrwt_i(g,x)\varphi_1(g)^2=\mathsf a^{(1)}_i(g)$ is invariant under $\pfrR^{(j)}$ by (K3) and vanishes off the points near at $j$ by (C2).
Lemma~\ref{pfr:invweight} (head $\bar p$ of length $j$, $b'=j$) gives $I_1\le\prod_{a<j}\pfsc_a(p_a^{-1})\,C_\Gamma\int\mathsf a^{(1)}_i\,d\theta_{[0,j)}$.

\emph{The second integral.} The function $u$ is invariant under $\pfrG^{(j)}$: for $\sigma\in\pfrG^{(j)}$, $w\sigma h=\sigma wh$, $\pfrwt_j(\sigma h,x)=
\pfrwt_j(h,x)$ (C3), $\pfrwt^-_i(w_{i,p}\sigma h)=\pfrwt^-_i(\sigma w_{i,p}h)=\pfrwt^-_i(w_{i,p}h)$, and $\varphi_2(\sigma h)=\varphi_2(h)$ by (K2) for the
kernel $\mathsf k^d$ at floor $j$. By Lemma~\ref{pfr:weylinv}, $I_2=\int u\,d\pfrm_j$. By Lemma~\ref{pfr:rearr} and convention (ii),
$\pfrm_j=\pfsA_j*\theta_{[0,j)}=\eta_i*\nu$ with $\nu=\pfsA_i*\theta_{[0,j)}$. For $n\in\pfrN_i$: $\pfrwt_j(nh,x)=\pfrwt_j(h,x)$ ($n\in\pfrR^{(j)}$),
and $\pfrwt^-_i(w_{i,p}nh)=\pfrwt^-_i(nw_{i,p}h)=\pfrwt^-_i(w_{i,p}h)$ ($n\in\pfrG^{(i)}$ commutes with $w_{i,p}$ and lies in $\pfrR^{(i)}$). Hence, by
Lemma~\ref{pfs:secbound} and the first claim of Lemma~\ref{pfr:floorsq} (at floor $j$, for $\mathsf k^d$),
\[
\int u(nh)\,d\eta_i(n)\le K\,\pfsc_i(p_i^{-1})\,C_\Gamma\,\pfrwt^-_i(w_{i,p}h)\,\mathsf b^{(d)}_j(h).
\]
The function $\mathsf b^{(d)}_j$ is invariant under $\pfrR^{(j)}$ and vanishes off the points near at $j$ (Remark~\ref{pfs:rem-corr}(a)). The
proof of Lemma~\ref{pfr:invweight} for the head $p$ of length $i$ and $b'=j$, with $\theta_{[0,j)}$ in place of $\theta_{[0,i)}$ in its Step~2
(Step~2 uses only that $\Phi(\nu y)=\Phi(y)$ for $\nu\in\pfrU_{(<i)}$ and the pointwise bound of Step~1), gives
$\int\mathsf b^{(d)}_j(h)\pfrwt^-_i(w_{i,p}h)\,d\nu(h)\le\prod_{a<i}\pfsc_a(p_a^{-1})\,C_\Gamma\,\mathsf S_j(\mathsf k^d;x)$. So
$I_2\le K\prod_{a<j}\pfsc_a(p_a^{-1})\,C_\Gamma^2\,\mathsf S_j(\mathsf k^d;x)$.

\emph{Conclusion.} $\int\mathsf a^{(1)}_i\,d\theta_{[0,j)}\le\mathsf S_0(\mathsf k^1;x)$ (unfold the recursion, as in the proof of
Proposition~\ref{pfs:frame}). Hence $\pfrF_j(x;1)^2\le K\bigl(\prod_{a<j}\pfsc_a(p_a^{-1})\bigr)^2C_\Gamma^3\,\mathsf S_0(\mathsf k^1;x)\,\mathsf S_j(\mathsf k^d;x)=
K\,\mathsf Q(x)$.
\end{proof}





\providecommand{\panot}[1]{\par\smallskip{\raggedright\noindent\textbf{[Not proved here:} #1\textbf{]}\par}\smallskip}
\providecommand{\pbdeg}{\operatorname{deg}}
\providecommand{\pbvol}{\operatorname{vol}}
\providecommand{\pbmTh}{\Theta}
\providecommand{\pfrG}{\mathsf G}
\providecommand{\pfrH}{\mathsf H}
\providecommand{\pfrU}{\mathsf U}
\providecommand{\pfrR}{\mathsf R}
\providecommand{\pfrM}{\mathsf M}
\providecommand{\pfrm}{\mathsf m}
\providecommand{\pfrK}{\mathsf K}
\providecommand{\pfrP}{\mathcal P}
\providecommand{\pfrCl}{\mathcal C}
\providecommand{\pfrI}{\mathcal I}
\providecommand{\pfrF}{\mathsf F}
\providecommand{\pfrB}{\mathsf B}
\providecommand{\pfrdep}{\operatorname{dp}}
\providecommand{\pfrkap}{\kappa}
\providecommand{\pfrN}{\mathsf N}
\providecommand{\pfrX}{\Xi}
\providecommand{\pfrip}[2]{\langle#1,#2\rangle}
\providecommand{\pfrell}{\ell}
\providecommand{\prpD}{\mathsf D}
\providecommand{\prpvolU}[1]{\pbvol\auU(#1)}
\providecommand{\prpdens}{\mathsf h}

\section{The real place: the local statements}\label{prp:sec}

This section proves the statements of the real place that \S\ref{pfr:ss-realforms} states without
proof: Lemmas~\ref{pfr:haar}, \ref{pfr:realgauge}, \ref{pfr:deepreal} and
Propositions~\ref{pfr:conv}, \ref{pfr:top}, \ref{pfr:formbound}, \ref{pfr:remote}. We use the notation
of \S\ref{pfr:sec} throughout, and its results with their proofs; of the statements of
\S\ref{pfr:ss-shared} that have no proof there, we use Theorems~\ref{pfr:oneframe}, \ref{pfr:gain-char}
and~\ref{pfr:gain-heads} as stated (they are shared with the finite places and are proved elsewhere).

\emph{What is proved and what is not.}
\begin{itemize}
\item Lemma~\ref{pfr:realgauge} (the gauge axioms (G0)--(G5) and the formula for the depth),
Lemma~\ref{pfr:deepreal}, Propositions~\ref{pfr:conv}, \ref{pfr:top}, \ref{pfr:formbound}
and~\ref{pfr:remote} are proved in full (from the shared theorems just named).
\item Lemma~\ref{pfr:haar} is proved in full; the comparison with Haar measure (Lemma~\ref{prp:fibre})
uses the Gram fibre normalisation {\cite[Lemma~8.80]{P51}} and {\cite[Lemma~8.81]{P51}} of {\cite[\S8.4.4]{P51}}.
\item Consequently Theorem~\ref{pfr:R} depends only on Theorems~\ref{pfr:oneframe}, \ref{pfr:gain-char},
\ref{pfr:gain-heads}, on results proved elsewhere, and on the hypotheses (H0)--(H5)
(\S\ref{prp:ss-R}).
\S\ref{prp:ss-H} says on what each of (H1)--(H5) rests for the Bessel functions of
{\cite[Theorem~3.22]{P51}}.
\end{itemize}

Throughout, $\pfrG=G_\tau$ in the normalised basis, which is orthonormal for the Euclidean norm
$|\cdot|$; $\|\cdot\|$ is the operator norm. The Gram matrix $J$ of the basis is a signed permutation
matrix ($J=J^{-1}=J^*$), so $g^{-1}=Jg^*J$ and $\|g^{-1}\|=\|g\|$ for $g\in\pfrG$; elements of $\pfrH$ are
Euclidean isometries. The place datum at $\tau$ is that of \S\ref{pfr:ss-realforms}, with box
$\pfrK=\{\|g\|\le\rho\}$, close pairs $|\omega_1|^2+|\omega_2|^2<1$, level $M=1$, and $\Gamma_L$ is the
filtration $\Gamma_L(i)=\{c:|c|\le L^i\}$ with loss $L^2$, for $L\ge\rho$. We write
$P_l(p)=p_0\cdots p_{l-1}$. By (H0), $\rho\ge8r\ge8$, $k-5=s_0+3n+1$, and
$r=2\varrho+k+1=s_0+5n+5$.

\subsection{The Haar ratio: proof of Lemma~\ref{pfr:haar}}\label{prp:ss-haar}

\subsubsection*{The Heisenberg measures}
At the real place the Heisenberg group $\pfrN_a$ of floor $a<\varrho$ consists of the elements
$n(v,\beta)$, $v\in V_{a+1}$ (the span of the basis vectors of floor $>a$), $\beta\in i\R$, with
$n(v,\beta)e_a^*=e_a^*+v+(\beta-\tfrac12h(v,v))e_a$, $n(v,\beta)e_a=e_a$, $n(v,\beta)w=w-h(w,v)e_a$ for
$w\in V_{a+1}$, and the identity on the basis vectors of floor $<a$; $\eta_a$ is the image of
Lebesgue measure $dv\,dt$ under $(v,t)\mapsto n(v,it)$. Put $\pfrU_{(<l)}=\pfrN_0\cdots\pfrN_{l-1}$
(so $\pfrU_{(<\varrho)}=\pfrU$, Lemma~\ref{pfr:rearr}) and $\eta_\pfrU=\eta_0*\eta_1*\dots*\eta_{\varrho-1}$,
that is $\int f\,d\eta_\pfrU=\int f(n_0n_1\cdots n_{\varrho-1})\,d\eta_0(n_0)\cdots d\eta_{\varrho-1}(n_{\varrho-1})$.
For $l\le\varrho$ and $g\in\pfrG$ let
\[
G_l(g)=\exp\Bigl(-\pi\sum_{b<l}\Bigl(\sum_{i:\ \text{floor}(i)>b}|g_{i,e_b^*}|^2+(\Iml g_{e_b,e_b^*})^2\Bigr)\Bigr),
\]
the Gaussian in the \emph{free entries} of the columns $e_b^*$, $b<l$, of $g$.

\begin{lemma}\label{prp:heis}
\textup{(a)} For $\nu\in\pfrU_{(<a)}$, $a<\varrho$, and $n=n(v,it)\in\pfrN_a$:
$G_{a+1}(\nu n)=e^{-\pi(|v|^2+t^2)}\,G_a(\nu)$.
\textup{(b)} $\int G_\varrho\,d\eta_\pfrU=1$.
\textup{(c)} The map $(n_0,\dots,n_{\varrho-1})\mapsto n_0\cdots n_{\varrho-1}$ is a homeomorphism of
$\pfrN_0\times\dots\times\pfrN_{\varrho-1}$ onto $\pfrU$, and $\eta_\pfrU$ is a Haar measure of $\pfrU$.
\end{lemma}
\begin{proof}
(a) Every element of $\pfrN_m$, $m<a$, maps a vector $w\in V_a$ to $w$ plus a multiple of $e_m$ (as
$V_a\subseteq V_{m+1}$); hence $\nu w-w\in\langle e_b:b<a\rangle$ for $w\in V_a$. The column $e_a^*$ of $\nu n$ is
$\nu(ne_a^*)$ with $ne_a^*\in V_a$; so its entries in the rows of floor $\ge a$ are those of $ne_a^*$: the
vector $v$ in the rows of floor $>a$, and $\beta-\frac12h(v,v)$, of imaginary part $t$ ($h(v,v)$ is
real), in the row $e_a$. For $b<a$, $n$ fixes $e_b^*$, so the column $e_b^*$ of $\nu n$ is that of $\nu$.

(b) Induction on $l\le\varrho$ for $\int G_l\,d(\eta_0*\dots*\eta_{l-1})=1$; for $l=0$ both sides are
$1$. By (a), Fubini's theorem and $\int_{\C^d\times\R}e^{-\pi(|v|^2+t^2)}dv\,dt=1$, the integral for $l=a+1$
equals the integral for $l=a$ (the measure $\eta_0*\dots*\eta_{a-1}$ is carried by $\pfrU_{(<a)}$).

(c) By the proof of (a) (applied to $\nu=n_0\cdots n_{a-1}$, and since $n_m$, $m>a$, fixes $e_a^*$), the free
entries of the column $e_a^*$ of $u=n_0\cdots n_{\varrho-1}$ are the coordinates $(v_a,t_a)$ of $n_a$. So
the product map is injective, with the continuous inverse $u\mapsto(n(v_a,it_a))_a$ read off the columns
of $u$; it is onto $\pfrU$ by Lemma~\ref{pfr:rearr}. Hence $\eta_\pfrU$ is the image of a Lebesgue measure under
a homeomorphism onto $\pfrU$: it is a nonzero Radon measure on $\pfrU$. It is right invariant under each
$\pfrN_l$: for $n\in\pfrN_l$ write $n_0\cdots n_{\varrho-1}n=n_0\cdots n_l\,n'\,n_{l+1}\cdots n_{\varrho-1}$ with
$n'=gng^{-1}$, $g=n_{l+1}\cdots n_{\varrho-1}\in\pfrG^{(l+1)}$; by (N2) of \S\ref{pfr:ss-floors}, $n'\in\pfrN_l$,
and integrating first in $n_l$ (the measure $\eta_l$ is right invariant under $\pfrN_l$) removes $n'$. As
the $\pfrN_l$ generate $\pfrU$, $\eta_\pfrU$ is a right Haar measure of $\pfrU$; $\pfrU$ is nilpotent, hence
unimodular, so $\eta_\pfrU$ is a Haar measure.
\end{proof}

\emph{Convention on $du$.} We take for the Haar measure $du$ of $\pfrU$ in {\cite[\S3.2.1]{P51}} the measure
$\eta_\pfrU$; equivalently (by Lemma~\ref{prp:heis}(b) and the uniqueness of Haar measure, $\pfrU$ being
unimodular so that $\int G_\varrho(u^{-1})\,du=\int G_\varrho\,du$), $du$ is the Haar measure with
$\int_\pfrU G_\varrho(u^{-1})\,du=1$. Since $u^{-1}=Ju^*J$, $G_\varrho(u^{-1})$ is the standard Gaussian in the
free entries of the rows $e_a$ of $u$ (the complex entries above the diagonal block that are not fixed by
unitarity, and the imaginary parts of the entries $(e_a,e_a^*)$); so $du$ is Lebesgue measure in these
entries ({\cite[Lemma~3.6]{P51}}). This is the measure $du$ fixed in {\cite[\S2.1]{P51}}. The invariant measure on $\pfrU\backslash\pfrG$ is
normalised by $\int_\pfrG f\,dx=\int_{\pfrU\backslash\pfrG}\int_\pfrU f(uy)\,du\,dy$ ($f\ge0$ measurable).

\subsubsection*{The comparison with Haar measure}

\begin{lemma}[Fibre measures of frames]\label{prp:fibre}
For every Borel set $B\subseteq\pfrG$, $\pfrM(\{g:g^{-1}\in B\})\le\dfrac{\prpvolU r}{\prpvolU k}\,\pbvol(B)$,
where $\pbvol$ is the Haar measure $dx$ of {\cite[\S3.2.1]{P51}} and $\prpvolU N$ is the volume of
$\auU(N)$ for the measure compatible with the compact dual \textup{(}Remark~\ref{pfr:rem-match}\textup{)}.
\end{lemma}
\begin{proof}
We use the fibre measures of {\cite[\S8.4.4]{P51}}: for a tuple $E_0$ the relation $\mu\preceq\lambda_{E_0}$,
{\cite[Lemma~8.81(a)]{P51}} (appending vectors) and {\cite[Lemma~8.80]{P51}} (if $\mu\preceq\lambda_1$, $1$ the tuple of
the basis, then $\mu\le\frac{\prpvolU r}{2^r}dg$; the order of the basis vectors in the tuple is irrelevant, a
permutation of the columns preserving Lebesgue measure and $dT$). Write $\nu^{\rm inv}$ for the image of $\nu$
under $g\mapsto g^{-1}$. Since $\pfrM=(\eta_{\varrho+1}*\theta_{\varrho+1})*\dots*(\eta_0*\theta_0)$ with
$\theta_{\varrho+1}=\mu_\pfrH$, $\eta_\varrho=\eta_{\varrho+1}=\delta_1$, and $\eta_a^{\rm inv}=\eta_a$
($n(v,it)^{-1}=n(-v,-it)$), $\mu_\pfrH^{\rm inv}=\mu_\pfrH$,
\[
\pfrM^{\rm inv}=\theta_0^{\rm inv}*\eta_0*\theta_1^{\rm inv}*\eta_1*\dots*\theta_{\varrho-1}^{\rm inv}*\eta_{\varrho-1}*
\theta_\varrho^{\rm inv}*\mu_\pfrH .
\]
Let $E_a=(e_0,e_0^*,\dots,e_{a-1},e_{a-1}^*)$ and $\mu'_a=\theta_0^{\rm inv}*\eta_0*\dots*\theta_{a-1}^{\rm inv}*\eta_{a-1}$.
We show $4^{-a}\mu'_a\preceq\lambda_{E_a}$ by induction ($a=0$: $\delta_1$ and the empty tuple). Write $V_a$ for the
span of the basis vectors of floor $\ge a$; the conditions $E_a^*\mathsf J_1X=0$ say $X\in V_a$, and they are $2a$
orthonormal complex coordinates of $X$, so the fibre measure of $X\mapsto(E_a^*\mathsf J_1X,\Phi(X))$ over
$(0,t)$ is the disintegration of Lebesgue measure of $V_a$ along $h(X,X)$ at $t$.

\emph{Appending $e_a$.} $\nu=\frac12\theta_a^{\rm inv}$ is carried by $\pfrG^{(a)}$, which fixes $E_a$; the image of
$\theta_a^{\rm inv}$ under $u\mapsto ue_a$ is the image of $\theta_a$ under $\sigma\mapsto\sigma^{-1}e_a$, which is
$c_a\mu_a=2\mu_a$ with $\mu_a$ the quadric measure of $e_a$ in $V_a$ (Lemma~\ref{pfr:floorquad}, $\kappa_\C=2$), and $\mu_a$
is the fibre measure just described at $t=0$. So the image of $\nu$ is that fibre measure, and by
{\cite[Lemma~8.81(a)]{P51}}, $4^{-a}\cdot\frac12\,\mu'_a*\theta_a^{\rm inv}\preceq\lambda_{(E_a,e_a)}$.

\emph{Appending $e_a^*$.} $\nu=\frac12\eta_a$ is carried by $\pfrN_a$, which fixes $E_a$ and $e_a$. The tuples $X$
with $h(X,e_b)=h(X,e_b^*)=0$ ($b<a$), $h(X,e_a)=1$ are $X=e_a^*+w+\mu e_a$, $w\in V_{a+1}$, and
$\Phi(X)=h(w,w)+2\Rel\mu$; in the coordinates $(w,\Rel\mu,\Iml\mu)$ Lebesgue measure disintegrates along $\Phi$ with
fibre measure $\frac12\,dw\,d\Iml\mu$ (as $d\Phi=2\,d\Rel\mu$). The image of $\eta_a$ under $n\mapsto ne_a^*=
e_a^*+v+(it-\frac12h(v,v))e_a$ is $dv\,dt$ in these coordinates. So the image of $\frac12\eta_a$ is the fibre measure,
and {\cite[Lemma~8.81(a)]{P51}} gives $4^{-a-1}\mu'_{a+1}\preceq\lambda_{E_{a+1}}$.

\emph{Appending $z$, then the $x_i$.} $\theta_\varrho^{\rm inv}$ is carried by $\pfrG^{(\varrho)}$, fixing $E_\varrho$, and
its image under $u\mapsto uz$ is the image of $\theta_\varrho$ under $\sigma\mapsto\sigma^{-1}z$, which is $c_\varrho\mu_\varrho=
2\mu_\varrho$ with $\mu_\varrho$ the quadric measure of $z$ in $V_\varrho$ (at $h=-1$; \S\ref{pfr:ss-realforms}), the fibre
measure for appending $z$; so $4^{-\varrho}\cdot\frac12\,\mu'_\varrho*\theta_\varrho^{\rm inv}\preceq\lambda_{(E_\varrho,z)}$. Finally,
as in the proof of {\cite[Lemma~8.81(b)]{P51}}, the image of $\frac{\prpvolU k}{2^k}\mu_\pfrH$ under $h\mapsto(hx_i)_i$ is the fibre
measure of $X\mapsto X^*X$ on $k$-tuples of $X_0$ over $1$ ({\cite[Lemma~8.80]{P51}} for $\auU(k)$), and the tuples
orthogonal to $(E_\varrho,z)$ are the tuples of $X_0$, in orthonormal coordinates. Hence
$2^{-2\varrho-1}\frac{\prpvolU k}{2^k}\pfrM^{\rm inv}\preceq\lambda_1$, and by {\cite[Lemma~8.80]{P51}}
\[
\pfrM^{\rm inv}\le2^{2\varrho+1+k}\frac{\prpvolU r}{2^r\prpvolU k}\,dg=\frac{\prpvolU r}{\prpvolU k}\,dg,
\]
since $r=k+2\varrho+1$. Here $dg$ is the Haar measure of \textup{(DS1)(iii)}, the measure $dx$ of
{\cite[\S3.2.1]{P51}}.
\end{proof}

\begin{proof}[Proof of Lemma~\ref{pfr:haar}]
\emph{$\pfrM$ is right invariant.} This is Lemma~\ref{pfr:upperinv} for $j=0$: $\rho_0=\pfrM$ and
$\pfrM*\delta_t=\pfrM$ for $t\in\pfrG^{(0)}=\pfrG$. (Its proof uses only the structure of the floors, not that
$\pfrM$ is a Haar measure.)

\emph{$\pfrM=\lambda\,dx$ with $0<\lambda<\infty$.} $\pfrM\neq0$, since each factor $\eta_a$, $\theta_a$ is a
nonzero measure carried by a group. By Lemma~\ref{prp:fibre} applied to $B^{-1}$, $\pfrM(B)\le C\pbvol(B^{-1})
=C\pbvol(B)$ with $C=\prpvolU r/\prpvolU k$ ($\pfrG$ is unimodular, so $dx$ is invariant under inversion);
hence $\pfrM$ is finite on compact sets, and, as $\pfrG$ is second countable and locally compact, it is a
right invariant Radon measure. By the uniqueness of right Haar measure, $\pfrM=\lambda\,dx$ (the right Haar
measures of the unimodular group $\pfrG$ are the multiples of $dx$), with $\lambda>0$ as $\pfrM\ne0$ and
$\lambda\le C<\infty$ by the inequality just written for a compact $B$ with $0<\pbvol(B)$. In particular
$\lambda\,\prpvolU k/\prpvolU r\le1$.

\emph{The quotient formula.} By Lemma~\ref{pfr:rearr} for $j=\varrho+2$ (the factors $\eta_l$, $l\ge\varrho$, are
$\delta_1$), $\int f\,d\pfrM=\iint f(uy)\,d\eta_\pfrU(u)\,d\theta_{\ge0}(y)$ for $f\ge0$ measurable, and
$\eta_\pfrU=du$ by our convention. Let $b(g)=e^{-\|g\|^2}$. For $y\in\pfrG$, $\|uy\|\ge\|u\|/\|y^{-1}\|$, and
$\|u\|$ is at least the modulus of each coordinate $v_a,t_a$ of the chart of Lemma~\ref{prp:heis}(c) (these
are entries of $u$); so $B(y):=\int_\pfrU b(uy)\,du$ is finite, positive, measurable in $y$, and
$B(u'y)=B(y)$ for $u'\in\pfrU$. Put $\beta=b/B$, so that $\int_\pfrU\beta(uy)\,du=1$ for all $y$. Let $\phi\ge0$ be measurable and
left invariant under $\pfrU$. Then
\[
\int\phi\,d\theta_{\ge0}=\iint\beta(uy)\phi(uy)\,du\,d\theta_{\ge0}(y)=\int\beta\phi\,d\pfrM=\lambda\int_\pfrG\beta\phi\,dx
=\lambda\int_{\pfrU\backslash\pfrG}\phi(y)\int_\pfrU\beta(uy)\,du\,dy=\lambda\int_{\pfrU\backslash\pfrG}\phi.\qedhere
\]
\end{proof}

\subsection{The gauge at $\C$: proof of Lemma~\ref{pfr:realgauge}}\label{prp:ss-gauge}

Let $L\ge\rho$. Then $L\ge8r\ge8$ and $L\ge r+1$. The space $X'$ is spanned by the $x_i$ other than $x_{i_0}$
and two blocks of two of them, so $\operatorname{rank}X'=k-5=s_0+3n+1$, $m=0$, and the volume exponent is
$s=\operatorname{rank}X'-(3n+1)=s_0$; by (H0), $s_0=k-3n-6\ge n^7-3n-6\ge10^{10}n^4$ (as $n>10^8$).

\begin{lemma}[A disc of products]\label{prp:disc}
Let $g,y\in\pfrG$ and let $v$ be a family of at most $n$ linear forms on $X'$. There are orthonormal
vectors $w_1,\dots,w_{s+1}$ of $X'$, killed by $v$, such that $|g^{-1}w|=|y^{-1}w|=|w|$ for every $w$ in their
span $O$. For $u\in X'$ put $q(u)=\sum_i|h(w_i,u)|^2$ \textup{(}on $X'$, $h(w,u)=u^*w$ in the coordinates of
the orthonormal basis, by (Cv1)\textup{)}. Then every $t\in\C$ with $|t|\le\sqrt{q(u)}/2$ is $h(w,u)$ for some
$w\in O$ with $|w|\le1/2$.
\end{lemma}
\begin{proof}
The vectors $w$ with $Jw=w$ are those with $w_{e_a}=w_{e_a^*}$ for all $a$ and $w_z=0$, the kernel of
$\varrho+1$ linear forms $N$; for them $h(w,w)=w^*Jw=|w|^2$. Let $C$ be the family of linear forms on $X'$
made of $v$, $N\circ g^{-1}$ and $N\circ y^{-1}$: at most $n+2(\varrho+1)\le3n$ forms. Its kernel has
dimension at least $\operatorname{rank}X'-3n=s+1$; let $w_1,\dots,w_{s+1}$ be orthonormal in it. For $w\in O$,
$vw=0$, and $g^{-1}w$ is fixed by $J$, so $|g^{-1}w|^2=h(g^{-1}w,g^{-1}w)=h(w,w)=|w|^2$ ($w\in X'$, on which $h$
is the Euclidean form); likewise for $y^{-1}$. Given $t$ with $|t|\le\sqrt{q}/2$, $q=q(u)>0$, put
$w=(t/q)\sum_ih(u,w_i)\,w_i$. Then, $h$ being linear in the first variable, $h(w,u)=(t/q)\sum_ih(u,w_i)h(w_i,u)=
(t/q)\sum_i|h(w_i,u)|^2=t$, and $|w|^2=|t|^2\sum_i|h(u,w_i)|^2/q^2=|t|^2/q\le\frac14$. (If
$q=0$ only $t=0$ occurs, with $w=0$.)
\end{proof}

\begin{proof}[Proof of Lemma~\ref{pfr:realgauge}]
\emph{(G0).} $|cc'|\le L^{i+i'}$; a sum of at most $r+1$ numbers of modulus $\le L^i$ has modulus
$\le(r+1)L^i\le L^{i+1}$; $\Gamma_L(i)$ is stable under $c\mapsto-c,\bar c$ and contains $0,1$; $J=J^{-1}$ has
entries $0,\pm1$. The consequences listed in (G0) follow from these rules, as stated there.

\emph{(G1).} $m=0\le9n$; $P=\pi^2\ge1$; $M=1$; at $\C$ the fixed element of trace one is $\tau_0=\frac12$ and
the skew unit is $\delta=\frac i2$, so $\tau_0\delta^{-1}=-i$; the basis vectors have entries of modulus at
most $2\le L$; an element of $\pfrK$ has entries of modulus at most $\|g\|\le\rho\le L$.

\emph{(G2).} If $|\omega_1|^2+|\omega_2|^2<1$, the entries of $\omega_1$ have modulus $\le1$ and those of
$\omega_1-\omega_2$ modulus $\le2\le L$. Given $i$, let $c=L^{-(i+2)}$, so $c^{-1}\in\Gamma_L(i+2)$; if
$\omega_1,\omega_2$ have entries of modulus $\le L^i$, then $|c\omega_j|^2\le rL^{2i}L^{-2i-4}$, and
$|c\omega_1|^2+|c\omega_2|^2\le2r/L^4<1$.

\emph{(G3).} For $|u|\le L^i$, $\operatorname{mod}(u)=|u|^2\le L^{2i}=\Lambda^i$; $\pbvol\Gamma_L(i)=\pi L^{2i}\le
L^{2i+2}=\Lambda^{i+1}$.

\emph{(G4).} Here $2m+2=2$, and the standard form on the block is $|y_1|^2+|y_2|^2$. In polar coordinates,
Lebesgue measure on $\C^2$ is $\int_0^\infty\kappa_t\,dt$, where $\kappa_t$ is the rotation-invariant
measure on the sphere $|y|^2=t$ of total mass $\pi^2t$ (the derivative of $\pbvol\{|y|^2\le t\}=\pi^2t^2/2$);
put $\kappa_t=0$ for $t\le0$. If $\kappa_t$ charges $\Gamma_L(i)^2$ then $t\le2L^{2i}$, so
$\kappa_t(\Gamma_L(i)^2)\le2\pi^2L^{2i}\le L^4L^{4i}=\Lambda^{2(i+1)}$.

\emph{(G5).} Let $g\in\pfrK$, $x\in\pfrG$, $v$ a family of $\varrho$ linear forms on $X'$, and $O$, $q$ as in
Lemma~\ref{prp:disc} for $g$ and $y=gx$. Put $\phi(u)=36/q(u)$ ($=\infty$ if $q(u)=0$); $\phi$ is measurable.
Each $w\in O$ with $|w|\le\frac12$ has $vw=0$ and $|g^{-1}w|^2+|(gx)^{-1}w|^2=2|w|^2<1$, so it is in
$\mathcal W$. Hence, by Lemma~\ref{prp:disc}, $\{h(\omega,u):\omega\in\mathcal W\}$ contains the disc of radius
$\sqrt q/2$, of area $\pi q/4\ge q/36$: the first condition. For the second, a product of two close pairs
satisfies $|h(\omega_1,\omega_1')|\le|\omega_1||\omega_1'|\le\frac12(|\omega_1|^2+|\omega_1'|^2)$, so (as $M=1$) every
element $h(\omega_1,\omega_1')-h(\omega_2,\omega_2')$ of $\mathcal D$ has modulus
$<\frac12(|\omega_1|^2+|\omega_2|^2)+\frac12(|\omega'_1|^2+|\omega'_2|^2)<1$. If $q>0$, take $\omega\in\mathcal W$ with
$h(\omega,u)=\sqrt q/2$; if $h(\omega,cu)=\bar c\,h(\omega,u)\in\mathcal D$ then $|c|\sqrt q/2<1$. So the set of these $c$ lies in
the disc of radius $2/\sqrt q$, of area $4\pi/q\le36/q$.

For the integral, complete $(w_i)$ to an orthonormal basis of $X'$ and let $Q$ be the unitary change of
coordinates; $u\mapsto Q^*u$ preserves Lebesgue measure, $q(u)=|y'|^2$ where $y'\in\C^{s+1}$ are the first
coordinates of $Q^*u$, and the other $3n$ coordinates $y''$ form the rest. If the entries of $u$ have
modulus $\le L^i$, those of $Q^*u$ have modulus $\le(\operatorname{rank}X')L^i\le L^{i+1}$; so $|y'|^2\le R:=
L\cdot L^{2i+2}$ and $y''$ lies in a polydisc of volume $(\pi L^{2i+2})^{3n}$. The law of $|y'|^2$ under
Lebesgue measure on $\C^{s+1}$ has density $\pi^{s+1}t^s/s!$ on $(0,\infty)$; hence
\[
\int_{\Gamma_L(i)^{X'}}\phi^s\le\int_0^R\Bigl(\frac{36}t\Bigr)^s\frac{\pi^{s+1}t^s}{s!}\,dt\cdot(\pi
L^{2i+2})^{3n}=\pi\frac{(36\pi)^s}{s!}R\,(\pi L^{2i+2})^{3n}\le4e^{144}R\,(\pi L^{2i+2})^{3n},
\]
using $\pi\le4$ and $(144)^s/s!\le e^{144}$. As $4,e,\pi\le L$, this is at most
$L^{1+144+(2i+3)+(2i+3)3n}=L^{145+(2i+3)(3n+1)}\le L^{2\cdot99n(i+1)}=\Lambda^{99n(i+1)}$, since
$198n(i+1)-145-(2i+3)(3n+1)=192ni-2i+189n-148\ge0$ for $n\ge1$.

\emph{The depth.} Let $a<\varrho$ and $D=\inf(|g^{-1}\zeta|^2+|(gx)^{-1}\zeta'|^2)$ over the $\zeta\perp e_b$
($b<a$) with $h(\zeta,e_a)=1$ and the $\zeta'\equiv\zeta$ modulo $\langle e_b:b<a\rangle$, that is, the
infimum of $|\zeta|_g^2+|\zeta|_{gx}^2$ over the duals of $e_a$ in $Y^{(a)}$ for the induced lengths. We
show $\mathfrak d_a(g,x)=\{c:|c|^2D<1\}$. The $e_a^*$-coordinate of $\omega$ is $h(\omega,e_a)$. If $c=0$,
take $\omega=0$. If $c\ne0$ and $c=h(\omega,e_a)$ with $\omega$ stable, $(g^{-1}\omega,(gx)^{-1}\omega')$ close, then
$\zeta=c^{-1}\omega$, $\zeta'=c^{-1}\omega'$ give $D\le|c|^{-2}(|g^{-1}\omega|^2+|(gx)^{-1}\omega'|^2)<|c|^{-2}$.
Conversely, if $|c|^2D<1$ choose $\zeta,\zeta'$ with $|c|^2(|g^{-1}\zeta|^2+|(gx)^{-1}\zeta'|^2)<1$; then
$\omega=c\zeta$ is stable with $h(\omega,e_a)=c$. So $\mathfrak d_a$ is the open disc of radius $D^{-1/2}$ (all of
$\C$ if $D=0$), of area $\pi/D$, and $\pfrdep_a(g,x)=D/\pi$.
\end{proof}

\begin{lemma}[The place axiom (P) at $\C$]\label{prp:planch}
The real place $(\C,\ dc\text{ Lebesgue},\ \psi(c)=e^{c-\bar c},\ \operatorname{mod}(c)=|c|^2)$ satisfies the axiom
\textup{(P)} of \S\ref{pfs:ss-planch} with $P=\pi^2$, with equality $\int_\C|\widehat G|=\pi^2$ for $G(v)=e^{-\pi|v|^2}$.
\end{lemma}
\begin{proof}
$G$ is continuous, integrable, $G(0)=1$. For $v=v_1+iv_2$, $z=z_1+iz_2$: $\bar vz-v\bar z=2i\Iml(\bar vz)=
2i(v_1z_2-v_2z_1)$, so $\widehat G(z)=\int_{\R^2}e^{-\pi|v|^2}e^{i\,k\cdot v}dv$ with $k=2(z_2,-z_1)$, $|k|=2|z|$; the
Gaussian integral gives $\widehat G(z)=e^{-|k|^2/(4\pi)}=e^{-|z|^2/\pi}>0$, and $\int_\C e^{-|z|^2/\pi}dz=\pi\cdot\pi=\pi^2$.
Moreover $\operatorname{mod}(\bar c)=|c|^2=\operatorname{mod}(c)$, $\operatorname{mod}$ is multiplicative, $d(cv)=|c|^2dv$
(the real Jacobian of $v\mapsto cv$ on $\R^2$), and the non-units $\{0\}$ form a null set.
\end{proof}

\emph{Conventions.} The convention of \S\ref{pfs:sec} ($h$ linear in its first and conjugate linear in its second
variable; $\pfrip vw=\sum\bar v_iw_i$) is the one used in this section: the $e_a^*$-coordinate of $\omega$ is
$h(\omega,e_a)$ and the $e_a$-coordinate of $v\in V_a$ is $h(v,e_a^*)$, both linear in $\omega$, $v$; the products
$h(w,u)=u^*w$ ($u,w\in X'$) of Lemma~\ref{prp:disc} and of (G5) are linear in $w$ and conjugate linear in $u$, and the
verification of (G5) above is written in this convention; the other uses of $h$ (isometries,
$h(\zeta,u^{-1}e_a)$, $s_w=h(w,v^\#)$ in Lemma~\ref{prp:unip}) hold in either convention. Only moduli of the
leading minors are used (Lemmas~\ref{prp:minor}, \ref{prp:headsupport}), which do not depend on it.

\subsection{Minors, supports, and four auxiliary bounds}\label{prp:ss-aux}

For $\gamma\in\pfrG$ let $\delta_l(\gamma)$ be the leading $l\times l$ minor of the frame matrix
$Z(\gamma)=(h(\gamma f_b,f_a))_{a,b}$.

\begin{lemma}[Minors of near points]\label{prp:minor}
Let $l\le n$, $\gamma,y\in\pfrG$, and suppose $y$ and $\gamma^{-1}y$ are near at $l$. Then
$|\delta_j(\gamma)|\le\rho^{2j}\le\rho^{9n}$ for $j\le l$.
\end{lemma}
\begin{proof}
Let $\xi y=k_1$, $\xi'\gamma^{-1}y=k_2$ with $\xi,\xi'\in\pfrR^{(l)}$ and $k_1,k_2\in\pfrK$. Then
$\gamma=\xi^{-1}k\,\xi'$ with $k=k_1k_2^{-1}$, $\|k\|\le\rho^2$. For $b<l$, $\xi'f_b\in f_b+\langle f_c:c<b\rangle$
and $\xi f_a\in f_a+\langle f_c:c<a\rangle$ (also for the inverses), so
$h(\gamma f_b,f_a)=h(k\xi'f_b,\xi f_a)$ and the leading $j\times j$ block of $Z(\gamma)$ is $\mathsf L^*Z_j(k)\mathsf R$
with unitriangular $\mathsf L,\mathsf R$, where $Z_j(k)$ is the leading block of $Z(k)$. Its entries are
entries of $Jk$, a $j\times j$ submatrix of $Jk$, of operator norm $\le\rho^2$; so
$|\delta_j(\gamma)|=|\det Z_j(k)|\le\rho^{2j}$.
\end{proof}

\begin{lemma}[Heads of near points]\label{prp:headsupport}
Let $j\le\varrho$, $W$ a cut Bessel function for the place datum at $\tau$, and $p,q$ sequences of units.
If $\pfrF^+_j(x)\ne0$ for some $x$, then $|P_l(p)|\le\rho^{9n}$ and $|P_l(q)|\le\rho^{9n}$ for all $l\le j$.
\end{lemma}
\begin{proof}
Let $w=w_{j,p}$, $d=w_{j,q}^{-1}w$, and $g$ with $\pfrB^1_j(g,gx)\pfrB^d_j(w^{-1}g,w^{-1}gx)\ne0$. As $\theta_j$
is carried by $\pfrG^{(j)}$ and $\theta_{\ge j+1}$ by $\pfrG^{(j+1)}$, the first factor gives $\varsigma\in\pfrG^{(j)}$
with $W(\varsigma g)\ne0\ne W(\varsigma gx)$, and the second gives $\tau\in\pfrG^{(j)}$ with $W(\tau w^{-1}g)\ne0$
and $W(d\tau w^{-1}gx)=W(\tau w_{j,q}^{-1}gx)\ne0$ ($d$ commutes with $\pfrG^{(j)}$). A point $y$ with $W(y)\ne0$
is near at every $l$ (cut condition, and $\pfrU\subseteq\pfrR^{(l)}$), and nearness at $j$ is invariant under left
translation by $\pfrG^{(j)}\subseteq\pfrR^{(j)}$. Hence $y_1=\varsigma g$ and $w^{-1}y_1=\varsigma\tau^{-1}(\tau
w^{-1}g)$ are near at $j$; and $y_2=\varsigma gx$ and $w_{j,q}^{-1}y_2$ are near at $j$. The frame matrix of
$w_{j,p}$ has leading block $\operatorname{diag}(p_0,\dots,p_{j-1})$ up to complex conjugation of the
entries ($w_pe_b=p_be_b^*$ and $h(e_b^*,e_a)=\delta_{ab}$), so $|\delta_l(w_{j,p})|=|P_l(p)|$, and
Lemma~\ref{prp:minor} applies to $(w_{j,p},y_1)$ and $(w_{j,q},y_2)$.
\end{proof}

\begin{lemma}[Heads for one gauge]\label{prp:exgauge}
Let $j\le\varrho$, $K\ge0$ and $p,q$ sequences of units. Let $L=\max\bigl(\rho,\allowbreak \max_{l\le j}|P_l(p)|^{-1},
\allowbreak \max_{l\le j}|P_l(q)|^{-1}\bigr)$. Then $L\ge\rho$, $p$ and $q$ are heads of length $j$ for $\Gamma_L$, and
$L^K\le\rho^K\prod_{l=1}^j\max(1,|P_l(p)|^{-1})^K\max(1,|P_l(q)|^{-1})^K$. In particular, if
$|P_l(p)|,|P_l(q)|\le b^{60n^2}$ for $l\le j$, then $L^K\le\rho^K\Lambda^\infty_j(p)\Lambda^\infty_j(q)$ for the
weights of Theorem~\ref{pfr:R} with exponent $K$ in place of $n^6$.
\end{lemma}
\begin{proof}
A head of length $j$ for $\Gamma_L$ means $|P_l(p)|^{-1}\le L$ for $l\le j$. The bound holds since the maximum of
numbers $\ge1$ is at most their product, $\rho\ge1$ and $P_0=1$.
\end{proof}

\begin{lemma}[Schur's bound on the quotient]\label{prp:schur}
Let $W$ be a measurable Bessel function of irreducible type $\mathbb W$ \textup{(}as in
\S\ref{pfr:ss-realdata}\textup{)} with $\int_{\pfrU\backslash\pfrG}\|W\|^2\le\#I$. Then
$\int_{\pfrU\backslash\pfrG}|\pfrip v{W(y)}|^2dy\le|v|^2$ for every $v\in\mathbb W$.
\end{lemma}
\begin{proof}
The function $y\mapsto W(y)W(y)^*$ is left invariant under $\pfrU$ (the character factors cancel), and
integrable on $\pfrU\backslash\pfrG$; let $A=\int_{\pfrU\backslash\pfrG}W(y)W(y)^*dy$, a positive endomorphism of
$\mathbb W$ with $\operatorname{tr}A=\int\|W\|^2\le\#I$. For $h\in\pfrH$, $W(hy)=\pi(h)W(y)$ with $\pi$ the
representation, and $hy$ runs over $\pfrU\backslash\pfrG$ with the invariant measure when $y$ does: $h$
normalises $\pfrU$, conjugation by $h$ preserves $du$ (the modulus of the compact group $\pfrH$ acting on
$\pfrU$ is $1$), and $dx$ is left invariant. Hence $\pi(h)A\pi(h)^*=A$, and by Schur's lemma $A=c\cdot1$ with
$c\,\#I=\operatorname{tr}A\le\#I$. So $\int|\pfrip v{W}|^2=\pfrip v{Av}=c|v|^2\le|v|^2$.
\end{proof}

\begin{lemma}[Two bounds for the cut Bessel function]\label{prp:normW}
Under \textup{(H0)}--\textup{(H1)}, with $\lambda$ the ratio of Lemma~\ref{pfr:haar}:
\textup{(a)} $\mathsf N(W^\flat)\le\pi^{2(k-5)\varrho}\lambda^2\,\#I$;
\textup{(b)} for $j\le\varrho$ and $L\ge\rho$, the mass of floor $j$ is at most
$V:=(\pi\rho^{24})^k\rho^{9999n}$, and $V\ge1$.
\end{lemma}
\begin{proof}
(a) At $\C$, $M=1$, $P=\pi^2$, $\operatorname{rank}X'=k-5$, so $\mathsf N(W^\flat)=\pi^{2\varrho(k-5)}\iint
|\pfrip{W^\flat(g)}{W^\flat(y)}|^2d\theta_{\ge0}(g)\,d\theta_{\ge0}(y)$. For fixed $g$ the function
$y\mapsto|\pfrip{W^\flat(g)}{W^\flat(y)}|^2$ is left invariant under $\pfrU$, so by Lemma~\ref{pfr:haar} and
Lemma~\ref{prp:schur} (with (H1)) its $\theta_{\ge0}$-integral is $\lambda\int_{\pfrU\backslash\pfrG}\le\lambda
\|W^\flat(g)\|^2$; and $\int\|W^\flat\|^2d\theta_{\ge0}=\lambda\int_{\pfrU\backslash\pfrG}\|W^\flat\|^2\le\lambda\,\#I$.

(b) Lemmas~\ref{pfr:chart} and~\ref{pfr:quadricX} hold for every $a\le\varrho$, with the constant $c_a=2$ of the
real place (\S\ref{pfr:ss-realforms}). Apply them, for the gauge $\Gamma_\rho$ (loss $\Lambda=\rho^2$), with $c\equiv1$
(invariant), $\Psi_g\equiv1$ and $\phi\equiv1$, for which $\int_{\Gamma_\rho(11)^{X'}}1=(\pi\rho^{22})^{k-5}\le
\Lambda^{12(k-5)}$: the mass of floor $j$ is at most $2\Lambda^{99n\cdot13+12(k-5)}=2\rho^{2574n}\rho^{24(k-5)}\le V$
(as $2\le\pi^k$ and $\rho\ge1$). The mass does not depend on the gauge. $V\ge1$ is clear.
\end{proof}

\subsection{Far means deep: proof of Lemma~\ref{pfr:deepreal}}\label{prp:ss-deep}

For $a\le\varrho$ let $V_a$ be the span of the basis vectors of floor $\ge a$ (so $V_0=\C^r$, and $V_\varrho$ is the
span of $z$ and the $x_i$), $E_a=\langle e_b:b<a\rangle$, and $\pi_a$ the coordinate projection onto $V_a$ along
$E_a\oplus\langle e_b^*:b<a\rangle$; on $E_a^\perp=V_a\oplus E_a$ it is the projection along $E_a$, and
$|\pi_aw|\le|w|$.

\begin{lemma}[Bounded duals bound a unipotent element]\label{prp:unip}
Let $m\in\pfrU$, $n=m^{-1}$, and $R\ge1$. Suppose that for every $a<\varrho$ there are vectors
$\alpha_a,\beta_a\perp E_a$ with $h(\alpha_a,e_a)=h(\beta_a,e_a)=1$, $|\alpha_a|,|\beta_a|\le R$,
$m\alpha_a\equiv\beta_a$ and $n\beta_a\equiv\alpha_a$ modulo $E_a$. Then $\|m\|,\|n\|\le(10R^2)^\varrho$.
\end{lemma}
\begin{proof}
An element of $\pfrU$ preserves each $E_a$, hence $E_a^\perp$; so $\bar m_a=\pi_a\circ m|_{V_a}$ is well defined on
$V_a$, and $\pi_a(mw)=\bar m_a(\pi_aw)$ for $w\in E_a^\perp$; likewise $\bar n_a$, and $\bar n_a=\bar m_a^{-1}$. As $E_a$
is isotropic and orthogonal to $V_a$, $\bar m_a$ preserves $h$ on $V_a$. Moreover: $\bar m_a e_a=e_a$
($me_a\in e_a+E_a$); $\bar m_\varrho=1$ (the middle block of an element of $\pfrU$ is $1$); and for $w\in V_{a+1}$,
$\pi_{a+1}(\bar m_aw)=\bar m_{a+1}w$, so $\bar m_aw=\bar m_{a+1}w+s_w e_a$ with $s_w=h(\bar m_aw,e_a^*)$. Let
$D_a=\max(\|\bar m_a-1\|,\|\bar n_a-1\|)$ (operator norms on $V_a$); $D_\varrho=0$. We show
$1+D_a\le10R^2(1+D_{a+1})$; then $\|m\|=\|\bar m_0\|\le1+D_0\le(10R^2)^\varrho$, and likewise for $n$.

Write $\pi_a\alpha_a=e_a^*+w_\alpha+t_\alpha e_a$ and $\pi_a\beta_a=e_a^*+w_\beta+t_\beta e_a$ with
$w_\alpha,w_\beta\in V_{a+1}$ (the $e_a^*$-coordinate is $h(\cdot,e_a)=1$); all these pieces have norm $\le R$. From
$m\alpha_a\equiv\beta_a$: $\bar m_a(\pi_a\alpha_a)=\pi_a\beta_a$, and similarly $\bar n_a(\pi_a\beta_a)=\pi_a\alpha_a$.
Write $\bar m_ae_a^*=e_a^*+w^\#+s^\#e_a$ and $\bar n_ae_a^*=e_a^*+v^\#+t^\#e_a$ ($w^\#,v^\#\in V_{a+1}$; the
$e_a^*$-coordinate is $h(\bar m_ae_a^*,e_a)=h(e_a^*,\bar n_ae_a)=1$).

\emph{The coefficients $s_w$.} For $w\in V_{a+1}$, $s_w=h(\bar m_aw,e_a^*)=h(w,\bar n_ae_a^*)=h(w,v^\#)$, so
$|s_w|\le|w|\,|v^\#|$. Comparing the $V_{a+1}$-components of $\bar n_a(\pi_a\beta_a)=\pi_a\alpha_a$ gives
$v^\#+\bar n_{a+1}w_\beta=w_\alpha$, so $v^\#=(w_\alpha-w_\beta)-(\bar n_{a+1}-1)w_\beta$ and
$|v^\#|\le(2+D_{a+1})R$. Symmetrically $|w^\#|\le(2+D_{a+1})R$, and the coefficient of $e_a$ in $\bar n_aw$ has
modulus $\le|w||w^\#|$.

\emph{The column $e_a^*$.} Comparing the $e_a$-components of $\bar m_a(\pi_a\alpha_a)=\pi_a\beta_a$:
$s^\#+s_{w_\alpha}+t_\alpha=t_\beta$, so $|s^\#|\le2R+R(2+D_{a+1})R$. So
$|(\bar m_a-1)e_a^*|\le|w^\#|+|s^\#|\le(2+D_{a+1})(R+R^2)+2R$.

\emph{Conclusion.} On $V_a=\C e_a\oplus V_{a+1}\oplus\C e_a^*$ (orthogonal for the Euclidean norm),
$(\bar m_a-1)e_a=0$, $|(\bar m_a-1)w|\le|(\bar m_{a+1}-1)w|+|s_w|\le(D_{a+1}+(2+D_{a+1})R)|w|$, and the bound on
$e_a^*$ above; so $\|\bar m_a-1\|$ is at most the sum of the two coefficients, $\le(1+D_{a+1})(1+6R+2R^2)$. The
same holds for $\bar n_a$, and $1+D_a\le(1+D_{a+1})(2+6R+2R^2)\le10R^2(1+D_{a+1})$.
\end{proof}

\begin{proof}[Proof of Lemma~\ref{pfr:deepreal}]
Let $y$ be active at $x$: $\pfrip{W^\flat(y)}{W^\flat(yx)}\ne0$, and $w_p^{-1}y$, $w_q^{-1}yx$ are near at $j$.
Suppose that $\pfrdep_a(y,x)<\tau'^{-2}$ for every $a$ with $j\le a<\varrho$; we show that all entries of $x-1$
have modulus $\le T^{99n}$, which proves the lemma. By the cut condition (H1) there are $u,v\in\pfrU$ with
$k_1=uy\in\pfrK$ and $k_2=vyx\in\pfrK$. Put $m=vu^{-1}\in\pfrU$, $n=m^{-1}$. Then $x=y^{-1}(yx)=k_1^{-1}nk_2$, so
$\|x\|\le\rho^2\|n\|$.

\emph{Duals for $a\ge j$.} By Lemma~\ref{pfr:realgauge}, the infimum $D$ defining the depth is $\pi\,\pfrdep_a(y,x)<\pi\tau'^{-2}\le\pi T^2/4<T^2$
(as $\tau'\ge2/T$); this gives $\zeta\perp E_a$ with $h(\zeta,e_a)=1$ and $\zeta'\equiv\zeta$ modulo $E_a$ with
$|y^{-1}\zeta|^2+|(yx)^{-1}\zeta'|^2<T^2$.

\emph{Duals for $a<j$.} Take $\zeta=\zeta'=e_a^*$. By Lemma~\ref{pfr:headentries} (for $\Gamma_L$, the head $p$,
the point $y$, $s=u\in\pfrU\subseteq\pfrR^{(j)}$ and the element $\eta\in\pfrR^{(j)}$ given by the nearness of
$w_p^{-1}y$), the column $ue_a^*$ has entries of modulus $\le L^{10n}$, so $|y^{-1}e_a^*|=|k_1^{-1}ue_a^*|\le
\rho\sqrt rL^{10n}\le L^{10n+2}$; likewise, with $q$, $yx$ and $v$, $|(yx)^{-1}e_a^*|\le L^{10n+2}\le T$.

In both cases put $\alpha_a=u\zeta$ and $\beta_a=v\zeta'$. They are orthogonal to $E_a$ (as $u,v$ preserve $E_a$
and $E_a^\perp$), $h(\alpha_a,e_a)=h(\zeta,u^{-1}e_a)=h(\zeta,e_a)=1$ and likewise $h(\beta_a,e_a)=1$;
$|\alpha_a|=|k_1y^{-1}\zeta|\le\rho T$ and $|\beta_a|=|k_2(yx)^{-1}\zeta'|\le\rho T$; and modulo $E_a$,
$m\alpha_a=v\zeta\equiv v\zeta'=\beta_a$ and $n\beta_a=u\zeta'\equiv u\zeta=\alpha_a$. By Lemma~\ref{prp:unip} with
$R=\rho T$, $\|n\|\le(10\rho^2T^2)^\varrho$. Hence every entry of $x-1$ has modulus at most
$1+\rho^2(10\rho^2T^2)^\varrho\le2\cdot10^\varrho\rho^{2\varrho+2}T^{2\varrho}\le T^{1+\varrho+(2\varrho+2)+2\varrho}=T^{5\varrho+3}\le
T^{99n}$, using $2,10,\rho\le L\le T$.
\end{proof}

\begin{lemma}[Dyadic shells]\label{prp:dyadic}
Let $p,q$ be heads of length $j\le\varrho$ for $\Gamma_L$, $L\ge\rho$, and $W=W^\flat$. For $i\ge0$ the set
$S_i$ of the $x$ for which $2^i$ is the least power of $2$ bounding the moduli of all entries of $x-1$ is a
shell of height $2^i$ for $\Gamma_L,j,p,q,W^\flat$ (Definition~\ref{pfr:shell}); the $S_i$ cover $\pfrG$.
\end{lemma}
\begin{proof}
Every $x$ lies in exactly one $S_i$. (a) Let $H=2^i$ and $L'=L2^i\ge\rho$. The position ball of $\Gamma_{L'}$
($M=1$, all positions allowed) contains $S_i$, and $\Lambda_{\Gamma_{L'}}=L^24^i=\Lambda H^2$; one gauge
suffices, and $1\le2H^{9n}$. (b) Let $\varsigma\in(0,1)$ with $\Lambda^{1200n^2}\le\varsigma^{400n^2}H$; then $i\ge1$
(otherwise $L^{2400n^2}\le\varsigma^{400n^2}<1$). Put $T=L^{11n+1}/\varsigma\ge L^{11n}$; then
$T^{-1}\le\varsigma/L\le\varsigma/2$. For $x\in S_i$ some entry of $x-1$ exceeds $2^{i-1}=H/2$, and
\[
\varsigma^{400n^2}\cdot2T^{99n}\le\varsigma^{99n}\,L\,(L^{11n+1}/\varsigma)^{99n}=L^{1089n^2+99n+1}\le L^{2400n^2}
\le\varsigma^{400n^2}H,
\]
so $2T^{99n}\le H$ and some entry of $x-1$ exceeds $T^{99n}$. By Lemma~\ref{pfr:deepreal} with $\tau'=\varsigma$,
$x$ has a deep vector of depth $\varsigma^{-2}$.
\end{proof}

\subsection{The sum over the dyadic shells; Propositions~\ref{pfr:formbound}, \ref{pfr:conv}, \ref{pfr:remote}}\label{prp:ss-sum}

Recall $\pbmTh=\pi^{(2\varrho+1)k}\allowbreak \prpvolU k/\prpvolU r$ and $\mathsf c=\frac{\prpvolU k}{\prpvolU r}(\lambda\,\#I)^{-2}$
\eqref{pfr:eq-c}.

\begin{proposition}[The sum over the shells at $\C$]\label{prp:sum}
Assume \textup{(H0)}--\textup{(H1)}. Let $j\le\varrho$, $L\ge\rho$, $p,q$ heads of length $j$ for $\Gamma_L$,
$t\ge1$, and $f\ge0$ measurable with $f\le\pfrF^+_j$ everywhere \textup{(}for $W=W^\flat$\textup{)} and, if
$t>1$, $f(x)^2\le(V/t^{s_0})^2\mathsf Q_{\Gamma_L}(x)$ for all $x$, with $V$ as in Lemma~\ref{prp:normW}. Then
\[
\#I\cdot\mathsf c\,\pfrkap_j(p,q)\int f\,d\pfrM\le\pbmTh\,L^{10^9n^4}\rho^{24k}\,t^{10^8n^4}/t^{s_0}.
\]
\end{proposition}
\begin{proof}
Apply Theorem~\ref{pfr:shells} to the dyadic shells $S_i$ of Lemma~\ref{prp:dyadic}, with the mass bound $V\ge1$
of Lemma~\ref{prp:normW}(b), $e=t^{-1}$, and $r'=s_0-10^8n^4$; then $9\cdot10^7n^4+r'\le s_0=s$, and if $e<1$ the
gain $Ve^{s}=V/t^{s_0}$ is our hypothesis. So
$\pfrkap_j\int f\,d\pfrM\le V\,t^{-r'}(L^2)^{4\cdot10^8n^4}\mathsf N(W^\flat)$. By Lemma~\ref{prp:normW}(a),
$\#I\cdot\mathsf c\cdot\mathsf N(W^\flat)\le\frac{\prpvolU k}{\prpvolU r}\pi^{2(k-5)\varrho}$. Hence the left side is at most
\[
\frac{\prpvolU k}{\prpvolU r}\,\pi^{k+2(k-5)\varrho}\rho^{24k}\rho^{9999n}L^{8\cdot10^8n^4}\frac{t^{10^8n^4}}{t^{s_0}}
\le\pbmTh\,L^{10^9n^4}\rho^{24k}\frac{t^{10^8n^4}}{t^{s_0}},
\]
since $k+2(k-5)\varrho\le(2\varrho+1)k$, $\rho\le L$ and $9999n+8\cdot10^8n^4\le10^9n^4$.
\end{proof}

\begin{lemma}\label{prp:arith}
Assume \textup{(H0)}. \textup{(a)} If $0<\nu$, $1\le\rho^{24}\nu$, then $(\pi^2/\nu^2)^{k-5}\le\bigl(V/(\rho^{24}\nu)^{s_0}\bigr)^2$.
\textup{(b)} Let $\nu>0$, $L\ge\rho$, $H\ge0$ with $L^{n^6}\le\rho^{n^6}H$, and $t=\max(1,\rho^{24}\nu)$. Then
$\pbmTh L^{10^9n^4}\rho^{24k}t^{10^8n^4}/t^{s_0}\le\pbmTh\,b^{2000n^9}\pfrell_{k-n^7}(\nu^{-1})\,H$.
\end{lemma}
\begin{proof}
(a) $(\pi/\nu)^{k-5}=(\pi\rho^{24})^{k-5}/(\rho^{24}\nu)^{k-5}\le(\pi\rho^{24})^k\rho^{9999n}/(\rho^{24}\nu)^{s_0}$, as
$\pi\rho^{24}\ge1$, $\rho^{24}\nu\ge1$ and $s_0\le k-5$.

(b) Write $k=n^7+K$ ($K\ge0$ by (H0)). As $n>10^8$, $10^9n^4\le n^6$, so $L^{10^9n^4}\le L^{n^6}\le\rho^{n^6}H$; and
$s_0=k-3n-6\ge10^8n^4+K+6$ (as $n^7-3n-12\ge10^8n^4$), so $t^{10^8n^4}/t^{s_0}\le t^{-K-6}$. Hence the left side
is at most $\pbmTh\,\rho^{n^6+24n^7}H\,(\rho^{24}/t)^Kt^{-6}$. Now $\rho^{24}/t\le\nu^{-1}$ and $t^{-1}\le\min(1,\nu^{-1})$,
so $(\rho^{24}/t)^Kt^{-6}\le\pfrell_K(\nu^{-1})$; and $\rho^{n^6+24n^7}=b^{(6n+2)(n^6+24n^7)}\le b^{200n^8}\le b^{2000n^9}$.
\end{proof}

\begin{proof}[Proof of Proposition~\ref{pfr:formbound}]
Let $\omega$ be one of the two weights; $|\omega|\le1$. If $|P_l(p)|>b^{60n^2}$ or $|P_l(q)|>b^{60n^2}$ for some
$l\le i$, then, as $\rho^{9n}=b^{(6n+2)9n}\le b^{60n^2}$, Lemma~\ref{prp:headsupport} gives $\pfrF^+_i=0$, so
$\pfrI_i(p,q;\omega)=0$ by Lemma~\ref{pfr:L1}. Otherwise let $L$ be given by Lemma~\ref{prp:exgauge} with $K=n^6$, so that
$p,q$ are heads of length $i$ for $\Gamma_L$ and $L^{n^6}\le\rho^{n^6}H$ with $H=\Lambda^\infty_i(p)\Lambda^\infty_i(q)$.
Let $f(x)=\pfrF_i(x;\omega)\le\pfrF^+_i(x)$ (Lemma~\ref{pfr:L1}) and $t=\max(1,\rho^{24}|\nu|)$. If $t>1$, then
$t=\rho^{24}|\nu|$, and Theorem~\ref{pfr:gain-char} (for the character weight, with $\xi=\eta=\nu$) or
Theorem~\ref{pfr:gain-heads} (for $\omega=1$, $\nu=p_{i-1}^{-1}-q_{i-1}^{-1}$, $j=i$) gives, at $\C$ where
$P/\operatorname{mod}(\nu)=\pi^2/|\nu|^2$ and $\operatorname{rank}X'=k-5$,
$f(x)^2\le(\pi^2/|\nu|^2)^{k-5}\mathsf Q_{\Gamma_L}(x)\le(V/t^{s_0})^2\mathsf Q_{\Gamma_L}(x)$ by
Lemma~\ref{prp:arith}(a). By Lemma~\ref{pfr:L1} (the integrability is Proposition~\ref{pfr:conv}, proved below
without using this proposition), Proposition~\ref{prp:sum} and Lemma~\ref{prp:arith}(b),
\[
\#I\cdot\mathsf c\,|\pfrI_i(p,q;\omega)|\le\#I\cdot\mathsf c\,\pfrkap_i\int f\,d\pfrM\le\pbmTh\,b^{2000n^9}
\pfrell_{k-n^7}(|\nu|^{-1})\,\Lambda^\infty_i(p)\Lambda^\infty_i(q).\qedhere
\]
\end{proof}

\begin{proof}[Proof of Proposition~\ref{pfr:conv}]
The support claim is Lemma~\ref{prp:headsupport} (and Lemma~\ref{pfr:L1} for $\pfrI_j$). For the finiteness,
let $L$ be as in Lemma~\ref{prp:exgauge} (any $K$), and apply Proposition~\ref{prp:sum} with $t=1$ and $f=\pfrF^+_j$
(the gain hypothesis is void): $\#I\cdot\mathsf c\,\pfrkap_j(p,q)\int\pfrF^+_j\,d\pfrM<\infty$, and $\mathsf c>0$,
$\pfrkap_j(p,q)>0$.
\end{proof}

\begin{proof}[Proof of Proposition~\ref{pfr:remote}]
Let $L=\rho$, $t=b^{6r}\ge1$, $T=t\rho^{11n}$, so that $\mathcal B$ is the set of the $x$ all of whose entries of
$x-1$ have modulus $\le T^{99n}$. At floor $0$ all sequences are heads of length $0$. Let
$f=\mathbf 1_{\pfrG\smallsetminus\mathcal B}\,\pfrF^+_0\le\pfrF^+_0$. For $x\notin\mathcal B$, Lemma~\ref{pfr:deepreal}
with $\tau'=t^{-1}$ (indeed $T\ge\rho^{11n}$ and $T^{-1}\le t^{-1}/2$) shows that $x$ has a deep vector of depth
$t^2$, and Theorem~\ref{pfr:deep} with $\varsigma=t^{-2}$ and the mass bound $V$ gives
$f(x)^2\le t^{-2s_0}V^2\mathsf Q_{\Gamma_\rho}(x)$; for $x\in\mathcal B$, $f(x)=0$. Proposition~\ref{prp:sum} and
$\#I\ge1$, $\pfrkap_0=1$ give
$\mathsf c\int_{\pfrG\smallsetminus\mathcal B}\pfrF^+_0\,d\pfrM\le\pbmTh\rho^{10^9n^4}\rho^{24k}t^{10^8n^4}/t^{s_0}$, which
is the claim.
\end{proof}

\subsection{The top floor: proof of Proposition~\ref{pfr:top}}\label{prp:ss-top}

Write $J(\gamma)=\int_\pfrG W^\flat(y)\otimes\overline{W^\flat(\gamma^{-1}y)}\,dy$, so that
$\pfrN_\sigma(x)=N(x)J(w_x\sigma)$ with $N(x)=\frac{\sqrt{c_G}}{\#I}\bigl((|x_\varrho|^2-1)^+\bigr)^{(k-1)/2}\prod_{a<\varrho}
|x_a|^{d_a}$, $d_a=k+1+2(\varrho-a-1)$. (For $|x_\varrho|\le1$ the factor $(|x_\varrho|^2-1)^{(k-1)/2}$ of
\S\ref{pfr:ss-realdata} is $0$.) Let $\prpdens(c)=\frac{2\pi^k}{(k-1)!}((|c|^2-1)^+)^{k-1}$, the
density of $\theta_\varrho$, and recall $c_G=\frac{\prpvolU k}{\prpvolU r}\frac{2\pi^k}{(k-1)!}$
(Remark~\ref{pfr:rem-match}).

\begin{lemma}[Compact supports]\label{prp:support}
Let $x^0$ be a sequence of non-zero complex numbers. There are a neighbourhood $\mathcal N$ of $x^0$ and a
compact set $C\subseteq\pfrG$ such that for $x\in\mathcal N$, $\sigma\in\pfrG^{(\varrho)}$ and $y\in\pfrG$:
if $W^\flat(y)\ne0$ and $W^\flat((w_x\sigma)^{-1}y)\ne0$, then $y\in C$.
\end{lemma}
\begin{proof}
Let $\varepsilon=\frac12\min_{l\le\varrho}|P_l(x^0)|$ and $\mathcal N$ the set of the $x$ with $|P_l(x)|>\varepsilon$ for
$l\le\varrho$. Let $\gamma=w_x\sigma$; the leading block of order $\varrho$ of $Z(\gamma)$ is that of $Z(w_x)$ ($\sigma$ fixes
$e_0,\dots,e_{\varrho-1}$), so $|\delta_l(\gamma)|=|P_l(x)|>\varepsilon$ for $l\le\varrho$. By (H1) there are $u,u'\in\pfrU$
with $k_1=uy\in\pfrK$, $k_2=u'\gamma^{-1}y\in\pfrK$; then $\kappa:=k_1k_2^{-1}=u\,w_x\sigma\,u'^{-1}$ lies in the compact
set $C_0=\{\kappa\in\pfrK\pfrK^{-1}:|\delta_l(\kappa)|\ge\varepsilon,\ l\le\varrho\}$ ($\delta_l$ is invariant under
$\gamma\mapsto u\gamma u'^{-1}$, Lemma~\ref{prp:minor}'s proof). $C_0$ is contained in the open set where the $\delta_l$,
$l\le\varrho$, are invertible, on which {\cite[Proposition~3.41]{P51}} decomposes every element uniquely as
$u(\kappa)w_{p}\sigma'u'(\kappa)$ with $u(\cdot)$ continuous. By uniqueness $u=u(\kappa)$, which lies in the compact
set $u(C_0)$; so $y=u^{-1}k_1\in C:=u(C_0)^{-1}\pfrK$.
\end{proof}

\begin{lemma}[Descent on the top floor]\label{prp:topdescent}
Let $\varphi\colon\C\to\C$ be measurable with $|\varphi|\le1$, let $p,q$ be sequences of units with
$\int\pfrF^+_\varrho\,d\pfrM<\infty$, $w=w_{\varrho,p}$, and $\Pi$ as in \S\ref{pfr:ss-descent} with $a=\varrho$. Then
$c\mapsto\prpdens(c)\varphi(Z_{zz}(\sigma_c))\allowbreak \iint\Pi(x;w\sigma_c;g')\,d\pfrM(x)\,d\pfrM(g')$ is integrable on $\C$, and
$\pfrkap_\varrho(p,q)$ times its integral is $\pfrI_\varrho(p,q;\varphi\circ Z_{zz})$.
\end{lemma}
\begin{proof}
This is the proof of Lemma~\ref{pfr:descent} (\S\ref{pfr:ss-descent}) with $a=\varrho$ and the following changes. Fact
(iv) is replaced by the definition of $\theta_\varrho$: it is the image of $\prpdens(c)\,dc\times\mu_\pfrH$ under
$(c,h)\mapsto\sigma_c^{-1}h$, with $\mu_\pfrH$ carried by $\pfrH=\pfrR^{(\varrho+1)}\cap\pfrG^{(\varrho)}$ (by (N1), as
$\pfrN_\varrho=1$). Fact (iii) holds for $t\in\pfrH$: $tz=z$, so $Z_{zz}(tyt')=Z_{zz}(y)$, and the correlations from floor
$\varrho+1$ are averages over $\mu_\pfrH$, which is right invariant under $\pfrH$, and $d=w_{\varrho,q}^{-1}w$ commutes with
$\pfrH$. With $F(x,g,\sigma,\tau)=\varphi(Z_{zz}(\sigma\tau^{-1}))\Pi(x;\sigma w\tau^{-1};\sigma g)$ on
$\pfrM\otimes\pfrm_\varrho\otimes\theta_\varrho\otimes\theta_\varrho$, $\int|F|\le\int\pfrF^+_\varrho\,d\pfrM<\infty$, and
$\pfrkap_\varrho\int F=\pfrI_\varrho(p,q;\varphi\circ Z_{zz})$. For $\sigma\in\pfrG^{(\varrho)}$, Lemma~\ref{pfr:floorright}
with $\xi=\sigma$ and the description of $\theta_\varrho$ give $\int f(F(x,g,\sigma,\tau))\,d\theta_\varrho(\tau)=\iint
f(F(x,g,\sigma,\sigma_c^{-1}h\sigma))\,\prpdens(c)\,dc\,d\mu_\pfrH(h)$; and $\sigma w(\sigma_c^{-1}h\sigma)^{-1}=wh^{-1}\sigma_c$,
$Z_{zz}(h^{-1}\sigma_c)=Z_{zz}(\sigma_c)$, $\Pi(x;wh^{-1}\sigma_c;\sigma g)=\Pi(x;w\sigma_c;h\sigma g)$. So
$F(x,g,\sigma,\sigma_c^{-1}h\sigma)=E(x,h\sigma g,c)$ with $E(x,g',c)=\varphi(Z_{zz}(\sigma_c))\Pi(x;w\sigma_c;g')$. When
$\sigma\sim\theta_\varrho$, $g\sim\pfrm_\varrho$ and $h\sim\mu_\pfrH$, the point $h\sigma g$ has law
$\mu_\pfrH*\theta_\varrho*\pfrm_\varrho=\pfrM$ (the factors $\eta_\varrho,\eta_{\varrho+1}$ are trivial and
$\theta_{\varrho+1}=\mu_\pfrH$). Hence $F$ and $E$ (the latter on $\pfrM\otimes\pfrM\otimes\prpdens(c)dc$) have the same
law, and Fubini's theorem concludes as in Lemma~\ref{pfr:descent}.
\end{proof}

\begin{proof}[Proof of Proposition~\ref{pfr:top}]
\emph{Continuity.} Let $x^0$, $\mathcal N$, $C$ be as in Lemma~\ref{prp:support}. On $\mathcal N$, $x\mapsto N(x)$,
$x\mapsto w_x$ and $x\mapsto\sigma_x$ are continuous (the boost depends continuously on its parameter
$x_\varrho/\min(1,|x_\varrho|)$), and $\pfrX(x)=N(x)\int_CW^\flat(y)\otimes\overline{W^\flat((w_x\sigma_x)^{-1}y)}\,dy$, with a
jointly continuous integrand bounded by $(\sup\|W^\flat\|)^2$ on the compact set $C$. By dominated convergence $\pfrX$ is
continuous at $x^0$.

\emph{The equality.} Let $\sigma\in\pfrG^{(\varrho)}$ and $c=h(\sigma z,z)=x_\varrho=x'_\varrho$. Write $\sigma z=-cz+v$ with $v$ in the
span $X_0$ of the $x_i$ (the $z$-coefficient of a vector $w$ is $-h(w,z)$, since $h(z,z)=-1$); $h(\sigma z,\sigma z)=-1$ gives $|v|^2=|c|^2-1$. If $|c|\le1$ both sides vanish ($N(x)=N(x')=0$).
If $|c|>1$, also $\sigma_cz=-cz+v_c$ with $|v_c|=|v|>0$, and there is $h_1\in\pfrH=\auU(X_0)$ with $h_1v_c=v$; then
$h_1\sigma_cz=\sigma z$, so $h_2:=\sigma_c^{-1}h_1^{-1}\sigma$ fixes $z$ and the $e_a,e_a^*$, i.e. $h_2\in\pfrH$, and
$\sigma=h_1\sigma_ch_2$. Since $h_1$ commutes with $w_x$, $W^\flat(hy)=\pi(h)W^\flat(y)$ for $h\in\pfrH$, and $dy$ is left
invariant, the substitution $y\mapsto h_1y$ gives $J(w_x\sigma)=(\pi(h_1)\otimes\overline{\pi(h_2)^{-1}})\,J(w_x\sigma_c)$,
with a unitary operator that does not depend on $x$. Hence
$\pfrip{\pfrN_\sigma(x)}{\pfrN_\sigma(x')}=\pfrip{\pfrN_{\sigma_c}(x)}{\pfrN_{\sigma_c}(x')}=\pfrip{\pfrX(x)}{\pfrX(x')}$.

\emph{The descent.} Let $\varphi(c)=e^{\xi c-\overline{\xi c}}$. By Proposition~\ref{pfr:conv},
$\int\pfrF^+_\varrho\,d\pfrM<\infty$, so Lemma~\ref{prp:topdescent} applies. We identify its integrand. The correlations
from floor $\varrho+1$ are averages over $\mu_\pfrH$, and $\pfrH$ acts unitarily on the values of $W^\flat$ and commutes
with $d$; so $\mathsf k^1_{\varrho+1}(g,y)=\pfrip{W^\flat(g)}{W^\flat(y)}$ and $\mathsf k^d_{\varrho+1}(g,y)=
\pfrip{W^\flat(g)}{W^\flat(dy)}$, and, as $\sigma_c$ commutes with $w$ and $d\sigma_c^{-1}w^{-1}=\sigma_c^{-1}w_{\varrho,q}^{-1}$,
\[
\Pi(x;w\sigma_c;g')=\pfrip{W^\flat(g')}{W^\flat(g'x)}\,\overline{\pfrip{W^\flat((w_{\varrho,p}\sigma_c)^{-1}g')}{W^\flat((w_{\varrho,q}
\sigma_c)^{-1}g'x)}}.
\]
With $\pfrM=\lambda\,dx$ and the measure-preserving map $(g',x)\mapsto(g',g'x)$ of $\pfrG\times\pfrG$, and by
Lemma~\ref{prp:support} (compact supports, so everything is integrable), expanding the inner products coordinatewise,
$\iint\Pi\,d\pfrM\,d\pfrM=\lambda^2\pfrip{J(w_{\varrho,p}\sigma_c)}{J(w_{\varrho,q}\sigma_c)}$. For $|c|>1$,
$Z_{zz}(\sigma_c)=c$ and
\[
N(p^c)N(q^c)=\frac{c_G}{(\#I)^2}\bigl((|c|^2-1)^+\bigr)^{k-1}\pfrkap_\varrho(p,q)=\mathsf c\,\lambda^2\,\prpdens(c)\,\pfrkap_\varrho(p,q);
\]
for $|c|\le1$ both sides are $0$. Hence $\varphi(c)\pfrip{\pfrX(p^c)}{\pfrX(q^c)}=\mathsf c\,\pfrkap_\varrho(p,q)\,
\prpdens(c)\varphi(Z_{zz}(\sigma_c))\iint\Pi(x;w\sigma_c;g')\,d\pfrM\,d\pfrM$ for all $c\ne0$, and
Lemma~\ref{prp:topdescent} gives the integrability and the value $\mathsf c\,\pfrI_\varrho(p,q;\varphi\circ Z_{zz})$.
\end{proof}

\subsection{The hypotheses (H1)--(H5) for the Bessel functions of \texorpdfstring{{\cite[Theorem~3.22]{P51}}}{[Moment, Theorem~3.22]}}\label{prp:ss-H}

Let $W=W_\tau$, $W^\flat=W_\tau^\flat$ be the functions of {\cite[Proposition~3.18]{P51}} and {\cite[Theorem~3.22]{P51}},
$\mathbb W=\pbW_\tau$, $\#I=d_\tau=\dim\pbW_\tau$, $f_W=f$ (the exponent of {\cite[Proposition~3.18(e)]{P51}}). They are
continuous Bessel functions of type $\mathbb W$ by (B1), and $\mathbb W$ is irreducible under $\pfrH$
({\cite[\S3.2.1]{P51}}).

\emph{(H1).} $\int\|W\|^2=d_\tau$ is the normalisation of {\cite[Proposition~3.18]{P51}}; $\int\|W^\flat\|^2\le d_\tau$
and $\int\|W^\flat-W\|^2\le b^{-2r^3f}d_\tau$ are {\cite[Proposition~3.18(e)]{P51}}. If $W^\flat(y)\ne0$ then
$\inf_{s\in R_\tau}(\|sy\|+1)\le\rho$ ({\cite[Proposition~3.18(e)]{P51}}), so some $s=u'h\in R_\tau=\pfrU\pfrH$ has
$\|sy\|<\rho$; as $h$ is a Euclidean isometry normalising $\pfrU$, $\|uy\|=\|sy\|$ with $u=h^{-1}u'h\in\pfrU$.

\emph{(H2), with equality.} Let $\mathsf k(g)=\pfrip W{R(g)W}=\int_{\pfrU\backslash\pfrG}\pfrip{W(y)}{W(yg)}\,dy$. As
$\|W\|$ is bounded ({\cite[Proposition~3.18(a)]{P51}}) and integrable on $\pfrU\backslash\pfrG$ (B2), $\mathsf k$ is
bounded and measurable; by the right invariance of the measure of $\pfrU\backslash\pfrG$,
$\int\pfrip{W(ys)}{W(yt)}dy=\mathsf k(s^{-1}t)$, and $\mathsf k(g^{-1})=\overline{\mathsf k(g)}$. By the reproduction
property (B7), $W(y)=\int e(a)W(ya)\,da$ with $e=e_\tau=\pbdeg_\tau\overline{c_\tau}$ integrable; inserting it in the
second factor of $\mathsf k(g)=\int\pfrip{W(yg^{-1})}{W(y)}dy$ (Fubini: $e\in L^1$, $W$ bounded, $\|W\|$ integrable)
gives $\mathsf k(g)=\int e(a)\mathsf k(ga)\,da$; conjugating, $\mathsf k(g)=\int\overline{e(a)}\mathsf k(a^{-1}g)\,da$.
Hence, substituting $t'=s^{-1}gt$,
\[
\mathsf k(g)=\iint\overline{e(s)}e(t)\mathsf k(s^{-1}gt)\,dt\,ds=\int\mathsf k(t')\Bigl(\int\overline{e(s)}\,e(g^{-1}st')\,ds\Bigr)dt'.
\]
The inner integral is $\pbdeg_\tau^2\int c_\tau(s)\overline{c_\tau(g^{-1}st')}\,ds$; with $a=s^{-1}$ and
$c_\tau(t'^{-1}ag)=\pfrip{\pi(t')v_0}{\pi(a)\pi(g)v_0}$ it equals $\pbdeg_\tau^2\int\overline{c_\tau(a)}\,c_\tau(t'^{-1}ag)\,da$,
and Schur's orthogonality relations ({\cite[\S5.4(ii)]{P51}}, for the four vectors $\pi(t')v_0,\pi(g)v_0,v_0,v_0$)
give $\pbdeg_\tau\,c_\tau(t'^{-1})\,c_\tau(g)=\pbdeg_\tau\overline{c_\tau(t')}c_\tau(g)$. So
$\mathsf k(g)=c_\tau(g)\int e(t')\mathsf k(t')\,dt'=c_\tau(g)\mathsf k(1)=d_\tau\,c_\tau(g)$, and
$\int_\pfrG|\mathsf k|^2=d_\tau^2\int|c_\tau|^2=d_\tau^2/\pbdeg_\tau$ by \textup{(DS1)(ii)}. This is (H2) with equality.

\emph{(H3).} The entries of $\varsigma_\tau$ are $\pm(\frac k2-1),\dots$ (modulus $\le r\le b$) and
$-\kappa(\tau)-j\Delta$ with $j\Delta\le m\Delta\le\kappa(\tau)\le b$ ({\cite[Lemma~13.4\,(S5)]{P46}} and (S7)(g)); so $|\varsigma_{\tau,i}|\le2b$, every factor $(\varsigma_{\tau,i}-\varsigma_{\tau,j})/(j-i)$ is
$\le4b\le b^3$, and $\pbdeg_\tau\le b^{3r^2/2}$. Next, $\prpvolU N=\prod_{j=1}^N\frac{2\pi^j}{(j-1)!}$, so
$\prpvolU k/\prpvolU r=\prod_{j=k+1}^r\frac{(j-1)!}{2\pi^j}\le(r^r)^{r-k}$; with $r-k=2\varrho+1\le r/2$ (as $k\ge r/2$,
{\cite[\S3.2.1]{P51}} sizes) and $r\le b$, this is $\le b^{r^2/2}$; and $\pi^{(2\varrho+1)k}\le4^{(r-k)k}\le4^{r^2/4}\le
b^{r^2/2}$. Hence $\pbdeg_\tau\pbmTh\le b^{3r^2/2+r^2/2+r^2/2}\le b^{4r^2}$.

\emph{(H4)} is \textup{(DS3)} for the Haar measure of \textup{(DS1)}: since $\|g\|\le T$ implies $(32/\|g\|)^{9r^2}\ge
(32/T)^{9r^2}$, $\pbdeg_\tau\pbvol\{\|g\|\le T\}\le(T/32)^{9r^2}b^{99r^2}\le T^{9r^2}b^{99r^2}$.

\emph{(H5).} $f\ge4000\,n$, {\cite[Lemma~13.4\,(S7)(b)]{P46}}. \emph{(H0)} is
Remark~\ref{pfr:rem-match}\,(b).

\subsection{What Theorem~\ref{pfr:R} now rests on}\label{prp:ss-R}

With the proofs of this section, Theorem~\ref{pfr:R} (through \S\ref{pfr:ss-Rproof}, Proposition~\ref{pfr:mass} and
the propositions of \S\ref{pfr:ss-realforms}) depends only on: the shared Theorems~\ref{pfr:oneframe},
\ref{pfr:gain-char} and~\ref{pfr:gain-heads}; the results of \S\ref{pfr:sec} with their proofs (and the standard
facts used there: Witt's theorem, measurable sections of orbit maps); {\cite[Proposition~3.41]{P51}} (in
Lemma~\ref{prp:support}); {\cite[Lemmas~8.80 and~8.81]{P51}} (in Lemma~\ref{prp:fibre}); and the hypotheses
(H0)--(H5). For the Bessel functions of {\cite[Theorem~3.22]{P51}}, (H0)--(H5) hold by \S\ref{prp:ss-H}, using
{\cite[Proposition~3.18]{P51}} and {\cite[Theorem~3.22\,(B1)]{P51}}, (B2), (B7) (resting on \textup{(DS1)}--\textup{(DS4)},
{\cite[Theorem~3.16]{P51}}, with \textup{(DS4)} = {\cite[Theorem~7.2]{P51}}), the Schur relations of {\cite[\S5.4]{P51}},
\textup{(DS3)}, and {\cite[Lemma~13.4]{P46}}. Schur's lemma, the uniqueness of Haar measure and the polar formula are
used as standard facts.



\providecommand{\auO}{\mathcal O}
\providecommand{\auPl}{\Phi}
\providecommand{\auU}{\mathrm U}
\providecommand{\auGL}{\mathrm{GL}}
\providecommand{\auA}{\mathbb A}
\providecommand{\auF}{\mathbb F}
\providecommand{\aurd}{\mathrm{rd}_E}
\providecommand{\aubar}[1]{\overline{#1}}
\providecommand{\panot}[1]{\par\smallskip{\raggedright\noindent\textbf{[Not proved here:} #1\textbf{]}\par}\smallskip}
\providecommand{\pbR}{\overline R_0}
\providecommand{\pbUb}{\overline U}
\providecommand{\pbHb}{\overline H}
\providecommand{\pbvol}{\operatorname{vol}}
\providecommand{\pbFix}{\operatorname{Fix}}
\providecommand{\pbDmain}{\mathcal D_{\mathrm{main}}}
\providecommand{\pbmGf}{G^{\flat}}
\providecommand{\pbmXp}{X^{+}}
\providecommand{\auG}{\overline G}
\providecommand{\pbmOrb}{\mathcal O}
\providecommand{\pbmT}{\mathcal T}
\providecommand{\pbmRT}{R_T}
\providecommand{\pbmE}{\mathsf E}
\providecommand{\pbmn}{\mathsf n}
\providecommand{\pbmden}{\mathsf{den}}
\providecommand{\pbmvol}{\mathsf v}
\providecommand{\pbmaS}{\mathsf a_*}
\providecommand{\pbmrS}{\mathsf r_*}
\providecommand{\pbmcG}{c_G}
\providecommand{\pffET}{E_T}
\providecommand{\pffOT}{\auO_T}
\providecommand{\pfgOh}{\auO_{\mathrm h}}
\providecommand{\pfgQ}{\mathsf Q_N}
\providecommand{\pffG}{G_T}
\providecommand{\pffK}{K_T}
\providecommand{\pffKM}{K_T(M)}
\providecommand{\pffch}{\chi_M}
\providecommand{\pffW}{W_Z}
\providecommand{\pffm}{\varkappa}
\providecommand{\pffr}{\rho}
\providecommand{\pffl}{\mathsf f}
\providecommand{\pffI}{\mathcal I}
\providecommand{\pffZ}{\mathsf Z}
\providecommand{\pfrF}{\mathsf F}
\providecommand{\pfrB}{\mathsf B}
\providecommand{\pfrM}{\mathsf M}
\providecommand{\pfrm}{\mathsf m}
\providecommand{\pfrN}{\mathsf N}
\providecommand{\pfrkap}{\kappa}
\providecommand{\pfrH}{\mathsf H}
\providecommand{\pbmKl}{\mathsf{Kl}}
\providecommand{\pbmass}{\operatorname{mass}}
\providecommand{\pfrmod}{\operatorname{mod}}
\providecommand{\pfrI}{\mathcal I}
\providecommand{\pfrip}[2]{\langle#1,#2\rangle}

\section{The local systems at the finite places and Siegel's step}\label{pff:sec}

This section concerns the parts \textup{(F)} and \textup{(S)} of {\cite[Theorem~3.51]{P51}}. The
part \textup{(R)}, and the theory of forms on pairs of heads that is common to the real and the
finite places (floor measures, correlations, gauges and sizes, depth of frame vectors, the sum
over dyadic shells of positions), are treated in \S\ref{pfr:sec}; from there we quote, with
their labels, the statements that we need, and we do not prove them again.

\emph{The results of this section.} We list the results by the label of the statement.
\begin{itemize}
\item[(i)] \emph{The exact statements} (\S\ref{pff:ss-statements}, Theorems~\ref{pff:F}
and~\ref{pff:S}). We fix the forms $J^T_i$ at the finite places, and we check that the
constants of \textup{(F)} and \textup{(S)} of {\cite[Theorem~3.51]{P51}} are exactly those that
{\cite[Theorem~3.53]{P51}} and {\cite[Lemma~3.52]{P51}} use: the mass $(1+b^{-6})b^{-6k\#\mathsf A}$, the
weights $\Lambda^T$, the decay $\beta^T$, the place $(E_T,dx,\psi_T,w\mapsto M^{-1}w^{-1})$, and in
\textup{(S)} the factor $(1+b^{-6})\pbmaS$ against $\pbmrS^{\,2}$. One reading has to be made precise in
{\cite[Theorem~3.51]{P51}}: in \textup{(F)} the forms depend on a group $\pffZ$ that is
chosen together with them, and \textup{(S)} is stated with the same group. The requirement that $T$
contain the places at which the chosen basis of the lattice is not integral is already part of the
hypothesis of \textup{(F)}, as these places lie above $N$. This does not affect the
deduction of {\cite[Theorem~3.53]{P51}} (Remark~\ref{pff:rem-L2}).
\item[(ii)] \emph{Theorem~\ref{pff:F} (part \textup{(F)}) is proved}, using the statements of the
common theory of \S\ref{pfr:ss-shared} (Lemmas~\ref{pfr:L1}, \ref{pfr:descent},
Theorems~\ref{pfr:deep}, \ref{pfr:gain-char}, \ref{pfr:gain-heads}, \ref{pfr:shells}; proved in
\S\S\ref{pfr:sec}, \ref{pfs:sec}) and Lemma~\ref{pff:twogauges}, proved in \S\ref{pfg:ss-two}. What is specific to the finite
places is written here: the basis at $T$, the gauges, the depth ideals, the floor masses and the volume of
$\pffK$, the dyadic shells, the mass (with the gain $b^{-6k}$ for each fixed vector,
Theorem~\ref{pff:fixed}), the decay (Theorem~\ref{pff:decay}), and the density at the last entry
(Proposition~\ref{pff:top}).
\item[(iii)] \emph{Theorem~\ref{pff:S} (Siegel's step) is proved} (\S\ref{pff:ss-siegel}), modulo Weil's Siegel--Weil theorem, from which Theorem~\ref{pff:siegelformula} is derived in \S\ref{pfg:ss-siegel}, and the top forms at the real places of \S\ref{pfr:sec}.
The lemmas used here and proved in \S\ref{pfg:sec} are listed in \S\ref{pff:ss-pending}.
\end{itemize}

\subsection{The setting at a finite set of places}\label{pff:ss-setting}

We use the notation of {\cite[\S2.1]{P51}} and {\cite[\S3.3]{P51}}. Hermitian forms follow convention (Cv1) of
{\cite[\S2.1]{P51}} in this section, as in \S\ref{pfr:sec} (linear in the first variable); \S\ref{pfg:sec} writes them
in the opposite convention, as declared there. Let $T$ be a finite set of
finite places of $E$, stable under conjugation, containing the places above $MN$. Let
$\pffET=\prod_{w\in T}E_w$, with ring of integers $\pffOT=\prod_{w\in T}\auO_w$, the Haar measure
$dx$ for which $\pffOT$ has volume $1$, and the module $|\xi|_T=\prod_{w\in T}|\xi|_w$; for a unit
$c$ of $E_T$ the map $x\mapsto cx$ multiplies $dx$ by $|c|_T$. The involution of $E/E^+$ acts on
$\pffET$, and $\psi_T(x)=\psi_{E,f}(-x/\varpi)$ is the character of {\cite[\S3.3.3]{P51}}; it is
trivial on $\pffOT$ (as $\varpi$ generates the different), and $\psi_T(x+\aubar x)=1$ for all
$x\in\pffET$: with $y=-x/\varpi$ and $\aubar\varpi=-\varpi$ we have $-(x+\aubar x)/\varpi=y-\aubar
y$, whose trace to $E^+_T$ vanishes, and $\psi_E$ factors through the trace.

\emph{The lattice.} Let $Y_0$ be the hermitian space of {\cite[\S2.1]{P51}} (dimension
$r=2\varrho+k+1$, with the hyperbolic pairs $e_a,e_a^*$ and the middle space $\pbmXp=X_0\perp
Ez$), and $L\subset Y_0$ the lattice of {\cite[\S2.1]{P51}}, whose stabiliser is $K$; it is
unimodular, and we fix a basis of it over $\auO_E$ (the \emph{lattice basis}; it differs from the
basis $e_a,x_i,z,e_a^*$ of $Y_0$ only in the plane of $x_{k_0}$ and $z$, by a matrix that is integral
with integral inverse at the places not above $N$). Let $V_T=Y_0\otimes_EE_T$,
$L_T=L\otimes\pffOT$, $\pffG=\auU(V_T)$, and
\[
\pffK=\{g\in\pffG:gL_T=L_T\},\qquad\pffKM=\{g\in\pffK:(g-1)L_T\subseteq ML_T\}.
\]
Both are compact open subgroups, $\pffKM$ is normal in $\pffK$, and $\pffKM$ is the component at $T$
of the group $J$ ({\cite[\S2.1]{P51}}; $J$ is the kernel of the reduction of $K$ modulo $M$),
as $T$ contains the places above $M$; the component of $K$ at $T$ is $\pffK$. We write $U_T$,
$H_T$, $\pbmGf_T$ for the groups $U$, $H$, $\pbmGf$ over $\pffET$, and $\sigma(u)$ for the sum of
the entries just above the diagonal blocks of $u\in U_T$ (in the frame
$f_0,\dots,f_\varrho$ of {\cite[\S3.3.2]{P51}}), so that $u\mapsto\psi_T(M^{-1}\sigma(u))$ is the
component at $T$ of the character $\psi_f$ of {\cite[\S3.3.2]{P51}}.

\emph{The function $\pffch$.} For $u\in U_T\cap\pffKM$ we have $\sigma(u)\in M\pffOT$, so
$\psi_T(M^{-1}\sigma(u))=1$. Hence
\begin{equation}\label{pff:eq-chi}
\pffch(uk)=\psi_T\bigl(M^{-1}\sigma(u)\bigr)\quad(u\in U_T,\ k\in\pffKM),\qquad\pffch=0\text{ off }U_T\pffKM,
\end{equation}
is well defined; $\pffch$ is locally constant, $|\pffch|\le1$, $\pffch(ug)=\psi_T(M^{-1}\sigma(u))
\pffch(g)$ for $u\in U_T$, and $\pffch(gk)=\pffch(g)$ for $k\in\pffKM$. It is the function $\chi_M$
of {\cite[\S3.3.7]{P51}}.

\emph{Intermediate groups.} An \emph{intermediate group} is a subgroup $\pffZ$ with $\pffKM\subseteq
\pffZ\subseteq\pffK$; it is compact open, normalises $\pffKM$, and has finitely many cosets
$\pffKM a$. Choosing a representative $a_c\in\pffZ$ of each coset $c$, the \emph{Bessel vector} of
$\pffZ$ is
\[
\pffW\colon\pffG\to\C^{\pffKM\backslash\pffZ},\qquad\pffW(y)_c=\pffch(ya_c^{-1}).
\]
By the right invariance of $\pffch$ under $\pffKM$ and the normality of $\pffKM$ in $\pffK$, a
change of representatives multiplies no coordinate, so $\pffW$ does not depend on the choices. It
satisfies $\pffW(uy)=\psi_T(M^{-1}\sigma(u))\pffW(y)$ for $u\in U_T$, and $\pffW(y)\ne0$ only if
$uy\in\pffK$ for some $u\in U_T$: it is a \emph{cut Bessel function} in the sense of
\S\ref{pfr:ss-setting}, for the box $\pffK$ and the level $M$.

\emph{Measures.} We fix on $\pffG$ the Haar measure $\pbvol=\pfrM$ of \S\ref{pfr:ss-setting}
(the \emph{group measure}, built floor by floor from the frame), for the place $(\pffET,dx,\psi_T)$
and the lattice $L_T$. The group $\pffG$ is unimodular, so $\pbvol$ is left and right invariant.
We use the following properties of it, all from \S\ref{pfr:sec}.
\begin{enumerate}
\item[(M1)] \emph{(Unfolding along $U_T$.)} There is a Haar measure $du$ on $U_T$ with
$\int_{U_T\cap\pffK}du=1$, and a measure $\theta$ on $\pffG$ (standing for the quotient
$U_T\backslash\pffG$), such that $\int_{\pffG}F\,d\pbvol=\int_{\pffG}\int_{U_T}F(uy)\,du\,
d\theta(y)$ for $F\ge0$ measurable. (Here $\theta=\theta_{\ge0}$ and $du=\eta_{[0,\varrho+2)}$, which is
carried by $U_T$; see \S\ref{pff:ss-large}.)
\item[(M2)] \emph{(Cosets of $\pffKM$ in $U_T$.)} For $y\in\pffG$, the set $\{u\in U_T:uy\in
\pffKM\}$ is empty or a coset of $U_T\cap\pffKM$; we write
\[
\pffl=\int_{U_T\cap\pffKM}du ,
\]
the \emph{unipotent mass}. (This is the product of the floor moduli of $M$ in
\S\ref{pfr:sec}.)
\end{enumerate}
We write $\pffm(\pffZ)$ for the \emph{mass integral} and $\pffr(\pffZ)$ for the
\emph{normalised mass} of an intermediate group:
\begin{equation}\label{pff:eq-mass}
\pffm(\pffZ)=\int_{\pffG}\Bigl|\int_{\pffZ}\pffch(yxy^{-1})\,dy\Bigr|^2dx,\qquad
\pffr(\pffZ)=\frac{\pbvol(\pffZ)\,\pbvol(\pffKM)}{\pffl^{\,2}},
\end{equation}
both integrals for $\pbvol$. The function $x\mapsto\int_{\pffZ}\pffch(yxy^{-1})dy$ is the
function $\psi^{\natural}$ of {\cite[Lemma~3.33(c)]{P51}} for the $p$-adic group, up to the factor
$\pbvol(\pffZ)$; compare {\cite[Remark~3.40(c)]{P51}}.

\subsection{The Bessel vector and the mass integral}\label{pff:ss-bessel}

Throughout, $\pffZ$ is an intermediate group and $\pffW$ its Bessel vector.

\begin{lemma}[Unfolding the Bessel vector]\label{pff:unfold}
For $x,y\in\pffG$,
\[
\pffl\cdot\langle\pffW(y),\pffW(yx)\rangle=\int_{U_T}\mathbf 1_{\pffZ}(uy)\,\pffch\bigl(uy\,x\,(uy)^{-1}
\bigr)\,du ,
\]
where $\langle a,b\rangle=\sum_c\aubar{a_c}\,b_c$.
\end{lemma}
\begin{proof}
For a coset $c$ let $S_c=\{u\in U_T:uya_c^{-1}\in\pffKM\}$, and $\alpha_c=\aubar{\pffW(y)_c}\,
\pffW(yx)_c$. We claim that for every $u\in U_T$
\begin{equation}\label{pff:eq-unf}
\mathbf 1_{\pffZ}(uy)\,\pffch\bigl(uyx(uy)^{-1}\bigr)=\sum_c\mathbf 1_{S_c}(u)\,\alpha_c .
\end{equation}
If $uy\notin\pffZ$, then $u\notin S_c$ for all $c$, since $uya_c^{-1}\in\pffKM\subseteq\pffZ$ would
give $uy\in\pffZ$. If $uy\in\pffZ$, then $uy$ lies in exactly one coset $\pffKM a_d$, so $u\in S_c$
exactly for $c=d$. Put $k=uya_d^{-1}\in\pffKM$. Then $\pffW(y)_d=\pffch(u^{-1}k)=\psi_T(M^{-1}
\sigma(u^{-1}))$, and
\[
\pffW(yx)_d=\pffch(yxa_d^{-1})=\pffch\bigl(u^{-1}\cdot uyx(uy)^{-1}\cdot k\bigr)=\psi_T(M^{-1}\sigma(
u^{-1}))\,\pffch\bigl(uyx(uy)^{-1}\bigr),
\]
by the two invariance properties of $\pffch$. As $|\psi_T|=1$, $\alpha_d=\pffch(uyx(uy)^{-1})$, which
proves \eqref{pff:eq-unf}. By (M2), $S_c$ is empty or a coset of $U_T\cap\pffKM$, of measure
$\pffl$; and if $S_c$ is empty then $\pffW(y)_c=0$ (for $\pffW(y)_c\ne0$ means $ya_c^{-1}\in
U_T\pffKM$, that is $S_c\ne\emptyset$), so $\alpha_c=0$. Integrating \eqref{pff:eq-unf} over
$U_T$ gives $\sum_c\pffl\,\alpha_c=\pffl\langle\pffW(y),\pffW(yx)\rangle$.
\end{proof}

\begin{corollary}\label{pff:k0}
Put $k_0(x)=\int\langle\pffW(y),\pffW(yx)\rangle\,d\theta(y)$ (the \emph{correlation from floor
$0$}, an absolutely convergent integral). Then for all $x\in\pffG$
\[
\int_{\pffZ}\pffch(yxy^{-1})\,dy=\pffl\cdot k_0(x),\qquad\text{hence}\qquad\pffm(\pffZ)=\pffl^{\,2}
\int_{\pffG}|k_0(x)|^2\,dx .
\]
Moreover $\pffl\,\|\pffW(y)\|^2\le1$ for all $y$, and $\pffl\int\|\pffW(y)\|^2d\theta(y)=
\pbvol(\pffZ)$.
\end{corollary}
\begin{proof}
The function $F(z)=\mathbf 1_{\pffZ}(z)\pffch(zxz^{-1})$ is bounded with compact support. By (M1)
and Fubini's theorem, $\int_{\pffZ}\pffch(yxy^{-1})dy=\int\int_{U_T}F(uy)\,du\,d\theta(y)$, and the
inner integral is $\pffl\langle\pffW(y),\pffW(yx)\rangle$ by Lemma~\ref{pff:unfold}. For $x=1$,
$\pffch(1)=1$ and Lemma~\ref{pff:unfold} gives $\pffl\|\pffW(y)\|^2=\int_{U_T}\mathbf 1_{\pffZ}(uy)\,du$.
This is at most $1$: the set $\{u:uy\in\pffZ\}$ lies in $\{u:uy\in\pffK\}$, which is empty or a coset
of $U_T\cap\pffK$, of measure $1$ by (M1). Integrating over $\theta$ and using (M1) for
$F=\mathbf 1_{\pffZ}$ gives the last identity; it also shows that the integral defining $k_0$
converges absolutely, by the Cauchy--Schwarz inequality.
\end{proof}

\begin{lemma}[The positions in $\pffKM$]\label{pff:inKM}
$\displaystyle\int_{\pffKM}\Bigl|\int_{\pffZ}\pffch(yxy^{-1})\,dy\Bigr|^2dx=\pbvol(\pffZ)^2\,
\pbvol(\pffKM)$.
\end{lemma}
\begin{proof}
For $x\in\pffKM$ and $y\in\pffZ$ we have $yxy^{-1}\in\pffKM$, as $\pffZ\subseteq\pffK$ normalises
$\pffKM$; so $\pffch(yxy^{-1})=\pffch(1)=1$ and the inner integral is $\pbvol(\pffZ)$.
\end{proof}

\begin{lemma}[Volumes]\label{pff:vol}
Assume the conditions \textup{(L1)}--\textup{(L3)} of \S\ref{pff:ss-large} below. Then:
\begin{enumerate}
\item[(a)] $\pbvol(\pffK)\le1+2^{-n}$;
\item[(b)] if the non-units of $\pffOT$ have volume at most $\frac12$, then $\pbvol(\pffK)\ge
2^{-n}$;
\item[(c)] for an intermediate group $\pffZ$, $\pbvol(\pffZ)=[\pffZ:\pffKM]\,\pbvol(\pffKM)$.
\end{enumerate}
\end{lemma}
\begin{proof}
(c) is the invariance of $\pbvol$. (a) and (b) are proved in \S\ref{pff:ss-floors}
(Lemma~\ref{pff:floormass}); in outline: in the unfolding of $\pbvol$ floor by floor, the frame vector of floor
$l$ ranges over a quadric of $\pffOT$-points of rank $2(m'+1)+\#\iota$, whose volume is at most the
ideal zeta sum $\zeta_{\pffOT}(m')=\sum_{I}[\pffOT:I]^{-m'}$ (sum over the ideals of finite index)
and at least the volume of $\pffOT^\times$; and $\zeta_{\pffOT}(m')^n\le1+2^{-n}$ is (L2).
\end{proof}

\subsection{The datum at $T$; positions outside $\pffKM$}\label{pff:ss-large}

The group measure and the common theory of \S\ref{pfr:sec} are applied at $T$ to a basis of $L_T$
adapted to the frame. We record the conditions under which the statements of \S\ref{pfr:sec} are
quoted (\emph{large lattice datum}); $m'=3\varrho+5$ and $\iota=k-6\varrho-13$.
\begin{enumerate}
\item[(L1)] $L_T$ has an $\pffOT$-basis consisting of the frame and dual frame
$e_1,\dots,e_\varrho$, the hyperbolic pair $(u,u')=(-v,z-N\tau v)$ spanning the plane $\langle v,z\rangle$
of Lemma~\ref{pff:L1} (so the frame vector $z=u'-N\tau u$ lies in this plane and has coordinate $1$ on $u'$),
$m'+1$ further hyperbolic pairs, $\iota$ orthonormal vectors, and $e_1^*,\dots,e^*_\varrho$; so
$k+1=\iota+2(m'+1)+2$ and $r=2\varrho+2+2(m'+1)+\iota$.
\item[(L2)] $\zeta_{\pffOT}(m')^n\le1+2^{-n}$, where $\zeta_{\pffOT}(s)=\sum_I[\pffOT:I]^{-s}=
\prod_{w\in T}(1-q_w^{-s})^{-1}$, the sum over the ideals of finite index; and $m'<9n$, $53n\le
\iota$.
\item[(L3)] $M\in\pffOT$, $\aubar M=M$, $\psi_T$ is trivial on $\pffOT$ and on the elements
$x+\aubar x$, and the hermitian form is unimodular on $L_T$.
\end{enumerate}
For the data, (L3) was seen in \S\ref{pff:ss-setting}. In (L2), $m'<9n$ and $53n\le\iota$ hold as
$k\ge x^5/2$ and $n\le2\sqrt x$; and, as $\zeta(s)\le1+2^{2-s}$ for $s\ge2$ and the number of places
of $E$ above a rational prime is at most $2f$,
\[
\log\zeta_{\pffOT}(m')^n\le n\cdot2f\cdot2^{2-m'}\le2^{-\varrho}\cdot fn\,2^{-2\varrho-1}\le
\log(1+2^{-n}),
\]
because $fn\le x^{13}\le2^{2\varrho-1}$ for the data ($\varrho\ge\sqrt x/2-1$ and $(\log b)^{100}\le
x$). Condition (L1) is Lemma~\ref{pff:L1}.

\subsubsection*{The finite place as a place datum}
We now place ourselves in the framework of \S\ref{pfr:ss-setting}, with the finite place
$A=\pffET$ ($\pfrmod(c)=|c|_T$, Plancherel constant $P=1$: axiom (P) of \S\ref{pfs:ss-planch} holds with $P=1$, Lemma~\ref{pfg:P}), the basis of (L1), and the place datum
$\mathfrak F_T$ whose box is $\pffK$, whose close pairs are the pairs $(\omega_1,\omega_2)$ with $\omega_1\in
L_T$ and $M^{-1}(\omega_1-\omega_2)\in L_T$, whose positions are the elements $y^{-1}uyk$
($y\in\pffG$, $u\in U_T$, $k\in\pffKM$), whose level is $M$, and whose measures are as follows.
\begin{itemize}
\item $\mu_H$ is the right Haar measure of the closed subgroup $H_T$ normalised by $\mu_H(H_T\cap\pffK)=1$. For every $g$
the $h\in H_T$ with $uhg\in\pffK$ for some $u\in U_T$ lie in one right coset of $H_T\cap\pffK$: if $u_1h_1g,u_2h_2g\in\pffK$, then
$(u_1h_1g)(u_2h_2g)^{-1}=u'h_1h_2^{-1}\in\pffK$ with $u'=u_1(h_1h_2^{-1}u_2^{-1}h_2h_1^{-1})\in U_T$ ($H_T$ normalises $U_T$), and in the
basis (L1) an element $u'h$ ($u'\in U_T$, $h\in H_T$) is block upper triangular with the diagonal blocks of $h$, as is its
inverse $h^{-1}u'^{-1}$ with the diagonal blocks of $h^{-1}$; so $u'h\in\pffK$ forces $h,h^{-1}$ integral, that is $h\in\pffK$; so the condition $\mu_H\{h:\exists u,\ uhg\in\pffK\}\le1$ holds.
\item The floor measures $\theta_l$, $l<\varrho$, are those of \S\ref{pfr:ss-setting}: $\theta_l=\lambda_l*\eta_l$ with
$\lambda_l$ the image of $|c|_T^{d_l}dc$ under $c\mapsto w_{l,c}^{-1}$, and $\eta_l$ the image of $d\omega\otimes d\beta_{\rm sk}$,
where $d\omega$ gives the lattice of $V_{l+1}$ volume $1$ and $d\beta_{\rm sk}$ gives the integral skew elements volume $1$
(Lemma~\ref{pfg:duality}(d)). By Lemma~\ref{pfr:floorquad} ($\kappa_{E_T}=1$) the image of $\theta_l$ under
$\sigma\mapsto\sigma^{-1}f_l$ is exactly the quadric measure $\mu_l$: $c_l=1$. Here $\mu_l$ is the fibre measure at $0$ of
Proposition~\ref{pfg:fibres} for the unimodular lattice of $V_l$ spanned by the basis vectors of floor $\ge l$.
\item The quadric measure of the frame vector $z$ is $\mu_\varrho=\nu^+_{-N}$, the fibre measure at $h(z,z)=-N$ of
\S\ref{pfg:ss-fibre} for the lattice $L^+$ of $V_\varrho=\pbmXp_T$; it is carried by the quadric $\pfgQ$, which is the orbit
$\{\sigma^{-1}z:\sigma\in\pffG^{(\varrho)}\}$ (Lemma~\ref{pfg:orbit}). Let $s\colon\pfgQ\to\pffG^{(\varrho)}$ be a Borel map
with $s(e)^{-1}z=e$ (Lemma~\ref{pfs:section} for $b=\varrho$). The \emph{top floor measure} is
\[
\theta_\varrho=s_*\mu_\varrho,\qquad\int F\,d\theta_\varrho=\int_{\pfgQ} F(s(e))\,d\mu_\varrho(e)\quad(F\ge0\text{ Borel}).
\]
It is finite on compact sets ($s^{-1}(C)\subseteq\{\sigma^{-1}z:\sigma\in C\}$, a compact set), and its image under
$\sigma\mapsto\sigma^{-1}z$ is $\mu_\varrho$: $c_\varrho=1$. The stabiliser of $z$ in $\pffG^{(\varrho)}$ is $H_T=\pfrR^{(\varrho+1)}\cap
\pffG^{(\varrho)}$, and $s(\sigma^{-1}z)\in H_T\sigma$. Every function integrated against $\theta_\varrho$ below is invariant under
left translation by $H_T$ (the correlations from floor $\varrho+1$ are $\mu_H$-averages, $\mu_H$ being right invariant
under $H_T$; nearness at $\varrho+1$ and the weights are $\pfrR^{(\varrho+1)}$-invariant), and for such $F$,
$\int F\,d\theta_\varrho=\int F_0\,d\mu_\varrho$ with $F(\sigma)=F_0(\sigma^{-1}z)$; so these integrals do not depend on the
choice of $s$.
\end{itemize}
\emph{The group measure is a Haar measure.} The group measure $\pfrM$ of $\mathfrak F_T$ is right invariant by
Lemma~\ref{pfr:upperinv} for $j=0$ (whose proof uses only the structure of the floors and the invariance of the $\mu_l$);
$0<\pfrM(\pffK)<\infty$ by Lemma~\ref{pff:floormass}(b),(c), whose proofs do not use that $\pfrM$ is a Haar measure; a compact set
is covered by finitely many right cosets $\pffK g$, each of measure $\pfrM(\pffK)$; so $\pfrM$ is a non-zero right invariant
Radon measure, hence a right Haar measure, hence a Haar measure ($\pffG$ is unimodular). It is the measure $\pbvol$ of
\S\ref{pff:ss-setting}. (M1) is the decomposition $\pfrM=\pfrm\ast\theta_{\ge0}$
with $\pfrm=\eta_{[0,\varrho+2)}$ carried by $U_T$, and (M2) holds with $\pffl=\prod_{l<\varrho}
\eta_l(\pfrN_l\cap\pffKM)$ (each Heisenberg coordinate lies in $M\pffOT$ or $\pffOT$).
The Bessel vector $\pffW$ is a cut Bessel function for $\mathfrak F_T$ (\S\ref{pff:ss-setting}).

\begin{lemma}[Gauges at $T$]\label{pff:gauge}
Let $\alpha\in\pffOT$ be a unit of $\pffET$ with $\aubar\alpha=\alpha$. Put $\Gamma_\alpha(i)=\{c\in
\pffET:\alpha^ic\in\pffOT\}$ and $\Lambda_\alpha=2/|\alpha|_T$. Under (L1)--(L3), $\Gamma_\alpha$ with
the loss $\Lambda_\alpha$ is a gauge for $\mathfrak F_T$ (\S\ref{pfr:ss-gauge}, axioms (G0)--(G5)). For
$\alpha=1$ the sets $\Gamma_1(i)$ are all $\pffOT$ and the loss is $2$.
\end{lemma}
\begin{proof}
This is Lemma~\ref{pfg:gauge}, where (G0)--(G5) are verified in their present form, including the clause
$\tau_0\delta^{-1}\in\Gamma(1)$ of (G1) (with $\tau_0$, $\delta$ of Lemma~\ref{pfg:psi}), the count (G4) from the
Gauss sum bound of Lemma~\ref{pfg:gauss} and Proposition~\ref{pfg:fibres}, and the count (G5) from
Lemmas~\ref{pfg:common}, \ref{pfg:Dint}, \ref{pfg:content}.
\end{proof}

\begin{lemma}[Depth at $T$]\label{pff:dideal}
Let $a<\varrho$, $x,y\in\pffG$.
\begin{enumerate}
\item[(a)] $\mathfrak d_a(y,x)=c\,\pffOT$ for a unit $c$ of $\pffET$; if $y$ is near at $a+1$, then
$c\in\pffOT$ and $\operatorname{dp}_a(y,x)=|c|_T^{-1}=[\pffOT:c\pffOT]$.
\item[(b)] If $\langle\pffW(y),\pffW(z)\rangle\ne0$ for some $z$, then $y$ is near at every floor and
$\mathfrak d_a(y,x)$ contains a $c\in\pffOT$ with $\pbvol\mathfrak d_a(y,x)=|c|_T$.
\end{enumerate}
\end{lemma}
\begin{proof}
(a) The pairs $(w,d)$ with $w$ stable at $(y,yx)$ for $a$ and $w+d$ its second lift are the solutions
of $X(w,d)\in L_T\oplus L_T$ with $w\perp e_b$, $d\in\langle e_b:b<a\rangle$ ($b<a$), for the invertible
matrix $X=\bigl(\begin{smallmatrix}y^{-1}&0\\M^{-1}(y^{-1}-(yx)^{-1})&-M^{-1}(yx)^{-1}\end{smallmatrix}\bigr)$.
The set of the values of one coordinate (here the $e_a^*$-coordinate of $w$) on the solutions of an
integrality condition $X v\in\pffOT^{2r}$ with $X$ invertible, subject to the vanishing of the
coordinates of larger weight, is an $\pffOT$-module that is open and compact, i.e.\ $c\pffOT$ for a unit
$c$ ($\pffOT$ is a finite product of discrete valuation rings, hence a principal ideal ring). If
$\xi y\in\pffK$ with $\xi\in R^{(a+1)}$, then $h(\xi y\cdot y^{-1}w,e_a)$-type coordinates are
integral: the $e_a^*$-coordinate of $w$ equals that of $\xi w$ ($\xi$ moves $e_b$, $b\le a$, by
combinations of the $e_{b'}$, $b'<b$, and $w\perp e_b$ for $b<a$), and $\xi w=(\xi y)(y^{-1}w)$ is
integral. So $c\in\pffOT$, and the volume of $c\pffOT$ is $|c|_T$.
(b) By the definition of $\pffch$, $\pffW(y)\ne0$ gives $u\in U_T$ with $uy\in\pffK$, so $y$ is near
at every floor ($U_T\subseteq R^{(l)}$); (a) applies.
\end{proof}

\begin{lemma}[Two gauges]\label{pff:twogauges}
Let
$\Gamma$, $\Gamma'$ be gauges with $\Gamma(10n)\subseteq\Gamma'(1)$ and $\pffOT\subseteq\Gamma'(2)$ (for instance any gauges $\Gamma_\alpha$ of Lemma~\ref{pff:gauge}), $p,q$ heads of length $j$ for
$\Gamma$, $M^{-1}\in\Gamma(1)$ if $j\ne0$. Let $y$ be active at $x$ with $\langle\pffW(y),
\pffW(yx)\rangle\ne0$, and suppose that for every $a$ with $j\le a<\varrho$ the ideal $\mathfrak
d_a(y,x)$ contains a unit $d_a$ with $d_a^{-1}\in\Gamma'(1)$. Then $M^{-1}(x-1)$ has entries in
$\Gamma'(99n)$.
\end{lemma}
(Proved in Corollary~\ref{pfg:twogaugesT}, from Theorem~\ref{pfg:twogauges}.)

\begin{lemma}[Dyadic shells]\label{pff:dyadic}
For $x$ a position, let $\mathfrak J(x)\subseteq\pffOT$ be the ideal of the $\varepsilon\in\pffOT$ with
$\varepsilon M^{-1}(x-1)$ integral, and for $i\ge0$ let $S_i$ be the set of positions with
$\lceil\log_2[\pffOT:\mathfrak J(x)]\rceil=i$. Let $\Gamma=\Gamma_\alpha$, $p,q$ heads of length $j\le\varrho$
for $\Gamma$, and $M^{-1}\in\Gamma(1)$ if $j\ne0$. Then $S_i$ is a shell of height $2^i$ for
$\Gamma,j,p,q,\pffW$ (Definition~\ref{pfr:shell}).
\end{lemma}
\begin{proof}
(a) By (L2) the number of ideals of index at most $H=2^i$ is at most $\zeta_{\pffOT}(m')H^{m'}\le2H^{9n}$.
For each such ideal $\mathfrak J=c\pffOT$, the gauge $\Gamma_{\alpha\aubar cc}$ has loss
$\Lambda|c|_T^{-2}\le\Lambda H^2$, and contains in its position ball the positions with $\mathfrak J(x)=
\mathfrak J$, since $\aubar c\,c\,M^{-1}(x-1)$ is integral. (b) Let $x\in S_i$, $\mathfrak J(x)=c\pffOT$,
$[\pffOT:\mathfrak J(x)]\ge H/2$, and $\varsigma<1$ with $\Lambda^{1200n^2}\le\varsigma^{400n^2}H$; then,
as $2\le\Lambda$, $|c|_T\Lambda^{990n^2}\le\varsigma^{400n^2}$. Let $y$ be active at $x$ and suppose that
$\pbvol\mathfrak d_a(y,x)>\varsigma^2$ for all $a\ge j$. By Lemma~\ref{pff:dideal}(b) choose $d_a\in\pffOT
\cap\mathfrak d_a(y,x)$ with $|d_a|_T\ge\varsigma^2$, and let $C=\prod_{a\ge j}d_a\aubar d_a$, so that
$|C|_T\ge\varsigma^{4n}$, $\aubar C=C$ and $C/d_a\in\pffOT$. For $\Gamma'=\Gamma_{\alpha^{10n}C}$ we have
$\Gamma(10n)\subseteq\Gamma'(1)$ and $d_a^{-1}\in\Gamma'(1)$, so by Lemma~\ref{pff:twogauges}
$(\alpha^{10n}C)^{99n}$ lies in $\mathfrak J(x)=c\pffOT$. Taking modules, $|C|_T^{99n}\le|c|_T|\alpha|_T^{-990n^2}\le|c|_T\Lambda^{990n^2}\le\varsigma^{400n^2}<
\varsigma^{396n^2}\le|C|_T^{99n}$, a contradiction. So some $a\ge j$ has depth $\ge\varsigma^{-2}$.
\end{proof}

\begin{lemma}[A deep vector outside $\pffKM$]\label{pff:deep}
Let $Q\ge2$ be such that $|c|_T\le Q^{-1}$ for every $c\in\pffOT\cap\pffET^\times$ with
$c^{-1}\notin\pffOT$. If $x\notin\pffKM$, then for the heads of length $0$ every point $y$ active at $x$
with $\langle\pffW(y),\pffW(yx)\rangle\ne0$ has a frame vector of depth at least $Q$.
\end{lemma}
\begin{proof}
Otherwise, by Lemma~\ref{pff:dideal}, each $\mathfrak d_a(y,x)$ contains $c_a\in\pffOT$ with
$|c_a|_T=\pbvol\mathfrak d_a>Q^{-1}$, hence $c_a^{-1}\in\pffOT$. Apply Lemma~\ref{pff:twogauges} with
$\Gamma=\Gamma'=\Gamma_1$ (all sizes mean integral): $M^{-1}(x-1)$ is integral. As $x$ is unitary and the form
unimodular, $x^{-1}=\mathsf J^{-1}{}^t\aubar x\mathsf J$ is integral, so $x\in\pffKM$.
\end{proof}
Active points with $\langle\pffW(y),\pffW(yx)\rangle=0$ contribute nothing to $\pfrF_0^+(x)$, whose
integrand involves this pairing; so only the points covered by Lemma~\ref{pff:deep} matter below.

\begin{proposition}[The positions outside $\pffKM$]\label{pff:outside}
Let $Q$ be as in Lemma~\ref{pff:deep} and $t\ge0$ with $t+2\cdot10^9n^4\le k$. Then
\[
\int_{\pffG\smallsetminus\pffKM}\Bigl|\int_{\pffZ}\pffch(yxy^{-1})\,dy\Bigr|^2dx\le Q^{-t}\,\pbvol(
\pffZ)\,\pbvol(\pffKM).
\]
\end{proposition}
\begin{proof}
At floor $0$ with the heads $1,1$ the majorant is $\pfrF_0^+(x)=\pfrB^1_0(1,x)^2\ge|k_0(x)|^2$
(Lemma~\ref{pfr:L1} and the definition of $k_0$ in Corollary~\ref{pff:k0}). Let $f=\mathbf
1_{\pffG\smallsetminus\pffKM}\pfrF^+_0$. In the gauge $\Gamma_1$ (loss $2$) the mass of floor $0$ is at most
$V=2^{1287n}$ (Lemma~\ref{pff:floormass}), and by Lemma~\ref{pff:deep} and the last assertion of Theorem~\ref{pfr:deep}
(the case $j=0$) with $\varsigma=Q^{-1}$, $\pfrF^+_0(x)^2\le Q^{-2s}\mathsf Q(x)$ for $x\notin\pffKM$, so that $f$ has the
gain $Q^{-s}\le VQ^{-s}$ at every position, where $s\ge9\cdot10^7n^4+r'$ for $r'=t+10^9n^4$ (as $s=\iota-2(2m'+2)-
(3n+1+m')$ and $k$ is large). By Lemma~\ref{pff:dyadic} and Theorem~\ref{pfr:shells}
\[
\int f\,d\pbvol\le V\,Q^{-t-10^9n^4}\,2^{4\cdot10^8n^4}\,\mathsf N(\pffW)\le Q^{-t}\pffr(\pffZ),
\]
using $Q\ge2$ and $\mathsf N(\pffW)\le\pffr(\pffZ)$; the last inequality follows from
$\pffl\iint|\langle\pffW(g),\pffW(y)\rangle|^2d\theta\,d\theta\le\pffr(\pffZ)$ (for fixed $g$, the $y$
with a non-zero pairing lie in a set of $\theta$-mass $\pbvol(\pffKM)/\pffl$, and
$\pffl|\langle\pffW(g),\pffW(y)\rangle|^2\le\|\pffW(g)\|^2$ by Corollary~\ref{pff:k0}; then integrate in $g$)
and $\prod_{l<\varrho}\mathsf c_l(M)=\pffl$ at a finite place (the floor moduli). Multiply by
$\pffl^{\,2}$ and use Corollary~\ref{pff:k0}.
\end{proof}

\begin{theorem}[The mass integral]\label{pff:massZ}
Let $Q$, $t$ be as in Proposition~\ref{pff:outside}. For every intermediate group $\pffZ$,
\[
\pbvol(\pffZ)^2\,\pbvol(\pffKM)\ \le\ \pffm(\pffZ)\ \le\ \pbvol(\pffZ)\,\pbvol(\pffKM)\,\bigl(\pbvol(
\pffZ)+Q^{-t}\bigr).
\]
\end{theorem}
\begin{proof}
Split the outer integral of $\pffm(\pffZ)$ into $\pffKM$ and its complement, and apply
Lemma~\ref{pff:inKM} and Proposition~\ref{pff:outside}; the lower bound is Lemma~\ref{pff:inKM}
alone.
\end{proof}

\begin{corollary}\label{pff:massK}
If $b\ge1000$ and $2^{-\varrho}\le b^{-6}$, then $\pffm(\pffK)\le(1+b^{-6})\,\pbvol(\pffK)\,
\pbvol(\pffKM)$.
\end{corollary}
\begin{proof}
Take $Q=2$: a $c\in\pffOT$ with $c^{-1}\notin\pffOT$ generates a proper ideal, of index
$|c|_T^{-1}\ge2$. Take $t=n$; then $t+2\cdot10^9n^4\le k$. By Theorem~\ref{pff:massZ} and
Lemma~\ref{pff:vol}(a), $\pffm(\pffK)\le\pbvol(\pffK)\pbvol(\pffKM)(1+2^{-n}+2^{-n})$, and
$1+2^{1-n}=1+2^{-\varrho}\le1+b^{-6}$.
\end{proof}

For the data, $2^{-\varrho}\le b^{-6}$ holds with room: $\varrho\ge\sqrt x/2-1$ and $(\log b)^{100}
\le x$ ({\cite[Lemma~13.4]{P46}}).

\subsection{The basis at $T$, floor masses and volumes}\label{pff:ss-floors}

\begin{lemma}[The basis (L1)]\label{pff:L1}
For the data, $L_T$ has a basis as in \textup{(L1)}, consisting of integral vectors with integral
dual basis.
\end{lemma}
\begin{proof}
In the lattice basis the middle block $X_0\perp Ez$ has the vectors $x_i$ $(i\ne k_0)$, orthonormal,
and $v=\aubar\pi_N^{-1}x_{k_0}+N^{-1}z$, $z$, with $h(v,v)=0$, $h(v,z)=-1$, $h(z,z)=-N$. The field
$E/E^+$ is unramified at every finite place ({\cite[\S2.1]{P51}}, data of
{\cite[\S14.1]{P46}}), so at every $w\in T$ the norm map $\auO_w^\times\to\auO_{E^+,v}^\times$ and the
trace $\auO_w\to\auO_{E^+,v}$ are onto (for a split place both are clear); hence there are
$c,\tau\in\pffOT$ with $\aubar cc=-1$ and $\tau+\aubar\tau=1$. Then $-v$ and $z-N\tau v$ form a
hyperbolic pair ($h(-v,z-N\tau v)=1$, $h(z-N\tau v,z-N\tau v)=-N+N\tau+N\aubar\tau=0$), and the second
vector has coordinate $1$ at $z$; for two orthonormal vectors $x,x'$, the vectors $p=x+cx'$ and
$q=\aubar\tau x-c\tau\,x'$ form a hyperbolic pair: $h(p,p)=1+\aubar cc=0$, $h(q,q)=(1+c\aubar c)\tau\aubar\tau=0$,
$h(p,q)=\tau+\aubar\tau=1$, and the change of basis $(x,x')\mapsto(p,q)$ has the matrix
$\bigl(\begin{smallmatrix}1&\aubar\tau\\c&-c\tau\end{smallmatrix}\bigr)$ of determinant $-c$ (the explicit basis and its
Gram matrix are written out in {\cite[\S3.1.6(ii)]{P51}}). Using $2(m'+1)$ of the $x_i$ this way and keeping
$\iota$ of them, the change of basis is integral with determinant a unit, as $c,\tau\in\pffOT$ and
$c^{-1}=-\aubar c$; and $k+1=\iota+2(m'+1)+2$.
\end{proof}

The places where the lattice basis is not integral with integral inverse: the matrix from the basis
$e_a,x_i,z,e_a^*$ to the lattice basis has the entries $\aubar\pi_N^{-1}$ and $(\pi_N\aubar\pi_N)^{-1}=
N^{-1}$, its inverse the entries $\aubar\pi_N$ and $-\pi_N^{-1}$; so both are integral at the places
not dividing $N$. This verifies the claim of \S\ref{pff:ss-setting}.

\begin{lemma}[Floor masses and the volume of $\pffK$]\label{pff:floormass}
Assume (L1)--(L3) and let $\Gamma=\Gamma_1$.
\begin{enumerate}
\item[(a)] For $l\le\varrho$ and every $g$, $\theta_l\{\sigma:\sigma g\text{ near at }l+1\}\le\zeta_{\pffOT}(m')$;
in the notation of \S\ref{pfr:ss-gauge}, the mass of every floor is at most $V=2\le2^{1287n}$ (the mass involves only
the box $\pffK$, not the gauge, so this holds for every gauge $\Gamma_\alpha$).
\item[(b)] $\pbvol(\pffK)\le\zeta_{\pffOT}(m')^n\le1+2^{-n}$.
\item[(c)] $\pbvol(\pffK)\ge\pbvol(\pffOT^\times)^n$; in particular $\pbvol(\pffK)\ge2^{-n}$ if the
non-units have volume at most $\frac12$.
\end{enumerate}
\end{lemma}
\begin{proof}
(a) Apply Lemma~\ref{pfr:chart} with $a=l$, the gauge $\Gamma_1$ (for which $\Gamma_1(7)$ is the set of integral
vectors), $c\equiv1$ and $\Psi_g\equiv1$: the mass is at most $c_l\,\mu_l(L_l)$, where $L_l$ is the lattice of the integral
vectors of $V_l$. Here $c_l=1$ (\S\ref{pff:ss-large}), and by Proposition~\ref{pfg:fibres}(e) for the unimodular lattice
$L_l$ (of rank $\kappa\ge k+1\ge2m'+1$), with $\mu=m'$, $g=1$, $x_0=0$: $\mu_l(L_l)\le\zeta_{\pffOT}(m')$. Finally
$\zeta_{\pffOT}(m')\le2$ by (L2).
(b) If $g\in\pffK$ then $g$ is near at every floor. Unfolding $\pbvol=\pfrM$ floor by floor
($\pfrM=\eta_{[0,\varrho+2)}\ast\theta_{\ge0}$ and the measures $\eta_l$ give $\pfrN_l\cap\pffK$ mass $1$),
$\pbvol(\pffK)\le\theta_{\ge0}\{g:g\text{ near at }\varrho+2\}\le\prod_{l<n}\zeta_{\pffOT}(m')$ by (a)
applied at each floor, and $\zeta_{\pffOT}(m')^n\le1+2^{-n}$ by (L2).
(c) $\pffK$ is a group, so for measures $\nu,\nu'$ on $\pffG$, $(\nu\ast\nu')(\pffK)\ge\nu(\pffK)\nu'(\pffK)$. As
$\eta_\varrho=\eta_{\varrho+1}=\delta_1$ and $\theta_{\varrho+1}=\mu_H$,
\[
\pfrM=(\mu_H\ast\theta_\varrho)\ast\eta_{\varrho-1}\ast\theta_{\varrho-1}\ast\dots\ast\eta_0\ast\theta_0,\qquad
\pbvol(\pffK)\ge(\mu_H\ast\theta_\varrho)(\pffK)\prod_{l<\varrho}\eta_l(\pffK)\,\theta_l(\pffK).
\]
\emph{Lower floors.} $\eta_l(\pffK)=\eta_l(\pfrN_l\cap\pffK)=1$: the element $n(\omega,b-\tau_0h(\omega,\omega))$ and its inverse
$n(-\omega,\bar\beta)$ are integral if and only if $\omega$ and $b$ are integral ($\tau_0h(\omega,\omega)$ is then integral, and the
lattice of $V_{l+1}$ is unimodular), a set of $d\omega\otimes d\beta_{\rm sk}$-measure $1$. For a unit $c$ of $\pffOT$,
$w_{l,c}^{-1}\in\pffK$; so $\theta_l(\pffK)\ge\int_{\pffOT^\times}|c|_T^{d_l}\,\eta_l(\pffK)\,dc=\pbvol(\pffOT^\times)$.

\emph{The top floor.} $(\mu_H\ast\theta_\varrho)(\pffK)=\int_{\pfgQ}\mu_H\{k\in H_T:ks(e)\in\pffK\}\,d\mu_\varrho(e)$. Use the hyperbolic
pair $(u,u')$ of (L1), with $z=u'-N\tau u$, and the orthogonal decomposition $L^+=\pffOT u\oplus\pffOT u'\oplus L'$, $L'$ the
span of the remaining basis vectors of $L^+$ (unimodular). Let $E_1\subseteq L^+$ be the compact open set of the
$e=au+bu'+y'$ ($a\in\pffOT$, $b\in\pffOT^\times$, $y'\in L'$). (i) \emph{For $e\in E_1\cap\pfgQ$ there is $\sigma_e\in\pffK\cap
\pffG^{(\varrho)}$ with $\sigma_e^{-1}z=e$.} For $\omega\in L'$ and integral $\beta$ with $\beta+\bar\beta=-h(\omega,\omega)$ let
$n_{\omega,\beta}$ be the isometry of $\pbmXp_T$ fixing $u$, sending $x\mapsto x-h(x,\omega)u$ on $L'\otimes\pffET$ and
$u'\mapsto u'+\omega+\beta u$; for a unit $c$ of $\pffOT$ let $d_c$ fix $L'$ and send $u\mapsto cu$, $u'\mapsto\bar c^{-1}u'$. These
are isometries (as in \S\ref{pfr:ss-floors}), preserve $L^+$ together with their inverses, and extended by the identity on
the $e_a,e_a^*$ they lie in $\pffK\cap\pffG^{(\varrho)}$ (by (L1), $L_T$ is the orthogonal sum of the planes $\langle e_a,e_a^*\rangle$
and $L^+$). With $\omega=-b^{-1}y'$ and $\beta=-\tau_0h(\omega,\omega)$, $n_{\omega,\beta}e=a_1u+bu'$ for some $a_1\in\pffOT$; then
$d_{\bar b}(a_1u+bu')=a_2u+u'$ with $a_2=a_1\bar b$; as these maps are isometries, $a_2+\bar a_2=h(e,e)=-N$. With the skew
integral $\beta'=-N\tau-a_2$ ($\beta'+\bar\beta'=-N+N=0$), $n_{0,\beta'}(a_2u+u')=-N\tau u+u'=z$. So $\gamma=n_{0,\beta'}d_{\bar b}
n_{\omega,\beta}\in\pffK\cap\pffG^{(\varrho)}$ has $\gamma e=z$, and $\sigma_e=\gamma$ works. (ii) For such $e$, $s(e)\in H_T\sigma_e$, say
$s(e)=k_0\sigma_e$; then $\{k:ks(e)\in\pffK\}=(H_T\cap\pffK\sigma_e^{-1})k_0^{-1}=(H_T\cap\pffK)k_0^{-1}$, of $\mu_H$-measure $1$ by
right invariance. Hence $(\mu_H\ast\theta_\varrho)(\pffK)\ge\mu_\varrho(E_1)$. (iii) $\mu_\varrho(E_1)=\pbvol(\pffOT^\times)$: by
Proposition~\ref{pfg:fibres}(a), $\mu_\varrho(E_1)=\lim_\ell A_\ell(\mathbf 1_{E_1};-N)$, and in the coordinates $(a,b,y')$, in which
$dy=da\,db\,dy'$ and $h(e,e)=a\bar b+\bar ab+h(y',y')$,
\[
A_\ell(\mathbf 1_{E_1};-N)=\frac1{\pbvol(\varpi_1^\ell\pfgOh)}\int_{\pffOT^\times}\int_{L'}\pbvol\{a\in\pffOT:\
a\bar b+\overline{a\bar b}\in-N-h(y',y')+\varpi_1^\ell\pfgOh\}\,dy'\,db=\pbvol(\pffOT^\times),
\]
because $a\mapsto a\bar b$ preserves the Haar measure of $\pffOT$ ($b$ a unit of $\pffOT$), the trace $x\mapsto x+\bar x$ maps
$\pffOT$ onto $\pfgOh$ (Lemma~\ref{pfg:duality}(a)) and so carries the Haar probability measure of $\pffOT$ to that of $\pfgOh$, and
$-N-h(y',y')\in\pfgOh$. So $(\mu_H\ast\theta_\varrho)(\pffK)\ge\pbvol(\pffOT^\times)$.

Altogether $\pbvol(\pffK)\ge\pbvol(\pffOT^\times)^{\varrho+1}=\pbvol(\pffOT^\times)^n$. If the non-units have volume $\le\frac12$,
$\pbvol(\pffOT^\times)\ge\frac12$.
\end{proof}

\subsubsection*{The density at the last entry}
For $l\in\pbmGf_T$ and a sequence of units $p$ let
\[
\mathcal O_p(l)=\pfrkap_\varrho(p,1)\int_{\pffG}\pffW(y)\otimes\aubar{\pffW\bigl((w_pl^{-1})^{-1}y\bigr)}
\,dy\ \in\C^{C\times C}\qquad(C=\pffKM\backslash\pffZ),
\]
the \emph{scaled orbital integral} at the element $w_pl^{-1}$ (head $p$, link $l^{-1}$), and
\[
\Pi_{p,q}(l^{-1}z)=\int_{H_T}\langle\mathcal O_p(kl),\mathcal O_q(kl)\rangle\,d\mu_H(k),
\]
the \emph{fibre pairing}; it depends only on $e=l^{-1}z$, as two $l$ with the same $e$ differ on the left
by an element of $H_T$ and $\mu_H$ is right invariant. Put $\Pi_{p,q}=0$ off the orbit of $z$.

\begin{definition}\label{pff:lastentry}
Let $\pi_N$ be the component at $T$ of the idele with $\pi_N\aubar\pi_N=N$. For heads $p,q$ of length
$n$ with $p_\varrho=q_\varrho=t$, the \emph{density at the last entry} is
\[
\mathcal D(p,q)=|\pi_N|_T^{-1}\int_{X_{0,T}}\Pi_{p,q}\bigl(x-N^{-1}t\,z\bigr)\,d\nu_{\beta(t)}(x),
\qquad\beta(t)=\frac{t\aubar t}N-N,
\]
where $\nu_\beta$ is the measure $dx/d\langle x,x\rangle$ on the sphere $\{x\in X_{0,T}:h(x,x)=\beta\}$ of
$X_{0,T}=X_0\otimes\pffET$ (the fibre measures of the disintegration of $dx$ along $x\mapsto h(x,x)$, constructed in Proposition~\ref{pfg:fibres}),
and the vector $x-N^{-1}tz$ is read in the lattice basis.
\end{definition}

The vector $e=x-N^{-1}tz$ satisfies $h(e,e)=h(x,x)-t\aubar t/N=-N$ and, by (Cv1), $h(e,z)=-N^{-1}t\,h(z,z)=t$: it runs over the
points of the quadric $h(e,e)=-N$ of $\pbmXp\otimes\pffET$ with last entry $t$.

\begin{lemma}[Properties of the fibre measures and of $\Pi$]\label{pff:fibre}
Assume (L1)--(L3).
\begin{enumerate}
\item[(a)] (\emph{Coarea.}) For $g$ continuous with compact support on $\pbmXp\otimes\pffET$,
$\int g\,d\theta_\varrho^{\sharp}=\int_{\pffET}|\pi_N|_T^{-1}\int g(x-N^{-1}tz)\,d\nu_{\beta(t)}(x)\,dt$, where
$\theta_\varrho^\sharp$ is the measure of the top quadric (image of $\theta_\varrho$ under $\sigma\mapsto
\sigma^{-1}z$); and $t\mapsto\int g(x-N^{-1}tz)d\nu_{\beta(t)}(x)$ is continuous.
\item[(b)] (\emph{Step function.}) There is a function $\Phi_{p,q}$ on $\pbmXp\otimes\pffET$, locally constant with
compact support and invariant under $H_T\cap\pffKM$, that equals $\Pi_{p,q}$ on the quadric $h(e,e)=-N$; and
$\Pi_{p',q'}=\Pi_{p,q}$ for all $(p',q')$ in a neighbourhood of $(p,q)$ ($\Pi$ does not depend on the entries of
index $\varrho$).
\end{enumerate}
\end{lemma}
\begin{proof}
(a) is Lemma~\ref{pfg:coarea}, with the fibre measures of Proposition~\ref{pfg:fibres} (Remark~\ref{pfg:rem-fibres});
(b) is Proposition~\ref{pfg:fibrepairing}(c),(d).
\end{proof}

\begin{proposition}[Top floor]\label{pff:top}
Assume (L1)--(L3), $k\ge10^{10}n^4$. For heads $p,q$ and $\xi\in\pffET$, writing $p^t,q^t$ for $p,q$ with
the last entry replaced by $t$:
\begin{enumerate}
\item[(a)] $t\mapsto\psi_T(\xi t)\mathcal D(p^t,q^t)$ is integrable and
$\int\psi_T(\xi t)\mathcal D(p^t,q^t)\,dt=\pfrI_\varrho\bigl(p,q;\ g\mapsto\psi_T(\xi h(gz,z))\bigr)$;
\item[(b)] for heads $p^{(1)},\dots,p^{(\kappa)}$ with a common last entry, the matrix
$(\mathcal D(p^{(a)},p^{(b)}))_{a,b}$ is hermitian and positive semidefinite;
\item[(c)] $(p,q)\mapsto\mathcal D(p,q)$ is continuous.
\end{enumerate}
\end{proposition}
\begin{proof}
(a) Write $\omega(g)=\psi_T(\xi h(gz,z))$, $\nu=\mu_H\ast\theta_\varrho$, $w=w_p$, $d=w_q^{-1}w_p$. Three facts are used.
(i) $\omega$ is invariant under left and right translation by $H_T$: $H_Tz=z$ and $h(kgz,z)=h(gz,k^{-1}z)$ for $k\in H_T$.
(ii) $\nu$ is invariant under right translation by $\pffG^{(\varrho)}$: for Borel $F\ge0$ put $F^H(\sigma)=\int F(h\sigma)\,
d\mu_H(h)$; it is left $H_T$-invariant ($\mu_H$ is right invariant), so $F^H(\sigma)=F_0(\sigma^{-1}z)$ with $F_0$ on $\pfgQ$, and
for $g\in\pffG^{(\varrho)}$, $\int F(lg)\,d\nu(l)=\int F_0(g^{-1}s(e)^{-1}z)\,d\mu_\varrho(e)=\int F_0(g^{-1}e)\,d\mu_\varrho(e)=\int F_0\,d\mu_\varrho
=\int F\,d\nu$, by the $\pffG^{(\varrho)}$-invariance of $\mu_\varrho$ (\S\ref{pfr:ss-setting}). Moreover $\nu$ is finite on compact
sets $C$: $\nu(C)=\int\mu_H(H_T\cap Cs(e)^{-1})\,d\mu_\varrho(e)$, the integrand vanishes unless $e\in C^{-1}z$ (if $hs(e)\in C$ then
$e=(hs(e))^{-1}z$), and $H_T\cap Cs(e)^{-1}\subseteq(H_T\cap CC^{-1})h_0$ for any of its elements $h_0$, so
$\nu(C)\le\mu_H(H_T\cap CC^{-1})\,\mu_\varrho(C^{-1}z)<\infty$.
(iii) $\pfrM=\nu\ast\pfrm_\varrho$, that is $\int F\,d\pfrM=\iint F(sg)\,d\nu(s)\,d\pfrm_\varrho(g)$: this is the definition of
$\pfrM$ and $\pfrm_\varrho$ (\S\ref{pfr:ss-setting}), as $\eta_\varrho=\eta_{\varrho+1}=\delta_1$ and $\theta_{\varrho+1}=\mu_H$; and
$\theta_{\ge\varrho+1}=\mu_H$, so $\mathsf k^{d'}_{\varrho+1}(g,y)=\int\pfrip{\pffW(hg)}{\pffW(d'hy)}\,d\mu_H(h)$.
Also $\theta_\varrho$, $\mu_H$ are carried by $\pffG^{(\varrho)}$, whose elements commute with $w_p$, $w_q$ and $d$.

By the definitions of \S\ref{pfr:ss-setting} with $j=\varrho$, $\pfrI_\varrho(p,q;\omega)$ is $\pfrkap_\varrho(p,q)$ times the integral over
$g$ ($\pfrm_\varrho$), $x$ ($\pfrM$), $\sigma,\tau$ ($\theta_\varrho$) and $h,h'$ ($\mu_H$) of
\[
\omega(\sigma\tau^{-1})\,\pfrip{\pffW(h\sigma g)}{\pffW(h\sigma gx)}\,\overline{\pfrip{\pffW(h'\tau w^{-1}g)}{\pffW(dh'\tau w^{-1}gx)}}.
\]
We first transform the integral of the absolute value of this function by Tonelli's theorem; the same steps, applied to the
function itself, are then justified by Fubini's theorem. Put $s=h\sigma$, $L=h'\tau s^{-1}\in\pffG^{(\varrho)}$. By (i),
$\omega(\sigma\tau^{-1})=\omega(s(h'\tau)^{-1})=\omega(L^{-1})$. Substitute $y=sgx$ in the $x$-integral ($\pfrM$ is a Haar measure of
the unimodular group $\pffG$): then $h'\tau w^{-1}g=Lw^{-1}sg$ and $dh'\tau w^{-1}gx=dLw^{-1}y=Lw_q^{-1}y$. Next, for fixed $s$,
substitute $l=(h'\tau)s^{-1}$ in the $(h',\tau)$-integral, which is an integral against $\nu$; by (ii) it preserves $\nu$. Finally
integrate over $(h,\sigma,g)$ with $u=sg$, which by (iii) is an integral against $\pfrM$. The result is
\[
\pfrI_\varrho(p,q;\omega)=\pfrkap_\varrho(p,q)\int\omega(l^{-1})\iint\pfrip{\pffW(u)}{\pffW(y)}\,
\overline{\pfrip{\pffW(lw_p^{-1}u)}{\pffW(lw_q^{-1}y)}}\,du\,dy\,d\nu(l),
\]
and the same with absolute values inside for the absolute integral. The latter is finite: by Lemma~\ref{pfg:Op}(a),(b) (for $p$
and for $q$) the inner integrand vanishes unless $u\in S_p$, $y\in S_q$ and $l\in C_p$, compact sets, it is bounded by (W3),
and $\nu(C_p)<\infty$ by (ii). So $\pfrI_\varrho(p,q;\omega)$ is defined, and the identity holds. With the pairing
$\langle a\otimes\bar b,c\otimes\bar d\rangle=\pfrip ac\,\overline{\pfrip bd}$ on $\C^{C\times C}$, the inner double integral is
$\langle\int\pffW(u)\otimes\overline{\pffW(lw_p^{-1}u)}\,du,\int\pffW(y)\otimes\overline{\pffW(lw_q^{-1}y)}\,dy\rangle$, and
$\pfrkap_\varrho(p,q)=\pfrkap_\varrho(p,1)\pfrkap_\varrho(q,1)$; hence
\[
\pfrI_\varrho(p,q;\omega)=\int\omega(l^{-1})\,\langle\mathcal O_p(l),\mathcal O_q(l)\rangle\,d\nu(l).
\]
Thus the argument of the weight is $l^{-1}$, where $l$ is the argument of the scaled orbital integrals: no conjugation of $\xi$
occurs. Writing $l=hs(e)$ ($h\in H_T$, $e\in\pfgQ$), so that $l^{-1}z=e$ and $\omega(l^{-1})=\psi_T(\xi h(e,z))$, and integrating
first over $h$ gives $\int\psi_T(\xi h(e,z))\,\Pi_{p,q}(e)\,d\mu_\varrho(e)$, where $\mu_\varrho=\theta_\varrho^\sharp$ ($c_\varrho=1$).
Lemma~\ref{pff:fibre}(a) with $g(e)=\psi_T(\xi h(e,z))\Phi_{p,q}(e)$ (locally constant with compact support by
Lemma~\ref{pff:fibre}(b); on the quadric $h(x-N^{-1}tz,z)=t$) gives $\int\psi_T(\xi t)\mathcal D(p^t,q^t)\,dt$, the function
$t\mapsto\psi_T(\xi t)\mathcal D(p^t,q^t)$ being continuous with compact support, hence integrable ($\Pi$ does not depend on the
last entry of the heads).
(b) For $u\in\C^\kappa$, $\sum_{a,b}\aubar u_au_b\Pi_{p^{(a)},p^{(b)}}(e)=\int_{H_T}\|\sum_bu_b\mathcal
O_{p^{(b)}}(kl)\|^2d\mu_H\ge0$, and hermitian symmetry is clear; integrate against the positive
measure $\nu_\beta$.
(c) Near $(p,q)$, by Lemma~\ref{pff:fibre}(b), $\mathcal D(p',q')=|\pi_N|_T^{-1}\int\Phi_{p,q}(x-N^{-1}
p'_\varrho z)\,d\nu_{\beta(p'_\varrho)}(x)$, continuous in $p'_\varrho$ by Lemma~\ref{pff:fibre}(a).
\end{proof}

\subsubsection*{The decay at $T$}
For a sequence $r$ and $K\ge0$ let $\Lambda_K(r)=\prod_{l\le i}[\delta_l\in\pffOT]\,\pbmn(\delta_l)^{K/2}$
($\delta_l=r_0\cdots r_{l-1}$), so that $\Lambda^T_i=\Lambda_{n^6}$ for heads of length $i$.

\begin{lemma}[Heads of a non-zero form]\label{pff:heads}
If the majorant $\pfrF^+_j$ of $p,q,\pffW$ is not identically zero, then $\delta_l(p),\delta_l(q)\in
\pffOT$ for $l\le j$, and there is $\alpha\in\pffOT$, $\aubar\alpha=\alpha$, such that in $\Gamma_\alpha$:
$p,q$ are heads of length $j$; $M^{-1}\in\Gamma_\alpha(1)$ if $j\ne0$; $p_i^{-1}-q_i^{-1}\in
\Gamma_\alpha(1)$ for $i<j$; and $\Lambda_\alpha^E\le(2/|M|_T)^K\Lambda_K(p)\Lambda_K(q)$ whenever
$4E\le K$.
\end{lemma}
\begin{proof}
Where the integrand of the majorant is non-zero there are points $y$, $yx$ with
$w_p^{-1}y$, $w_q^{-1}yx$ near at $j$ (\S\ref{pfr:ss-gauge}); in the gauge $\Gamma_1$ the leading
minors of order $l\le j$ of the frame matrix of an element moved into $\pffK$ by $R^{(j)}$ are integral, and
these minors are the $\delta_l$. Let $\beta=\prod_{l\le j}\delta_l(p)$, $\gamma=\prod_{l\le j}\delta_l(q)$;
they are integral multiples of each $\delta_l$ and of each $p_i$, resp.\ $q_i$ ($p_i=\delta_{i+1}/\delta_i$),
and $\Lambda_K(p)=\pbmn(\beta)^{K/2}$, $\Lambda_K(q)=\pbmn(\gamma)^{K/2}$. Put $\alpha=M\,\aubar{\beta\gamma}\,
\beta\gamma$. Then $\alpha\delta_l(p)^{-1}$, $\alpha\delta_l(q)^{-1}$, $\alpha M^{-1}$, $\alpha(p_i^{-1}-q_i^{-1})$
are integral, and $|\alpha|_T=|M|_T(|\beta|_T|\gamma|_T)^2$, so $\Lambda_\alpha=2/|\alpha|_T=(2/|M|_T)(\pbmn
(\beta)\pbmn(\gamma))^2$; as each factor is $\ge1$ and $4E\le K$, $\Lambda_\alpha^E\le(2/|M|_T)^K\pbmn(\beta)^{K/2}
\pbmn(\gamma)^{K/2}$.
\end{proof}

\begin{theorem}[Decay at $T$]\label{pff:decay}
Assume (L1)--(L3), $j\le\varrho$, $t\ge0$ with $t+10^{10}n^4\le k$, and $K\ge10^{10}n^4$. Let $\xi$ be a
unit and either (i) $\omega(g)=\psi_T(\xi\,h(gf_j,f_j))$, or (ii) $\omega=1$, $j=i+1$ and
$\xi=M^{-1}(p_i^{-1}-q_i^{-1})$. Then
\[
|\pfrI_j(p,q;\omega)|\le\pffr(\pffZ)\,\pbmn(\xi)^{t/2}\,\pbmden(\xi)^{-6}\,(2/|M|_T)^K\,\Lambda_K(p)\Lambda_K(q).
\]
\end{theorem}
\begin{proof}
If $\pfrF^+_j\equiv0$ the form vanishes (Lemma~\ref{pfr:L1}). Otherwise take $\alpha$ of
Lemma~\ref{pff:heads}, $\Gamma=\Gamma_\alpha$, $\Lambda=\Lambda_\alpha$, and $\beta'=\alpha^{200n^2}$.

\emph{A unit $\eta$ dividing $\xi$.} In case (i), as $\pffOT$ is a finite product of discrete valuation
rings, there is a unit $\eta=\xi+c\beta'$ of $\pffET$, $c\in\pffOT$, with $\xi\eta^{-1}$ and $\eta^{-1}\beta'$
integral (at each component take $\eta$ to be $\xi$ or $\xi+\beta'$, whichever generates the ideal
$(\xi,\beta')$ there); then $\psi_T(\xi c')=\psi_T(\eta c')$ for $c'\in\Gamma(200n^2)$, because
$\psi_T$ is trivial on $\pffOT$ and $c\beta'c'\in\pffOT$. Put $w=\eta^{-1}\beta'$. In case (ii) put
$\eta=\xi$ and $w=\beta'$; here $\xi\in\Gamma(2)\subseteq\Gamma(200n^2)$ by Lemma~\ref{pff:heads}, so
$\xi w\in\pffOT$. In both cases $w$, $\xi w$, $\xi\eta^{-1}$ are integral, and $|\eta|_T^{-1}=|w|_T
|\beta'|_T^{-1}$ in case (i), while in case (ii) $|w|_T|\beta'|_T^{-1}=1$.

\emph{The gain.} Let $e=|\eta|_T^{-1/2}$. By Theorem~\ref{pfr:gain-char} (case (i), with $\eta$ in place
of the character's argument) or Theorem~\ref{pfr:gain-heads} (case (ii)), with $P=1$,
$\pfrF_j(x;\omega)^2\le e^{2\operatorname{rank}X'}\mathsf Q(x)$ at every position; if $e<1$, as the volume
exponent satisfies $s\le\operatorname{rank}X'$, $\pfrF_j(x;\omega)$ has the gain $e^{s}\le Ve^s$.

\emph{Sum over the shells.} By Lemma~\ref{pfr:L1}, Lemma~\ref{pff:dyadic} (the dyadic shells cover the
positions, outside of which $\pfrF_j^+$ vanishes since the integrand involves $\langle\pffW(y),\pffW(yx)
\rangle\ne0$, which forces $x$ to be a position), Lemma~\ref{pff:floormass} ($V=2$) and
Theorem~\ref{pfr:shells} with $r'=t+12$ (admissible as $s\ge9\cdot10^7n^4+t+12$ for $t+10^{10}n^4\le k$),
\[
|\pfrI_j(p,q;\omega)|\le\pfrkap_j\int\pfrF_j(x;\omega)\,d\pbvol\le2\min(1,e)^{t+12}\Lambda^{4\cdot10^8n^4}
\mathsf N(\pffW),\qquad\mathsf N(\pffW)\le\pffr(\pffZ)
\]
(the last inequality as in the proof of Proposition~\ref{pff:outside}).

\emph{From $\eta$ to $\xi$.} As $\xi\eta^{-1}$ is integral, $\pbmn(\xi)=\pbmn(\xi\eta^{-1})\pbmn(\eta)\ge
\pbmn(\eta)=e^2$, so $\min(1,e)^t\le\pbmn(\xi)^{t/2}$. As $w\pffOT\subseteq\pffOT\cap\xi^{-1}\pffOT$,
$|w|_T\le\pbmden(\xi)^{-1}$; and $|\beta'|_T^{-1}\le\Lambda^{200n^2}$. Moreover $\min(1,e)^2\le|w|_T|\beta'|_T^{-1}$: in case (i) because $e^2=|\eta|_T^{-1}=|w|_T|\beta'|_T^{-1}$, in case (ii) because the right side is $1$. Hence $\min(1,e)^{12}\le(\min(1,e)^2)^6
\le(|w|_T|\beta'|_T^{-1})^6\le\pbmden(\xi)^{-6}\Lambda^{1200n^2}$. Altogether
\[
|\pfrI_j(p,q;\omega)|\le\pffr(\pffZ)\,\pbmn(\xi)^{t/2}\pbmden(\xi)^{-6}\,\Lambda^{E},\qquad E=1+1200n^2+4\cdot10^8
n^4,
\]
and $4E\le10^{10}n^4\le K$, so Lemma~\ref{pff:heads} gives $\Lambda^E\le(2/|M|_T)^K\Lambda_K(p)\Lambda_K(q)$.
\end{proof}


\subsection{Fixed vectors: the gain $b^{-6k}$}\label{pff:ss-fixed}

Let $\mathsf A$ be a set of chosen places above $p_0$ and $\mathsf d=(\mathsf d_w)$ a family of
frame vectors ({\cite[\S2.1]{P51}}). Let $K^{\mathrm{ad}}\subset G_f$ be the stabiliser of the
lattice and $\mathrm{red}\colon K^{\mathrm{ad}}\to\auG$ the reduction modulo the chosen places; it
is onto (strong approximation, {\cite[\S2.1]{P51}}). The residue field at a chosen place above
$p_0$ is $\auF_{p_0}$.

\begin{lemma}[Few non-units]\label{pff:nonunits}
Let $T'$ be a set of places each of which is a chosen place or the conjugate of one, with
$\#T'\le b$. Then the non-units of $\auO_{T'}$ have volume at most $b^{-9}$.
\end{lemma}
\begin{proof}
An element $y\in\auO_{T'}$ that is not a unit has a component in the maximal ideal $\mathfrak m_w$
for some $w\in T'$. The set of such $y$, for fixed $w$, has volume $q_w^{-1}\le b^{-10}$, as
$q_w\in\{p_0,p_1\}$ and $p_0,p_1\ge b^{10}$. Sum over the at most $b$ places.
\end{proof}

For the data $\#T'\le4f\le b$ (the chosen places above $p_0$ and $p_1$ and their conjugates).

\begin{lemma}\label{pff:index}
Let $w\in\mathsf A$ and $H_w=\{\kappa\in K^{\mathrm{ad}}:\mathrm{red}(\kappa)\in\pbFix(\mathsf d_w)\}$.
Then $[K^{\mathrm{ad}}:H_w]=p_0^{\,r}-1$.
\end{lemma}
\begin{proof}
As $\mathrm{red}$ is onto, $[K^{\mathrm{ad}}:H_w]=[\auG:\pbFix(\mathsf d_w)]$. The map $s\mapsto
s_w$ (or $s\mapsto s_w^\vee$, if $\mathsf d_w\in F_{\aubar w}$) maps $\auG$ onto $\auGL_r(\auF_{p_0})$,
and $\pbFix(\mathsf d_w)$ is the preimage of the stabiliser of the non-zero vector $\mathsf d_w$; as
$\auGL_r(\auF_{p_0})$ acts transitively on the $p_0^r-1$ non-zero vectors, the index is
$p_0^r-1$.
\end{proof}

\begin{proposition}[One place]\label{pff:oneplace}
Let $w\in\mathsf A$, let $T'=T\cap\{w,\aubar w\}$, and use the notation of \S\ref{pff:ss-setting}
with $T'$ in place of $T$. Assume (L1)--(L3) at $T'$, $b\ge1000$, $k\ge3\cdot10^{10}n^4$ and $r\ge
k+2$. Let $\pffZ_w$ be the image in $G_{T'}$ of $H_w$. Then $\pffZ_w$ is an intermediate group, it
contains the components at $T'$ of all $\kappa\in K^{\mathrm{ad}}$ with $\mathrm{red}(\kappa)\in
\pbFix(\mathsf A,\mathsf d)$, and
\[
\pffm(\pffZ_w)\,\pbvol(K_{T'})\,b^{6k}\le\pffm(K_{T'})\,\pbvol(\pffZ_w).
\]
\end{proposition}
\begin{proof}
Let $p\colon K^{\mathrm{ad}}\to G_{T'}$ be the projection; its image is $K_{T'}$ (strong
approximation for the lattice: a point of $K_{T'}$ is the component of an element of $G_f$ that is
integral elsewhere, which lies in $K^{\mathrm{ad}}$). If $p(\kappa)\in K_{T'}(M)$, then $\kappa-1$
vanishes modulo $w$, so $\mathrm{red}(\kappa)_w=1$ and $\kappa\in H_w$. Hence $p^{-1}(K_{T'}(M))
\subseteq H_w$, so $K_{T'}(M)\subseteq p(H_w)=\pffZ_w\subseteq K_{T'}$, and $\ker p\subseteq H_w$
gives $[K_{T'}:\pffZ_w]=[K^{\mathrm{ad}}:H_w]=p_0^r-1$ (Lemma~\ref{pff:index}). The second claim
is clear, as $\pbFix(\mathsf A,\mathsf d)\subseteq\pbFix(\mathsf d_w)$.

Put $B=b^9$; by Lemma~\ref{pff:nonunits}, the non-units of $\auO_{T'}$ have volume at most
$B^{-1}\le\frac12$. Let $t=k-10^{10}n^4$; then $t+1+n+2\cdot10^9n^4\le k$, $6k\le9t$ and $t+1<r$.

\emph{Small volume.} As $B\le b^{10}\le p_0$ and $t+1<r$, $B^{t+1}\le p_0^{t+1}\le p_0^r-1=
[K_{T'}:\pffZ_w]$, so by Lemma~\ref{pff:vol}(c)
\begin{equation}\label{pff:eq-rare}
B^{t+1}\,\pbvol(\pffZ_w)\le\pbvol(K_{T'}).
\end{equation}
\emph{Upper bound for $\pffZ_w$.} A $c\in\auO_{T'}$ with $c^{-1}\notin\auO_{T'}$ generates an ideal
of non-units, so $|c|_{T'}=\pbvol(c\auO_{T'})\le B^{-1}$. By Theorem~\ref{pff:massZ} with $Q=B$
and the exponent $t+1+n$, by $B^{-n}\le2^{-n}$ and by Lemma~\ref{pff:vol}(b),
\[
\pffm(\pffZ_w)\le\pbvol(\pffZ_w)\pbvol(K_{T'}(M))\bigl(\pbvol(\pffZ_w)+B^{-t-1}\,\pbvol(K_{T'})\bigr).
\]
\emph{Lower bound for $K_{T'}$.} By Theorem~\ref{pff:massZ}, $\pbvol(K_{T'})^2\pbvol(K_{T'}(M))\le
\pffm(K_{T'})$.

\emph{Conclusion.} By \eqref{pff:eq-rare} and $B\ge2$, $\bigl(\pbvol(\pffZ_w)+B^{-t-1}\pbvol(K_{T'})
\bigr)B^t\le2B^{-1}\pbvol(K_{T'})\le\pbvol(K_{T'})$ (each term times $B^{t+1}$ is at most
$\pbvol(K_{T'})$). As $b^{6k}\le b^{9t}=B^t$,
\[
\pffm(\pffZ_w)\pbvol(K_{T'})b^{6k}\le\pbvol(\pffZ_w)\pbvol(K_{T'}(M))\pbvol(K_{T'})^2\le\pffm(K_{T'})
\pbvol(\pffZ_w).\qedhere
\]
\end{proof}

\begin{lemma}[Products]\label{pff:products}
Let $T=T_1\sqcup\dots\sqcup T_s$ be a partition into sets stable under conjugation, and assume
(L1)--(L3) at $T$ and at each $T_i$. Then $G_T=\prod_iG_{T_i}$, and there is a constant $\lambda>0$
such that for all closed subgroups $Y_i\subseteq K_{T_i}$ containing $K_{T_i}(M)$ and $Y=\prod Y_i$
\[
\pbvol(Y)=\lambda\prod_i\pbvol(Y_i),\qquad\pffm(Y)=\lambda^3\prod_i\pffm(Y_i).
\]
\end{lemma}
\begin{proof}
The decomposition of $E_T$ as $\prod E_{T_i}$ gives $G_T=\prod_iG_{T_i}$, with $K_T=\prod K_{T_i}$,
$K_T(M)=\prod K_{T_i}(M)$, $U_T=\prod U_{T_i}$, and $\psi_T=\prod\psi_{T_i}$, $\sigma$ being additive;
hence $\chi_{M}^{(T)}(g)=\prod_i\chi_M^{(T_i)}(g_i)$. Both $\pbvol_T$ and $\prod_i\pbvol_{T_i}$ are Haar
measures on $G_T$, so $\pbvol_T=\lambda\prod_i\pbvol_{T_i}$. Then $\int_Y\chi_M(yxy^{-1})dy=\lambda
\prod_i\int_{Y_i}\chi_M^{(T_i)}(y_ix_iy_i^{-1})dy_i$, and integrating the square of its absolute
value over $x$ gives another factor $\lambda$ by Fubini's theorem.
\end{proof}

\begin{theorem}[The gain for fixed vectors]\label{pff:fixed}
Assume (L1)--(L3) at $T$ and at its subsets stable under conjugation, $b\ge1000$, $4f\le b$,
$k\ge3\cdot10^{10}n^4$, and that $T$ contains the places above $M$. For every set $\mathsf A$ of
chosen places above $p_0$ and every family $\mathsf d$ of frame vectors there is an intermediate
group $\pffZ=\pffZ(\mathsf A,\mathsf d)$ such that
\begin{enumerate}
\item[(a)] $\pffZ$ contains the component at $T$ of every $\kappa\in K^{\mathrm{ad}}$ with
$\mathrm{red}(\kappa)\in\pbFix(\mathsf A,\mathsf d)$;
\item[(b)] $\pffm(\pffZ)\,\pbvol(\pffK)\,b^{6k\#\mathsf A}\le\pffm(\pffK)\,\pbvol(\pffZ)$.
\end{enumerate}
\end{theorem}
\begin{proof}
A place of $T$ lies in $\{w,\aubar w\}$ for at most one $w\in\mathsf A$ (the conjugate of a chosen
place is not chosen). Partition $T$ into the sets $T_w=T\cap\{w,\aubar w\}$, $w\in\mathsf A$, which
are non-empty as $T$ contains the places above $M$, and the rest $T_0$. On $T_w$ take $\pffZ_w$ of
Proposition~\ref{pff:oneplace}, on $T_0$ take $K_{T_0}$ (for which (b) holds with the factor $1$),
and let $\pffZ=\prod_w\pffZ_w\times K_{T_0}$, an intermediate group. (a) holds componentwise. By
Lemma~\ref{pff:products} applied to $Y=\pffZ$ and to $Y=\pffK$, the inequality (b) is the product of
the inequalities of Proposition~\ref{pff:oneplace}, the constants $\lambda^3\cdot\lambda$ being the
same on both sides.
\end{proof}

\subsection{The forms at $T$ and the exact statements of \textup{(F)} and \textup{(S)}}\label{pff:ss-statements}

\subsubsection*{The forms}
Let $\pffZ$ be an intermediate group and $\pffW$ its Bessel vector. For heads $p,q$ of length $i$
(tuples of units of $\pffET$) and $0\le i\le\varrho$, let $\pffI_i(p,q)=\pffI_i(p,q;\pffW,1)$ be the
\emph{form of floor $i$} of the cut Bessel function $\pffW$, with weight $1$
(\S\ref{pfr:ss-setting}); for heads of length $n$ with a common last entry, let
$\mathcal D(p,q)=\mathcal D_{\pi_N}(p,q;\pffZ)$ be the \emph{density at the last entry}
(Definition~\ref{pff:lastentry}): the integral, over the points of the quadric $\{y\in\pbmXp\otimes
\pffET:h(y,y)=-N\}$ with last entry $h(y,z)=p_\varrho$, of the pairing of the scaled orbital
integrals of $\pffW$ at $w_{p'}\sigma_y$ and $w_{q'}\sigma_y$ ($p',q'$ the first $\varrho$ entries),
for the unit $\pi_N$ of $\pffET$ with $\pi_N\aubar\pi_N=N$. The \emph{forms at $T$} are
\begin{equation}\label{pff:eq-forms}
J^T_i(p,q)=\frac{\pffI_i(p,q)}{\pffr(\pffZ)}\quad(0\le i\le\varrho),\qquad J^T_n(p,q)=
\frac{\mathcal D(p,q)}{\pffr(\pffZ)} .
\end{equation}
They depend on $T$ and on $\pffZ$; in \textup{(F)} and \textup{(S)} the group is
$\pffZ=\pffZ(\mathsf A,\mathsf d)$ of Theorem~\ref{pff:fixed}. This is the description of
{\cite[\S3.3.7]{P51}}, item \textup{(F)}: the vector of translates of $\chi_M$ by the cosets of
$\pffKM$ in $\pffZ$, divided by $\pbvol(\pffZ)\pbvol(\pffKM)/\pffl^{\,2}$, and $\pffl^{-1}=\#\pbUb$
for the data (the volume of $U_T\cap\pffKM$ for the measure that gives $U_T\cap\pffK$ volume $1$ is
$[U_T\cap\pffK:U_T\cap\pffKM]^{-1}$, and reduction identifies this quotient with $\pbUb$).

\begin{lemma}[The form of floor $0$]\label{pff:J0}
$\pffI_0(p,q)=\int_{\pffG}|k_0(x)|^2dx=\pffm(\pffZ)/\pffl^{\,2}$ for all $p,q$; hence
$0\le J^T_0=\pffm(\pffZ)/(\pbvol(\pffZ)\pbvol(\pffKM))$.
\end{lemma}
\begin{proof}
At floor $0$ the head density is $1$, $\pfrm_0=\delta_1$, and $\theta_{\ge0}=\theta_{\ge1}\ast
\theta_0$, so $\Phi_0(x;1;1)=\bigl|\int\mathsf k^1_1(\sigma,\sigma x)\,d\theta_0(\sigma)\bigr|^2=|k_0(x)|^2$
and $\pfrI_0(p,q)=\int|k_0|^2\,d\pbvol$; the second is
Corollary~\ref{pff:k0}; the third follows from \eqref{pff:eq-mass}.
\end{proof}

\begin{theorem}[The mass in \textup{(F)}]\label{pff:mass}
Under the hypotheses of Theorem~\ref{pff:fixed} and $2^{-\varrho}\le b^{-6}$, for
$\pffZ=\pffZ(\mathsf A,\mathsf d)$ and all heads $p,q$,
\[
0\le J^T_0(p,q)\le(1+b^{-6})\,b^{-6k\#\mathsf A}.
\]
\end{theorem}
\begin{proof}
By Lemma~\ref{pff:J0}, Theorem~\ref{pff:fixed}(b) and Corollary~\ref{pff:massK},
\[
J^T_0=\frac{\pffm(\pffZ)}{\pbvol(\pffZ)\pbvol(\pffKM)}\le\frac{b^{-6k\#\mathsf A}\,\pffm(\pffK)}{
\pbvol(\pffK)\pbvol(\pffKM)}\le(1+b^{-6})\,b^{-6k\#\mathsf A}.\qedhere
\]
\end{proof}

\begin{remark}\label{pff:rem-E4}
For $\mathsf A=\emptyset$ the theorem says that the mass of $\pffK$ exceeds its main part
$\pbvol(\pffK)^2\pbvol(\pffKM)$, which comes from the positions in $\pffKM$, by at most the factor
$1+b^{-6}$ (up to $\pbvol(\pffK)\le1+2^{-n}$). The positions in $\pffK\smallsetminus\pffKM$ are among
those treated by Proposition~\ref{pff:outside}; on them the statement is of the same nature as
{\cite[Theorem~3.39]{P51}}, which is however not used here. The gain for fixed vectors does not come from
a bound for the sum of $f_{\mathsf D}$ over a stabiliser as suggested in {\cite[Remark~3.40(b)]{P51}},
but from the index $p_0^r-1$ of the stabiliser of a vector, which makes $\pbvol(\pffZ)$ small, while
the mass integral of $\pffZ$ is controlled by its main part $\pbvol(\pffZ)^2\pbvol(\pffKM)$ and an
error that is small relative to $\pbvol(\pffK)$.
\end{remark}

\subsubsection*{The exact statements}
Recall $\Lambda^T$, $\beta^T$ of {\cite[(38)]{P51}}, and the place $\mathfrak P_T=(\pffET,dx,\psi_T,
w\mapsto M^{-1}w^{-1})$ of {\cite[Theorem~3.50]{P51}}. A head of length $i$ is a tuple $p=(p_0,\dots,
p_{i-1})$ of units of $\pffET$, and $(p,c)$ is $p$ followed by $c$.

\begin{theorem}[Local system at $T$; {\cite[Theorem~3.51]{P51}}\,\textup{(F)}]\label{pff:F}
Assume the hypotheses on the data of {\cite[Theorem~3.53]{P51}}. Let $T$ be a finite set of finite places
of $E$, stable under conjugation, containing the places above $MN$. Let $\mathsf A$ be a set of
chosen places above $p_0$ and $\mathsf d$ a family of frame vectors, and let $\pffZ=\pffZ(\mathsf A,
\mathsf d)$ be the intermediate group of Theorem~\ref{pff:fixed}. Then the forms
\eqref{pff:eq-forms} for $\pffZ$ satisfy:
\begin{enumerate}
\item[(a)] $\pffZ$ contains the component at $T$ of every $\kappa\in K^{\mathrm{ad}}$ with
$\mathrm{red}(\kappa)\in\pbFix(\mathsf A,\mathsf d)$;
\item[(b)] $(J^T_i)_{0\le i\le n}$ is a sieve system at $\mathfrak P_T$
({\cite[Definition~3.44]{P51}}) with mass $(1+b^{-6})\,b^{-6k\#\mathsf A}$, weights $\Lambda^T$ and
decay $\beta^T$.
\end{enumerate}
\end{theorem}

\begin{proof}[Proof of Theorem~\ref{pff:F}]
(a) is Theorem~\ref{pff:fixed}(a). For (b): \emph{(S4)} is Theorem~\ref{pff:mass}.
\emph{(S2), (S3)} are Proposition~\ref{pff:top}(b), (c), divided by $\pffr(\pffZ)>0$.
\emph{(S1).} For $i<\varrho$ the identity and the integrability are Lemma~\ref{pfr:descent}, which
applies as $\int\pfrF^+_{i+1}d\pbvol<\infty$ (Theorem~\ref{pfr:shells} with $e=1$, the right side being
finite); for $i=\varrho$ they are Proposition~\ref{pff:top}(a). In both cases the twisted integral is the
form of floor $i$ with the weight $\psi_T(\xi h(gf_i,f_i))$, and for a unit $\xi$
Theorem~\ref{pff:decay}(i) with $t=k-n^7$, $K=n^6$ (admissible: $10^{10}n^4\le n^6$ and $n^7+10^{10}n^4\le
k$ for the data) bounds it by $\pffr(\pffZ)\pbmn(\xi)^{(k-n^7)/2}\pbmden(\xi)^{-6}(2/|M|_T)^{n^6}
\Lambda^T_i(p)\Lambda^T_i(q)$. \emph{(S5)} is the same with Theorem~\ref{pff:decay}(ii), $j=i$ (the floor
$i$ form for heads of length $i$ with $\xi=M^{-1}(p_{i-1}^{-1}-q_{i-1}^{-1})$). Finally $\pbmden(\xi)^{-6}\le
\pbmden(\xi)^{-3}$ and $(2/|M|_T)^{n^6}=(2M^{2f})^{n^6}\le(2b^{44f})^{n^6}\le b^{n^9f}$ (as $M\le b^{22}$,
$2\le b$ and $44f+1\le n^3f$), which gives the decay $\beta^T$; dividing by $\pffr(\pffZ)$ gives
\eqref{pff:eq-forms}.
\end{proof}

\begin{theorem}[Siegel's step; {\cite[Theorem~3.51]{P51}}\,\textup{(S)}]\label{pff:S}
Assume the hypotheses of Theorem~\ref{pff:F}, and in addition that $T$ contains the places at which
an element of the set $\mathsf B$ of {\cite[\S3.3.2]{P51}} is not a unit. Let $J^\tau_n$ be the top
forms at the real places of Theorem~\ref{pfr:R}. Let $(c_j,c'_j,s_j,\sigma_j)_{j\in\mathsf I}$ be a
finite family of pairwise distinct quadruples of classes $c_j,c'_j$, elements $s_j\in\pbFix(\mathsf
A,\mathsf d)$ and $\sigma_j\in\pbmGf(E^+)$ such that $t_j=t(\sigma_j)\in\pbmT$, $|\tau(t_j)|>N$
for all $\tau$, and $h_{c_j}^{-1}\sigma_jh_{c'_j}$ stabilises the lattice at every finite place
outside $T$. Then
\[
\pbmrS^{\,2}\sum_{j\in\mathsf I}\Bigl\|\sum_{p\in\pbmT^\varrho}\pbmOrb_{c_j,s_j,c'_j}(w_p\sigma_j)
\Bigr\|^2\le(1+b^{-6})\,\pbmaS\,\Rel\sum_{t\in\{t_j\}}\ \sum_{p,q\in\pbmT^\varrho}\prod_\tau J^\tau_n
\bigl([p,t]_\tau,[q,t]_\tau\bigr)\,J^T_n\bigl((p,t),(q,t)\bigr),
\]
with $J^T_n$ the top form \eqref{pff:eq-forms} for the group $\pffZ(\mathsf A,\mathsf d)$ of
Theorem~\ref{pff:F}.
\end{theorem}

\begin{remark}[Comparison with {\cite[Theorem~3.51]{P51}}, and the constants]\label{pff:rem-L2}
(1) \emph{The constants of \textup{(F)}.} The mass $(1+b^{-6})b^{-6k\#\mathsf A}$, the weights
$\Lambda^T_i(p)=\prod_{l\le i}[\delta_l\in\pffOT]\,\pbmn(\delta_l)^{n^6/2}$, the decay
$\beta^T(\xi)=\pbmn(\xi)^{(k-n^7)/2}\pbmden(\xi)^{-3}b^{n^9f}$, and the place $\mathfrak P_T$ (with
$\iota(w)=M^{-1}w^{-1}$, so that the condition (S5) concerns $\xi=M^{-1}(p_{i-1}^{-1}-q_{i-1}^{-1})$)
are exactly those of {\cite[Theorem~3.51]{P51}}\,\textup{(F)} and of {\cite[Theorem~3.50]{P51}}, which
{\cite[Theorem~3.53]{P51}} applies. Here $\pbmn(\xi)=|\xi|_T^{-1}$ is the generalised index $[\pffOT:\xi
\pffOT]$ and $\pbmden(\xi)=[\pffOT:\pffOT\cap\xi^{-1}\pffOT]$. The decay of Theorem~\ref{pff:decay} carries the
factor $(2/|M|_T)^{n^6}=(2M^{2f})^{n^6}$ (as $|M|_T=M^{-2f}$), which is bounded by $b^{n^9f}$ because
$M\le b^{22}$ and $44f+1\le n^3f$; this is the factor $b^{n^9f}$ of $\beta^T$, and the power
$\pbmden(\xi)^{-6}$ of Theorem~\ref{pff:decay} is weakened to $\pbmden(\xi)^{-3}$.

(2) \emph{The constant of \textup{(S)}.} With $\#\pbR=\#\pbHb\,\#\pbUb$ and $\#\pbUb=M^{f\varrho(2k+
2\varrho+1)}$ (the unipotent radical has dimension $\varrho(2k+2\varrho+1)$, and reduction modulo $M$ gives
$[U_T\cap\pffK:U_T\cap\pffKM]=|\pffOT^+/M\pffOT^+|^{\dim U}$ with $|\pffOT^+/M\pffOT^+|=M^f$, $\pffOT^+$ the
integers of $E^+_T$),
\[
\pbmaS=\#\pbR\cdot\#\pbUb\cdot\#\mathcal C\cdot\Bigl(\frac{\pi^k}{(k-1)!\,N^k}\Bigr)^f\aurd^{\,f/2}
\pbmvol^{\,k},\qquad\aurd^{\,f/2}=|\operatorname{disc}E|^{1/4},\qquad\pbmvol=(2/\aurd)^f=\pbvol(
\widehat\auO_E),
\]
where $\pbvol(\widehat\auO_E)$ is the volume of $\widehat\auO_E$ for the Haar measure of $\auA_{E,f}$ which,
together with Lebesgue measure on each $E_\tau=\C$, gives $E\backslash\auA_E$ volume $1$
(Lemma~\ref{pfg:tamagawa}(a)); up to the factor $1+b^{-6}$ it is
(number of classes) $\times\#\pbR\#\pbUb\times$ (Siegel's constant). The
normalisation is $\pbmrS=d^{-1}\prod_\tau\bigl(\sqrt{\pbmcG}\prod_as_{\tau,a}^{k-1+2(\varrho-a)}\bigr)$. {\cite[Lemma~3.52]{P51}} uses exactly $\pbmaS$, $\pbmrS$ and the mass
$(1+b^{-6})b^{-6k\#\mathsf A}$; so the constants transcribed in {\cite[Theorem~3.51]{P51}} are right.

(3) \emph{Two additions.} (a) In {\cite[Theorem~3.51]{P51}}, the functions $J^T_i$ are said to
depend on $T,\mathsf A,\mathsf d$; they do so through the group $\pffZ(\mathsf A,\mathsf d)$, and
\textup{(S)} holds for the top form of \emph{that} group: the elements $s_j\in\pbFix(\mathsf A,
\mathsf d)$ are reductions of elements $\kappa_j\in K^{\mathrm{ad}}$, whose components at $T$ lie in
$\pffZ$ by (a) of Theorem~\ref{pff:F}, and this is what Siegel's step uses. (b) $T$ must contain
the places at which the lattice basis is not integral with integral inverse; these lie above $N$
(\S\ref{pff:ss-setting}), so the hypothesis of {\cite[Theorem~3.51]{P51}}\,\textup{(F)} already
contains it. With these readings the proof of {\cite[Theorem~3.53]{P51}} applies verbatim: it uses
\textup{(F)} and \textup{(S)} for the same $T,\mathsf A,\mathsf d$.
\end{remark}

\subsection{Siegel's step: proof of Theorem~\ref{pff:S}}\label{pff:ss-siegel}

Throughout, $T$, $\mathsf A$, $\mathsf d$, $\pffZ=\pffZ(\mathsf A,\mathsf d)$ are as in
Theorem~\ref{pff:S}, $C=\pffKM\backslash\pffZ$ is the set of cosets, and for a class $c$ we write
$h_c\in H(\auA_f)$ for its representative. For a rational $\gamma$ we write $\gamma_T$ for its component
at $T$ and put $g(c,c',\sigma)=(h_c^{-1}\sigma h_{c'})_T\in\pbmGf_T$.

\subsubsection*{Siegel's formula and the density at the finite places}
Let $\mathcal C=H(E^+)\backslash H(\auA_f)/J_X$ be the set of classes ({\cite[\S2.1]{P51}}); by
{\cite[Lemma~14.5]{P46}} the stabilisers $H(E^+)\cap hJ_Xh^{-1}$ are trivial. On $\auA_{E,f}^k$ we use the
Haar measure for which $\widehat\auO_E$ has volume $\pbvol(\widehat\auO_E)=\pbmvol=(2/\aurd)^f$ in each
coordinate (the measure of Lemma~\ref{pfg:tamagawa}(a)), and on the self-adjoint finite adeles $\mathrm{Herm}_f$
the Haar measure for which the codifferent $\{z:\psi_{E,f}(z\widehat\auO_E)=1\}\cap\mathrm{Herm}_f$ has
volume $1$. For $\Phi\ge0$ locally constant with compact support on $\auA_{E,f}^k$ and $\alpha\in
\auA_{E,f}$ self-adjoint put
\[
\delta(\Phi,\alpha)=\int_{\mathrm{Herm}_f}\int_{\auA_{E,f}^k}\Phi(x)\,\psi_{E,f}\Bigl(\tfrac12z\,(x^*x-\alpha)
\Bigr)\,dx\,dz ,
\]
whenever the outer integral converges absolutely, and $0$ otherwise.

\begin{theorem}[Siegel's formula]\label{pff:siegelformula}
Let $k\ge3$, let $\Phi\ge0$ be locally constant with compact support on $\auA_{E,f}^k$ and invariant under
$J_X$, and let $\alpha\in E^+$ be totally positive. Then
\[
\sum_{c\in\mathcal C}\ \sum_{v\in E^k,\ v^*v=\alpha}\Phi(h_c^{-1}v)=\#\mathcal C\cdot|\operatorname{disc}E|^{1/4}
\prod_{\tau\in\auPl}\frac{\pi^k\,\tau(\alpha)^{k-1}}{(k-1)!}\cdot\delta(\Phi,\alpha),
\]
the inner sums being finite.
\end{theorem}
This is Siegel's formula for the representations of a totally positive number by the genus of $x^*x$; it is
derived, with every constant, from Weil's Siegel--Weil theorem (Theorem~\ref{pfg:weil}) in
Proposition~\ref{pfg:siegel}: the measures are Tamagawa measures (Lemma~\ref{pfg:tamagawa}), the archimedean
density is computed, and the $z$-integral defining $\delta(\Phi,\alpha)$ converges absolutely (classically: for quadratic forms over $\Q$
\cite[\S7]{pff-Siegel35}, for hermitian forms \cite{pff-Braun41}).

\begin{lemma}[The density off $T$ is at most $(1+b^{-6})\pbmvol^k$]\label{pff:offT}
Let $\Psi\ge0$ be locally constant with compact support on $X_{0,T}=\pffET^k$, let $\Phi=\Psi\circ
\mathrm{pr}_T\cdot\mathbf 1_{\prod_{v\notin T}\auO_v^k}$, and let $\alpha\in E^+$ (identified with its image in
$\auA_{E,f}$). Then, for the data,
\[
|\delta(\Phi,\alpha)|\le(1+b^{-6})\,\pbmvol^{\,k}\int\Psi\,d\nu_{\alpha,T},
\]
where $\nu_{\alpha,T}$ is the fibre measure at $\alpha$ of $x\mapsto x^*x$ on $\pffET^k$
(Definition~\ref{pff:lastentry}).
\end{lemma}
\begin{proof}
\emph{Factorisation.} Write $\auA_{E,f}=\pffET\times\auA^T$, $x=(x_T,x^T)$, $z=(z_T,z^T)$. Because
$\Phi(x)$ depends on $x_T$ and on the condition $x^T$ integral, and because $\psi_{E,f}(\frac12z^T(x^{T*}x^T-
\alpha))$ is $1$ when $z^T$ is in the codifferent and $x^T$ integral, an exact computation (translate $x$
by the integral vectors off $T$ and $z$ by the self-adjoint elements that are in the codifferent off $T$;
the measures are translation invariant and the codifferent has volume $1$) gives
\[
\delta(\Phi,\alpha)=\Bigl(\int_{\mathrm{Herm}^{(T)}}\ \int_{\widehat\auO_E^k}\psi_{E,f}\bigl(\tfrac12z(w^*w-\alpha^{(T)})
\bigr)\,dw\,dz\Bigr)\cdot\int_{\mathrm{Skew}_T}\int_{\pffET^k}\Psi(y)\,\psi_T\bigl(\tau b(y^*y-\alpha)\bigr)\,dy\,db,
\]
where $\mathrm{Herm}^{(T)}$ is the set of self-adjoint $z$ whose components at $T$, multiplied by $-\varpi$,
are integral, $\alpha^{(T)}$ is $\alpha$ with its components at $T$ removed, $\tau\in\pffOT$ has
$\tau+\aubar\tau=1$, and $db$ is the Haar measure of the skew elements of $\pffET$ that gives the skew
integers volume $1$ (the substitution $z_T=-\varpi^{-1}b$ turns $\psi_{E,f}(\frac12 z_T\cdot)$ into
$\psi_T(\tau b\,\cdot)$ on the relevant set).

\emph{The Gauss sums.} For self-adjoint $z$ let $I_z\subseteq\widehat\auO_E$ be the largest ideal on which
$v\mapsto\psi_{E,f}(zv)$ is trivial, and $G(z)=\int_{\widehat\auO_E^k}\psi_{E,f}(\frac12z(w^*w-\alpha^{(T)}))dw$. By
Weyl differencing, $|G(z)|^2=\int_{\widehat\auO^k}\int_{\widehat\auO^k}\aubar{F(w)}F(w+d)\,dw\,dd$ with $F$ the
integrand; the inner integral vanishes unless every $d_i$ lies in $K_z=\{r:\psi_{E,f}(zr\widehat\auO)=1\}$
(translate $w$ by $\aubar v e_i$, $v\in\widehat\auO$, which multiplies it by $\psi_{E,f}(zd_iv)$), so
$|G(z)|^2\le(\pbvol(K_z\cap\widehat\auO)\pbmvol)^k=([\widehat\auO:I_z]^{-1}\pbmvol^2)^k$. The $z$ with $I_z=I$
lie in at most $[\widehat\auO:I]$ cosets of the codifferent (one for each character of $\widehat\auO/I$), each
of volume $1$. Hence, for $\mu$ with $2(\mu+1)\le k$,
\[
\int|G(z)|\,dz\le\sum_I[\widehat\auO:I]\cdot[\widehat\auO:I]^{-k/2}\pbmvol^k\le\zeta_{\widehat\auO}(\mu)\,
\pbmvol^k,\qquad\zeta_{\widehat\auO}(\mu)=\sum_I[\widehat\auO:I]^{-\mu}.
\]
Take $\mu=3\varrho+4$ ($2(\mu+1)\le k$ for the data). As $\zeta_{\widehat\auO}(\mu)\le\zeta_E(\mu)$ and
$\log\zeta_E(\mu)\le2f\cdot2^{2-\mu}$, we get $\zeta_{\widehat\auO}(\mu)\le1+2^{-\varrho}\le1+b^{-6}$ for
the data (as in \S\ref{pff:ss-large}).

\emph{The density at $T$.} For $\Psi\ge0$ locally constant with compact support,
$\bigl|\int\!\int\Psi(y)\psi_T(\tau b(y^*y-\alpha))\,dy\,db\bigr|\le\int\Psi\,d\nu_{\alpha,T}$. Indeed, if the
outer integral does not converge absolutely the left side is $0$ by convention. Otherwise, cut the
frequency to $b\in u\aubar u\,\mathcal O$ for a unit $u$ ($\mathcal O$ the skew integers): by orthogonality of
characters the cut integral equals $\pbvol(u\aubar u\mathcal O)$ times the mass, for $dy$, of the slab
$\{y:y^*y-\alpha\in(u\aubar u)^{-1}\mathcal O_{E^+,T}\}$ weighted by $\Psi$, where $\mathcal O_{E^+,T}$ is the lattice of self-adjoint integers, dual to $\mathcal O$ under
$(b,X)\mapsto\psi_T(\tau bX)$ since $E/E^+$ is unramified (so the two volumes multiply to $1$); this mass, by the disintegration, is the
mean of $\int\Psi\,d\nu_{\alpha+\beta,T}$ over $\beta$ in that slab. As $u\to0$ the cut integrals tend
to the full integral (dominated convergence) and the means tend to $\int\Psi\,d\nu_{\alpha,T}$
(continuity of $\beta\mapsto\int\Psi\,d\nu_{\beta,T}$, Lemma~\ref{pff:fibre}(a)); so the integral equals
$\int\Psi\,d\nu_{\alpha,T}\ge0$.

Multiplying the two bounds gives the lemma.
\end{proof}

\begin{proposition}[Siegel's formula as an upper bound]\label{pff:siegelsum}
Let $t\in E$ with $|\tau(t)|>N$ for all $\tau$, $\alpha=(t\aubar t-N^2)/N$, and let $\Psi\ge0$ on
$V_T=Y_0\otimes\pffET$ be locally constant with compact support and invariant under $H_T\cap\pffKM$.
For a class $c$ and $v\in E^k$ let $x(c,v)=h_c^{-1}v$. Then for every finite set $\mathcal S$ of pairs
$(c,v)$ with $Nv^*v=t\aubar t-N^2$ and $x(c,v)$ integral outside $T$,
\[
\sum_{(c,v)\in\mathcal S}\Psi\bigl(x(c,v)_T-N^{-1}tz\bigr)\le(1+b^{-6})\,\#\mathcal C\,\mathsf c_S\,G_t\cdot
|\pi_N|_T^{-1}\int\Psi\bigl(y-N^{-1}tz\bigr)\,d\nu_{\alpha,T}(y),
\]
with $\mathsf c_S=(\pi^k/((k-1)!N^k))^f|\operatorname{disc}E|^{1/4}\pbmvol^k$ and $G_t=\prod_\tau(|\tau(t)|^2-
N^2)^{k-1}$.
\end{proposition}
\begin{proof}
Let $\Phi(x)=\Psi(x_T-N^{-1}tz)\,\mathbf 1[x\text{ integral off }T]$ on $\auA_{E,f}^k$. It is locally
constant with compact support; it is invariant under $J_X$: an element $\kappa\in J_X$ is integral off $T$
(so preserves the integrality condition), and at $T$ it acts on the vectors $x-N^{-1}tz$ as an element of
$H_T\cap\pffKM$ ({\cite[Lemma~2.1]{P51}}: $K_H\cap J=J_X$), under which $\Psi$ is invariant. The pairs of
$\mathcal S$ are among the terms of the left side of Theorem~\ref{pff:siegelformula} (for each class, the
$v$ with $v^*v=\alpha$ and $\Phi(h_c^{-1}v)=\Psi(\cdots)$), and all terms are $\ge0$; so the left side is at
most $\#\mathcal C|\operatorname{disc}E|^{1/4}\prod_\tau\pi^k\tau(\alpha)^{k-1}/(k-1)!\cdot|\delta(\Phi,\alpha)|$.
By Lemma~\ref{pff:offT} (with $\Psi(\cdot-N^{-1}tz)$), $|\delta|\le(1+b^{-6})\pbmvol^k\int\Psi(y-N^{-1}tz)
d\nu_{\alpha,T}$. Finally $\prod_\tau\tau(\alpha)^{k-1}=N^{-(k-1)f}G_t$, and the product formula for the
$T$-unit $N$ ($N^{-1}$ is integral off $T$) gives $|\pi_N|_T^{-1}=|N|_T^{-1/2}=N^f$; so
$N^{-(k-1)f}=N^{-kf}\cdot|\pi_N|_T^{-1}$, which is the stated constant.
\end{proof}

\subsubsection*{From class sums to the fibre pairing}

\begin{lemma}[The real places]\label{pff:topreal}
Let $\tau\in\auPl$, $\sigma\in\pbmGf_\tau$ and $x$ a head with $h(\sigma z,z)=x_\varrho$ (in the normalised
basis). Then $\pfrip{\pfrN_\sigma(x)}{\pfrN_\sigma(x')}=J^\tau_n(x,x')$ whenever $x_\varrho=x'_\varrho$.
\end{lemma}
\begin{proof}
$\sigma^{-1}z$ and $\sigma_x^{-1}z$ have the same norm and the same pairing with $z$, so their components in
the span of the $x_i$ have the same norm, and by Witt's theorem there is $h\in\pfrH$ with $\sigma^{-1}z=
h\sigma_x^{-1}z$. Then $h_2=\sigma_xh^{-1}\sigma^{-1}$ fixes $z$ and the $e_a,e_a^*$, so $h_2\in\pfrH$ and
$\sigma=h_2^{-1}\sigma_xh^{-1}$. As $w_x$ commutes with $\pfrH$ and $W^\flat(h'y)=\pi(h')W^\flat(y)$ for
$h'\in\pfrH$, the substitution $y\mapsto h_2^{-1}y$ gives $\pfrN_\sigma(x)=(\pi(h_2)^{-1}\otimes\aubar{\pi(h)})
\pfrN_{\sigma_x}(x)$, with a unitary operator that does not depend on $x$. Take inner products and use
the definition of $J^\tau_n$ in \eqref{pfr:eq-c}.
\end{proof}

\begin{lemma}[Orbital integrals at $T$ are Kloosterman sums]\label{pff:locKl}
Let $\sigma\in\pbmGf(E^+)$, $c,c'$ classes, $p$ a head of units of $\pbmRT$ (entries and inverses integral
off $T$), $\gamma=w_p\sigma$ with $h_c^{-1}\gamma h_{c'}$ integral off $T$, $l=g(c,c',\sigma)$, and let
$\kappa\in K$ and cosets $a,b\in C$ with $\kappa_T^{-1}a_{\vphantom{1}}b^{-1}\in\pffKM$. Then
\[
\mathcal O_p(l^{-1})_{(a,b)}\cdot O(\gamma)=\pfrkap_\varrho(p,1)\,\pbvol(\pffKM)\,\pbmOrb_{c,\bar\kappa,c'}(\gamma),
\]
where $\mathcal O_p$ is the scaled orbital integral of \S\ref{pff:ss-floors}, $O(\gamma)$ the orbital integral
at the real places ({\cite[\S3.3.2]{P51}}), and $\bar\kappa$ the reduction of $\kappa$.
\end{lemma}
\begin{proof}
By {\cite[Proposition~3.42]{P51}}, $\pbmOrb_{c,\bar\kappa,c'}(\gamma)=\pbmKl_{c,\bar\kappa,c'}(\gamma)O(\gamma)$, and by
the remark after it $\pbmKl$ is the sum of $\psi_f(u)\aubar{\psi_f(u')}$ over $(u,u')\in(U(\auA_f)/U(\auA_f)\cap
h_cJh_c^{-1})\times(\cdots)$ with $h_c^{-1}u^{-1}\gamma u'h_{c'}\in K$ of reduction $\bar\kappa$. Off $T$ all
data are integral and the characters trivial, so the sum is over $U_T/(U_T\cap\pffKM)$ in both variables,
with the condition $u^{-1}(w_pl)u'\in\kappa_T\pffKM$ at $T$ (the classes being represented in $H$, which
normalises $U$). On the other hand, by definition,
\[
\int_{\pffG}\pffW(y)_a\aubar{\pffW((w_pl)^{-1}y)_b}\,dy=\int_{\pffG}\pffch(ya^{-1})\aubar{\pffch(l^{-1}w_p^{-1}yb^{-1})}dy .
\]
The first factor vanishes unless $ya^{-1}=uk$ ($u\in U_T$, $k\in\pffKM$); write $y=uka$ and integrate over
$u\in U_T/(U_T\cap\pffKM)$ and $k\in\pffKM$. The second factor is non-zero iff $l^{-1}w_p^{-1}uka b^{-1}\in
u'\pffKM$, i.e.\ $u'^{-1}(w_pl)^{-1}u\in\pffKM\,ba^{-1}=(\kappa_T\pffKM)^{-1}$ (using $\kappa_T^{-1}ab^{-1}\in\pffKM$
and normality), and then the integrand is $\psi_T(M^{-1}\sigma(u))\aubar{\psi_T(M^{-1}\sigma(u'))}$. The
$k$-integral gives $\pbvol(\pffKM)$; so the integral is $\pbvol(\pffKM)$ times the Kloosterman sum above (the
inversion $(u,u')\leftrightarrow$ the pairs of the global sum being the identity on the condition). Multiply by
$\pfrkap_\varrho(p,1)$.
\end{proof}

\begin{lemma}[One position]\label{pff:headsum}
Let $c,c'$ be classes and $\sigma\in\pbmGf(E^+)$ with $h_c^{-1}\sigma h_{c'}$ integral off $T$, let $\mathcal P$ be
a finite set of heads of units of $\pbmRT$, $\rho_p>0$ with $\rho_p\pfrkap_\varrho(p,1)=w$ for all $p$, and let
$R$ be the Gram matrix of the vectors $\rho_pO(w_p\sigma)$. Let $\kappa_i\in K$ ($i\in I$) have distinct
reductions $s_i$ and components $\kappa_{i,T}\in\pffZ$. Then, with $l=g(c,c',\sigma)$,
\[
\pbvol(\pffKM)^2\,\#C\sum_{i\in I}w^2\Bigl\|\sum_{p\in\mathcal P}\pbmOrb_{c,s_i,c'}(w_p\sigma)\Bigr\|^2\le
\Rel\sum_{p,q}R_{pq}\pfrip{\mathcal O_p(l^{-1})}{\mathcal O_q(l^{-1})}=:\mathcal H(l).
\]
\end{lemma}
\begin{proof}
$\mathcal H(l)=\sum_{\mathbf x\in C\times C}v_{\mathbf x}^*Rv_{\mathbf x}$ with $v_{\mathbf x}=(\mathcal
O_p(l^{-1})_{\mathbf x})_p$, a sum of non-negative terms ($R$ is a Gram matrix). For $i\in I$ and $b\in C$ let
$A_i(b)$ be the coset with $\kappa_{i,T}^{-1}bA_i(b)^{-1}\in\pffKM$ (it exists as $\kappa_{i,T}\in\pffZ$). For
$\mathbf x=(b,A_i(b))$, Lemma~\ref{pff:locKl} gives $\sum_pv_{\mathbf x,p}\rho_pO(w_p\sigma)=w\,\pbvol(\pffKM)\sum_p
\pbmOrb_{c,s_i,c'}(w_p\sigma)$, so $v_{\mathbf x}^*Rv_{\mathbf x}=w^2\pbvol(\pffKM)^2\|\sum_p\pbmOrb\|^2$. The pairs
$(b,A_i(b))$ are distinct for distinct $(i,b)$: if $A_i(b)=A_{i'}(b)$ then $\kappa_{i,T}^{-1}\kappa_{i',T}\in
\pffKM$, so (as both are integral off $T$, hence in $K$) $\kappa_i^{-1}\kappa_{i'}\in J$ and $s_i=s_{i'}$. Summing
over $i$ and the $\#C$ cosets $b$ gives the claim.
\end{proof}

\begin{lemma}[Separation]\label{pff:sep}
Let $(c,c_1',\sigma_1)$, $(c,c_2',\sigma_2)$ have the same left class $c$, with $\sigma_i\in\pbmGf(E^+)$,
$h(\sigma_1z,z)=h(\sigma_2z,z)$, the same $X_0$-component of $\sigma_iz$, $h_c^{-1}\sigma_ih_{c_i'}$ integral off
$T$, and $g(c,c_1',\sigma_1)^{-1}g(c,c_2',\sigma_2)\in\pffKM$. Then $c_1'=c_2'$ and $\sigma_1=\sigma_2$.
\end{lemma}
\begin{proof}
$\sigma_1z=\sigma_2z$ (same $X_0$-part and same pairing with $z$; the $z$-coordinate is determined by the
latter), so $\nu=\sigma_1^{-1}\sigma_2\in\pbmGf(E^+)$ fixes $z$: $\nu\in H(E^+)$. Moreover $h_{c_1'}^{-1}\nu
h_{c_2'}=(h_c^{-1}\sigma_1h_{c_1'})^{-1}(h_c^{-1}\sigma_2h_{c_2'})$ is integral off $T$ and in $\pffKM$ at $T$, so in
$J\cap H(\auA_f)=J_X$ ({\cite[Lemma~2.1]{P51}}). Hence $c_1'=c_2'$ in $\mathcal C$, and then $\nu\in H(E^+)\cap
h_{c_1'}J_Xh_{c_1'}^{-1}=1$ by {\cite[Lemma~14.5]{P46}}.
\end{proof}

\begin{lemma}[Fibre pairing]\label{pff:fibresum}
Let $R$ be positive semidefinite, $\mathcal P$ a finite set of heads, and $l_\theta\in\pbmGf_T$ ($\theta\in\Theta$) with
the same $e=l_\theta^{-1}z$ and in distinct right cosets of $\pffKM$ ($l_\theta l_{\theta'}^{-1}\in\pffKM\Rightarrow
\theta=\theta'$). Then $\mu_H(H_T\cap\pffKM)\sum_\theta\mathcal H(l_\theta^{-1})^{\vphantom{1}}\le\Rel\sum_{p,q}
R_{pq}\Pi_{p,q}(e)$, with $\mathcal H$ formed with the $\mathcal O_p(\cdot)$ as in Lemma~\ref{pff:headsum}.
\end{lemma}
\begin{proof}
Fix $l_0$ with $l_0^{-1}z=e$; then $l_\theta=k_\theta l_0$ with $k_\theta\in H_T$. By definition
$\sum R_{pq}\Pi_{p,q}(e)=\int_{H_T}\phi(k)\,d\mu_H(k)$ with $\phi(k)=\sum_{p,q}R_{pq}\pfrip{\mathcal O_p(kl_0)}
{\mathcal O_q(kl_0)}\ge0$. For $k'\in H_T\cap\pffKM$, $\mathcal O_p(k'kl_0)=\mathcal O_p(kl_0)$: by
definition $\mathcal O_p(k'kl_0)=\pfrkap_\varrho(p,1)\int_{\pffG}\pffW(y)\otimes\overline{\pffW(k'kl_0w_p^{-1}y)}\,dy$, and
$\pffW(k'x)=\pffW(x)$ for $k'\in H_T\cap\pffKM$ by \textup{(W2)} in the setting of Proposition~\ref{pfg:fibrepairing}
($H_T$ normalises $U_T$ and fixes the frame). So $\phi$ is constant on the cosets
$(H_T\cap\pffKM)k_\theta$, where it equals $\mathcal H(l_\theta^{-1})$. These cosets are disjoint, and each has
measure $\mu_H(H_T\cap\pffKM)$ by right invariance.
\end{proof}

\begin{lemma}[The normalised mass]\label{pff:rhoD}
Put $D=\#\pbHb\cdot\#\pbUb^{\,2}$. Then $\pffr(\pffZ)\le\mu_H(H_T\cap\pffKM)\,\pbvol(\pffKM)^2\,\#C\cdot D$, and
$\pbmaS=D\cdot\#\mathcal C\cdot\mathsf c_S$.
\end{lemma}
\begin{proof}
$\pbvol(\pffZ)=\#C\,\pbvol(\pffKM)$ (Lemma~\ref{pff:vol}(c)); $\pffl^{-1}=\#\pbUb$ (\S\ref{pff:ss-statements}); and
$\mu_H(H_T\cap\pffK)=1$ with $[H_T\cap\pffK:H_T\cap\pffKM]=\#\pbHb$ (reduction, {\cite[Lemma~2.1]{P51}}), so
$\mu_H(H_T\cap\pffKM)\#\pbHb=1$. Hence $\pffr(\pffZ)=\#C\,\pbvol(\pffKM)^2\#\pbUb^2=\mu_H(H_T\cap\pffKM)\pbvol(\pffKM)^2\#C
\cdot D$. The second identity is $\#\pbR=\#\pbHb\#\pbUb$ and Remark~\ref{pff:rem-L2}(2).
\end{proof}

\begin{lemma}[The weights]\label{pff:weights}
For a head $p\in\pbmT^\varrho$ with last entry $t$, let
\[
\rho_p=\prod_\tau\frac{\sqrt{\pbmcG}(|x_{\tau,\varrho}|^2-1)^{(k-1)/2}}
{d_\tau}\prod_{a<\varrho}|x_{\tau,a}|^{k+1+2(\varrho-a-1)}
\]
with $x_\tau=[p,t]_\tau$ and $d=\prod_\tau d_\tau$. Then
$\rho_p\,\pfrkap_\varrho(p,1)=\pbmrS\sqrt{G_t}$.
\end{lemma}
\begin{proof}
Put $d_a=k+1+2(\varrho-a-1)=k-1+2(\varrho-a)$. Then $\pfrkap_\varrho(p,1)=\prod_{a<\varrho}|p_a|_T^{d_a/2}$, and $|x_{\tau,a}|=
s_{\tau,a}|\tau(p_a)|$. The product formula for the $T$-units $p_a$ gives $|p_a|_T\prod_\tau|\tau(p_a)|^2=1$, so
$\prod_\tau\prod_{a<\varrho}|x_{\tau,a}|^{d_a}\cdot\pfrkap_\varrho(p,1)=\prod_\tau\prod_{a<\varrho}s_{\tau,a}^{d_a}$. And $s_{\tau,\varrho}=
N^{-1}$, so $|x_{\tau,\varrho}|^2-1=(|\tau(t)|^2-N^2)/N^2$ and $\prod_\tau(|x_{\tau,\varrho}|^2-1)^{(k-1)/2}=\sqrt{G_t}
\prod_\tau s_{\tau,\varrho}^{k-1}$. Comparing with the definition of $\pbmrS$ in {\cite[(40)]{P51}} gives the
claim.
\end{proof}

\begin{proof}[Proof of Theorem~\ref{pff:S}]
Group the family by the last entry $t\in\{t_j\}$; it suffices to prove the inequality for the subfamily
$\mathsf I_t$ with last entry $t$, the right side being additive in $t$. Fix $t$, let $\mathcal P=\pbmT^\varrho$
and, for $p\in\mathcal P$, let $\tilde p=(p,t)$; put $R_{pq}=\prod_\tau J^\tau_n([p,t]_\tau,[q,t]_\tau)$. By
Lemma~\ref{pff:topreal} and the multiplicativity of inner products of tensors over $\tau$, for every $j\in
\mathsf I_t$ the matrix $R$ is the Gram matrix of the vectors $\rho_pO(w_p\sigma_j)$ ($O$ the orbital integral
of the cut Bessel function at the real places, with $\rho_p$ of Lemma~\ref{pff:weights}); in particular $R$ is
positive semidefinite. For $j\in\mathsf I_t$ choose $\kappa_j\in K$ with reduction $s_j$; by
Theorem~\ref{pff:F}(a), $\kappa_{j,T}\in\pffZ$.

Let $v_j\in E^k$ be the $X_0$-component of $\sigma_jz$; then $Nv_j^*v_j=t\aubar t-N^2$ {\cite[(32)]{P51}}, and
$g_j=g(c_j,c'_j,\sigma_j)$ satisfies $g_jz=x(c_j,v_j)_T-N^{-1}tz=:e_j$, read in the lattice basis; $x(c_j,
v_j)$ is integral off $T$ because $h_{c_j}^{-1}\sigma_jh_{c_j'}$ and the lattice basis are. Group $\mathsf I_t$ by
the pair $(c_j,v_j)$. For a fixed pair, group further by $(c'_j,\sigma_j)$: by Lemma~\ref{pff:sep} the elements
$l=g^{-1}$ for distinct $(c',\sigma)$ lie in distinct right cosets of $\pffKM$ ($l_1l_2^{-1}=g_1^{-1}g_2\in\pffKM$
forces equality), and all satisfy $l^{-1}z=e$; and for fixed $(c,c',\sigma)$ the $s_j$ are distinct.
By Lemma~\ref{pff:headsum} with $w=\pbmrS\sqrt{G_t}$ (Lemma~\ref{pff:weights}) and Lemma~\ref{pff:fibresum},
\[
\mu_H(H_T\cap\pffKM)\pbvol(\pffKM)^2\#C\sum_{j:\,(c_j,v_j)=(c,v)}\pbmrS^{\,2}G_t\Bigl\|\sum_p\pbmOrb_j\Bigr\|^2\le\Rel\sum_{p,q}
R_{pq}\Pi_{\tilde p,\tilde q}(e)=:\Psi_t(e),
\]
where $\pbmOrb_j=\pbmOrb_{c_j,s_j,c_j'}(w_p\sigma_j)$. By Lemma~\ref{pff:fibre}(b), $\Psi_t$ agrees on the quadric with a
function $\Psi\ge0$ (the positive part of the real part of $\sum R_{pq}\Phi_{\tilde p,\tilde q}$) that is locally
constant with compact support and invariant under $H_T\cap\pffKM$. By Lemma~\ref{pff:rhoD}, the left side is at
least $\pffr(\pffZ)D^{-1}\pbmrS^{\,2}G_t\sum_{j}\|\sum_p\pbmOrb_j\|^2$. Summing over the pairs $(c,v)$ and applying
Proposition~\ref{pff:siegelsum},
\[
\begin{aligned}\frac{\pffr(\pffZ)}D\pbmrS^{\,2}G_t\sum_{j\in\mathsf I_t}\Bigl\|\sum_p\pbmOrb_j\Bigr\|^2&\le(1+b^{-6})\#\mathcal C\mathsf c_SG_t\,
|\pi_N|_T^{-1}\int\Psi(y-N^{-1}tz)\,d\nu_{\alpha,T}(y)\\&=(1+b^{-6})\#\mathcal C\mathsf c_SG_t\,\Rel\sum_{p,q}R_{pq}\mathcal
D(\tilde p,\tilde q),\end{aligned}
\]
the last equality because $\Rel\sum R_{pq}\Pi_{\tilde p\tilde q}(e)\ge0$ on the quadric (the entrywise product of the positive semidefinite matrices $R$ and $(\Pi_{pq}(e))$ is positive semidefinite, as in the proof of Proposition~\ref{pff:top}(b)), so that there $\Psi$ is this function itself, and because almost every point $y-N^{-1}tz$ of the sphere lies on the quadric
(Definition~\ref{pff:lastentry}). Divide by $G_t\pffr(\pffZ)/D$ and use $D\#\mathcal C\mathsf c_S=\pbmaS$
(Lemma~\ref{pff:rhoD}) and $J^T_n=\mathcal D/\pffr(\pffZ)$.
\end{proof}
\subsection{The inputs of \textup{(F)} and \textup{(S)} proved in \S\ref{pfg:sec}}\label{pff:ss-pending}

The points used above that are proved in \S\ref{pfg:sec} are: Lemma~\ref{pff:twogauges}
(Corollary~\ref{pfg:twogaugesT}), the axioms (G0)--(G5) at $T$ (Lemma~\ref{pfg:gauge}), Lemma~\ref{pff:fibre}
(Lemma~\ref{pfg:coarea}, Proposition~\ref{pfg:fibrepairing}), and Theorem~\ref{pff:siegelformula}, derived from Weil's
theorem (Proposition~\ref{pfg:siegel}). The published results used are Theorem~\ref{pfg:weil}, Witt's theorem for hermitian forms
(in the form of {\cite[Lemma~7.12]{P51}}) and the Smith normal form.



\providecommand{\auO}{\mathcal O}
\providecommand{\auA}{\mathbb A}
\providecommand{\auU}{\mathrm U}
\providecommand{\auGL}{\mathrm{GL}}
\providecommand{\auPl}{\Phi}
\providecommand{\aurd}{\mathrm{rd}_E}
\providecommand{\aubar}[1]{\overline{#1}}
\providecommand{\panot}[1]{\par\smallskip{\raggedright\noindent\textbf{[Not proved here:} #1\textbf{]}\par}\smallskip}
\providecommand{\pbvol}{\operatorname{vol}}
\providecommand{\pbmvol}{\mathsf v}
\providecommand{\pffET}{E_T}
\providecommand{\pffOT}{\auO_T}
\providecommand{\pffG}{G_T}
\providecommand{\pffK}{K_T}
\providecommand{\pffKM}{K_T(M)}
\providecommand{\pffch}{\chi_M}
\providecommand{\pffW}{W_Z}
\providecommand{\pffZ}{\mathsf Z}
\providecommand{\pfrK}{\mathsf K}
\providecommand{\pfrG}{\mathsf G}
\providecommand{\pfrU}{\mathsf U}
\providecommand{\pfrR}{\mathsf R}
\providecommand{\pfrCl}{\mathcal C}
\providecommand{\pfrkap}{\kappa}
\providecommand{\pfrdep}{\operatorname{dp}}
\providecommand{\pfrip}[2]{\langle#1,#2\rangle}
\providecommand{\pbmXp}{X^{+}}
\providecommand{\pfgH}{\mathrm{Herm}_T}
\providecommand{\pfgS}{\mathrm{Skew}_T}
\providecommand{\pfgOh}{\auO_{\mathrm h}}
\providecommand{\pfgOs}{\auO_{\mathrm s}}
\providecommand{\pfgE}{\mathsf e}
\providecommand{\pfgD}{\mathsf D}
\providecommand{\pfgEb}{\mathcal E}
\providecommand{\pfgQ}{\mathsf Q_N}
\providecommand{\pfgz}{\zeta_{\pffOT}}

\section{The finite places: gauges at the top, fibres and Siegel's formula}\label{pfg:sec}

This section proves the lemmas used in \S\ref{pff:sec} and listed in \S\ref{pff:ss-pending}. We use the
notation of \S\ref{pff:sec} and of \S\ref{pfr:sec} (place data, gauges, the axioms (G0)--(G5) of
\S\ref{pfr:ss-gauge} in their present form, heads, depth ideals $\mathfrak d_a$, nearness), and cite their
statements by label.

\emph{Contents of this section.}
\begin{itemize}
\item[(i)] \emph{Two gauges} (\S\ref{pfg:ss-two}): Theorem~\ref{pfg:twogauges}, a statement for every place
datum and every pair of gauges, and its case at $T$, Corollary~\ref{pfg:twogaugesT}, which is
Lemma~\ref{pff:twogauges} exactly as stated there. The proof is complete: the comparison of lengths is
carried down from $Y^{(\varrho)}$ to the whole space through short duals, and below the length $j$ of the
heads the short duals come from units of $\mathfrak d_a$ produced by Cramer's rule
(Lemmas~\ref{pfg:cramer}, \ref{pfg:headunits}). The off-diagonal block $h(se_b^*,e_c)$ is not needed.
\item[(ii)] \emph{Fourier analysis at $T$} (\S\ref{pfg:ss-fourier}): the duality of self-adjoint and skew
elements, a bound for Gauss sums of residue classes (Lemma~\ref{pfg:gauss}), and the construction of
the fibre measures of a unimodular hermitian lattice, with their continuity, the disintegration and a
bound for the mass of residue classes (Proposition~\ref{pfg:fibres}).
\item[(iii)] \emph{The gauges at $T$} (\S\ref{pfg:ss-gauge}): Lemma~\ref{pfg:gauge}, a complete proof of
Lemma~\ref{pff:gauge}, including the axioms (G0) and the clause $\tau_0\delta^{-1}\in\Gamma(1)$ of (G1),
and the counts (G4) and (G5) written out in full.
\item[(iv)] \emph{The fibre measures at the last entry and the step function} (\S\ref{pfg:ss-fibre}):
Lemma~\ref{pfg:coarea} (Lemma~\ref{pff:fibre}(a)) and Proposition~\ref{pfg:fibrepairing} (Lemma~\ref{pff:fibre}(b)); what
Proposition~\ref{pff:top}(c) uses is that $\Pi_{p',q'}=\Pi_{p,q}$ for \emph{all} $(p',q')$ near $(p,q)$.
\item[(v)] \emph{Siegel's formula} (\S\ref{pfg:ss-siegel}): we state Weil's Siegel--Weil theorem for the
dual pair $(\auU(1,1),\auU(V))$ in the range of absolute convergence (Theorem~\ref{pfg:weil}, cited) and
derive from it Theorem~\ref{pff:siegelformula} with every constant (Proposition~\ref{pfg:siegel}): the
measures are checked to be Tamagawa measures, and the archimedean density is computed. So the
normalisation of Theorem~\ref{pff:siegelformula} is verified directly.
\end{itemize}
The only results cited without proof are Theorem~\ref{pfg:weil} and the Smith normal form over a principal
ideal domain; Witt's theorem is used in the form of {\cite[Lemma~7.12]{P51}}, proved there.

\subsection{Conventions}\label{pfg:ss-conv}
The hermitian form is written $h(u,v)=u^*\mathsf Jv$, conjugate linear in $u$ (the opposite of
convention (Cv1) of {\cite[\S2.1]{P51}}; by Remark~\ref{pfg:rem-conv} the statements of this section
hold for either), with Gram matrix $\mathsf J$ of the
basis of \S\ref{pfr:ss-setting}; vectors are columns of coordinates in that basis, and matrices act on them.
For a vector $\omega$ we write $\omega_{e^*_a}$ for its coordinate on $e^*_a$ and $\omega_{e_a}$ for its
coordinate on $e_a$; since $e_a$ is orthogonal to every basis vector except $e^*_a$ and $h(e_a,e^*_a)=1$,
\[
\omega_{e^*_a}=h(e_a,\omega).
\]
Put $\pfgEb_{<l}=\langle e_b:b<l\rangle$. A vector $\omega$ is \emph{orthogonal below $l$} if $h(e_b,\omega)=0$
for $b<l$, i.e.\ $\omega_{e_b^*}=0$ for $b<l$; and $\omega\equiv\omega'\ (\mathrm{mod}\ \pfgEb_{<l})$ means
$\omega-\omega'\in\pfgEb_{<l}$. The module $Y^{(l)}$ of \S\ref{pfr:ss-gauge} is the quotient of the vectors
orthogonal below $l$ by $\pfgEb_{<l}$. Sizes refer to a gauge (\S\ref{pfr:ss-gauge}); we use the
consequences of (G0) listed there: a sum of at most $r+1$ terms of size $i$ has size $i+1$, a product of
numbers of sizes $i,i'$ has size $i+i'$, a product of matrices (or a matrix and a vector) of sizes $i,i'$ has
size $i+i'+1$, $h(u,v)$ has size $i+i'+3$ for $u,v$ of sizes $i,i'$ (as $\mathsf J$ has size $1$), the
inverse of an isometry of size $i$ has size $i+4$, and an $l\times l$ determinant ($l\le r+1$) of entries of
size $i$ has size $l(i+1)$. Sizes are monotone: size $i$ implies size $i'$ for $i'\ge i$ (as $1\in\Gamma(i'-i)$).

We use two properties of the groups of \S\ref{pfr:ss-setting}.
\begin{enumerate}
\item[(R)] An element $\xi\in\pfrR^{(l)}$ satisfies $\xi e_a\in e_a+\pfgEb_{<a}$ for $a<l$, and so does
$\xi^{-1}$. Hence $\xi$ preserves each $\pfgEb_{<a}$ ($a\le l$), maps vectors orthogonal below $a$ to
vectors orthogonal below $a$ for $a\le l$ (as $h(e_b,\xi\omega)=h(\xi^{-1}e_b,\omega)$), and does not
change the coordinate on $e^*_a$ of a vector orthogonal below $a$, for $a<l$.
\item[(U)] The Bessel unipotent group $\pfrU$ is contained in $\pfrR^{(\varrho)}$, and every $u\in\pfrU$ acts
trivially on $Y^{(\varrho)}$: $u\omega\equiv\omega\ (\mathrm{mod}\ \pfgEb_{<\varrho})$ for every $\omega$
orthogonal below $\varrho$. (The group $\pfrU$ is block unitriangular for the flag $\pfgEb_{<1}\subset\dots
\subset\pfgEb_{<\varrho}\subset\pfgEb_{<\varrho}^\perp\subset\dots$, with middle block $1$.) In coordinates:
for a vector $w$ orthogonal below $\varrho$, $(u-1)w$ has non-zero coordinates only on the $e_a$, $a<\varrho$.
\end{enumerate}

\subsection{Two gauges}\label{pfg:ss-two}

Throughout, $\mathfrak F$ is a place datum (\S\ref{pfr:ss-setting}) with level $M$, box $\pfrK$ and close pairs
$\pfrCl$, and sizes refer to a gauge $\Gamma$ for $\mathfrak F$ unless another gauge is named.

\begin{definition}\label{pfg:def}
Let $g,g'\in\pfrG$, let $r$ be a unit with $\bar r=r$ and $r\in\Gamma(1)$, $l\le\varrho$, $a<\varrho$ and
$t,c\ge0$.
\begin{enumerate}
\item[(a)] \emph{The lengths at $g$ and $g'$ differ by $t$ modulo $r$ on $Y^{(l)}$}, written $\pfgD_r(g,g';l,t)$,
if for every $i\ge0$ and every $\omega$ orthogonal below $l$ with $g^{-1}\omega$ of size $i$ there is
$\omega'\equiv\omega\ (\mathrm{mod}\ \pfgEb_{<l})$ such that $r^{-1}(g^{-1}\omega-g'^{-1}\omega')$ has size
$i+t$.
\item[(b)] $e_a$ has a \emph{short lift of size $c$ at $g$} if some $f\equiv e_a\ (\mathrm{mod}\
\pfgEb_{<a})$ has $g^{-1}f$ of size $c$.
\item[(c)] $e_a$ has a \emph{short dual of size $c$ for $(g,g')$ modulo $r$} if there are $\zeta$ orthogonal
below $a$ with $\zeta_{e^*_a}=1$ and $\zeta'\equiv\zeta\ (\mathrm{mod}\ \pfgEb_{<a})$ such that $g^{-1}\zeta$ and
$r^{-1}(g^{-1}\zeta-g'^{-1}\zeta')$ have size $c$.
\end{enumerate}
\end{definition}

\begin{lemma}[The pairing with a dual]\label{pfg:pairing}
Let $\zeta$ be orthogonal below $a$ with $\zeta_{e^*_a}=1$, $\zeta'\equiv\zeta\ (\mathrm{mod}\ \pfgEb_{<a})$, and
let $\omega$ be orthogonal below $a+1$ and $\omega'\equiv\omega\ (\mathrm{mod}\ \pfgEb_{<a+1})$. Then
$h(\zeta',\omega')-h(\zeta,\omega)=\mu$, the coordinate of $\omega'-\omega$ on $e_a$.
\end{lemma}
\begin{proof}
$\omega'-\mu e_a\equiv\omega\ (\mathrm{mod}\ \pfgEb_{<a})$. The vector $\zeta'$ is orthogonal below $a$ (as
$\zeta$ is and the $e_b$ are isotropic and mutually orthogonal), so $h(\zeta',\omega'-\mu e_a)=h(\zeta',\omega)$;
and $h(\zeta',\omega)=h(\zeta,\omega)$ as $\zeta'-\zeta\in\pfgEb_{<a}$ and $\omega$ is orthogonal below $a$.
Finally $h(\zeta',e_a)=h(\zeta,e_a)=\overline{\zeta_{e^*_a}}=1$ (again because $\zeta'-\zeta\in\pfgEb_{<a}$).
Hence $h(\zeta',\omega')-\mu=h(\zeta,\omega)$.
\end{proof}

\begin{lemma}[One floor down]\label{pfg:step}
Let $a<\varrho$, $r$ as in Definition~\ref{pfg:def}, and suppose that $e_a$ has short lifts of size $c$ at
$g$ and at $g'$ and a short dual of size $c$ for $(g,g')$ modulo $r$. Then
$\pfgD_r(g,g';a+1,t)$ implies $\pfgD_r(g,g';a,t+4c+12)$.
\end{lemma}
\begin{proof}
Let $f,f'$ be the short lifts at $g,g'$ and $\zeta,\zeta'$ the short dual. Let $\omega$ be orthogonal below $a$
with $g^{-1}\omega$ of size $i$. Put $\alpha=\omega_{e^*_a}$. As $f\equiv e_a\ (\mathrm{mod}\ \pfgEb_{<a})$ and
$\omega$ is orthogonal below $a$, $\alpha=h(e_a,\omega)=h(f,\omega)=h(g^{-1}f,g^{-1}\omega)$, of size $c+i+3$.
Put $\omega_0=\omega-\alpha\zeta$; it is orthogonal below $a+1$ (below $a$ since $\omega,\zeta$ are, and
$(\omega_0)_{e^*_a}=\alpha-\alpha=0$), and $g^{-1}\omega_0=g^{-1}\omega-\alpha g^{-1}\zeta$ has size
$i_0:=2c+i+4$. By $\pfgD_r(g,g';a+1,t)$ there is $\omega_0'\equiv\omega_0\ (\mathrm{mod}\ \pfgEb_{<a+1})$ with
$r^{-1}(g^{-1}\omega_0-g'^{-1}\omega_0')$ of size $i_0+t$. Let $\mu$ be the coordinate of $\omega_0'-\omega_0$ on
$e_a$. By Lemma~\ref{pfg:pairing} and the invariance of $h$ under $g,g'$,
\[
\mu=h(g'^{-1}\zeta',g'^{-1}\omega'_0)-h(g^{-1}\zeta,g^{-1}\omega_0).
\]
Write $v=g^{-1}\zeta$, $v'=g'^{-1}\zeta'$, $w=g^{-1}\omega_0$, $w'=g'^{-1}\omega_0'$; then $v$, $r^{-1}(v-v')$ have
size $c$ and $w$, $r^{-1}(w-w')$ have size $i_0+t$. Since $\bar r=r$,
$r^{-1}(h(v',w')-h(v,w))=-h(v',r^{-1}(w-w'))-h(r^{-1}(v-v'),w)$, and $v'=v-r\cdot r^{-1}(v-v')$ has size $c+2$; so
$r^{-1}\mu$ has size $c+i_0+t+6$. Put
\[
\omega'=\alpha\zeta'+\omega_0'-\mu f' .
\]
Then $\omega'\equiv\omega\ (\mathrm{mod}\ \pfgEb_{<a})$: indeed $\omega_0'-\mu f'\equiv\omega_0'-\mu e_a\equiv\omega_0$
(the coordinates of $\omega_0'-\omega_0\in\pfgEb_{<a+1}$ on $e_b$, $b<a$, do not matter, and the one on $e_a$ is
$\mu$), and $\zeta'\equiv\zeta$. Finally
\[
r^{-1}(g^{-1}\omega-g'^{-1}\omega')=r^{-1}(g^{-1}\omega_0-g'^{-1}\omega'_0)+\alpha\,r^{-1}(g^{-1}\zeta-g'^{-1}\zeta')
+(r^{-1}\mu)\,g'^{-1}f' ,
\]
a sum of three vectors of sizes $i_0+t$, $(c+i+3)+c$ and $(c+i_0+t+6)+c$, each at most
$i+t+4c+10$; so the sum has size $i+t+4c+11\le i+t+(4c+12)$.
\end{proof}

\begin{corollary}[Down to the whole space]\label{pfg:whole}
Let $j\le\varrho$, and suppose that for every $a<j$ the vector $e_a$ has short lifts of size $c$ at $g$ and at
$g'$ and a short dual of size $c$ for $(g,g')$ modulo $r$. Then $\pfgD_r(g,g';j,t)$ implies
$\pfgD_r(g,g';0,t+(4c+12)j)$.
\end{corollary}
\begin{proof}
Apply Lemma~\ref{pfg:step} for $a=j-1,j-2,\dots,0$.
\end{proof}

\begin{lemma}[From lengths to the matrix]\label{pfg:matrix}
If $\pfgD_r(g,gx;0,i)$, then $r^{-1}(x-1)$ has size $i+4$.
\end{lemma}
\begin{proof}
For a basis vector $e$ put $\omega=ge$; then $g^{-1}\omega=e$ has size $0$. On $Y^{(0)}$ the only lift of $\omega$
is $\omega$ itself, so $r^{-1}(e-x^{-1}e)$ has size $i$: the matrix $N=r^{-1}(x^{-1}-1)$ has size $i$. As $x$ is
an isometry, $x^{-1}=\mathsf J^{-1}x^*\mathsf J$, and applying this to $x^{-1}$ gives
$x=\mathsf J^{-1}(x^{-1})^*\mathsf J$; with $\bar r=r$ we get $r^{-1}(x-1)=\mathsf J^{-1}N^*\mathsf J$. As
$\mathsf J,\mathsf J^{-1}$ have size $1$ (G0), $N^*\mathsf J$ has size $i+2$ and $\mathsf J^{-1}N^*\mathsf J$ size
$i+4$.
\end{proof}

\begin{lemma}[The last quotient]\label{pfg:top}
Let $u,v\in\pfrU$ with $uy\in\pfrK$, $vyx\in\pfrK$, and $M^{-1}(vyx-uy)$ of size $2$. Then
$\pfgD_M(y,yx;\varrho,9)$.
\end{lemma}
\begin{proof}
$M$ is a unit with $\bar M=M$ and $M\in\Gamma(1)$ (\S\ref{pfr:ss-setting}, (G1)). Let $\omega$ be orthogonal below
$\varrho$ with $y^{-1}\omega$ of size $i$, and put $\omega'=v^{-1}u\,\omega$; by (U),
$\omega'\equiv\omega\ (\mathrm{mod}\ \pfgEb_{<\varrho})$. Then
\[
M^{-1}\bigl(y^{-1}\omega-(yx)^{-1}\omega'\bigr)=(vyx)^{-1}\cdot M^{-1}(vyx-uy)\cdot y^{-1}\omega ,
\]
because $(vyx)^{-1}(vyx-uy)y^{-1}=y^{-1}-(yx)^{-1}v^{-1}u$. Here $(vyx)^{-1}$ has size $1+4=5$ by (G1), so the
right side has size $(5+2+1)+i+1=i+9$.
\end{proof}

\begin{lemma}[Short lifts at a near point]\label{pfg:lift}
If $g$ is near at $j$, then for every $a<j$ the vector $e_a$ has a short lift of size $5$ at $g$.
\end{lemma}
\begin{proof}
Let $\xi\in\pfrR^{(j)}$ with $\xi g\in\pfrK$, and $f=\xi^{-1}e_a$; by (R), $f\equiv e_a\ (\mathrm{mod}\
\pfgEb_{<a})$. And $g^{-1}f=(\xi g)^{-1}e_a$ is a column of $(\xi g)^{-1}$, of size $5$ by (G1), (G0).
\end{proof}

\begin{lemma}[A unit of the depth ideal gives a short dual]\label{pfg:dualunit}
Let $c$ be a unit with $c\in\mathfrak d_a(g,x)$ and $c^{-1}\in\Gamma(1)$. Then $e_a$ has a short dual of size
$2$ for $(g,gx)$ modulo $M$.
\end{lemma}
\begin{proof}
By the definition of $\mathfrak d_a$ there is $\omega$ orthogonal below $a$ with $\omega_{e^*_a}=c$, and
$\omega'\equiv\omega\ (\mathrm{mod}\ \pfgEb_{<a})$ with $(g^{-1}\omega,(gx)^{-1}\omega')\in\pfrCl$. By (G2),
$g^{-1}\omega$ and $M^{-1}(g^{-1}\omega-(gx)^{-1}\omega')$ have size $1$. Put $\zeta=c^{-1}\omega$, $\zeta'=c^{-1}\omega'$:
$\zeta$ is orthogonal below $a$, $\zeta_{e^*_a}=1$, $\zeta'\equiv\zeta$, and $g^{-1}\zeta$,
$M^{-1}(g^{-1}\zeta-(gx)^{-1}\zeta')$ have size $1+1=2$.
\end{proof}

\begin{lemma}[Cramer's rule at a head]\label{pfg:cramer}
Let $j\le\varrho$, $r$ a head of length $j$ for $\Gamma$, $z\in\pfrG$, and $s,\eta\in\pfrR^{(j)}$ with $sz\in\pfrK$
and $\eta w_r^{-1}z\in\pfrK$. Then for $a<j$ the column $se^*_a$ has size $8\varrho+1$.
\end{lemma}
\begin{proof}
This is the column part of Lemma~\ref{pfr:headentries}, with the size read off its proof; we repeat the
argument. $\gamma:=(sz)(\eta w_r^{-1}z)^{-1}=sw_r\eta^{-1}$ has size $1+5+1=7$. Fix a basis index $i$ and let
$Z'$ be the $(a+1)\times(a+1)$ block of $\gamma$ with rows $e_0^*,\dots,e^*_{a-1},i$ and columns
$e_0,\dots,e_a$. As $\eta^{-1}e_d\in e_d+\pfgEb_{<d}$ (R), $Z'=\mathsf L'\mathsf C$ with $\mathsf C$ the block
of $\eta^{-1}$ on $e_0,\dots,e_a$, unitriangular, and $\mathsf L'$ the block of $sw_r$ with the same rows and
columns $e_0,\dots,e_a$; as $w_re_d=r_de_d^*$, $\mathsf L'_{bd}=s_{R(b),e^*_d}r_d$ ($R(b)$ the $b$-th row index).
For $b\le d$, $s_{e^*_b,e^*_d}=h(e_b,se^*_d)=h(s^{-1}e_b,e^*_d)=\delta_{bd}$ by (R). So $\mathsf L'$ is lower
triangular with diagonal $r_0,\dots,r_{a-1},s_{i,e^*_a}r_a$, and $\det Z'=P_{a+1}s_{i,e^*_a}$ with
$P_{a+1}=r_0\cdots r_a$. By (G0), $\det Z'$ has size $8(a+1)$; as $P_{a+1}^{-1}\in\Gamma(1)$ (head),
$s_{i,e^*_a}=P_{a+1}^{-1}\det Z'$ has size $8(a+1)+1\le8\varrho+1$.
\end{proof}

\begin{lemma}[Units of the depth ideals below a head]\label{pfg:headunits}
Let $j\le\varrho$, $j\ne0$, $p,q$ heads of length $j$ for $\Gamma$, $M^{-1}\in\Gamma(1)$, and let $y$, $yx$,
$w_p^{-1}y$, $w_q^{-1}yx$ be near at $j$. Then for every $a<j$ the ideal $\mathfrak d_a(y,x)$ contains a unit $c$
with $c^{-1}\in\Gamma(8\varrho+9)\subseteq\Gamma(10n)$.
\end{lemma}
\begin{proof}
Let $z\in\{y,yx\}$ and $r\in\{p,q\}$ the corresponding head. Let $s\in\pfrR^{(j)}$ with $sz\in\pfrK$; by
Lemma~\ref{pfg:cramer} (with $\eta$ given by the nearness of $w_r^{-1}z$) the column $se^*_a$ has size
$8\varrho+1$, so $z^{-1}e^*_a=(sz)^{-1}(se^*_a)$ has size $5+(8\varrho+1)+1=8\varrho+7$. By (G2) (as
$M^{-1}\in\Gamma(1)$) there is a unit $c$ with $c^{-1}\in\Gamma(8\varrho+9)$ such that $(c\omega_1,c\omega_2)\in\pfrCl$
for all vectors $\omega_1,\omega_2$ of size $8\varrho+7$. Put $\omega=ce^*_a$. It is orthogonal below $a$
($h(e_b,e^*_a)=0$ for $b\ne a$), and $(y^{-1}\omega,(yx)^{-1}\omega)=(c\,y^{-1}e^*_a,c\,(yx)^{-1}e^*_a)\in\pfrCl$; so
$\omega$ is stable at $(y,yx)$ for $a$ (with $\omega'=\omega$), and $c=\omega_{e^*_a}\in\mathfrak d_a(y,x)$. Finally
$8\varrho+9\le10\varrho+10=10n$.
\end{proof}

\begin{theorem}[Two gauges]\label{pfg:twogauges}
Let $\Gamma,\Gamma'$ be gauges for $\mathfrak F$, $j\le\varrho$, $p,q$ heads of length $j$ for $\Gamma$, and if
$j\ne0$ assume $M^{-1}\in\Gamma(1)$ and $\Gamma(10n)\subseteq\Gamma'(1)$. Let $y,x\in\pfrG$ and suppose:
\begin{enumerate}
\item[(a)] there are $u,v\in\pfrU$ with $uy\in\pfrK$, $vyx\in\pfrK$ and $M^{-1}(vyx-uy)$ of size $2$ for $\Gamma'$;
\item[(b)] $w_p^{-1}y$ and $w_q^{-1}yx$ are near at $j$;
\item[(c)] for every $a$ with $j\le a<\varrho$ the ideal $\mathfrak d_a(y,x)$ contains a unit $d_a$ with
$d_a^{-1}\in\Gamma'(1)$.
\end{enumerate}
Then $M^{-1}(x-1)$ has entries in $\Gamma'(99n)$.
\end{theorem}
\begin{proof}
All sizes are for $\Gamma'$ except where $\Gamma$ is named. \emph{Units for every $a<\varrho$.} For $a\ge j$ take
$d_a$ of (c). For $a<j$ (so $j\ne0$): $y$ and $yx$ are near at $j$ by (a), since $\pfrU\subseteq\pfrR^{(\varrho)}
\subseteq\pfrR^{(j)}$; with (b), Lemma~\ref{pfg:headunits} for $\Gamma$ gives a unit $d_a\in\mathfrak d_a(y,x)$ with
$d_a^{-1}\in\Gamma(10n)\subseteq\Gamma'(1)$. \emph{Short lifts and duals.} By Lemma~\ref{pfg:dualunit} (for
$\Gamma'$), each $e_a$, $a<\varrho$, has a short dual of size $2\le5$ for $(y,yx)$ modulo $M$; by (a) $y$ and $yx$
are near at $\varrho$, so by Lemma~\ref{pfg:lift} each $e_a$ has short lifts of size $5$ at $y$ and at $yx$.
\emph{Descent.} By Lemma~\ref{pfg:top}, $\pfgD_M(y,yx;\varrho,9)$; by Corollary~\ref{pfg:whole} with $c=5$,
$\pfgD_M(y,yx;0,9+32\varrho)$; by Lemma~\ref{pfg:matrix}, $M^{-1}(x-1)$ has size $13+32\varrho\le99n$.
\end{proof}

\begin{corollary}[Lemma~\ref{pff:twogauges}]\label{pfg:twogaugesT}
At the finite places, with the place datum $\mathfrak F_T$ of \S\ref{pff:ss-large}, the statement of
Lemma~\ref{pff:twogauges} holds for all gauges $\Gamma,\Gamma'$ for $\mathfrak F_T$ with $\pffOT\subseteq\Gamma'(2)$
(in particular for all gauges $\Gamma_\alpha$ of Lemma~\ref{pff:gauge}, whose sets $\Gamma_\alpha(0)$ contain
$\pffOT$).
\end{corollary}
\begin{proof}
We verify (a)--(c) of Theorem~\ref{pfg:twogauges}. (b) is part of ``$y$ active at $x$''; (c) is assumed. For
(a): $\langle\pffW(y),\pffW(yx)\rangle\ne0$ gives a coset $\pffKM a_c$ with $\pffch(ya_c^{-1})\ne0$ and
$\pffch(yxa_c^{-1})\ne0$, so $ya_c^{-1}=u^{-1}k_1$ and $yxa_c^{-1}=v^{-1}k_2$ with $u,v\in U_T$ and
$k_1,k_2\in\pffKM$ \eqref{pff:eq-chi}. Then $uy=k_1a_c$ and $vyx=k_2a_c$ lie in $\pffK$, and
$M^{-1}(vyx-uy)=k_1\cdot M^{-1}(k_1^{-1}k_2-1)\cdot a_c$ is integral, because $k_1^{-1}k_2\in\pffKM$ and
$k_1,a_c\in\pffK$ are integral. So its entries lie in $\pffOT\subseteq\Gamma'(2)$.
\end{proof}

We also record the bound for the unipotent factor and its consequence, used in \S\ref{pfg:ss-fibre}.

\begin{lemma}[Inverse of a unipotent factor]\label{pfg:unipinv}
Let $s\in\pfrU$ and suppose that the columns $se^*_a$, $a<\varrho$, have size $d$. Then $s^{-1}$ has size
$(4d+16)\varrho+5$.
\end{lemma}
\begin{proof}
Take $g=1$, $g'=s^{-1}$, $r=1$ (a unit with $\bar r=r$, $1\in\Gamma(1)$). \emph{Top.} $\pfgD_1(1,s^{-1};\varrho,0)$:
for $\omega$ orthogonal below $\varrho$, $\omega':=s^{-1}\omega\equiv\omega$ by (U), and $\omega-(s^{-1})^{-1}\omega'=0$.
\emph{Lifts.} $f=e_a$ at $1$, and $f'=s^{-1}e_a\equiv e_a\ (\mathrm{mod}\ \pfgEb_{<a})$ (R) at $s^{-1}$, with
$(s^{-1})^{-1}f'=e_a$; both of size $0$. \emph{Duals.} $\zeta=\zeta'=e^*_a$: orthogonal below $a$, coordinate $1$
on $e^*_a$, and $e^*_a-se^*_a$ has size $d+1$. By Corollary~\ref{pfg:whole} with $c=d+1$ and $j=\varrho$,
$\pfgD_1(1,s^{-1};0,(4d+16)\varrho)$, and by Lemma~\ref{pfg:matrix} (with $x=s^{-1}$) $s^{-1}-1$ has size
$(4d+16)\varrho+4$; so $s^{-1}=1+(s^{-1}-1)$ has size $(4d+16)\varrho+5$.
\end{proof}

\begin{corollary}[Size of a point under two cut conditions]\label{pfg:sizehead}
Let $W$ be a cut Bessel function for $\mathfrak F$, $r$ a head of length $\varrho$ for $\Gamma$, $l\in\pfrG^{(\varrho)}$,
and $y\in\pfrG$ with $W(y)\ne0$ and $W(lw_r^{-1}y)\ne0$. Then $y$ has size $60n^2$.
\end{corollary}
\begin{proof}
By the cut condition there are $s,u'\in\pfrU$ with $sy\in\pfrK$ and $u'lw_r^{-1}y\in\pfrK$; $\eta:=u'l$ lies in
$\pfrR^{(\varrho)}$ (which contains $\pfrU$ by (U) and $\pfrG^{(\varrho)}$, whose elements fix the $e_a$). By
Lemma~\ref{pfg:cramer} with $j=\varrho$, the columns $se^*_a$ have size $8\varrho+1$; by Lemma~\ref{pfg:unipinv},
$s^{-1}$ has size $(32\varrho+20)\varrho+5$; so $y=s^{-1}(sy)$ has size $32\varrho^2+20\varrho+7\le60n^2$.
\end{proof}

\subsection{Fourier analysis on hermitian lattices at $T$}\label{pfg:ss-fourier}

We keep the setting of \S\ref{pff:ss-setting}: $\pffET=\prod_{w\in T}E_w$, $T$ stable under conjugation,
$E/E^+$ unramified at every finite place, $|\cdot|_T$ the module, $dx$ the Haar measure with
$\pbvol(\pffOT)=1$, $\psi_T(x)=\psi_{E,f}(-x/\varpi)$.

\begin{lemma}[The character and two elements]\label{pfg:psi}
\begin{enumerate}
\item[(a)] $\psi_T(y+\bar y)=1$ for all $y$, and for $c\in\pffET$: $\psi_T(c\pffOT)=1$ if and only if $c\in\pffOT$.
\item[(b)] There are $\tau_0\in\pffOT$ with $\tau_0+\bar\tau_0=1$ and a unit $\delta$ of $\pffET$ with
$\bar\delta=-\delta$ and $\delta,\delta^{-1}\in\pffOT$; hence $\tau_0\delta^{-1}\in\pffOT$.
\item[(c)] There is $\varpi_1\in\pffOT$ with $\bar\varpi_1=\varpi_1$, a unit of $\pffET$, with $|\varpi_1|_w<1$ for
every $w\in T$.
\end{enumerate}
\end{lemma}
\begin{proof}
(a) The first claim is in \S\ref{pff:ss-setting}. The standard character $\psi_{E,f}=\psi_{\Q,f}\circ\mathrm{Tr}$ has
$\psi_{E,f}(y\widehat\auO_E)=1$ exactly for $y$ in the inverse different; at $T$ the different is generated by
$\varpi$, so $\psi_T(c\pffOT)=\psi_{E,f}(-c\varpi^{-1}\pffOT)=1$ iff $c\varpi^{-1}\in\varpi^{-1}\pffOT$, i.e.
$c\in\pffOT$. (b) $\tau_0$ is the element $\tau$ of Lemma~\ref{pff:L1}. For $\delta$ we work at a pair
$\{w,\bar w\}\subseteq T$. If $w\ne\bar w$, the components at $w,\bar w$ are exchanged by the involution, and
$\delta=(1,-1)$ works. If $w=\bar w$, $E_w/E^+_v$ is an unramified quadratic extension, the involution induces
the non-trivial automorphism of the residue extension, and for a unit $\epsilon$ of $\auO_w$ whose residue is
not in the residue field of $E^+_v$, $\delta=\epsilon-\bar\epsilon$ is skew and has non-zero residue, so it is a
unit of $\auO_w$. (c) At a pair $w\ne\bar w$ take $(\varpi_v,\varpi_v)$ for a uniformiser $\varpi_v$ of $E^+_v$; at
$w=\bar w$ a uniformiser of $E^+_v$, which is one of $E_w$ (unramified).
\end{proof}

The trace-one element $\tau_0$ and the skew unit $\delta$ are fixed from now on; they are the elements of axiom
(G1). Put $\pfgH=\{\beta\in\pffET:\bar\beta=\beta\}$, $\pfgS=\{b:\bar b=-b\}$, $\pfgOh=\pfgH\cap\pffOT$,
$\pfgOs=\pfgS\cap\pffOT$, with the Haar measures $d\beta$, $db$ giving $\pfgOh$, $\pfgOs$ volume $1$, and the
pairing
\[
\pfgE(b,\beta)=\psi_T(\tau_0b\beta)\qquad(b\in\pfgS,\ \beta\in\pfgH).
\]

\begin{lemma}[Duality]\label{pfg:duality}
\begin{enumerate}
\item[(a)] Every $x\in\pffOT$ can be written $x=\tau_0h+s$ with $h=x+\bar x\in\pfgOh$, $s\in\pfgOs$, and
$x=\tau_0\delta h_1+h_2$ with $h_1=\delta^{-1}(x-\bar x)\in\pfgOh$, $h_2\in\pfgOh$. $\psi_T(\gamma)=1$ for
$\gamma\in\pfgH$.
\item[(b)] $\pfgE$ is a continuous bicharacter, trivial on $\pfgOs\times\pfgOh$; for $b\in\pfgS$, $\pfgE(b,\cdot)\equiv1$
on $\pfgOh$ implies $b\in\pfgOs$; for $\beta\in\pfgH$, $\pfgE(\cdot,\beta)\equiv1$ on $\pfgOs$ implies $\beta\in
\pfgOh$. For $c\in\pfgH$, $\pfgE(b,c\beta)=\pfgE(cb,\beta)$.
\item[(c)] For $b\in\pfgS$ and $z\in\pffET$: $\pfgE(b,z+\bar z)=\psi_T(bz)$.
\item[(d)] The map $\pfgS\times\pfgS\to\pffET$, $(b_1,b_2)\mapsto\tau_0\delta^{-1}b_1+b_2$, is an isomorphism of
topological groups carrying $\pfgOs\times\pfgOs$ onto $\pffOT$, and so carries $db\,db$ to $dx$. Consequently, for
a unit $u$ of $\pffET$, $\pbvol\{b\in\pfgS:ub\in\pffOT\}^2\le|u|_T^{-1}$; and for a unit $c\in\pfgH$,
multiplication by $c$ multiplies $db$ and $d\beta$ by $|c|_T^{1/2}$.
\end{enumerate}
\end{lemma}
\begin{proof}
(a) $s=x-\tau_0(x+\bar x)=\bar\tau_0x-\tau_0\bar x$ is skew and integral. $x-\bar x=\tau_0\delta h_1-
\overline{\tau_0\delta h_1}=(\tau_0+\bar\tau_0)\delta h_1=\delta h_1$, and $h_2=x-\tau_0(x-\bar x)=\bar\tau_0x+\tau_0\bar x$
is self-adjoint and integral. For $\gamma\in\pfgH$, $\gamma=\tau_0\gamma+\overline{\tau_0\gamma}$ is a trace, so
$\psi_T(\gamma)=1$. (b) Continuity and multiplicativity are clear, and $\tau_0b\beta\in\pffOT$ for integral $b,\beta$.
If $\pfgE(b,\cdot)\equiv1$ on $\pfgOh$, then for $x=\tau_0h+s\in\pffOT$ as in (a), $\psi_T(bx)=\psi_T(\tau_0bh)
\psi_T(bs)=1$, as $bs\in\pfgH$; so $b\in\pffOT$ by Lemma~\ref{pfg:psi}(a). If $\pfgE(\cdot,\beta)\equiv1$ on $\pfgOs$, then
for $x=\tau_0\delta h_1+h_2$, $\psi_T(x\beta)=\pfgE(\delta h_1,\beta)\psi_T(h_2\beta)=1$ (as $\delta h_1\in\pfgOs$ and
$h_2\beta\in\pfgH$); so $\beta\in\pffOT$. The last identity is clear. (c) $\overline{\tau_0b\bar z}=-\bar\tau_0bz$, and
$\psi_T(y)=\psi_T(-\bar y)$ by (a) of Lemma~\ref{pfg:psi}; so $\psi_T(\tau_0b\bar z)=\psi_T(\bar\tau_0bz)$ and
$\pfgE(b,z+\bar z)=\psi_T((\tau_0+\bar\tau_0)bz)=\psi_T(bz)$. (d) The inverse is $x\mapsto(\delta(x+\bar x),x-\tau_0(x+
\bar x))$: indeed $x+\bar x=(\tau_0+\bar\tau_0)\delta^{-1}b_1=\delta^{-1}b_1$ for $x=\tau_0\delta^{-1}b_1+b_2$. Both maps
preserve integrality, so the compact open subgroups correspond and the Haar measures, both normalised to give them
volume $1$, correspond. If $S=\{b:ub\in\pffOT\}$, the image of $S\times S$ lies in $\{x:ux\in\pffOT\}$ (as
$\tau_0\delta^{-1}$ is integral), of volume $|u|_T^{-1}$. Multiplication by $c\in\pfgH$ on $\pfgS\times\pfgS$
corresponds to multiplication by $c$ on $\pffET$, of module $|c|_T$, and acts the same way on both factors; and
$\beta\mapsto\delta\beta$ is a measure preserving isomorphism $\pfgH\to\pfgS$ commuting with $c$.
\end{proof}

\begin{lemma}[Orthogonality]\label{pfg:orth}
Let $\varpi_1$ be as in Lemma~\ref{pfg:psi}(c), $\ell\ge0$, $B_\ell=\varpi_1^{-\ell}\pfgOs$. For $y\in\pfgH$,
\[
\int_{B_\ell}\pfgE(b,y)\,db=\pbvol(B_\ell)\,\mathbf 1[y\in\varpi_1^{\ell}\pfgOh],\qquad
\pbvol(B_\ell)=|\varpi_1|_T^{-\ell/2}=\pbvol(\varpi_1^\ell\pfgOh)^{-1}.
\]
\end{lemma}
\begin{proof}
$b\mapsto\pfgE(b,y)$ is a character of the compact group $B_\ell$; it is trivial iff $\pfgE(b',\varpi_1^{-\ell}y)=1$
for all $b'\in\pfgOs$, i.e.\ (Lemma~\ref{pfg:duality}(b)) iff $\varpi_1^{-\ell}y\in\pfgOh$. The volumes follow from
Lemma~\ref{pfg:duality}(d).
\end{proof}

Let $\pfgz(\mu)=\sum_I[\pffOT:I]^{-\mu}=\prod_{w\in T}(1-q_w^{-\mu})^{-1}$, the sum over the ideals of finite index
(\S\ref{pff:ss-large}); it is finite for every $\mu>0$. Every ideal of finite index is $w\pffOT$ for some
$w\in\pffOT\cap\pffET^\times$, with $[\pffOT:w\pffOT]=|w|_T^{-1}$ ($\pffOT$ is a finite product of discrete valuation
rings).

\emph{Setting for the rest of this subsection.} $V$ is a free $\pffET$-module of rank $\kappa$ with a hermitian form
$h$ and a lattice $L$ (a free $\pffOT$-module spanning $V$) on which $h$ is \emph{unimodular}: in a basis of $L$ the
Gram matrix $\mathsf J$ satisfies $\mathsf J^*=\mathsf J$, $\mathsf J,\mathsf J^{-1}$ integral. Coordinates and the
Haar measure $dx$ of $V$ ($\pbvol(L)=1$) are taken in that basis; $h(x,x)\in\pfgH$.

\begin{lemma}[Gauss sums of residue classes]\label{pfg:gauss}
Let $g$ be a unit of $\pffET$, $x_0\in V$, $\beta\in\pfgH$, and for $b\in\pfgS$
\[
G(b)=\int_{x_0+gL}\pfgE\bigl(b,h(x,x)-\beta\bigr)\,dx .
\]
Then $|G(b)|^2\le|g|_T^{2\kappa}[\pffOT:I_b]^{-\kappa}$, where $I_b=\{c\in\pffOT:bg\bar gc\in\pffOT\}$; and for every
$\mu>0$ with $2\mu+1\le\kappa$,
\[
\int_{\pfgS}|G(b)|\,db\le\pfgz(\mu)\,|g|_T^{\kappa-1}.
\]
\end{lemma}
\begin{proof}
\emph{Weyl differencing.} Writing the second point of $C=x_0+gL$ as $x+d$, $d\in gL$,
\[
|G(b)|^2=\int_C\int_{gL}\overline{\pfgE(b,h(x,x)-\beta)}\,\pfgE(b,h(x+d,x+d)-\beta)\,dd\,dx=\int_{gL}\pfgE(b,h(d,d))
\int_C\psi_T\bigl(b\,h(x,d)\bigr)\,dx\,dd ,
\]
since $h(x+d,x+d)-h(x,x)=h(d,d)+z+\bar z$ with $z=h(x,d)$, and Lemma~\ref{pfg:duality}(c). In the inner integral put
$x=x_0+gm$, $m\in L$: it is $\psi_T(bh(x_0,d))$ times the integral over $gL$ of the character $gm\mapsto\psi_T(b\bar g
h(m,d))$, which is $\pbvol(gL)=|g|_T^\kappa$ if the character is trivial and $0$ otherwise. The values $h(m,d)=m^*\mathsf
Jd$, $m\in L$, form the $\pffOT$-module generated by the coordinates of $\mathsf Jd$; by Lemma~\ref{pfg:psi}(a) the
character is trivial iff $b\bar g\,\mathsf Jd$ is integral, i.e.\ (unimodularity) iff $b\bar gd\in L$. Writing
$d=gm'$, this says $bg\bar gm'\in L$, i.e.\ every coordinate of $m'$ lies in $I_b$. The ideal $I_b$ has finite index
(at each $w\in T$ it is $\auO_w$ or $(bg\bar g)_w^{-1}\auO_w$), so $I_b=w_b\pffOT$ and the set of such $d$ is $gw_bL$,
of volume $|g|_T^\kappa|w_b|_T^\kappa$. Hence $|G(b)|^2\le|g|_T^{2\kappa}|w_b|_T^\kappa$.

\emph{Sum over the ideals.} For an ideal $I=w\pffOT$ of finite index let $S_I=\{b:I_b=I\}$. As $w\in I_b$ for $b\in S_I$,
$S_I\subseteq\{b:bg\bar gw\in\pffOT\}$, and by Lemma~\ref{pfg:duality}(d) $\pbvol(S_I)\le|g\bar gw|_T^{-1/2}=|g|_T^{-1}
[\pffOT:I]^{1/2}$. The sets $S_I$ cover $\pfgS$ and are measurable: $S_I$ is the set $\{b:bg\bar gw\in\pffOT\}$ minus the union of the
countably many sets $\{b:bg\bar gw'\in\pffOT\}$ with $w'\pffOT\supsetneq I$. So
\[
\int|G|\le\sum_I\pbvol(S_I)\,|g|_T^\kappa[\pffOT:I]^{-\kappa/2}\le|g|_T^{\kappa-1}\sum_I[\pffOT:I]^{-(\kappa-1)/2}\le
|g|_T^{\kappa-1}\pfgz(\mu),
\]
as $(\kappa-1)/2\ge\mu$ and every index is $\ge1$.
\end{proof}

\begin{proposition}[Fibre measures]\label{pfg:fibres}
Assume $\kappa\ge3$. For $\phi\colon V\to\C$ locally constant with compact support and $\beta\in\pfgH$ put
$F_\phi(b;\beta)=\int_V\phi(x)\,\pfgE(b,h(x,x)-\beta)\,dx$ and
\[
A_\ell(\phi;\beta)=\pbvol(\varpi_1^\ell\pfgOh)^{-1}\int_{\{x:\,h(x,x)-\beta\in\varpi_1^\ell\pfgOh\}}\phi(x)\,dx\qquad(\ell\ge0).
\]
\begin{enumerate}
\item[(a)] $b\mapsto F_\phi(b;\beta)$ is integrable on $\pfgS$, with $\int|F_\phi(b;\beta)|db=:C_\phi$ independent of
$\beta$, and
\[
\nu_\beta(\phi):=\int_{\pfgS}F_\phi(b;\beta)\,db=\lim_{\ell\to\infty}A_\ell(\phi;\beta),\qquad|A_\ell(\phi;\beta)|\le C_\phi .
\]
\item[(b)] If $\phi\ge0$ then $\nu_\beta(\phi)\ge0$. Hence $\nu_\beta$ is a positive Radon measure on $V$ (the
\emph{fibre measure} at $\beta$), and it is carried by the sphere $S_\beta=\{x:h(x,x)=\beta\}$.
\item[(c)] $\beta\mapsto\nu_\beta(\phi)$ is continuous.
\item[(d)] (\emph{Disintegration.}) $\int_{\pfgH}\nu_\beta(\phi)\,d\beta=\int_V\phi\,dx$; hence $\int_Vf\,dx=\int_{\pfgH}
\int f\,d\nu_\beta\,d\beta$ for every Borel $f\ge0$, $\beta\mapsto\int f\,d\nu_\beta$ being measurable.
\item[(e)] For every unit $g$, every $x_0$ and every $\mu>0$ with $2\mu+1\le\kappa$:
$\nu_\beta(x_0+gL)\le\pfgz(\mu)|g|_T^{\kappa-1}$.
\item[(f)] (\emph{Uniqueness.}) If $(\nu'_\beta)$ is a family of Radon measures, $\nu'_\beta$ carried by $S_\beta$, with
$\int\nu'_\beta(\phi)\,d\beta=\int\phi\,dx$ and $\beta\mapsto\nu'_\beta(\phi)$ continuous for every such $\phi$, then
$\nu'_\beta=\nu_\beta$ for every $\beta$.
\end{enumerate}
\end{proposition}
\begin{proof}
Choose a unit $g$ such that $\phi$ is constant on the classes $x_0+gL$ and vanishes off finitely many of them; then
$\phi$ is a finite combination of indicators of such classes.

(a) For an indicator, $F$ is the Gauss sum $G$ of Lemma~\ref{pfg:gauss}, integrable (take $\mu=1$; $\kappa\ge3$).
As $\pfgE(b,h(x,x)-\beta)=\pfgE(b,h(x,x))\pfgE(b,-\beta)$, $|F_\phi(b;\beta)|=|F_\phi(b;0)|$. By Fubini's theorem on the
compact set $B_\ell$ and Lemma~\ref{pfg:orth}, $\int_{B_\ell}F_\phi(b;\beta)\,db=\pbvol(B_\ell)\int\phi(x)\mathbf 1[h(x,x)-\beta
\in\varpi_1^\ell\pfgOh]\,dx=A_\ell(\phi;\beta)$. So $|A_\ell|\le C_\phi$, and as the $B_\ell$ increase to $\pfgS$ (as
$|\varpi_1|_w<1$ at every $w$), dominated convergence gives the limit.

(b) $A_\ell(\phi;\beta)\ge0$ for $\phi\ge0$. A positive linear functional on the locally constant compactly supported
functions extends uniquely to $C_c(V)$ (these functions are dense in $C_c(V)$ for the uniform norm with supports in a
fixed compact open set, and $|\nu_\beta(\phi)|\le\nu_\beta(\mathbf 1_U)\sup|\phi|$ for $\phi$ supported in $U$), hence is a
Radon measure (Riesz). If $\operatorname{supp}\phi\cap S_\beta=\emptyset$, the compact set $\{h(x,x)-\beta:x\in
\operatorname{supp}\phi\}$ does not contain $0$, so it misses $\varpi_1^\ell\pfgOh$ for large $\ell$, and
$A_\ell(\phi;\beta)=0$; hence $\nu_\beta(V\smallsetminus S_\beta)=0$.

(c) $\nu_\beta(\phi)=\int F_\phi(b;0)\pfgE(b,-\beta)\,db$ with $F_\phi(\cdot;0)$ integrable; dominated convergence.

(d) By Fubini, $\int_{\pfgH}A_\ell(\phi;\beta)\,d\beta=\int\phi(x)\,\pbvol(\varpi_1^\ell\pfgOh)^{-1}\pbvol\{\beta:h(x,x)-\beta\in
\varpi_1^\ell\pfgOh\}\,dx=\int\phi\,dx$. For $\ell\ge0$, $A_\ell(\phi;\beta)=0$ unless $\beta$ lies in the compact set
$\{h(x,x):x\in\operatorname{supp}\phi\}+\pfgOh$, and $|A_\ell|\le C_\phi$; dominated convergence gives the identity. Both $\phi\mapsto\int\phi\,dx$ and
$\phi\mapsto\int\nu_\beta(\phi)d\beta$ are then Radon measures agreeing on a dense class, so they are equal; the
extension to Borel $f\ge0$ and the measurability follow by the monotone class theorem.

(e) For $\phi=\mathbf 1_{x_0+gL}$, $\nu_\beta(\phi)=\int G\le\int|G|$ and Lemma~\ref{pfg:gauss}.

(f) Let $\phi_\ell(x)=\phi(x)\mathbf 1[h(x,x)-\beta\in\varpi_1^\ell\pfgOh]$, again locally constant with compact support.
On $S_{\beta'}$ it equals $\phi\cdot\mathbf 1[\beta'-\beta\in\varpi_1^\ell\pfgOh]$, so, $\nu'_{\beta'}$ being carried by
$S_{\beta'}$, the disintegration gives $A_\ell(\phi;\beta)=\pbvol(\varpi_1^\ell\pfgOh)^{-1}\int_{\beta+\varpi_1^\ell\pfgOh}
\nu'_{\beta'}(\phi)\,d\beta'$, a mean of a continuous function over a shrinking neighbourhood of $\beta$. Hence
$\nu'_\beta(\phi)=\lim_\ell A_\ell(\phi;\beta)=\nu_\beta(\phi)$ by (a).
\end{proof}

\begin{remark}\label{pfg:rem-fibres}
By (f), the fibre measures of any continuous disintegration (for instance those built from a flat chart of the quadric) coincide with these. In particular the measures $\nu_\beta$ of Definition~\ref{pff:lastentry}
(``$dx/d\langle x,x\rangle$'' on $X_{0,T}$) and the measure of the top quadric are those of
Proposition~\ref{pfg:fibres}, and the formula $\nu_\alpha(\Psi)=\int_{\pfgS}\int\Psi(y)\psi_T(\tau_0b(y^*y-\alpha))\,dy\,db$
used at the end of the proof of Lemma~\ref{pff:offT} is Proposition~\ref{pfg:fibres}(a) (the integral over $b$
converges absolutely, so the convention ``$0$ otherwise'' there is never invoked).
\end{remark}

\subsection{The gauges at $T$: the axioms in full}\label{pfg:ss-gauge}

We use the place datum $\mathfrak F_T$ of \S\ref{pff:ss-large} under (L1)--(L3), with the parameter $m$ of
\S\ref{pfr:ss-gauge} taken to be $m'$ (it enters only through (G1), the rank $2m+2$ of the blocks of (G4) and the
exponent $s$), and $X'$ as in
\S\ref{pfr:ss-gauge}: the span of the orthonormal vectors of (L1) left after two blocks of $2m'+2$ of them are set
aside; $d':=\operatorname{rank}X'$ and the volume exponent is $s=d'-(3n+1+m')$. Coordinates are taken in the basis of
(L1), which is a basis of $L_T$ with unimodular Gram matrix. For $u\in\pffET^{X'}$ (a vector of coordinates on the
orthonormal basis $x'_1,\dots,x'_{d'}$ of $X'$) and $\omega\in X'$ put $\langle u,\omega\rangle=\sum_i\bar u_i\omega_i$,
$\omega_i$ the coordinates of $\omega$.

\begin{lemma}[Products of close pairs are integral]\label{pfg:Dint}
For close pairs $(\omega_1,\omega_2)$, $(\omega'_1,\omega'_2)$ of $\mathfrak F_T$, the number
$M^{-1}(h(\omega_1,\omega'_1)-h(\omega_2,\omega'_2))$ is integral.
\end{lemma}
\begin{proof}
Write $\omega_2=\omega_1-Mb$, $\omega'_2=\omega'_1-Md$ with $b,d\in L_T$. As $\bar M=M$,
$M^{-1}(h(\omega_1,\omega'_1)-h(\omega_2,\omega'_2))=h(\omega_1,d)+h(b,\omega'_1)-M\,h(b,d)$, and $h$ is integral on
$L_T$ (unimodular), $M\in\pffOT$.
\end{proof}

\begin{lemma}[Common vectors]\label{pfg:common}
Let $g\in\pffK$, $x=y^{-1}uyk$ a position ($y\in\pffG$, $u\in U_T$, $k\in\pffKM$), and $v$ a $\varrho\times d'$ matrix
(the coefficients of $\varrho$ linear forms on $X'$). Put $d=d'-3\varrho$. There is $V\in\auGL_{d'}(\pffOT)$ such that for
every $t\in\pffOT^{d}$ the vector $\omega=\sum_i\tau_ix'_i$, $\tau=V\binom0t$, lies in the set $\mathcal W$ of (G5):
$v\tau=0$ and $(g^{-1}\omega,(gx)^{-1}\omega)\in\pfrCl$.
\end{lemma}
\begin{proof}
Let $B$ be the $r\times d'$ matrix whose columns are the coordinates of the $x'_i$ (integral, by (L1)), $N=yg^{-1}B$,
and $C$ the $3\varrho\times d'$ matrix whose rows are: the $\varrho$ rows of $v$; the rows of $N$ of index $e^*_a$
($a<\varrho$); the rows of $(u^{-1}-1)N$ of index $e_a$ ($a<\varrho$). \emph{Column reduction.} At each $w\in T$ (the
components of $\pffET$ are the fields $E_w$), after multiplying $C_w$ by a power of a uniformiser to make it integral,
the Smith normal form over the principal ideal domain $\auO_w$ gives $C_w=P_wD_wQ_w$ with $P_w\in\auGL_{3\varrho}(\auO_w)$,
$Q_w\in\auGL_{d'}(\auO_w)$ and $D_w$ diagonal of size $3\varrho\times d'$, whose columns of index $>3\varrho$ vanish. Let
$V=(Q_w^{-1})_w\in\auGL_{d'}(\pffOT)$; then the last $d$ columns of $CV$ vanish. \emph{Verification.} For
$\tau=V\binom0t$ we have $C\tau=0$: so $v\tau=0$; $N\tau$ has zero coordinates on the $e^*_a$, i.e.\ is orthogonal
below $\varrho$; and $(u^{-1}-1)N\tau$ has zero coordinates on the $e_a$. By (U) applied to $u^{-1}\in U_T$ and the
vector $N\tau$, the other coordinates of $(u^{-1}-1)N\tau$ vanish too; so $u^{-1}N\tau=N\tau$. As $N\tau=yg^{-1}\omega$,
\[
(gx)^{-1}\omega=k^{-1}y^{-1}u^{-1}yg^{-1}\omega=k^{-1}y^{-1}u^{-1}N\tau=k^{-1}y^{-1}N\tau=k^{-1}g^{-1}\omega .
\]
Now $g^{-1}\omega=g^{-1}B\tau$ is integral ($g\in\pffK$, $B$, $\tau$ integral), and $M^{-1}(g^{-1}\omega-k^{-1}g^{-1}\omega)
=M^{-1}(1-k^{-1})g^{-1}\omega$ is integral as $k^{-1}\in\pffKM$. So the pair is close.
\end{proof}

\begin{lemma}[Integral of a content index]\label{pfg:content}
Let $d'=a+d$, $V\in\auGL_{d'}(\pffOT)$, and for $u\in\pffET^{d'}$ let $w(u)\in\pffET^d$ be the last $d$ coordinates of
$V^*u$ and $\phi(u)=\pbvol\{c\in\pffET:c\,w(u)_j\in\pffOT\ \forall j\}\in(0,\infty]$. Let $\beta\in\pffOT$ be a unit of
$\pffET$ and $0\le s'<d$. Then
\[
\int_{\{u:\,\beta u\in\pffOT^{d'}\}}\phi(u)^{s'}\,du\le\pfgz(d-s')\,|\beta|_T^{-a-(d-s')} .
\]
\end{lemma}
\begin{proof}
$u\mapsto V^*u$ preserves $du$ and the set $\{\beta u\in\pffOT^{d'}\}$ ($V^*\in\auGL_{d'}(\pffOT)$). So the integral is
$\pbvol(\beta^{-1}\pffOT)^a\allowbreak \int_{\{\beta w\in\pffOT^d\}}\phi_0(w)^{s'}dw=|\beta|_T^{-a}\int\phi_0^{s'}$, with
$\phi_0(w)=\pbvol\{c:cw\in\pffOT^d\}$. For almost every $w$, at every $w'\in T$ some coordinate of $w$ has non-zero
component, and the content $\mathfrak c(w)=\sum_j\pffOT w_j$ is $t\pffOT$ for a unit $t$; then $\phi_0(w)=|t|_T^{-1}$, and if
$\beta w$ is integral, $I(w):=\beta t\pffOT$ is an ideal of finite index with $[\pffOT:I(w)]=|\beta t|_T^{-1}$. Given $I$,
$t$ is determined up to a unit of $\pffOT$, and the $w$ with $I(w)=I$ lie in $(t\pffOT)^d$, of volume $|t|_T^d$. Hence
\[
\int\phi_0^{s'}\le\sum_I|t|_T^{-s'}|t|_T^{d}=\sum_I\bigl(|\beta|_T[\pffOT:I]\bigr)^{-(d-s')}=|\beta|_T^{-(d-s')}\pfgz(d-s').\qedhere
\]
\end{proof}

\begin{lemma}[Gauges at $T$; Lemma~\ref{pff:gauge}]\label{pfg:gauge}
Let $\alpha\in\pffOT$ be a unit of $\pffET$ with $\bar\alpha=\alpha$, $\Gamma_\alpha(i)=\{c:\alpha^ic\in\pffOT\}$ and
$\Lambda_\alpha=2/|\alpha|_T$. Under (L1)--(L3), $\Gamma_\alpha$ with the loss $\Lambda_\alpha$ is a gauge for $\mathfrak F_T$:
it satisfies (G0)--(G5) of \S\ref{pfr:ss-gauge} in their present form, with $\tau_0,\delta$ of Lemma~\ref{pfg:psi}(b).
\end{lemma}
\begin{proof}
$\Gamma_\alpha(i)=\alpha^{-i}\pffOT$ is compact and increasing in $i$ ($\alpha$ integral), and $\Lambda_\alpha\ge2$ as
$|\alpha|_T\le1$.

(G0) $\alpha^{i+i'}cc'=(\alpha^ic)(\alpha^{i'}c')$; for a sum $\sigma$ of elements of $\Gamma_\alpha(i)$, $\alpha^i\sigma$ is
integral, so $\alpha^{i+1}\sigma$ is; $\pffOT$ is stable under $-$ and conjugation and $\bar\alpha=\alpha$; $0,1\in
\Gamma_\alpha(i)$; the Gram matrix of the basis (L1) and its inverse are integral (unimodularity, (L3)), hence in
$\Gamma_\alpha(0)\subseteq\Gamma_\alpha(1)$.

(G1) $m'\le9n$ by (L2); $P=1$; $M\in\pffOT$; $\tau_0\delta^{-1}\in\pffOT$ (Lemma~\ref{pfg:psi}(b)); the basis vectors are
integral (Lemma~\ref{pff:L1}) and the elements of $\pffK$ are integral matrices in that basis. All these lie in
$\Gamma_\alpha(1)$.

(G2) For a close pair, $\omega_1$ and $M^{-1}(\omega_1-\omega_2)$ are integral. If $M^{-1}\in\Gamma_\alpha(1)$ and $i\ge0$,
take $c=\alpha^{i+1}$: $\alpha^{i+2}c^{-1}=\alpha$ is integral, so $c^{-1}\in\Gamma_\alpha(i+2)$; and for $\omega_1,\omega_2$ of
size $i$, $c\omega_1=\alpha(\alpha^i\omega_1)$ is integral and $M^{-1}(c\omega_1-c\omega_2)=(\alpha M^{-1})\alpha^i(\omega_1-\omega_2)$
is integral; so $(c\omega_1,c\omega_2)$ is close.

(G3) For a unit $u$ with $\alpha^iu\in\pffOT$, $|\alpha|_T^i|u|_T\le1$, so $|u|_T\le|\alpha|_T^{-i}\le\Lambda_\alpha^i$; and
$\pbvol\Gamma_\alpha(i)=|\alpha|_T^{-i}\le\Lambda_\alpha^{i+1}$.

(G4) A block of $2m'+2$ of the orthonormal vectors of (L1) spans a free module of rank $\kappa=2m'+2\ge3$ on which
$h$ is the standard form $\sum_i\bar y_iy_i$, with the unimodular lattice $\pffOT^{2m'+2}$, and its Haar measure is the one of \S\ref{pfg:ss-fourier}. By Proposition~\ref{pfg:fibres} its
measure disintegrates along the form into the fibre measures $\nu_\beta$, and by (e) there with $\mu=m'$
($2m'+1\le\kappa$), $g=\alpha^{-i}$, $x_0=0$:
\[
\nu_\beta\bigl(\Gamma_\alpha(i)^{2m'+2}\bigr)\le\pfgz(m')\,|\alpha|_T^{-i(2m'+1)}\le2\Bigl(\frac{\Lambda_\alpha}2\Bigr)^{i(2m'+1)}
\le\Lambda_\alpha^{(2m'+2)(i+1)},
\]
using $\pfgz(m')\le2$ ((L2): $\pfgz(m')^n\le1+2^{-n}$).

(G5) Let $g\in\pffK$, $x$ a position, $v$ as there, and $V$ as in Lemma~\ref{pfg:common}; put $d=d'-3\varrho$,
$w(u)$ and $\phi(u)$ as in Lemma~\ref{pfg:content} (with $a=3\varrho$). $\phi$ is measurable. Axiom (G5) is stated with
the form of convention (Cv1), for which, for $u,\omega\in X'$ with coordinates $u_i,\omega_i$ on the orthonormal basis
$x'_i$, the product is $h_{\rm (Cv1)}(\omega,u)=\sum_i\omega_i\bar u_i=u^*\omega=\langle u,\omega\rangle$ (in the convention of
\S\ref{pfg:ss-conv} this is $h(u,\omega)$; the number is the same), and $h_{\rm (Cv1)}(\omega,cu)=\bar c\langle u,\omega\rangle$. We
write the verification with $\langle u,\omega\rangle=u^*\omega$. \emph{First condition.} By Lemma~\ref{pfg:common}, $\mathcal W$
contains the vectors $\omega$ with coordinates $\tau=V\binom0t$, $t\in\pffOT^d$, and for them
$\langle u,\omega\rangle=u^*V\binom0t=\sum_j\overline{w(u)_j}\,t_j$. These values fill the $\pffOT$-module generated by the
$\overline{w(u)_j}$, which is $\overline{\mathfrak c(w(u))}$; if $\mathfrak c(w(u))=t_0\pffOT$ with $t_0$ a unit, this is $\bar t_0\pffOT$,
of volume $|\bar t_0|_T=|t_0|_T=\phi(u)^{-1}$ ($T$ is stable under conjugation), and otherwise $\phi(u)=\infty$. So
$\pbvol\{h_{\rm (Cv1)}(\omega,u):\omega\in\mathcal W\}^{-1}\le\phi(u)$. \emph{Second condition.} The set $\mathcal D$ of (G5) consists of
the numbers $M^{-1}(h_{\rm (Cv1)}(\omega_1,\omega_1')-h_{\rm (Cv1)}(\omega_2,\omega_2'))$, the conjugates of the numbers of
Lemma~\ref{pfg:Dint}, so $\mathcal D\subseteq\pffOT$ by that lemma ($\pffOT$ is stable under conjugation). If
$h_{\rm (Cv1)}(\omega,cu)\in\mathcal D$ for all $\omega\in\mathcal W$, then $\bar c\,\overline{w(u)_j}=h_{\rm (Cv1)}(\omega_j,cu)\in\pffOT$ for the
vectors $\omega_j$ with $t=e_j$; conjugating, $c\,w(u)_j\in\pffOT$ for all $j$, and $c$ lies in the set whose volume is $\phi(u)$. \emph{The integral.} By
Lemma~\ref{pfg:content} with $\beta=\alpha^i$, $a=3\varrho$, $s'=s$ and $d-s=d'-3\varrho-s=m'+4$:
\[
\int_{\Gamma_\alpha(i)^{X'}}\phi^s\le\pfgz(m'+4)\,|\alpha|_T^{-i(3\varrho+m'+4)}\le2\Bigl(\frac{\Lambda_\alpha}2\Bigr)^{i(3\varrho+m'+4)}
\le\Lambda_\alpha^{99n(i+1)},
\]
as $\pfgz(m'+4)\le\pfgz(m')\le2$ and $3\varrho+m'+4\le3n+9n\le99n$.
\end{proof}

For $\alpha=1$ all the sets $\Gamma_1(i)$ are $\pffOT$ and the loss is $2$, as stated in Lemma~\ref{pff:gauge}.

\subsection{The fibre measures at the last entry, and the step function}\label{pfg:ss-fibre}

\emph{Notation.} $\pbmXp_T=X_{0,T}\perp\pffET z$; its \emph{lattice basis} is $x_i$ ($i\ne k_0$), $v=\bar\pi_N^{-1}x_{k_0}+N^{-1}z$,
$z$ (\S\ref{pff:ss-floors}), with Gram matrix unimodular ($x_i$ orthonormal and orthogonal to $v,z$; $h(v,v)=0$,
$h(v,z)=-1$, $h(z,z)=-N$). Let $dy$ be the Haar measure of $\pbmXp_T$ giving the $\pffOT$-span $L^+$ of the lattice basis
volume $1$, $dx$ that of $X_{0,T}$ giving $\bigoplus\pffOT x_i$ volume $1$, and $\nu^+_\beta$, $\nu^0_\beta$ the fibre
measures of Proposition~\ref{pfg:fibres} for these two unimodular lattices (ranks $k+1$ and $k$, both $\ge3$). As in
Remark~\ref{pfg:rem-fibres}, $\nu^0_\beta$ is the measure $\nu_\beta$ of Definition~\ref{pff:lastentry}, and the measure
$\theta^\sharp_\varrho$ of the top quadric (the image of the top floor measure $\theta_\varrho$ of $\mathfrak F_T$ under
$\sigma\mapsto\sigma^{-1}z$) is normalised to be $\nu^+_{-N}$, with $dy$ as just fixed (the continuous disintegration of $dy$ along $y\mapsto h(y,y)$, at
the value $h(z,z)=-N$). Put $\pfgQ=\{e\in\pbmXp_T:h(e,e)=-N\}$.

\begin{lemma}[Coarea; Lemma~\ref{pff:fibre}(a)]\label{pfg:coarea}
For $\phi$ locally constant with compact support on $\pbmXp_T$ and $\beta\in\pfgH$, put $\phi_t(x)=\phi(x-N^{-1}tz)$ on
$X_{0,T}$. Then $t\mapsto\nu^0_{\beta+t\bar t/N}(\phi_t)$ is continuous with compact support, and
\[
\nu^+_\beta(\phi)=|\pi_N|_T^{-1}\int_{\pffET}\nu^0_{\beta+t\bar t/N}(\phi_t)\,dt .
\]
For $\beta=-N$ this is Lemma~\ref{pff:fibre}(a), with $\beta(t)=t\bar t/N-N$; the case of a continuous $g$ with compact
support follows by uniform approximation by locally constant functions with supports in a fixed compact open set (both
sides are finite measures on it).
\end{lemma}
\begin{proof}
\emph{Measures.} Let $L_0=\bigoplus\pffOT x_i\oplus\pffOT z$ and $dy_0$ the Haar measure with $\pbvol(L_0)=1$. The matrix taking
the basis $(x_1,\dots,x_k,z)$ to the lattice basis has determinant $\bar\pi_N^{-1}$, so $\pbvol_{y_0}(L^+)=|\pi_N|_T^{-1}$ and
$dy=|\pi_N|_T\,dy_0$. The map $\Theta(x,t)=x-N^{-1}tz$ carries $dx\,dt$ to $|N|_T\,dy_0=|\pi_N|_T\,dy$ (the coordinate on $z$
is $-N^{-1}t$, and $|N|_T=|\pi_N|_T^2$). So $\int\phi\,dy=|\pi_N|_T^{-1}\iint\phi_t(x)\,dx\,dt$. Moreover
$h(\Theta(x,t),\Theta(x,t))=h(x,x)-t\bar t/N$, as $h(z,z)=-N$, $\bar N=N$, $x\perp z$.

\emph{Slabs.} With $A_\ell$ as in Proposition~\ref{pfg:fibres}, these two facts give
$A^+_\ell(\phi;\beta)=|\pi_N|_T^{-1}\int A^0_\ell(\phi_t;\beta+t\bar t/N)\,dt$. Choose $a$ with $\phi$ invariant under translation
by $\varpi_1^aL^+$; as $z\in L^+$, $\phi_t$ depends only on $t$ modulo $\varpi_1^aN\pffOT$. As $t=h(z,y)$ for
$y=\Theta(x,t)$, $\phi_t=0$ unless $t$ lies in the compact set $h(z,\operatorname{supp}\phi)$. So only finitely many
functions $\phi_t$ occur, and $|A^0_\ell(\phi_t;\cdot)|\le C:=\max_tC_{\phi_t}$. By Proposition~\ref{pfg:fibres}(a) for both
lattices and dominated convergence in $t$, letting $\ell\to\infty$ gives the formula. Continuity of the integrand: $\phi_t$
is locally constant in $t$, and $\beta'\mapsto\nu^0_{\beta'}(\phi_t)$ is continuous (Proposition~\ref{pfg:fibres}(c)).
\end{proof}

\begin{lemma}[The top orbit]\label{pfg:orbit}
$\pffG^{(\varrho)}=\auU(\pbmXp_T)$ (the elements fixing the $e_a,e^*_a$) acts transitively on $\pfgQ$, and for $e_0\in\pfgQ$ the
map $\sigma\mapsto\sigma^{-1}e_0$ from $\pffG^{(\varrho)}$ to $\pfgQ$ is open.
\end{lemma}
\begin{proof}
\emph{Transitivity.} It suffices to work at one pair $\{w,\bar w\}\subseteq T$. Let $e,e'\in\pfgQ$. As $h(e,e)=-N$ is a unit,
$\pbmXp=\pffET e\perp e^\perp=\pffET e'\perp e'^\perp$. If $w=\bar w$, $E_w$ is a field with non-trivial involution, and by
{\cite[Lemma~7.12(a)]{P51}} (characteristic $0$, $h(e,e)=h(e',e')\ne0$; see also \cite[Chapter~7, Theorem~9.1]{Scharlau}) there is an element of
$\auU(\pbmXp_w)$ taking $e$ to $e'$. If $w\ne\bar w$, the components at $w,\bar w$ are $X_1,X_2$ (vector spaces over
$E_w\cong E^+_v$), $h$ is given by a non-degenerate bilinear pairing $B\colon X_1\times X_2\to E_w$, the unitary group is
$\auGL(X_1)$ acting on $X_2$ contragrediently, and $e=(e_1,e_2)$, $e'=(e'_1,e'_2)$ with $B(e_1,e_2)=B(e'_1,e'_2)=c\ne0$. An
element $\gamma\in\auGL(X_1)$ with $\gamma e_1=e'_1$ and $\gamma(\ker B(\cdot,e_2))=\ker B(\cdot,e'_2)$ exists (both kernels are
complements of the lines), and its contragredient takes $e_2$ to $e'_2$ (both are characterised by their values on
$e'_1$ and on $\ker B(\cdot,e'_2)$).

\emph{Openness.} $\pffG^{(\varrho)}$ is locally compact and second countable, $\pfgQ$ is closed in $\pbmXp_T$, hence locally
compact. Let $U$ be a neighbourhood of $1$ and $W$ a compact neighbourhood of $1$ with $W^{-1}W\subseteq U$. Countably many
translates $g_nW$ cover $\pffG^{(\varrho)}$, so $\pfgQ=\bigcup_ng_nWe_0$ with compact $g_nWe_0$; by Baire's theorem some $g_nWe_0$,
hence $We_0$, has an interior point $we_0$ ($w\in W$), and $e_0$ is interior to $w^{-1}We_0\subseteq Ue_0$. So $\sigma\mapsto
\sigma e_0$ is open at $1$, hence everywhere, and so is $\sigma\mapsto\sigma^{-1}e_0$.
\end{proof}

We use three properties of the Bessel vector $\pffW$ of an intermediate group $\pffZ$ (\S\ref{pff:ss-setting}):
\begin{enumerate}
\item[(W1)] $\pffW(gk)=\pffW(g)$ for $k\in\pffKM$: $\pffch(gka_c^{-1})=\pffch(ga_c^{-1}\cdot a_cka_c^{-1})$ and $a_cka_c^{-1}\in\pffKM$;
\item[(W2)] $\pffW(kg)=\pffW(g)$ for $k\in H_T\cap\pffKM$: $H_T$ normalises $U_T$ and fixes the frame, so $\sigma(kuk^{-1})=\sigma(u)$;
if $ga_c^{-1}=uk'$ ($u\in U_T$, $k'\in\pffKM$) then $kga_c^{-1}=(kuk^{-1})(kk')$ and $\pffch(kga_c^{-1})=\pffch(ga_c^{-1})$, and if
$ga_c^{-1}\notin U_T\pffKM$ the same holds for $kga_c^{-1}$;
\item[(W3)] $\pffW$ is a cut Bessel function with $|\pffW(g)_c|\le1$.
\end{enumerate}
Recall $\mathcal O_p(l)=\pfrkap_\varrho(p,1)\int_{\pffG}\pffW(y)\otimes\overline{\pffW(lw_p^{-1}y)}\,dy$ for $l\in\pffG^{(\varrho)}$
(as $(w_pl^{-1})^{-1}=lw_p^{-1}$); it involves only the entries $p_a$, $a<\varrho$.

\begin{lemma}[A gauge for heads near $p$]\label{pfg:headgauge}
Let $p$ be a sequence of units and $\mathcal N'_p=\{r:r_bp_b^{-1}\in\pffOT^\times\ (b<\varrho)\}$, an open neighbourhood of $p$.
There is $\alpha$ as in Lemma~\ref{pfg:gauge} such that every $r\in\mathcal N'_p$ and its conjugate $\bar r$ are heads of length
$\varrho$ for $\Gamma_\alpha$; and $\pfrkap_\varrho(r,1)=\pfrkap_\varrho(p,1)$ for $r\in\mathcal N'_p$.
\end{lemma}
\begin{proof}
Take $\alpha=\varpi_1^a$ with $a$ so large that $\varpi_1^a(p_0\cdots p_{l-1})^{-1}\in\pffOT$ for $l\le\varrho$. For $r\in\mathcal N'_p$,
$(r_0\cdots r_{l-1})^{-1}$ is $(p_0\cdots p_{l-1})^{-1}$ times a unit of $\pffOT$, so has size $1$, and so has its conjugate
($\bar\alpha=\alpha$). And $|r_b|_T=|p_b|_T$, so the head densities agree.
\end{proof}

\begin{lemma}[Support and local constancy of $\mathcal O_p$]\label{pfg:Op}
Let $p$ be a sequence of units, $\alpha$ as in Lemma~\ref{pfg:headgauge}, and $S=\{y\in\pffG:y\text{ has size }60n^2\text{ for }
\Gamma_\alpha\}$, a compact set.
\begin{enumerate}
\item[(a)] For $r\in\mathcal N'_p$ and $l\in\pffG^{(\varrho)}$, the integrand of $\mathcal O_r(l)$ vanishes off $S$; so
$\mathcal O_r(l)$ converges absolutely and $|\mathcal O_r(l)|\le\pfrkap_\varrho(p,1)\,\#C\,\pbvol(S)$.
\item[(b)] If $\mathcal O_p(l)\ne0$ ($l\in\pffG^{(\varrho)}$), then $l\in C_p:=\{y'y^{-1}w_p:y,y'\in S\}$, a compact set.
\item[(c)] There is a neighbourhood $\mathcal N_p$ of $(p,1)$ in $\mathcal N'_p\times\pffG^{(\varrho)}$ such that
$\mathcal O_r(ls)=\mathcal O_p(l)$ for all $(r,s)\in\mathcal N_p$ and all $l\in\pffG^{(\varrho)}$.
\end{enumerate}
\end{lemma}
\begin{proof}
(a) If $\pffW(y)\ne0$ and $\pffW(lw_r^{-1}y)\ne0$, then $y\in S$ by Corollary~\ref{pfg:sizehead} ($r$ a head for $\Gamma_\alpha$,
$l\in\pffG^{(\varrho)}$). The bound follows from (W3). (b) Some $y$ has $\pffW(y)\ne0$ and $\pffW(y')\ne0$ for $y'=lw_p^{-1}y$; so
$y\in S$. As $l$ commutes with $w_p$ and $w_p^{-1}=w_{\bar p}$ (\S\ref{pfr:ss-gainsproof}), $l^{-1}w_{\bar p}^{-1}y'=
l^{-1}w_plw_p^{-1}y=y$; Corollary~\ref{pfg:sizehead} for the head $\bar p$ and $l^{-1}\in\pffG^{(\varrho)}$ gives $y'\in S$. And
$l=y'y^{-1}w_p$. (c) The map $(r,s,y)\mapsto y^{-1}w_psw_r^{-1}y$ is continuous and equals $1$ at $(p,1,y)$; as $\pffKM$ is open
and $S$ compact, there is a neighbourhood $\mathcal N_p\subseteq\mathcal N'_p\times\pffG^{(\varrho)}$ of $(p,1)$ with
$y^{-1}w_psw_r^{-1}y\in\pffKM$ for $(r,s)\in\mathcal N_p$, $y\in S$. Let $(r,s)\in\mathcal N_p$, $l\in\pffG^{(\varrho)}$, and
$y$ with $\pffW(y)\ne0$. If $\pffW(lsw_r^{-1}y)\ne0$ or $\pffW(lw_p^{-1}y)\ne0$, then $y\in S$ (by (a) for $r$ and $ls$, or for $p$
and $l$), and $lsw_r^{-1}y=lw_p^{-1}y\cdot(y^{-1}w_psw_r^{-1}y)$, so the two values agree by (W1). Hence the integrands of
$\mathcal O_r(ls)$ and $\mathcal O_p(l)$ coincide; and $\pfrkap_\varrho(r,1)=\pfrkap_\varrho(p,1)$.
\end{proof}

\begin{proposition}[The fibre pairing; Lemma~\ref{pff:fibre}(b)]\label{pfg:fibrepairing}
Assume (L1)--(L3). Let $p,q$ be heads. Then $\Pi_{p,q}$ is defined on all of $\pfgQ$, depends only on the entries of index
$<\varrho$ of $p,q$, and:
\begin{enumerate}
\item[(a)] the integral defining $\Pi_{p,q}(e)$ converges absolutely, and $\Pi_{p,q}(ke)=\Pi_{p,q}(e)$ for $k\in H_T\cap\pffKM$;
\item[(b)] $\Pi_{p,q}$ vanishes outside a compact subset of $\pfgQ$, and is locally constant on $\pfgQ$;
\item[(c)] there is a neighbourhood $\mathcal N$ of $(p,q)$ such that $\Pi_{p',q'}=\Pi_{p,q}$ on $\pfgQ$ for $(p',q')\in\mathcal N$;
\item[(d)] there is a function $\Phi_{p,q}$ on $\pbmXp_T$, locally constant with compact support and invariant under $H_T\cap
\pffKM$, with $\Phi_{p,q}=\Pi_{p,q}$ on $\pfgQ$.
\end{enumerate}
\end{proposition}
\begin{proof}
By Lemma~\ref{pfg:orbit}, $\pfgQ$ is the orbit of $z$, so $\Pi_{p,q}$ is defined everywhere on it. Fix $e=l^{-1}z$,
$l\in\pffG^{(\varrho)}$.

(a) By Lemma~\ref{pfg:Op}(a),(b), $k\mapsto\langle\mathcal O_p(kl),\mathcal O_q(kl)\rangle$ is bounded and vanishes unless
$k\in H_T\cap C_pl^{-1}$, a compact set, of finite $\mu_H$-measure. For $k\in H_T\cap\pffKM$, $ke=(lk^{-1})^{-1}z$, and
$\mathcal O_p(glk^{-1})=\mathcal O_p(gl)$: substitute $y\mapsto ky$ (Haar measure is invariant), use $\pffW(ky)=\pffW(y)$ (W2) and
$k^{-1}w_p^{-1}k=w_p^{-1}$ ($k\in\pffG^{(\varrho)}$ commutes with $w_p$).

(b) If $\Pi_{p,q}(e)\ne0$, some $k\in H_T$ has $\mathcal O_p(kl)\ne0$, so $kl\in C_p$ and $e=(kl)^{-1}z$ (as $k$ fixes $z$) lies in
the compact set $\{c^{-1}z:c\in C_p\cap\pffG^{(\varrho)}\}$. Local constancy at $e_0=l_0^{-1}z$: by Lemma~\ref{pfg:Op}(c) there
is a neighbourhood $U$ of $1$ in $\pffG^{(\varrho)}$ with $(p,\sigma)\in\mathcal N_p$ and $(q,\sigma)\in\mathcal N_q$ for $\sigma\in U$. By
Lemma~\ref{pfg:orbit}, every $e\in\pfgQ$ near $e_0$ is $\sigma^{-1}e_0=(l_0\sigma)^{-1}z$ with $\sigma\in U$; then
$\mathcal O_p(kl_0\sigma)=\mathcal O_p(kl_0)$ and likewise for $q$, for all $k\in H_T$, so $\Pi_{p,q}(e)=\Pi_{p,q}(e_0)$.

(c) Let $\mathcal N=\{(p',q'):(p',1)\in\mathcal N_p,\ (q',1)\in\mathcal N_q\}$. By Lemma~\ref{pfg:Op}(c) with $s=1$,
$\mathcal O_{p'}=\mathcal O_p$ and $\mathcal O_{q'}=\mathcal O_q$ on $\pffG^{(\varrho)}$.

(d) Let $C\subseteq\pfgQ$ be the compact set of (b) and $L^\sharp$ the lattice spanned by a basis of $\pbmXp_T$ in which
$H_T\cap\pffKM$ acts by integral matrices (the lattice basis; $\pffKM\subseteq\pffK$ is integral, and $H_T$ preserves $\pbmXp_T$).
\emph{A uniform modulus.} For each $y\in C$ choose an integral unit $u_y$ with $\Pi(x)=\Pi(y)$ for $x\in\pfgQ\cap(y+u_yL^\sharp)$
(by (b)); cover $C$ by finitely many $y_i+\varpi_1u_{y_i}L^\sharp$, and put $u=\varpi_1\prod_iu_{y_i}$. If $y\in C$, $x\in\pfgQ$ and
$x-y\in uL^\sharp$, then for $i$ with $y\in y_i+\varpi_1u_{y_i}L^\sharp$ both $x,y$ lie in $y_i+u_{y_i}L^\sharp$, so
$\Pi(x)=\Pi(y_i)=\Pi(y)$. \emph{The extension.} Put $\Phi(x)=\Pi(y)$ if there is $y\in\pfgQ$ with $\Pi(y)\ne0$ and
$x-y\in uL^\sharp$, and $\Phi(x)=0$ otherwise. It is well defined (two such $y,y'$ lie in $C$ and $y'-y\in uL^\sharp$, so
$\Pi(y')=\Pi(y)$), constant on the classes $x+uL^\sharp$, and supported in the compact set $C+uL^\sharp$. On $\pfgQ$:
if $\Pi(x)\ne0$ take $y=x$; if $\Pi(x)=0$ and some $y$ as above existed, then $\Pi(x)=\Pi(y)\ne0$, a contradiction; so
$\Phi=\Pi$ on $\pfgQ$. For $k\in H_T\cap\pffKM$: $k$ preserves $\pfgQ$, $uL^\sharp$ and the values of $\Pi$ by (a), so it maps the
set of admissible $y$ for $x$ onto that for $kx$, and $\Phi(kx)=\Phi(x)$.
\end{proof}

\begin{remark}\label{pfg:rem-top}
Proposition~\ref{pfg:fibrepairing}(c),(d) is what the proof of Proposition~\ref{pff:top}(c) uses (for $(p',q')$ near $(p,q)$,
$\mathcal D(p',q')=|\pi_N|_T^{-1}\int\Phi_{p,q}(x-N^{-1}p'_\varrho z)\,d\nu_{\beta(p'_\varrho)}(x)$, continuous in $p'_\varrho$ by
Lemma~\ref{pfg:coarea}); (d) and Lemma~\ref{pfg:coarea} are what the proofs of Proposition~\ref{pff:top}(a) and of
Theorem~\ref{pff:S} use (there $\Psi\ge0$ is the positive part of $\Rel\sum R_{pq}\Phi_{\tilde p,\tilde q}$, again locally
constant, compactly supported and $H_T\cap\pffKM$-invariant).
\end{remark}

\subsection{The place axiom (P) at $T$, and conventions of the shared file}\label{pfg:ss-P}

\begin{lemma}[Axiom (P) at $T$]\label{pfg:P}
The place $(\pffET,dx,\psi_T)$ (with modulus $|\cdot|_T$) of Theorem~\ref{pff:F} satisfies axiom (P) of
\S\ref{pfs:ss-planch} with $P=1$. (The map $w\mapsto M^{-1}w^{-1}$ of the place $\mathfrak P_T$ plays no role in (P).)
\end{lemma}
\begin{proof}
Take $G=\mathbf 1_{\pffOT}$: continuous (as $\pffOT$ is open and closed), integrable, $G(0)=1$. Then
$\widehat G(z)=\int_{\pffOT}\psi_T(\bar vz)\,dv$ is the integral of a character of the compact group $\pffOT$
($\bar v$ runs over $\pffOT$), equal to $1$ if $\psi_T(z\pffOT)=1$ and $0$ otherwise; by Lemma~\ref{pfg:psi}(a),
$\widehat G=\mathbf 1_{\pffOT}$, so $\int|\widehat G|=1$. The module $|c|_T=\prod_{w\in T}|c_w|_w$ is multiplicative,
$|\bar c|_T=|c|_T$ because conjugation permutes the places of the conjugation-stable set $T$ and maps $E_w$
isometrically onto $E_{\bar w}$, and $d(cv)=|c|_Tdv$ by definition of the module. A non-unit of $\pffET$ has some
component $0$, and $\{x:x_w=0\}$ is null for each of the finitely many $w$.
\end{proof}

\begin{remark}[Conventions of \S\ref{pfs:ss-conv}]\label{pfg:rem-conv}
(1) The shared proofs take $h$ linear in its first variable (convention (Cv1) of {\cite[\S2.1]{P51}}), so that the coordinate on $e^*_a$ is $h(\omega,e_a)$;
\S\ref{pfg:ss-conv} wrote $h(u,v)=u^*\mathsf Jv$, linear in the second variable. The two readings differ by
replacing $h(u,v)$ by $h(v,u)=\overline{h(u,v)}$. Every statement of this section is about sizes, integrality,
orthogonality, volumes, or characters evaluated on traces, and each of these is unchanged under conjugation of the
values of $h$ (the sets $\Gamma(i)$, $\pffOT$, $\pfgOh$, $\pfgOs$ are stable under conjugation, conjugation preserves
$dx$, and $\psi_T(\bar y)=\psi_T(-y)$); in the proofs, $h(e_a,\omega)$ becomes $h(\omega,e_a)$, $h(gm,d)=\bar g\,h(m,d)$
becomes $g\,h(m,d)$, with the same conclusions; the verification of (G5) in Lemma~\ref{pfg:gauge} is written directly
with the form $h_{\rm (Cv1)}$ in which (G5) is stated. In particular
Lemma~\ref{pff:gauge}, in the form Lemma~\ref{pfg:gauge}, holds verbatim under the convention of \S\ref{pfs:ss-conv}.
(2) Fact (C4) is used in \S\ref{pff:sec} and here only through Theorems~\ref{pfr:deep}, \ref{pfr:shells}, i.e.\ in the
indexing recalled in Remark~\ref{pfs:rem-corr}; nothing in Lemma~\ref{pff:gauge} or in this
section depends on the indexing. (3) ``Near at $\varrho+2$'' with $\pfrR^{(\varrho+2)}=\pfrU$: in the proof of
Lemma~\ref{pff:floormass}(b) an element $g\in\pffK$ is near at $\varrho+2$ (take $\xi=1\in\pfrU$), which is all that is
used; and $\pfrU\subseteq\pfrR^{(j)}$ for $j\le\varrho$ is property (U) of \S\ref{pfg:ss-conv}. So the axioms verified in
Lemma~\ref{pfg:gauge} and the use of nearness in \S\ref{pff:sec} are consistent with (C4) and with the definition of nearness in \S\ref{pfr:sec} (Remark~\ref{pfs:rem-corr}).
\end{remark}

\subsection{Siegel's formula: the published theorem and the normalisation}\label{pfg:ss-siegel}

\emph{Data.} $E^+$ is totally real of degree $f$, $E$ is CM, unramified over $E^+$ at every finite place, $\psi_E=
\psi_\Q\circ\mathrm{Tr}_{E/\Q}$ the standard character of $\auA_E/E$ (its components at the infinite places are
$x\mapsto e^{\mp2\pi i(x+\bar x)}$), and $\psi_{E^+}(y):=\psi_E(\tfrac12y)=\psi_\Q(\mathrm{Tr}_{E^+/\Q}y)$ for
$y\in\auA_{E^+}$, the standard character of $\auA_{E^+}/E^+$. $V=E^k$ with $h(x,y)=x^*y$, $H=\auU(V)$, and $\mathcal C$,
$h_c$, $J_X$ as in \S\ref{pff:ss-siegel}. $G'=\auU(1,1)$ is the unitary group of the split skew-hermitian plane over $E$,
$P=MN$ its Borel subgroup, $n(b)\in N(\auA)$ ($b\in\auA_{E^+}$) and $w$ the Weyl element.

\emph{The Weil representation.} Fix a splitting character $\chi$ of $\auA_E^\times/E^\times$ (with
$\chi|_{\auA_{E^+}^\times}=\varepsilon_{E/E^+}^k$). The Weil representation $\omega=\omega_{\psi,\chi}$ of
$G'(\auA)\times H(\auA)$ on the Schwartz--Bruhat space $\mathcal S(V(\auA_E))$ (Schr\"odinger model) satisfies
\[
\omega(1,h)\Phi(x)=\Phi(h^{-1}x),\quad\omega(n(b))\Phi(x)=\psi_{E^+}(b\,h(x,x))\Phi(x),\quad
\omega(w)\Phi(0)=\gamma\int_{V(\auA_E)}\Phi(x)\,dx ,
\]
where $dx$ is the self-dual measure for the pairing $\psi_E(h(x,y))$, i.e.\ the Haar measure for which $V(E)\backslash
V(\auA_E)$ has volume $1$, and $\gamma=\prod_v\gamma_v=1$ is the product of the local Weil indices; and
$\omega(m(a))\Phi(0)=\chi(a)|a|^{k/2}\Phi(0)$, which is $\Phi(0)$ for $a\in E^\times$. Put $\theta(g,h;\Phi)=\sum_{\xi\in V(E)}
\omega(g,h)\Phi(\xi)$ and $E(g,\Phi)=\sum_{\gamma\in P(E^+)\backslash G'(E^+)}\omega(\gamma g)\Phi(0)$.

\begin{theorem}[Siegel--Weil formula, convergent anisotropic case; Weil]\label{pfg:weil}
Let $V$ be anisotropic of dimension $k\ge3$. Then $\sum_\gamma|\omega(\gamma g)\Phi(0)|$ converges and is locally bounded in $g$,
and for every $\Phi\in\mathcal S(V(\auA_E))$ and $g\in G'(\auA)$
\[
\int_{H(E^+)\backslash H(\auA)}\theta(g,h;\Phi)\,dh=E(g,\Phi),
\]
$dh$ the Haar measure of total mass $1$ on the compact quotient.
\end{theorem}
This is Weil's theorem \cite[Th\'eor\`eme~5]{WeilSiegel} for the dual pair $(\auU(1,1),\auU(V))$; the condition $k>2$ is
the range of absolute convergence of the Eisenstein series on $\auU(1,1)$ at the point given by $V$ ($k>2n$ for
$\auU(n,n)$), in which Weil's hypotheses hold for anisotropic $V$ (condition (B) of \cite[no.~38]{WeilSiegel} in the hermitian case
$\varepsilon=\frac12$ is $k>2n$ with $n=1$, and the theta integral converges since $[\auU(V)]$ is compact, \cite[Proposition~8]{WeilSiegel},
$r=0$); the statement in this notation for unitary groups is recalled in \cite[Introduction, p.~242]{IchinoSW}. We use
it as stated.

\begin{lemma}[The measures are Tamagawa measures]\label{pfg:tamagawa}
\begin{enumerate}
\item[(a)] The product of Lebesgue measure $dx\,dy$ on each $E_\tau=\C$ and the Haar measure on $\auA_{E,f}$ giving
$\widehat\auO_E$ volume $\pbmvol=(2/\aurd)^f$ gives $E\backslash\auA_E$ volume $1$.
\item[(b)] The product of $|\operatorname{disc}E|^{1/4}$, Lebesgue measure on each $E^+_\tau=\R$ and the Haar measure $dz$
on $\mathrm{Herm}_f=\auA_{E^+,f}$ of \S\ref{pff:ss-siegel} (codifferent of volume $1$) gives $E^+\backslash\auA_{E^+}$
volume $1$.
\end{enumerate}
\end{lemma}
\begin{proof}
(a) $\auA_E=E+(E_\infty\times\widehat\auO_E)$ with $E\cap(E_\infty\times\widehat\auO_E)=\auO_E$, so the volume is
$\operatorname{covol}(\auO_E\subset\C^f)\cdot\pbmvol=2^{-f}|\operatorname{disc}E|^{1/2}\cdot2^f\aurd^{-f}=1$, as
$\aurd^{f}=|\operatorname{disc}E|^{1/2}$. (b) For $z\in\auA_{E^+,f}$, $\psi_{E,f}(z\widehat\auO_E)=\psi_{E^+,f}(z\,
\mathrm{Tr}_{E/E^+}\widehat\auO_E)$ and $\mathrm{Tr}_{E/E^+}\widehat\auO_E=\widehat\auO_{E^+}$ (unramified), so the
self-adjoint codifferent is the inverse different $\widehat{\mathfrak d}^{-1}_{E^+}$, and $dz(\widehat\auO_{E^+})=
N(\mathfrak d_{E^+})^{-1}=|\operatorname{disc}E^+|^{-1}$. The volume is $|\operatorname{disc}E|^{1/4}\cdot
|\operatorname{disc}E^+|^{1/2}\cdot|\operatorname{disc}E^+|^{-1}=1$, since $|\operatorname{disc}E|=|\operatorname{disc}E^+|^2$
($E/E^+$ unramified at all finite places).
\end{proof}

\begin{proposition}[Theorem~\ref{pff:siegelformula}]\label{pfg:siegel}
Theorem~\ref{pff:siegelformula} holds as stated, with the measures and the normalisation of \S\ref{pff:ss-siegel}; moreover
the $z$-integral defining $\delta(\Phi,\alpha)$ converges absolutely.
\end{proposition}
\begin{proof}
$V=E^k$ with $h(x,y)=x^*y$ is anisotropic: for $0\ne x\in E^k$ and every embedding $\tau$,
$\tau(x^*x)=\sum_i|\tau(x_i)|^2>0$; so Theorem~\ref{pfg:weil} applies.
Let $\Phi_f$ be the function $\Phi$ of Theorem~\ref{pff:siegelformula}, $\Phi_\infty(x)=\exp(-\sum_\tau x_\tau^*x_\tau)$, and
$\Phi=\Phi_\infty\otimes\Phi_f$. Measures: $dx$ as in Lemma~\ref{pfg:tamagawa}(a) on each coordinate (the self-dual
measure: covolume $1$), and $db$ the measure of Lemma~\ref{pfg:tamagawa}(b) on $\auA_{E^+}$.

\emph{Coefficient of the theta side.} For $h\in H(\auA)$, $\theta(n(b),h;\Phi)=\sum_\xi\psi_{E^+}(b\,\xi^*\xi)\Phi(h^{-1}\xi)$,
absolutely and uniformly convergent on compacta. Since $E^+\backslash\auA_{E^+}$ has volume $1$ and $\psi_{E^+}$ is
non-trivial and trivial on $E^+$, $\int_{E^+\backslash\auA_{E^+}}\theta(n(b),h;\Phi)\psi_{E^+}(-\alpha b)\,db=\sum_{\xi^*\xi=\alpha}
\Phi(h^{-1}\xi)$. For such $\xi$, $\Phi_\infty(h_\infty^{-1}\xi)=e^{-\sum_\tau\tau(\alpha)}$. Integrating over the compact
$H(E^+)\backslash H(\auA)$ (Fubini), the $\alpha$-coefficient of the left side of Theorem~\ref{pfg:weil} at $g=1$ is
$e^{-\sum\tau(\alpha)}\int\Sigma(h_f)\,dh$ with $\Sigma(h_f)=\sum_{\xi^*\xi=\alpha}\Phi_f(h_f^{-1}\xi)$ (a finite sum: $\Phi_f$ has compact
support and the $\xi$ lie in a lattice and on a compact sphere at infinity). $H(\auA)=\bigsqcup_cH(E^+)h_cH(\R)J_X$, the image of
each piece in the quotient is $(H(\R)\times J_X)$ divided by $H(E^+)\cap h_cJ_Xh_c^{-1}=1$ ({\cite[Lemma~14.5]{P46}}), so each class
has mass $1/\#\mathcal C$; $\Sigma$ is right $J_X$-invariant. Hence the coefficient is
$e^{-\sum\tau(\alpha)}\#\mathcal C^{-1}\sum_c\Sigma(h_c)$.

\emph{Coefficient of the Eisenstein side.} $P(E^+)\backslash G'(E^+)=\{1\}\sqcup\{wn(\beta):\beta\in E^+\}$ (Bruhat). The
identity term contributes $\Phi(0)\int\psi_{E^+}(-\alpha b)db=0$. By the absolute convergence of Theorem~\ref{pfg:weil} and
Tonelli, $\int_{E^+\backslash\auA}\sum_\beta|\omega(wn(\beta+b))\Phi(0)|\,db=\int_{\auA_{E^+}}|\omega(wn(b))\Phi(0)|\,db<\infty$, so
the coefficient is
\[
\mathcal E_\alpha=\int_{\auA_{E^+}}\mathcal W(b)\,db,\qquad\mathcal W(b)=\int_{V(\auA_E)}\Phi(x)\,\psi_{E^+}\bigl(b(x^*x-\alpha)\bigr)\,dx ,
\]
with $\gamma=1$. The integrand factorises: $\mathcal W(b)=\prod_\tau\mathcal W_\tau(b_\tau)\cdot\mathcal W_f(b_f)$, where
$\mathcal W_f(z)=\int\Phi_f(x)\psi_{E,f}(\tfrac12z(x^*x-\alpha))dx$ and, with $a=\tau(\alpha)>0$ and $t=b_\tau\in\R$,
\[
\mathcal W_\tau(t)=\int_{\C^k}e^{-|x|^2}e^{\mp2\pi it(|x|^2-a)}\,dx=e^{\pm2\pi ita}\Bigl(\frac\pi{1\pm2\pi it}\Bigr)^k .
\]
$\mathcal W_\tau$ is integrable ($k\ge2$) and not identically $0$, and $\int|\mathcal W|<\infty$; so $\mathcal W_f$ is integrable for
$dz$ (this is the absolute convergence of $\delta$), and $\mathcal E_\alpha=|\operatorname{disc}E|^{1/4}\prod_\tau\int_\R
\mathcal W_\tau\cdot\delta(\Phi_f,\alpha)$ by Lemma~\ref{pfg:tamagawa}(b). \emph{The archimedean density.} The image of Lebesgue
measure of $\C^k=\R^{2k}$ under $x\mapsto|x|^2$ is $\sigma(r)dr$ with $\sigma(r)=\pi^kr^{k-1}/(k-1)!$ (the ball of radius
$\sqrt r$ has volume $\pi^kr^k/k!$). So $\mathcal W_\tau(t)=\int_0^\infty F(r)e^{\mp2\pi it(r-a)}dr$ with $F(r)=e^{-r}\sigma(r)$,
continuous at $a$ and integrable, with integrable Fourier transform; by Fourier inversion
$\int_\R\mathcal W_\tau(t)\,dt=F(a)=e^{-\tau(\alpha)}\pi^k\tau(\alpha)^{k-1}/(k-1)!$.

\emph{Conclusion.} Theorem~\ref{pfg:weil} equates the two coefficients:
$e^{-\sum\tau(\alpha)}\#\mathcal C^{-1}\sum_c\Sigma(h_c)=|\operatorname{disc}E|^{1/4}\allowbreak \prod_\tau e^{-\tau(\alpha)}\allowbreak \pi^k\tau(\alpha)^{k-1}
/(k-1)!\cdot\delta(\Phi_f,\alpha)$. Cancelling $e^{-\sum\tau(\alpha)}$ and multiplying by $\#\mathcal C$ gives
Theorem~\ref{pff:siegelformula}.
\end{proof}

\begin{remark}\label{pfg:rem-siegel}
Every constant of Theorem~\ref{pff:siegelformula} is now accounted for: $\#\mathcal C$ from the class decomposition and the
freeness of $J_X$; $|\operatorname{disc}E|^{1/4}$ from Lemma~\ref{pfg:tamagawa}(b); $\pi^k\tau(\alpha)^{k-1}/(k-1)!$ from the
archimedean density; $\pbmvol^k$ enters through the measure $dx$ in $\delta$, fixed by Lemma~\ref{pfg:tamagawa}(a). The check
through {\cite[Lemma~2.21]{P51}} and \cite{GHY} in \S\ref{pff:ss-siegel} is not needed.
\end{remark}


\appendix

\section{Statements used from \texorpdfstring{\cite{P46}}{[P46]}, \texorpdfstring{\cite{P51}}{[P51]}}\label{B2:quoted}

For the convenience of the reader we reproduce here, without proofs, the statements of other papers that are used in this one, in the notation of the paper quoted. Each is headed by its place in that paper.

\begin{quotedthm}{\cite[Lemma~13.4]{P46}}{}
For admissible data the following hold.
\begin{enumerate}
\item[(S1)] $x^{10^{6}}\le x_1$, $x_1^{10^{6}}\le b$,
$4\log x_1\le(\log x_1)^{6}\le\log b$,
$(\log b)^{100}\le x$, $\log b\ge128$, and
$\log x\ge10^{14}(q_0+2d+|\aneu|)$.
\item[(S2)] $x^{11}\le2^{s}\le2^{22}x^{11}$; $l_0(s)\le2592\log x$;
$2m\le r$, so $r/2\le k\le r$.
\item[(S3)] $x^{11}\le T(s)=\#\anPa\le f\le q_1\le x^{12}$, and
$f=\varphi(q_0)\,\#\anPa$; moreover $1\le\#\anPa/\#\anPb\le\log b$.
\item[(S4)] $\anrd_E=\sqrt d\;q_1\prod_{p\mid q_0}p^{-1/(p-1)}$, and
$x^{10}\le\anrd_E\le x^{12}$.
\item[(S5)] $1\le xb/x_1\le\kappa(\tau)\le b\le a\le b^{2}$ for all $\tau$;
$b^{20}\le M\le b^{22}$;
$20\log b\le\log(M\anrd_E/2\pi)\le23\log b$; and
$m\log b\le\log N\le25\,m\log b$.
\item[(S6)] With $\Omega:=e^{\aneu}\cdot2\pi/\anrd_E$: $0<\Omega\le1$ and
$20\log b\le\log(M/\Omega)\le23\log b$.
\item[(S7)] With $n=m/2$ and $\varrho=n-1$ (so that $r=k+2\varrho+1$): $r$ is odd; every prime
dividing $q_1$ divides $q_0$; $q_0^{\,q_0}\le x$; $\varrho\ge10^{5}$;
$k\ge512000\,n^{9}$ and $k\ge1000\,(2\varrho+2)^{9}$; $k\ge2\varrho^{2}+3$; $18\varrho f\le b^{8}$;
$f\le b/1000$; $\anrd_E\le b$; and $N\le b^{60n}$. Moreover, with $\Delta=2r^{8}$:
\begin{itemize}
\item[(a)] $\varrho\ge10^{8}$ (so $n\ge100$); $n^{7}\le k$; $12n^{3}+500\le k$; $k-3n-6>10^{9}n^{3}$;
\item[(b)] $r\ge100$ (so $2r^{8}\ge500r^{3}+10r^{5}$); $3r^{3}f\le r^{7}$; $8nr\le b$; $kf\ge40$; $f\ge4000\,n$;
\item[(c)] $2f\le2^{2\varrho}$ and $2^{-\varrho}\le b^{-6}$;
\item[(d)] $(200\,f\,k\,\anrd_E)^{1000}\le b$ and $(25\,n^{3}f\,\anrd_E)^{1000}\le b$;
\item[(e)] $\sqrt d\le\anrd_E$, and $n,r,f,\sqrt d\le b$; $M\le b^{23}$;
\item[(f)] $\binom{a+k-1}{k-1}^{f}\le b^{2kf}$ (this bounds the dimension $d_\infty=\dim W_{\lambda^X}$ of
{\cite[\S14.3]{P46}}, a product of at most $f$ such binomial coefficients);
\item[(g)] $m\Delta+r\le\kappa(\tau)$ for all $\tau$; $k+2b+2m\Delta\le b^{2}$ (this bounds every
entry of $2\varsigma_\tau$, so that $L_\tau:=\max_i|2\varsigma_{\tau,i}|\le b^{2}$); and
$(r+1)^{4}(b^{2}+r+1+2n^{2})\le b^{5}$.
\end{itemize}
\end{enumerate}
\end{quotedthm}

\begin{quotedthm}{\cite[Lemma~14.5]{P46}}{ (Free levels)}
Let $n\le2n_m$ and let $J\subseteq K^\circ$ consist of elements congruent to $1$ modulo every
chosen place. Then $\Gamma\cap gJg^{-1}=1$ for all $g\in G_f$.
\end{quotedthm}

\begin{quotedthm}{\cite[Lemma~2.1]{P51}}{}
$K_R=K_U\rtimes K_H$. The group $K_H$ consists of the elements of the stabiliser $K^\circ$ of
$\auO_E^k$ in $H(\auA_{E^+,f})$ whose last row is $\equiv(0,\dots,0,1)$ modulo $\pi_N$, and
$K_H\cap J=J_X$.
\end{quotedthm}

\begin{quoteddef}{\cite[Definition~2.16]{P51}}{ (Hypothesis on the Bessel functions)}
We say that \emph{Bessel functions exist} for the data if there are a point $g_0\in G_\infty$, a
linear form $\mathsf t$ on $W_{\lambda^X}$, a constant $c_W>0$ and continuous functions
$W_\tau,W_\tau^\flat\colon G_\tau\to\pbW_\tau$ $(\tau\in\auPl)$ with the following properties, where
$W(y)=\bigotimes_\tau W_\tau(y_\tau)$ and $W^\flat(y)=\bigotimes_\tau W^\flat_\tau(y_\tau)$ for
$y\in G_\infty$.
\begin{enumerate}
\item[(B1)] $W_\tau(uhy)=\psi_\tau(u)\,\tau_\tau(h)\,W_\tau(y)$ for $u\in U_\tau$, $h\in H_\tau$,
where $\psi_\tau$ is the component of $\psi_\infty$; and the same for $W_\tau^\flat$.
\item[(B2)] $\|W_\tau\|$ is integrable and square-integrable on $U_\tau\backslash G_\tau$, and
$\int_{U_\tau\backslash G_\tau}\|W_\tau\|^2=\dim\pbW_\tau$.
\item[(B3)] For every class $c$ and $i\in I$ the series $\sum_{\gamma\in U_c\backslash\Gamma}
\bigl|W_i(\gamma_\infty g_\infty)\bigr|\,\pbone_{h_cJ}(\gamma g_f)$ is bounded on $G(\auA)$, where
$U_c=U(E^+)\cap h_cJh_c^{-1}$; and the function $P_{c,i}(g)=\sum_{\gamma\in U_c\backslash\Gamma}W_i(
\gamma_\infty g_\infty)\pbone_{h_cJ}(\gamma g_f)$ belongs to $\auAz$.
\item[(B4)] For every $\varphi\in\auS$, as functions on $G(\auA)$,
\[
\auP\varphi(g)=c_W\sum_{c\in\mathcal C}\ \sum_{\gamma\in U_c\backslash\Gamma}\pbip{W(\gamma_\infty
g_\infty)}{\varphi(h_c)}\,\pbone_{h_cJ}(\gamma g_f),
\]
where $\auP$ is defined by $g_0$, $\mathsf t$ and {\cite[(4)]{P51}}. (Replacing $\varphi(h_c)$ by
the value at $h_c$ of the translate of $\varphi$ by a fixed central element of $H(\auA_f)$ composes
$\auP$ with a unitary automorphism of $\auS$ that commutes with $\pbrpsi$, and changes nothing below.)
\item[(B5)] $W^\flat_\tau=\chi_\tau W_\tau$ with
$\chi_\tau(y)=\max\bigl(0,\min(1,2-2\ell_\tau(y)/b^{6n+2})\bigr)$, where
$\ell_\tau(y)=\inf_{s\in R_\tau}(\|sy\|_\tau+1)$ and $\|\cdot\|_\tau$ is the operator norm of
{\cite[Definition~3.14]{P51}}. So $\chi_\tau=1$ where $\ell_\tau\le\frac12b^{6n+2}$, $\chi_\tau=0$
where $\ell_\tau\ge b^{6n+2}$, and $\chi_\tau$ is continuous, $0\le\chi_\tau\le1$, and left
invariant under $R_\tau$; in particular the support of $W^\flat_\tau$ is compact modulo $R_\tau$.
\item[(B5$'$)] For each $\tau$ there is a compact set $C_\tau\subset G_\tau$ with
$\operatorname{supp}W^\flat_\tau\subseteq U_\tau C_\tau$. (This follows from (B5): $\chi_\tau(y)\ne0$ gives
$s=uh\in R_\tau$ with $\|sy\|_\tau<b^{6n+2}$, so $y\in h^{-1}u^{-1}B=u'h^{-1}B$ with $u'=h^{-1}u^{-1}h\in
U_\tau$ and $B$ the closed ball of radius $b^{6n+2}$ in $G_\tau$; take $C_\tau=H_\tau B$, compact as
$H_\tau$ is compact.)
\end{enumerate}
\end{quoteddef}

\begin{quotedthm}{\cite[Lemma~2.21]{P51}}{ (Ratio of masses)}
For $l\ge1$,
\[
\frac{\pbmass_{l+1}}{\pbmass_l}=\Bigl(\frac{\aurd^{\,l+1/2}\;l!}{(2\pi)^{l+1}}\Bigr)^{f}\cdot L\bigl(
l+1,\chi_{E/E^+}^{\,l+1}\bigr),
\]
where $\chi_{E/E^+}$ is the quadratic character of $E/E^+$ and $L(s,\chi^j_{E/E^+})$ means
$\zeta_{E^+}(s)$ for even $j$. In particular $\bigl|\log L(l+1,\cdot)\bigr|\le f\,2^{1-l}$ for
$l\ge2$.
\end{quotedthm}

\begin{quotedthm}{\cite[Lemma~3.6]{P51}}{ (The measure of \textup{(Cv3)} is a Haar measure)}
For every $\tau$ and every $x\in U_\tau$ the left translation $u\mapsto xu$ preserves $du_\tau$. So $du_\tau$
is a Haar measure of $U_\tau$ (a nilpotent Lie group, hence unimodular), and it is the Haar measure of
Lemma~\ref{prp:heis}\,\textup{(c)}: the two descriptions of $du_\tau$ in \textup{(Cv3)} agree.
\end{quotedthm}

\begin{quotedthm}{\cite[Theorem~3.16]{P51}}{ (Properties of the discrete series coefficient)}
For every $\tau\in\auPl$:
\begin{enumerate}
\item[(DS1)] \textup{(i)} $c_\tau$ is continuous, $c_\tau(g^{-1})=\aubar{c_\tau(g)}$ and $|c_\tau|
\le1$. \textup{(ii)} If $c_\tau$ is integrable, then $\int_{G_\tau}e_\tau(a)\,c_\tau(ga)\,da=c_\tau(g)$
for all $g\in G_\tau$, for the Haar measure $da$ for which $\int_{G_\tau}|c_\tau|^2=\pbdeg_\tau^{-1}$.
\textup{(iii)} This Haar measure is the one that corresponds to the measure of mass $1$ on the
compact form $\auU(r)$.
\item[(DS2)] $|c_\tau(g)|\le\bigl(2/\|g\|\bigr)^{\,j+\pbeq}$ for all $g\in G_\tau$ and all integers
$0\le j\le r^7$.
\item[(DS3)] $\pbdeg_\tau\cdot\int_{G_\tau}\bigl(32/\|g\|\bigr)^{\pbeq}\,dg\le b^{99r^2}$.
\item[(DS4)] There are a point $x_\tau\in G_\tau$ with $\pbel_\tau(x_\tau)\le b^{5n}$ and a unit
vector $v_\tau\in\pbW_\tau$ such that
\[
\Bigl|\int_{R_\tau}c_\tau(x_\tau^{-1}s\,x_\tau)\,\pbip{v_\tau}{\pi_R(s)^{-1}v_\tau}\,ds\Bigr|\ \ge\
b^{-60(r+k)^2}\,\pbvol(\pbeB_{x_\tau}).
\]
\end{enumerate}
\end{quotedthm}

\begin{quotedthm}{\cite[Proposition~3.18]{P51}}{ (The local Bessel function)}
Assume {\cite[Theorem~3.16]{P51}}. For $y\in G_\tau$ the integral
\[
W^\circ(y)=\int_{R_\tau}c_\tau(x^{-1}sy)\,\pi_R(s)^{-1}v\,ds\ \in\pbW_\tau
\]
converges absolutely, and:
\begin{enumerate}
\item[(a)] $W^\circ$ is continuous, $\|W^\circ(y)\|\le16^{\pbeq}\pbvol(\pbeB_x)$, and $W^\circ(sy)=
\pi_R(s)W^\circ(y)$ for $s\in R_\tau$.
\item[(b)] $c_\tau$ is integrable on $G_\tau$, $\int_{G_\tau}e_\tau(a)W^\circ(ya)\,da=W^\circ(y)$ for
all $y$, and $\int_{U_\tau\backslash G_\tau}\|W^\circ\|\le\int_{G_\tau}|c_\tau|$.
\item[(c)] \emph{(The norm is a period.)} $\displaystyle\pbdeg_\tau\int_{U_\tau\backslash G_\tau}\|W^
\circ(y)\|^2dy=\pbip v{W^\circ(x)}=\int_{R_\tau}c_\tau(x^{-1}sx)\pbip v{\pi_R(s)^{-1}v}\,ds$. In
particular this number is real and, by \textup{(DS4)}, at least $b^{-60(r+k)^2}\pbvol(\pbeB_x)>0$.
\end{enumerate}
Let $c_0>0$ be such that $W_\tau:=c_0W^\circ$ has $\int_{U_\tau\backslash G_\tau}\|W_\tau\|^2=d_\tau$,
and $W_\tau^\flat:=\chi_\tau W_\tau$ with $\chi_\tau$ as in {\cite[(23)]{P51}}. Put
$\mathsf A_\tau(y)=c_0\int_{R_\tau}(2/\|x^{-1}sy\|)^{\pbeq}ds$ and $\mathsf a_\tau=c_0\,16^{\pbeq}\pbvol(
\pbeB_x)$. Then:
\begin{enumerate}
\item[(d)] $\mathsf A_\tau$ is continuous and left invariant under $R_\tau$, $\mathsf A_\tau\le\mathsf
a_\tau$, and
\[
\pbels_\tau(y)^{3r^3f}\,\|W_\tau(y)\|\le\mathsf A_\tau(y)\quad(y\in G_\tau),\qquad\mathsf a_\tau\int_{
U_\tau\backslash G_\tau}\mathsf A_\tau\le b^{400r^2}\,d_\tau .
\]
\item[(e)] $\|W^\flat_\tau\|\le\|W_\tau\|$, $W^\flat_\tau(y)=0$ unless $\pbel_\tau(y)\le b^{6n+2}$, and
$\displaystyle\int_{U_\tau\backslash G_\tau}\|W_\tau^\flat-W_\tau\|^2\le b^{-2r^3f}\,d_\tau$.
\end{enumerate}
\end{quotedthm}

\begin{quotedthm}{\cite[Theorem~3.22]{P51}}{ ({\cite[Hypothesis~2.16]{P51}} holds)}
Assume {\cite[Theorem~3.16]{P51}}. Let $g_0=(x_\tau)_\tau$, let $\mathsf t=\bigotimes_\tau v_\tau\in\pbW$,
regarded as a linear form on $W_{\lambda^X}$, and let $W_\tau$, $W_\tau^\flat$, $\chi_\tau$ be as in
{\cite[Proposition~3.18]{P51}}. Then \textup{(B1)--(B4)} of {\cite[Definition~2.16]{P51}} hold, with
\[
c_W=\frac{\pbvol(R_0(\auA_f)\cap J)}{\pbvol(J)}\prod_\tau\frac{\pbdeg_\tau}{c_{0,\tau}}>0,
\]
and \textup{(B5)} holds in the following exact form:
\begin{enumerate}
\item[(B5)] $W_\tau^\flat=\chi_\tau W_\tau$, where $\chi_\tau(y)=\max\bigl(0,\min(1,2-2\pbel_\tau(y)/
b^{6n+2})\bigr)$ and $\pbel_\tau(y)=\inf_{s\in R_\tau}(\|sy\|_\tau+1)$ is the length of
{\cite[Definition~3.14]{P51}}.
\end{enumerate}
Moreover:
\begin{enumerate}
\item[(B6)] \emph{(Decay.)} {\cite[(25)]{P51}} holds, with $\mathsf A$ continuous and left
invariant under $R_\infty$.
\item[(B7)] \emph{(Reproduction.)} $\int_{G_\tau}e_\tau(a)W_\tau(ya)\,da=W_\tau(y)$ for all $\tau$
and $y\in G_\tau$.
\end{enumerate}
\end{quotedthm}

\begin{quotedthm}{\cite[Lemma~3.33]{P51}}{}
Let $\pbmD$ be a chain datum, $U=U_\pbmD$, $\psi=\psi_\pbmD$, $f=f_\pbmD$ and $\mathsf G=\auGL(V)$.
\begin{enumerate}
\item[(a)] $f(h)$ is $\#(U\cap h^{-1}Uh)$ or $0$; in particular $f\ge0$. Moreover $f(uhu')=f(h)$
for $u,u'\in U$ and $f(h^{-1})=f(h)$.
\item[(b)] $f$ is of positive type on $\mathsf G$ (in the sense of {\cite[\S2.3]{P51}}).
\item[(c)] Let $\tilde\psi$ be $\psi$ on $U$ and $0$ on $\mathsf G\smallsetminus U$, and
$\psi^{\pbnat}(s)=(\#\mathsf G)^{-1}\sum_{g\in\mathsf G}\tilde\psi(gsg^{-1})$. Then
$\sum_{s\in\mathsf G}|\psi^{\pbnat}(s)|^2=\pbmPhi(\pbmD)$.
\end{enumerate}
\end{quotedthm}

\begin{quotedthm}{\cite[Theorem~3.39]{P51}}{ ({\cite[Theorem~2.4]{P51}}\,\textup{(E4)})}
Assume $b\ge3$, $k\ge2\varrho^2+3$, $\auNm w\ge b^{10}$ for the chosen places $w$, and $18\,\varrho f\le b^8$
(these hold for the data of {\cite[\S2.1]{P51}}: $\varrho<2\sqrt x$, $k\ge x^5/2$, $f\le x^{12}$, $p_i\ge b^{10}$, and $b\ge1000$ by {\cite[Lemma~13.4(S1)]{P46}}). Then
$\sum_{s\in\auG}|\psi^{\pbnat}(s)|^2\le2$.
\end{quotedthm}

\begin{quotedthm}{\cite[Proposition~3.41]{P51}}{ (The big cell)}
Let $\gamma$ be a point of $G$ with values in $E^+$ or in a completion $E^+_v$; its matrix has
entries in $F=E$ or $F=E\otimes_{E^+}E^+_v$. Suppose that $\delta_l(\gamma)$ is invertible in $F$ for
$1\le l\le\varrho$. Then there are unique $u,u'\in U$, a
head $p$ and $\sigma\in\pbmGf$ with $\gamma=u\,w_p\,\sigma\,u'$; moreover $p_a=\delta_{a+1}(\gamma)/
\delta_a(\gamma)$ and $t(\sigma)=\delta_n(\gamma)/\delta_\varrho(\gamma)$. In particular $U\cap\gamma
U\gamma^{-1}=1$. The maps $\gamma\mapsto u,u',p,\sigma$ are continuous on the open set where the
$\delta_l$, $l\le\varrho$, are invertible. If $v$ is a finite place, $\gamma$ stabilises the lattice
of $Y_0$ at $v$ and the $\delta_l(\gamma)$, $l\le\varrho$, are units at $v$, then $u$, $u'$,
$\sigma$ stabilise the lattice at $v$.
\end{quotedthm}

\begin{quotedthm}{\cite[Proposition~3.42]{P51}}{ (Orbit sums are Kloosterman sums)}
Let $\gamma_0$ be in the big cell and
\[
\pbmKl_{c,s,c'}(\gamma_0)=\sum_{\substack{u\in U(E^+)/U_c,\ u'\in U(E^+)/U_{c'}\\ u^{-1}\gamma_0u'\in\Gamma_{c,s,
c'}}}\psi_f(u)\,\aubar{\psi_f(u')} .
\]
If $O(\gamma_0)\ne0$ the sum is finite and $\pbmOrb_{c,s,c'}(\gamma_0)=\pbmKl_{c,s,c'}(\gamma_0)\cdot
O(\gamma_0)$; if $O(\gamma_0)=0$ then $\pbmOrb_{c,s,c'}(\gamma_0)=0$.
\end{quotedthm}

\begin{quoteddef}{\cite[Definition~3.44]{P51}}{ (Sieve system)}
Let $\mathfrak P$ be a place, $\pbmE\ge0$, $\beta\colon C^\times\to[0,\infty)$ and $\Lambda_i\colon
(C^\times)^i\to[0,\infty)$ for $0\le i\le\varrho$. A \emph{sieve system} at $\mathfrak P$ with mass
$\pbmE$, weights $\Lambda$ and decay $\beta$ is a family of functions $J_i\colon(C^\times)^i\times
(C^\times)^i\to\C$, $0\le i\le n=\varrho+1$, such that:
\begin{enumerate}
\item[(S1)] for $0\le i\le\varrho$, heads $p,q$ of length $i$ and $\xi\in C$, the function $c\mapsto
\chi(\xi c)J_{i+1}((p,c),(q,c))$ is integrable on $C$; $\int_CJ_{i+1}((p,c),(q,c))\,dc=J_i(p,q)$; and
$\bigl|\int_C\chi(\xi c)J_{i+1}((p,c),(q,c))\,dc\bigr|\le\beta(\xi)\Lambda_i(p)\Lambda_i(q)$ for
$\xi\in C^\times$;
\item[(S2)] for heads $p^{(1)},\dots,p^{(\kappa)}$ of length $\varrho$ and $t\in C^\times$ the matrix
$\bigl(J_n((p^{(a)},t),(p^{(b)},t))\bigr)_{a,b}$ is hermitian and positive semidefinite;
\item[(S3)] $J_n$ is continuous;
\item[(S4)] $|J_0|\le\pbmE$;
\item[(S5)] for $1\le i\le\varrho$ and heads $p,q$ of length $i$ such that $\xi=\iota(p_{i-1})-\iota(
q_{i-1})$ is a unit, $|J_i(p,q)|\le\beta(\xi)\Lambda_i(p)\Lambda_i(q)$.
\end{enumerate}
\end{quoteddef}

\begin{quotedthm}{\cite[Theorem~3.50]{P51}}{ (The sieve)}
Assume $n\ge3$, $b\ge1000$, $2n\le b$, $1\le|\kappa_\tau|\le b^{23}$ for all $\tau$, $N\le b^{60n}$
and $\aurd\le b$. Let $T$ and $\pbmT$ be as in {\cite[Theorem~3.48]{P51}}, with $M^{-1}\in\pbmRT$. For
each $\tau\in\auPl$ let $(J^\tau_i)$ be a sieve system at $(\C,\text{Lebesgue},e^{c-\aubar c},w\mapsto
w^{-1})$ with mass $\pbmE_\tau$, weights $\Lambda^\infty$ and decay $\beta^\infty$, and let $(J^T_i)$
be a sieve system at $(E_T,dx,\psi_T,w\mapsto M^{-1}w^{-1})$ with mass $\pbmE_T$, weights $\Lambda^T$
and decay $\beta^T$. Then for every $\mathsf U\subseteq\pbmT$
\[
\sum_{t\in\mathsf U}\ \sum_{p,q\in\pbmT^\varrho}\Rel\Bigl(\prod_\tau J^\tau_n\bigl([p,t]_\tau,[q,t]_\tau
\bigr)\cdot J^T_n\bigl((p,t),(q,t)\bigr)\Bigr)\le\prod_{a=0}^{\varrho}c_a\cdot\Bigl(\prod_\tau\pbmE_\tau
\cdot\pbmE_T+b^{2403\,n^9f}\,\pbmdec\Bigr).
\]
\end{quotedthm}

\begin{quotedthm}{\cite[Theorem~3.51]{P51}}{ (Local systems and Siegel's step)}
Bessel functions exist for the data ({\cite[Definition~2.16]{P51}}), with \textup{(B5$'$)}, and they
can be chosen, together with functions $J^\tau_i$ $(\tau\in\auPl)$ and $J^T_i$ $(0\le i\le n$; the
$J^T_i$ depend on $T$, $\mathsf A$, $\mathsf d$ as in \textup{(F)}$)$, such that the following
holds, for the entry set $\pbmT$ of {\cite[\S3.3.2]{P51}}.
\begin{enumerate}
\item[(R)] \emph{(Real places.)} For every $\tau\in\auPl$, $(J^\tau_i)_{0\le i
\le n}$ is a sieve system at the place $(\C,\text{Lebesgue},e^{c-\aubar c},w\mapsto w^{-1})$ with mass $(1+b^{-10})/
\pbdeg_\tau$, weights $\Lambda^\infty$ and decay $\beta^\infty$; and its top form is the Gram form
$J^\tau_n(x,x')=\langle\mathsf N(x),\mathsf N(x')\rangle$ of the normalised orbital integrals of
{\cite[\S3.3.7]{P51}} (Theorem~\ref{pfr:R}\,(R-b)).
\item[(F)] \emph{(Finite places.)} Let $T$ be a finite set of finite places of $E$, stable under
conjugation, that contains the places above $MN$ and those at which some element of the set
$\mathsf B$ of {\cite[\S3.3.2]{P51}} is not a unit. For every set $\mathsf A$ of chosen places above
$p_0$ and every family $\mathsf d$ of frame vectors, $(J^{T}_i)_{0\le i\le
n}$ is a sieve system at the place $(E_T,dx,\psi_T,w\mapsto M^{-1}w^{-1})$ with mass $(1+b^{-6})\,b^{-6k\#\mathsf A}$,
weights $\Lambda^T$ and decay $\beta^T$.
\item[(S)] \emph{(Siegel's step.)} Let $T$, $\mathsf A$, $\mathsf d$ be as in \textup{(F)}, and let $J^T_n$
be the top form of \textup{(F)} for the same $T$, $\mathsf A$, $\mathsf d$ (it is built from a group
$\mathsf Z(\mathsf A,\mathsf d)$ that contains the components at $T$ of the integral elements whose
reductions lie in $\pbFix(\mathsf A,\mathsf d)$; see Theorem~\ref{pff:F}). Let
$(c_j,c_j',s_j,\sigma_j)_{j\in\mathsf I}$ be a finite family of pairwise distinct quadruples of
classes $c_j,c_j'$, elements $s_j\in\pbFix(\mathsf A,\mathsf d)$ and $\sigma_j\in\pbmGf(E^+)$ such
that the last entry $t_j=t(\sigma_j)$ lies in $\pbmT$, $|\tau(t_j)|>N$ for all $\tau$, and
$h_{c_j}^{-1}\sigma_jh_{c_j'}$ stabilises the lattice at every finite place outside $T$. Then
\[
\pbmrS^{\,2}\sum_{j\in\mathsf I}\Bigl\|\sum_{p\in\pbmT^\varrho}\pbmOrb_{c_j,s_j,c_j'}(w_p\sigma_j)\Bigr\|^2\le
(1+b^{-6})\,\pbmaS\sum_{t\in\{t_j\}}\ \sum_{p,q\in\pbmT^\varrho}\Rel\Bigl(\prod_\tau J^\tau_n\bigl([p,t]_\tau,
[q,t]_\tau\bigr)\,J^T_n\bigl((p,t),(q,t)\bigr)\Bigr).
\]
\end{enumerate}
\end{quotedthm}

\begin{quotedthm}{\cite[Lemma~3.52]{P51}}{ (Constants)}
Assume $k\ge512000\,n^9$, $n\ge100$, $b\ge1000$, $f\le b/1000$, $1\le|\kappa_\tau|$ for all
$\tau$, $M\le b^{23}$, $N\le b^{60n}$, $d\le b^{2kf}$ and $\prod_\tau|\kappa_\tau|=(M\aurd/2\pi)^f$.
Then:
\begin{enumerate}
\item[(C0)] $\displaystyle\pbmaS\prod_{a=0}^{\varrho}c_a\le(1+b^{-9})\,\pbmrS^{\,2}\,\pbdeg\cdot\pbDmain\,
\pbmnuS^{\,2}$;
\item[(C1)] $\displaystyle\pbmaS\prod_{a=0}^{\varrho}c_a\cdot b^{2403\,n^9f}\,\pbmdec\le b^{-10}\,b^{-6kf}\,
\pbmrS^{\,2}\,\pbDmain\,\pbmnuS^{\,2}$.
\end{enumerate}
\end{quotedthm}

\begin{quotedthm}{\cite[Theorem~3.53]{P51}}{ ({\cite[Theorem~2.23]{P51}}\,\textup{(L2)})}
Assume {\cite[Theorem~3.51]{P51}} and the inequalities on the data listed in {\cite[Theorem~3.50]{P51}}
and {\cite[Lemma~3.52]{P51}}. Then, for the Bessel functions of {\cite[Theorem~3.51]{P51}}, every set
$\mathsf A$ of chosen places above $p_0$ and every family $\mathsf d$ of frame vectors,
\[
\pbLM\bigl(\{s\in\pbFix(\mathsf A,\mathsf d):s\text{ open}\}\bigr)\le(1+b^{-4})\,b^{-6k\#\mathsf A}\,\pbDmain
\,\nu^2 .
\]
\end{quotedthm}

\begin{quotedthm}{\cite[Theorem~7.2]{P51}}{ (\textup{(DS4)})}
There are $x_\tau\in G_\tau$ with $\pbel_\tau(x_\tau)\le b^{5n}$ and a unit vector $v_\tau\in
\pbW_\tau$ such that
\[
\pbvol(\pbeB_{x_\tau})\le b^{60(r+k)^2}\,\Bigl|\int_{R_\tau}c_\tau(x_\tau^{-1}s\,x_\tau)\,
\pbip{v_\tau}{\pi_R(s)^{-1}v_\tau}\,ds\Bigr| .
\]
\end{quotedthm}

\begin{quotedthm}{\cite[Lemma~7.12]{P51}}{ (Witt's theorem in the form used)}
Let $F$ be a field of characteristic $\ne2$ with an involution $a\mapsto\aubar a$ \textup(possibly
trivial\textup), $F_0$ its fixed field, and $V,V'$ finite-dimensional $F$-spaces with non-degenerate
hermitian forms $h,h'$ \textup(linear in the first variable\textup) that are isometric.
\textup{(a)} If $e\in V$, $e'\in V$ and $h(e,e)=h(e',e')\ne0$, some isometry $g$ of $V$ has $ge=e'$.
\textup{(b)} Every injective linear isometry $\phi\colon S\to V'$ from a subspace $S\subseteq V$ extends to
a bijective isometry $V\to V'$.
\end{quotedthm}

\begin{quotedthm}{\cite[Lemma~8.80]{P51}}{ (Gram fibre normalisation \textup{(GF)})}
Let $dg$ be the Haar measure of \textup{(DS1)(iii)} and $\sigma$ a real-analytic section of $\Phi$ near
$\mathsf J_1$ with $\sigma(\mathsf J_1)=1$ and invertible values (e.g.\ $\sigma(T)=(\mathsf J_1T)^{1/2}$, the
power series square root near $1$). Then $(g,T)\mapsto g\sigma(T)$ carries $\frac{\pboV(r)}{2^r}\,dg\,dT$ to
Lebesgue measure $dY$ on a neighbourhood of $G=\{\Phi=\mathsf J_1\}$. Consequently every measure
$\mu\preceq\lambda_1$ satisfies $\mu\le\frac{\pboV(r)}{2^r}dg$.
\end{quotedthm}

\begin{quotedthm}{\cite[Lemma~8.81]{P51}}{ (Appending vectors)}
\begin{enumerate}
\item[(a)] Let $E=(E_1|E_2)$ be a tuple split into two parts, let $\mu\preceq\lambda_{E_1}$, and let
$\nu$ be a measure on the stabiliser $G_{E_1}$ of $E_1$ such that the image of $\nu$ under $u\mapsto
uE_2$ is at most the fibre measure, over $(E_1^*\mathsf J_1E_2,\Phi(E_2))$, of the map $X\mapsto(E_1^*
\mathsf J_1X,\Phi(X))$ on tuples $X$ of the size of $E_2$. Then $\mu*\nu\preceq\lambda_E$.
\item[(b)] Let $\mu\preceq\lambda_{\pboF}$ and let $ds=du^0\,dh$ be the measure on $\pboR$, $du^0$ the
Haar measure on $\pboU$ that is the image of Lebesgue measure on the parameters $q=(\beta,\sigma)$
($\beta\in M_{\varrho\times k}(\C)$, $\sigma$ real skew-hermitian coordinates)  $q$ of
$u_q$ in {\cite[(81)]{P51}} below, and $dh$ of mass one. Then $\mu*ds\le\frac{\pboV(r)}{2^{\varrho+1}\pboV(k)}\,dg$.
\end{enumerate}
\end{quotedthm}

\begin{thebibliography}{P46}
\bibitem[P46]{P46} D.~A.~Trejo Pizzo, \emph{The sign of the radial intersection form on the idele class group}, preprint, 2026, \url{https://github.com/dtrejopizzo/riemann/tree/main/04-papers/46-row-d-sign-radial-intersection-form}.
\bibitem[P51]{P51} D.~A.~Trejo Pizzo, \emph{The Spanning Theorem: the second moment of a Poincar\'e map}, preprint, 2026, \url{https://github.com/dtrejopizzo/riemann/tree/main/04-papers/51-spanning-theorem-second-moment-poincare-map}.
\bibitem{GHY} W.~T.~Gan, J.~P.~Hanke and J.-K.~Yu, \emph{On an exact mass formula of Shimura},
Duke Math.\ J.\ \textbf{107} (2001), 103--133.
\bibitem{Scharlau} W.~Scharlau, \emph{Quadratic and Hermitian Forms}, Grundlehren Math.\ Wiss.\
\textbf{270}, Springer, 1985.

\bibitem{IchinoSW} A.~Ichino, \emph{A regularized Siegel--Weil formula for unitary groups}, Math.\
Z.\ \textbf{247} (2004), 241--277.
\bibitem{WeilSiegel} A.~Weil, \emph{Sur la formule de Siegel dans la th\'eorie des groupes
classiques}, Acta Math.\ \textbf{113} (1965), 1--87.
\bibitem{pff-Siegel35} C. L. Siegel, \emph{\"Uber die analytische Theorie der quadratischen Formen}, Ann.\ of Math.\ (2) \textbf{36} (1935), 527--606.
\bibitem{pff-Braun41} H. Braun, \emph{Zur Theorie der hermitischen Formen}, Abh.\ Math.\ Sem.\ Hansischen Univ.\ \textbf{14} (1941), 61--150.

\end{thebibliography}

\end{document}
